\documentclass[11pt,a4paper]{article}

\usepackage[T1]{fontenc}
\usepackage{lmodern}
\usepackage[utf8]{inputenc}
\usepackage[spanish,es-nodecimaldot]{babel}
\usepackage[margin=2cm]{geometry}
\usepackage{amsmath,amssymb,amsfonts,amsthm,mathtools}
\usepackage{booktabs}
\usepackage{enumitem}
\usepackage[most]{tcolorbox}
\usepackage{xcolor}
\usepackage{fancyhdr}
\usepackage{hyperref}

% Definición de Colores Académicos
\definecolor{uniblue}{RGB}{15, 45, 95}
\definecolor{primaryblue}{RGB}{30, 64, 175}
\definecolor{darkslate}{RGB}{30, 41, 59}
\definecolor{accentgreen}{RGB}{22, 101, 52}
\definecolor{accentamber}{RGB}{180, 83, 9}
\definecolor{accentrose}{RGB}{190, 18, 60}
\definecolor{boxbg}{RGB}{248, 250, 252}

\hypersetup{
    colorlinks=true,
    linkcolor=primaryblue,
    citecolor=primaryblue,
    urlcolor=primaryblue
}

% Configuración de Encabezados y Pies de Página
\pagestyle{fancy}
\fancyhf{}
\fancyhead[L]{\small\textsc{UNI - Facultad de Ciencias} $\bullet$ \textit{Cálculo Avanzado CM2A1}}
\fancyhead[R]{\small\textbf{Práctica Calificada 4}}
\fancyfoot[C]{\thepage}
\renewcommand{\headrulewidth}{0.5pt}
\renewcommand{\footrulewidth}{0.4pt}

% Estilo de Cajas Teóricas
\newtcolorbox{theorempage}[1]{
    enhanced,
    colback=boxbg,
    colframe=uniblue,
    coltitle=white,
    fonttitle=\bfseries\sffamily,
    title={#1},
    arc=3mm,
    boxrule=1pt,
    left=4mm, right=4mm, top=3mm, bottom=3mm
}

\newtcolorbox{cautionbox}[1]{
    enhanced,
    colback=yellow!5,
    colframe=accentamber,
    coltitle=white,
    fonttitle=\bfseries\sffamily,
    title={#1},
    arc=2mm,
    boxrule=0.8pt,
    left=3mm, right=3mm, top=2mm, bottom=2mm
}

\newtcolorbox{trapbox}{
    enhanced,
    colback=red!5,
    colframe=accentrose,
    coltitle=white,
    fonttitle=\bfseries\sffamily\small,
    title={[!] Errores Frecuentes y Trampas UNI},
    arc=1.5mm,
    boxrule=0.6pt,
    left=3mm, right=3mm, top=2mm, bottom=2mm
}

\newtcolorbox{intuitionbox}[1]{
    enhanced,
    colback=cyan!4!white,
    colframe=cyan!60!blue,
    coltitle=white,
    fonttitle=\bfseries\sffamily,
    title={Explicación Sencilla y Detallada $\bullet$ #1},
    arc=2mm,
    boxrule=0.8pt,
    left=3.5mm, right=3.5mm, top=2.5mm, bottom=2.5mm
}

\begin{document}

% ==============================================================================
% PORTADA
% ==============================================================================
\begin{titlepage}
    \centering
    \vspace*{1cm}
    {\scshape\LARGE Universidad Nacional de Ingeniería \par}
    \vspace{0.4cm}
    {\scshape\Large Facultad de Ciencias $\bullet$ Escuela Profesional de Matemática \par}
    \vspace{1.5cm}
    
    \rule{\linewidth}{1.2pt}\\[0.4cm]
    {\huge\bfseries\color{uniblue} CÁLCULO DIFERENCIAL E INTEGRAL AVANZADO (CM2A1)\par}
    \vspace{0.3cm}
    {\Large\textbf{Compendio Teórico y Banco Maestro de Problemas Resueltos}\par}
    \vspace{0.2cm}
    {\large\textit{Preparación Integral para la Práctica Calificada 4 (Ciclo 2026-2)}\par}
    \rule{\linewidth}{1.2pt}\\[1.5cm]
    
    \begin{minipage}{0.9\textwidth}
        \small
        \textbf{Temario Dosificado Cubierto (Sesiones 20 a 24):}
        \begin{itemize}[leftmargin=*, noitemsep]
            \item \textbf{Sesión 20:} Teorema de la Función Inversa, Difeomorfismos locales vs globales y Singularidades.
            \item \textbf{Sesión 21:} Teorema de la Función Implícita, Sistemas acoplados no lineales y derivadas matriciales.
            \item \textbf{Sesión 22:} Integrales en Bloques: Rectángulos, Partición, Sumas de Darboux y Límites de Riemann.
            \item \textbf{Sesión 23:} Integrabilidad de Riemann, Continuidad uniforme y Criterio de oscilación.
            \item \textbf{Sesión 24:} Teorema de Fubini en Bloques, Funciones con Valor Absoluto y Teorema del Valor Medio.
            \item \textbf{Sesión 25 (Bonus):} Generalización a Regiones Planas de Tipo 1 y 2, y Cálculo de Áreas.
        \end{itemize}
    \end{minipage}
    
    \vfill
    {\large Basado en los exámenes pasados oficiales de la UNI (2024-II, 2025-II y 2026-I)\par}
    \vspace{0.5cm}
    {\small Lima, Perú $\bullet$ 2026\par}
\end{titlepage}

\tableofcontents
\newpage

% ==============================================================================
% PARTE I: TEORÍA RIGUROSA
% ==============================================================================
\part{Fundamentos Teóricos Rigurosos}

\section{Subtema 1: Teorema de la Función Inversa (Sesión 20)}

El Teorema de la Función Inversa establece las condiciones bajo las cuales una función diferenciable admite localmente una inversa que hereda la misma regularidad que la función directa.

\begin{theorempage}{Teorema de la Función Inversa en $\mathbb{R}^n$}
Sea $U \subseteq \mathbb{R}^n$ un conjunto abierto no vacío y sea $f: U \to \mathbb{R}^n$ una aplicación de clase $C^k$ con $k \ge 1$. Sea $x_0 \in U$ un punto tal que la matriz jacobiana $Df(x_0)$ es invertible, es decir:
\[
\det Df(x_0) \neq 0.
\]
Entonces:
\begin{enumerate}[label=\textbf{\alph*)}]
    \item Existen un entorno abierto $V \subseteq U$ de $x_0$ y un entorno abierto $W \subseteq \mathbb{R}^n$ de $y_0 = f(x_0)$ tales que la restricción $f|_V: V \to W$ es una biyección.
    \item La función inversa $g = (f|_V)^{-1}: W \to V$ es de clase $C^k$.
    \item Para todo $y \in W$, si $x = g(y)$, la matriz diferencial de la inversa es exactamente la inversa algebraica de la matriz jacobiana directa:
    \[
    Dg(y) = \left[ Df(x) \right]^{-1}.
    \]
\end{enumerate}
\end{theorempage}

\begin{cautionbox}{Trampa Conceptual Crítica: Invertibilidad Local vs Invertibilidad Global}
Que una función cumpla $\det Df(x) \neq 0$ en \textbf{todos} los puntos de su dominio garantiza que sea un difeomorfismo local alrededor de cualquier punto, pero \textbf{NO} garantiza que sea globalmente inyectiva o invertible.
\begin{itemize}[noitemsep]
    \item \textbf{Contraejemplo UNI:} $F(x,y) = (e^x\cos y, e^x\sin y)$ satisface $\det DF(x,y) = e^{2x} > 0$ en todo $\mathbb{R}^2$, pero tiene periodo $2\pi$ en $y$, de modo que $F(x, y + 2k\pi) = F(x,y)$, perdiendo la inyectividad global.
\end{itemize}
\end{cautionbox}

\begin{intuitionbox}{Intuición Geométrica de la Deformación y el Jacobiano}
\textbf{¿Qué hace una función multivariable?} Imagina que tienes una lámina de goma elástica en el plano $\mathbb{R}^2$ (o un bloque elástico en $\mathbb{R}^n$). Una función $f(x,y)$ toma esa lámina y la estira, la gira o la deforma espacialmente. La \textbf{matriz jacobiana} $Df(x_0)$ representa la «lupa matemática»: si te acercas muchísimo a un punto $x_0$, la deformación no lineal curva se ve indistinguible de una transformación lineal plana (la mejor aproximación afín local).

\textbf{¿Por qué el determinante jacobiano es la clave?} En una dimensión, $f'(x_0) \neq 0$ significa que la curva no tiene tangente horizontal (no se estaciona), por lo que crece o decrece estrictamente y podemos invertirla localmente. En varias dimensiones, $\det Df(x_0)$ mide cómo cambia el \textit{área orientada} (en 2D) o el \textit{volumen} (en 3D) de un elemento infinitesimal tras la transformación:
\begin{itemize}[noitemsep]
    \item Si $\det Df(x_0) \neq 0$, un rectángulo diminuto se transforma en un paralelogramo con área no nula: \textbf{no se aplasta a una línea recta ni colapsa a un punto}. Al preservarse la dimensión, no hay pérdida de información geométrica y podemos «deshacer» el movimiento de retorno de manera única, suave y diferenciable.
    \item Si $\det Df(x_0) = 0$, la transformación aplasta el volumen a cero (configuración singular). Es imposible recuperar el punto original porque infinitos puntos se han comprimido en el mismo sitio.
\end{itemize}

\textbf{La trampa entre Local y Global:} El teorema garantiza la existencia de la inversa únicamente en una bola abierta pequeña alrededor de $x_0$. Piensa en enrollar un pliego de papel en forma de cilindro: cualquier trozo pequeño del pliego sigue siendo plano y no se interseca consigo mismo (es un difeomorfismo local en todo punto), pero a escala global el papel se solapa una y otra vez perdiendo la inyectividad global.
\end{intuitionbox}

\section{Subtema 2: Teorema de la Función Implícita (Sesión 21)}

El Teorema de la Función Implícita (TFI) resuelve el problema de parametrizar localmente conjuntos de nivel $F(x,y) = 0$ como la gráfica de una función $y = \varphi(x)$.

\begin{theorempage}{Teorema de la Función Implícita en Bloques Matriciales}
Sea $U \subseteq \mathbb{R}^n \times \mathbb{R}^m$ un abierto y sea $F: U \to \mathbb{R}^m$ una función de clase $C^k$ ($k \ge 1$). Denotemos las variables por $(x, y) \in \mathbb{R}^n \times \mathbb{R}^m$. Sea $(x_0, y_0) \in U$ tal que:
\[
F(x_0, y_0) = 0 \quad \text{y} \quad \det D_y F(x_0, y_0) \neq 0,
\]
donde $D_y F(x_0, y_0) \in \mathbb{R}^{m \times m}$ es la submatriz jacobiana formada por las derivadas parciales respecto a las variables dependientes $y$.

Entonces existen vecindades abiertas $V \subset \mathbb{R}^n$ de $x_0$ y $W \subset \mathbb{R}^m$ de $y_0$, y una única función $\varphi: V \to W$ de clase $C^k$ tal que:
\[
\varphi(x_0) = y_0 \quad \text{y} \quad F(x, \varphi(x)) = 0, \quad \forall x \in V.
\]
Además, la matriz jacobiana $D\varphi(x)$ satisface la relación matricial:
\[
D\varphi(x) = - \left[ D_y F(x, \varphi(x)) \right]^{-1} \cdot D_x F(x, \varphi(x)).
\]
\end{theorempage}

\begin{intuitionbox}{Intuición de Conjuntos de Nivel y la Prueba de la Línea Vertical}
\textbf{El dilema del despeje algebraico:} En matemáticas aplicadas y física, las leyes fundamentales suelen venir dadas como ecuaciones de equilibrio $F(x,y) = 0$ (por ejemplo, superficies equipotenciales o restricciones energéticas no lineales). Despejar algebraicamente la variable dependiente $y = \varphi(x)$ con papel y lápiz suele ser imposible. El Teorema de la Función Implícita resuelve este problema demostrando que \textbf{no hace falta despejar $y$} para saber que existe como función suave ni para calcular todas sus derivadas.

\textbf{¿Por qué la condición clave es $\det D_y F(x_0, y_0) \neq 0$?} Geométricamente en el plano, la ecuación $F(x,y) = 0$ describe una curva de nivel cuyo vector gradiente $\nabla F = (\frac{\partial F}{\partial x}, \frac{\partial F}{\partial y})$ es ortogonal a la curva:
\begin{itemize}[noitemsep]
    \item Si $\frac{\partial F}{\partial y}(x_0, y_0) \neq 0$, el vector tangente a la curva en ese punto \textbf{no es vertical}. Al no tener pendiente infinita, la curva no se dobla hacia atrás sobre sí misma en ese instante: pasa con éxito la \textit{prueba de la línea vertical} del cálculo elemental, garantizando que localmente existe una única curva función $y = \varphi(x)$.
    \item Si $\frac{\partial F}{\partial y}(x_0, y_0) = 0$, la tangente es exactamente vertical (como en los extremos $(\pm 1, 0)$ de la circunferencia unitaria $x^2 + y^2 - 1 = 0$). En esos puntos, la curva se repliega sobre sí misma y es imposible definir $y$ unívocamente como función de $x$.
\end{itemize}

\textbf{La regla de la cadena implícita:} Al diferenciar la identidad $F(x, \varphi(x)) \equiv 0$ respecto a $x$, obtenemos $D_x F + D_y F \cdot D\varphi = 0$. Al ser $D_y F$ una matriz cuadrada invertible ($\det D_y F \neq 0$), multiplicamos por su inversa a la izquierda y deducimos directamente la fórmula matricial:
\[
D\varphi(x) = - \left[ D_y F(x, \varphi(x)) \right]^{-1} \cdot D_x F(x, \varphi(x)).
\]
\end{intuitionbox}

\section{Subtema 3: Particiones y Sumas de Darboux y Riemann (Sesión 22)}

La formulación analítica de la integral en $\mathbb{R}^2$ parte de un rectángulo cerrado $R = [a,b] \times [c,d]$.

\begin{theorempage}{Partición, Sumas de Darboux y Criterio de Orden}
Una partición $P = P_1 \times P_2$ divide $R$ en subrectángulos $R_{ij} = [x_{i-1}, x_i] \times [y_{j-1}, y_j]$ de área $\text{Área}(R_{ij}) = \Delta x_i \Delta y_j$.
Para una función acotada $f: R \to \mathbb{R}$, con $m_{ij} = \inf_{R_{ij}} f$ y $M_{ij} = \sup_{R_{ij}} f$:
\[
L(f, P) = \sum_{i,j} m_{ij} \text{Área}(R_{ij}), \qquad U(f, P) = \sum_{i,j} M_{ij} \text{Área}(R_{ij}).
\]
\textbf{Propiedad de Refinamiento:} Si $P \subseteq P'$, entonces $L(f, P) \le L(f, P')$ y $U(f, P') \le U(f, P)$.
Para cualesquiera particiones $P_1, P_2$:
\[
L(f, P_1) \le U(f, P_2) \implies \sup_P L(f, P) \le \inf_P U(f, P).
\]
\end{theorempage}

\begin{intuitionbox}{Intuición de Topografía y Columnas Lego (Darboux y Riemann)}
\textbf{Construyendo el volumen bajo una superficie:} Imagina un terreno montañoso cuya elevación sobre el suelo rectangular $R = [a,b] \times [c,d]$ viene dada por $z = f(x,y) \ge 0$. ¿Cómo calculamos el volumen exacto de tierra bajo la montaña?
\begin{itemize}[noitemsep]
    \item \textbf{La cuadrícula (Partición $P$):} Trazamos una retícula de líneas paralelas a los ejes, dividiendo el suelo en pequeñas baldosas rectangulares $R_{ij}$.
    \item \textbf{Suma Inferior $L(f, P)$ (Columnas por defecto):} Sobre cada baldosa $R_{ij}$, construimos una columna prismática cuya altura es la elevación más baja del terreno en esa baldosa: $m_{ij} = \inf_{R_{ij}} f$. El conjunto de todas estas columnas queda totalmente \textbf{por debajo} del terreno montañoso real.
    \item \textbf{Suma Superior $U(f, P)$ (Columnas por exceso):} Sobre la misma baldosa, levantamos una columna cuya altura es la cumbre más alta del terreno en esa baldosa: $M_{ij} = \sup_{R_{ij}} f$. Estas columnas cubren totalmente la montaña por arriba.
    \item \textbf{El principio de refinamiento:} Si subdividimos la cuadrícula en baldosas más pequeñas ($P'$ más fina que $P$), las columnas inferiores ganan espacio libre y se elevan ($L(f, P) \le L(f, P')$), mientras que las columnas superiores se ajustan mejor al perfil y bajan ($U(f, P') \le U(f, P)$). Ambas sumas atrapan al volumen real en una prensa hidráulica geométrica.
    \item \textbf{Suma de Riemann $S(f, P)$:} Si en vez de tomar el ínfimo o el supremo tomamos una muestra cualquiera en cada baldosa $(x_{ij}^*, y_{ij}^*)$, la altura de la columna es $f(x_{ij}^*, y_{ij}^*)$. Como $m_{ij} \le f(x_{ij}^*, y_{ij}^*) \le M_{ij}$, la suma de Riemann siempre queda atrapada entre Darboux inferior y superior: $L(f, P) \le S(f, P) \le U(f, P)$.
\end{itemize}
\end{intuitionbox}

\section{Subtema 4: Criterios y Propiedades de Integrabilidad (Sesión 23)}

\begin{theorempage}{Criterio de Riemann y Continuidad Uniforme}
$f$ es integrable Riemann sobre $R$ si y sólo si:
\[
\forall \varepsilon > 0, \ \exists P \text{ tal que } U(f, P) - L(f, P) < \varepsilon.
\]
\textbf{Teorema:} Toda función $f: R \to \mathbb{R}$ continua sobre un rectángulo compacto es integrable Riemann.
\end{theorempage}

\begin{intuitionbox}{El Criterio del Sándwich de Darboux y el Teorema de Lebesgue-Vitali}
\textbf{¿Cuándo decimos que una función es integrable?} Decimos que una función acotada $f$ es integrable Riemann cuando la «brecha de incertidumbre» entre las columnas por exceso y las columnas por defecto puede hacerse arbitrariamente pequeña:
\[
U(f, P) - L(f, P) = \sum_{i,j} (M_{ij} - m_{ij}) \text{Área}(R_{ij}) < \varepsilon.
\]
La diferencia $M_{ij} - m_{ij} = \omega(f, R_{ij})$ mide la \textit{oscilación} de la función en cada subrectángulo.

\textbf{¿Por qué toda función continua sobre un rectángulo compacto es integrable?}
En un conjunto cerrado y acotado (compacto) de $\mathbb{R}^2$, el \textbf{Teorema de Heine-Cantor} asegura que toda función continua es \textit{uniformemente continua}. Esto significa que no existen oscilaciones infinitamente empinadas ni cambios bruscos imprevistos: fijado $\varepsilon > 0$, basta con tomar baldosas con diagonal menor a un $\delta > 0$ para garantizar que en cada subrectángulo la oscilación sea $M_{ij} - m_{ij} < \frac{\varepsilon}{\text{Área}(R)}$. Al sumar todas las baldosas:
\[
U(f, P) - L(f, P) < \frac{\varepsilon}{\text{Área}(R)} \sum_{i,j} \text{Área}(R_{ij}) = \varepsilon.
\]
\textbf{¿Cuándo falla la integrabilidad? (Criterio de Lebesgue-Vitali UNI):}
Una función acotada no necesita ser continua en todas partes para ser integrable: puede tener líneas de discontinuidad o saltos escalonados, siempre que el conjunto de sus discontinuidades tenga \textit{medida cero} (como curvas diferenciables o puntos aislados). En cambio, si las discontinuidades están densamente distribuidas por todo el rectángulo (como la función indicatriz de racionales en 2D), la oscilación es $M_{ij} - m_{ij} = 1$ en absolutamente cualquier subrectángulo por fino que sea, haciendo que $U(f, P) - L(f, P) = \text{Área}(R)$ nunca tienda a cero.
\end{intuitionbox}

\section{Subtema 5: Teorema de Fubini y Teorema del Valor Medio (Sesión 24)}

\begin{theorempage}{Teorema de Fubini en Rectángulos}
Si $f: [a,b] \times [c,d] \to \mathbb{R}$ es continua, entonces:
\[
\iint_R f(x,y) \, dA = \int_a^b \left( \int_c^d f(x,y) \, dy \right) dx = \int_c^d \left( \int_a^b f(x,y) \, dx \right) dy.
\]
\end{theorempage}

\begin{theorempage}{Teorema del Valor Medio para Integrales Dobles}
Si $R \subset \mathbb{R}^2$ es compacto y conexo y $f: R \to \mathbb{R}$ es continua, existe $(x_0, y_0) \in R$ tal que:
\[
\iint_R f(x,y) \, dA = f(x_0, y_0) \cdot \text{Área}(R).
\]
\end{theorempage}

\begin{intuitionbox}{El Pan de Molde de Fubini y el Lago Nivelado del Valor Medio}
\textbf{Intuición de Fubini (Rebanadas de pan de molde):} Imagina un pan de molde artesanal cuya superficie superior tiene una forma ondulada $z = f(x,y)$. ¿Cómo medimos su volumen total?
\begin{itemize}[noitemsep]
    \item Fijamos una posición $x$ en el eje longitudinal y pasamos un cuchillo transversal perpendicular a dicho eje. El corte genera una rebanada plana cuya área transversal bidimensional es $A(x) = \int_c^d f(x,y) \, dy$.
    \item Luego sumamos (integramos) todas las rebanadas desde el extremo inicial $x=a$ hasta el extremo final $x=b$:
    \[
    \text{Volumen} = \int_a^b A(x) \, dx = \int_a^b \left( \int_c^d f(x,y) \, dy \right) dx.
    \]
    \item El poder del Teorema de Fubini radica en que \textbf{la orientación del corte no altera el pan}: cortar longitudinalmente en $x$ o transversalmente en $y$ conduce exactamente al mismo volumen numérico.
\end{itemize}

\textbf{Intuición del Teorema del Valor Medio (El lago nivelado):} Imagina una piscina rectangular con olas congeladas de agua descritas por $z = f(x,y)$. Si descongelamos las olas, el agua se redistribuye por gravedad hasta formar una superficie perfectamente plana y horizontal a una altura de reposo $h = f(x_0, y_0)$. Como el volumen total no ha cambiado, se cumple:
\[
\text{Volumen} = \iint_R f(x,y) \, dA = f(x_0, y_0) \cdot \text{Área}(R).
\]
Al ser la función continua y el dominio conexo, el Teorema del Valor Intermedio garantiza que en algún punto $(x_0, y_0)$ de las olas originales la altura real era exactamente igual a la altura final de reposo.
\end{intuitionbox}

\section{Subtema 6 (Bonus): Integración en Regiones Planas de Tipo 1 y 2 (Sesión 25)}

Ampliación del concepto de integral sobre dominios acotados no rectangulares $D \subset \mathbb{R}^2$ mediante la función característica $\chi_D$: $\iint_D f \, dA = \iint_R f \chi_D \, dA$, donde $R$ es un rectángulo envolvente que contiene a $D$.

\begin{theorempage}{Clasificación de Regiones en $\mathbb{R}^2$}
\begin{itemize}[noitemsep]
    \item \textbf{Región de Tipo 1 (Verticalmente Simple):} Acotada por constantes en $x$ y funciones continuas en $y$:
    \[
    D = \left\{ (x,y) \in \mathbb{R}^2 : a \le x \le b, \ g_1(x) \le y \le g_2(x) \right\} \implies \iint_D f \, dA = \int_a^b \left( \int_{g_1(x)}^{g_2(x)} f(x,y) \, dy \right) dx.
    \]
    \item \textbf{Región de Tipo 2 (Horizontalmente Simple):} Acotada por constantes en $y$ y funciones continuas en $x$:
    \[
    D = \left\{ (x,y) \in \mathbb{R}^2 : c \le y \le d, \ h_1(y) \le x \le h_2(y) \right\} \implies \iint_D f \, dA = \int_c^d \left( \int_{h_1(y)}^{h_2(y)} f(x,y) \, dx \right) dy.
    \]
\end{itemize}
\end{theorempage}

\begin{intuitionbox}{El Escáner de Rayos y el Cambio de Orden de Integración}
\textbf{¿Por qué salir de los rectángulos?} En geometría, física e ingeniería, los recintos reales (círculos, sectores parabólicos, triángulos) casi nunca tienen bordes paralelos a los ejes cartesianos. Para integrar sobre $D$, se coloca una «valla eléctrica» matemática $\chi_D(x,y)$ (vale 1 si $(x,y) \in D$ y 0 fuera). Como la frontera de $D$ está formada por curvas continuas (medida de Jordan cero), la función $f \cdot \chi_D$ sigue siendo integrable.

\textbf{Cómo determinar los límites de integración sin fallar (El método del rayo):}
\begin{itemize}[noitemsep]
    \item \textbf{Tipo 1 (Disparo vertical):} Fijamos un $x \in [a,b]$ y lanzamos una flecha vertical hacia arriba. La flecha \textit{entra} a la región por el piso $y = g_1(x)$ y \textit{sale} por el techo $y = g_2(x)$. La integral interna va de $g_1(x)$ a $g_2(x)$, y la externa de constantes $a$ a $b$.
    \item \textbf{Tipo 2 (Disparo horizontal):} Fijamos una altura $y \in [c,d]$ y lanzamos una flecha horizontal hacia la derecha. La flecha \textit{entra} por la pared izquierda $x = h_1(y)$ y \textit{sale} por la pared derecha $x = h_2(y)$. La integral interna va de $h_1(y)$ a $h_2(y)$, y la externa de constantes $c$ a $d$.
\end{itemize}

\textbf{Técnica clave en exámenes UNI: Inversión del orden de integración.}
Muchas integrales dobles no pueden resolverse elementalmente en el orden en que son planteadas (por ejemplo $\int_0^1 \int_y^1 e^{x^2} dx dy$, dado que $e^{x^2}$ carece de primitiva elemental respecto a $x$). Al trazar la región y cambiar a Tipo 1 ($0 \le y \le x$ con $0 \le x \le 1$), la integral pasa a ser:
\[
\int_0^1 \left( \int_0^x e^{x^2} dy \right) dx = \int_0^1 x e^{x^2} dx = \left[ \frac{1}{2} e^{x^2} \right]_0^1 = \frac{e - 1}{2}.
\]
¡El cambio de orden genera automáticamente el factor $x$ necesario para integrar por sustitución directa!
\end{intuitionbox}

\newpage
\part{Compendio de 98 Problemas Resueltos por Niveles (Nivel 0 al 5)}

\section{Teorema de la Función Inversa}

\subsection*{Problema 1.0A: Derivada de la inversa en una dimensión (Conexión 1D) \hfill \color{teal}\small[Nivel 0]}
\addcontentsline{toc}{subsection}{Problema 1.0A: Derivada de la inversa en una dimensión (Conexión 1D) (Nivel 0)}

\begin{tcolorbox}[colback=white,colframe=gray!40,arc=1mm,boxrule=0.5pt,left=3mm,right=3mm,top=2mm,bottom=2mm]
\textbf{Enunciado:} Considere la función real $f: \mathbb{R} \to \mathbb{R}$ dada por $f(x) = x^3 + 2x$.
\begin{enumerate}[label=\textbf{\alph*)}]
\item Demuestre que $f'(x) > 0$ para todo $x \in \mathbb{R}$ y concluya que $f$ admite inversa diferenciable global $f^{-1}: \mathbb{R} \to \mathbb{R}$.
\item Sabiendo que $f(1) = 3$, calcule la derivada de la inversa $(f^{-1})'(3)$ sin necesidad de despejar la fórmula explícita de $f^{-1}(y)$.
\end{enumerate}
\end{tcolorbox}

\begin{tcolorbox}[colback=yellow!4,colframe=accentamber!70,arc=1mm,boxrule=0.5pt,left=3mm,right=3mm,top=1.5mm,bottom=1.5mm]
\textbf{Pista Conceptual:} En una dimensión, el Teorema de la Función Inversa establece que $(f^{-1})'(y_0) = \frac{1}{f'(x_0)}$, donde $y_0 = f(x_0)$.
\end{tcolorbox}

\paragraph{Solución Paso a Paso:}\par
\begin{enumerate}[label=\textbf{\arabic*.}]
\item \textbf{Cálculo de la derivada directa:}
\[ f'(x) = 3x^2 + 2. \]
Como $x^2 \ge 0$ para todo $x \in \mathbb{R}$, se cumple $3x^2 + 2 \ge 2 > 0$ en toda la recta real. Al ser $f'(x) > 0$, la función es estrictamente creciente en todo su dominio, lo que garantiza que es inyectiva globalmente.
\item \textbf{Evaluación de la derivada en el punto base:}
Para $x_0 = 1$, tenemos $y_0 = f(1) = 1^3 + 2(1) = 3$. La derivada en dicho punto es:
\[ f'(1) = 3(1)^2 + 2 = 5. \]
\item \textbf{Cálculo de la derivada de la inversa:}
Por el Teorema de la Función Inversa (caso 1D):
\[ (f^{-1})'(3) = \frac{1}{f'(1)} = \frac{1}{5}. \]
¡Observe que obtuvimos el valor exacto de la derivada de la función inversa sin haber tenido que resolver la compleja ecuación cúbica $y = x^3 + 2x$!
\end{enumerate}

\begin{trapbox}
Intentar despejar algebraicamente $x$ en función de $y$, o evaluar $f'(3)$ en lugar de $f'(1)$ en el denominador. Recuerde: el punto donde se evalúa $f'$ es la preimagen $x_0 = 1$, no la imagen $y_0 = 3$.
\end{trapbox}

\vspace{0.5cm}
\hrule
\vspace{0.5cm}
\subsection*{Problema 1.0B: Matriz jacobiana e inversión en un sistema desacoplado en R$^2$ \hfill \color{teal}\small[Nivel 0]}
\addcontentsline{toc}{subsection}{Problema 1.0B: Matriz jacobiana e inversión en un sistema desacoplado en R$^2$ (Nivel 0)}

\begin{tcolorbox}[colback=white,colframe=gray!40,arc=1mm,boxrule=0.5pt,left=3mm,right=3mm,top=2mm,bottom=2mm]
\textbf{Enunciado:} Sea $f: \mathbb{R}^2 \to \mathbb{R}^2$ la transformación definida por $f(x,y) = (3x - 1, \, 2y + 5)$.
\begin{enumerate}[label=\textbf{\alph*)}]
\item Halle la matriz jacobiana $Df(x,y)$ y calcule su determinante.
\item Verifique que $\det Df(x,y) \neq 0$ en todo $\mathbb{R}^2$ y determine la matriz jacobiana de la inversa $D(f^{-1})(u,v)$.
\item Despeje directamente la inversa explícita $f^{-1}(u,v)$ y corrobore su derivada.
\end{enumerate}
\end{tcolorbox}

\begin{tcolorbox}[colback=yellow!4,colframe=accentamber!70,arc=1mm,boxrule=0.5pt,left=3mm,right=3mm,top=1.5mm,bottom=1.5mm]
\textbf{Pista Conceptual:} En un sistema lineal o desacoplado, las derivadas parciales cruzadas son cero, facilitando el cálculo directo de la matriz jacobiana diagonal.
\end{tcolorbox}

\paragraph{Solución Paso a Paso:}\par
\begin{enumerate}[label=\textbf{\arabic*.}]
\item \textbf{Matriz jacobiana:}
Las funciones componentes son $f_1(x,y) = 3x - 1$ y $f_2(x,y) = 2y + 5$.
\[ Df(x,y) = \begin{pmatrix} \frac{\partial f_1}{\partial x} & \frac{\partial f_1}{\partial y} \\[4pt] \frac{\partial f_2}{\partial x} & \frac{\partial f_2}{\partial y} \end{pmatrix} = \begin{pmatrix} 3 & 0 \\[4pt] 0 & 2 \end{pmatrix}. \]
\item \textbf{Determinante jacobiano:}
\[ \det Df(x,y) = (3)(2) - (0)(0) = 6 \neq 0. \]
Como el determinante es una constante no nula ($6 > 0$) en todo el plano, $f$ es invertible localmente en cualquier punto.
\item \textbf{Jacobiano de la inversa mediante el teorema:}
\[ D(f^{-1}) = [Df]^{-1} = \begin{pmatrix} 3 & 0 \\[4pt] 0 & 2 \end{pmatrix}^{-1} = \begin{pmatrix} 1/3 & 0 \\[4pt] 0 & 1/2 \end{pmatrix}. \]
\item \textbf{Verificación por despeje directo:}
Igualando $(u,v) = (3x - 1, 2y + 5)$:
\[ u = 3x - 1 \implies x = \frac{u+1}{3}, \qquad v = 2y + 5 \implies y = \frac{v-5}{2}. \]
Por tanto, $f^{-1}(u,v) = \left( \frac{u+1}{3}, \frac{v-5}{2} \right)$. Derivando directamente:
\[ D(f^{-1})(u,v) = \begin{pmatrix} 1/3 & 0 \\[4pt] 0 & 1/2 \end{pmatrix}, \]
coincidiendo perfectamente con la fórmula matricial del Teorema de la Función Inversa.
\end{enumerate}

\begin{trapbox}
Confundir la matriz jacobiana con el gradiente de una sola función escalar, o invertir los elementos de la matriz en lugar de invertir la matriz como operador lineal.
\end{trapbox}

\vspace{0.5cm}
\hrule
\vspace{0.5cm}
\subsection*{Problema 1.0C: Punto singular y pérdida de invertibilidad local \hfill \color{teal}\small[Nivel 0]}
\addcontentsline{toc}{subsection}{Problema 1.0C: Punto singular y pérdida de invertibilidad local (Nivel 0)}

\begin{tcolorbox}[colback=white,colframe=gray!40,arc=1mm,boxrule=0.5pt,left=3mm,right=3mm,top=2mm,bottom=2mm]
\textbf{Enunciado:} Considere la función cuadrática $f: \mathbb{R}^2 \to \mathbb{R}^2$ dada por $f(x,y) = (x^2, \, y)$.
\begin{enumerate}[label=\textbf{\alph*)}]
\item Calcule la matriz jacobiana $Df(x,y)$ y determine todos los puntos donde el determinante jacobiano se anula.
\item En el punto $p_0 = (0,1)$, verifique que $\det Df(0,1) = 0$. Explique geométricamente por qué $f$ no puede tener función inversa en ninguna vecindad alrededor de $(0,1)$.
\end{enumerate}
\end{tcolorbox}

\begin{tcolorbox}[colback=yellow!4,colframe=accentamber!70,arc=1mm,boxrule=0.5pt,left=3mm,right=3mm,top=1.5mm,bottom=1.5mm]
\textbf{Pista Conceptual:} Analice qué ocurre con los puntos $(+\varepsilon, 1)$ y $(-\varepsilon, 1)$ bajo la acción de $f$ para valores pequeños de $\varepsilon > 0$.
\end{tcolorbox}

\paragraph{Solución Paso a Paso:}\par
\begin{enumerate}[label=\textbf{\arabic*.}]
\item \textbf{Cálculo del jacobiano y determinante:}
\[ Df(x,y) = \begin{pmatrix} 2x & 0 \\[4pt] 0 & 1 \end{pmatrix} \implies \det Df(x,y) = 2x(1) - 0 = 2x. \]
El determinante se anula exactamente cuando $2x = 0$, es decir, sobre todo el eje $y$ ($x = 0$).
\item \textbf{Evaluación en $p_0 = (0,1)$:}
\[ Df(0,1) = \begin{pmatrix} 0 & 0 \\[4pt] 0 & 1 \end{pmatrix} \implies \det Df(0,1) = 0. \]
\item \textbf{Falla de la inyectividad local:}
Cualquier bola abierta centrada en $(0,1)$ de radio $r > 0$ contiene pares de puntos de la forma $p_1 = (\varepsilon, 1)$ y $p_2 = (-\varepsilon, 1)$ con $0 < \varepsilon < r$. Evaluando $f$ en ambos puntos:
\[ f(\varepsilon, 1) = (\varepsilon^2, 1) = f(-\varepsilon, 1). \]
Como $p_1 \neq p_2$ pero tienen exactamente la misma imagen, $f$ \textbf{no es inyectiva} en ninguna vecindad de $(0,1)$, por pequeña que sea. Al no ser inyectiva, es imposible definir una función inversa local $f^{-1}$, demostrando por qué la condición $\det Df \neq 0$ es indispensable.
\end{enumerate}

\begin{trapbox}
Asumir que si una sola variable tiene derivada no nula (como $\frac{\partial f_2}{\partial y} = 1$), la función completa es invertible. La invertibilidad requiere que todo el determinante sea no nulo.
\end{trapbox}

\vspace{0.5cm}
\hrule
\vspace{0.5cm}


\subsection*{Problema 1.1: Inversión local elemental en R$^2$ \hfill \color{primaryblue}\small[Nivel 1]}
\addcontentsline{toc}{subsection}{Problema 1.1: \texorpdfstring{Inversión local elemental en R$^2$}{Inversión local elemental en R^2} (Nivel 1)}

\begin{tcolorbox}[colback=white,colframe=gray!40,arc=1mm,boxrule=0.5pt,left=3mm,right=3mm,top=2mm,bottom=2mm]
\textbf{Enunciado:} Sea $f: \mathbb{R}^2 \to \mathbb{R}^2$ dada por $f(x,y) = (x^2 - y^2, 2xy)$. Verifique que $f$ es localmente invertible en $p_0 = (1,1)$, halle $q_0 = f(p_0)$ y calcule la matriz jacobiana de la inversa $D(f^{-1})(q_0)$.
\end{tcolorbox}

\begin{tcolorbox}[colback=yellow!4,colframe=accentamber!70,arc=1mm,boxrule=0.5pt,left=3mm,right=3mm,top=1.5mm,bottom=1.5mm]
\textbf{Pista Conceptual:} Calcule la matriz jacobiana $Df(x,y)$, evalúe su determinante en $(1,1)$. Si $\det Df(1,1) \neq 0$, aplique la fórmula $D(f^{-1})(q_0) = [Df(p_0)]^{-1}$.
\end{tcolorbox}

\paragraph{Solución Paso a Paso:}\par
\begin{enumerate}[label=\textbf{\arabic*.}]
\item \textbf{Cálculo de la imagen:} Evaluamos $p_0 = (1,1)$:
\begin{equation*}
q_0 = f(1,1) = (1^2 - 1^2, 2(1)(1)) = (0,2).
\end{equation*}
\end{enumerate}

\begin{enumerate}[label=\textbf{\arabic*.}]
\item \textbf{Matriz jacobiana de $f$:}
\begin{equation*}
Df(x,y) = \begin{pmatrix} \frac{\partial f_1}{\partial x} & \frac{\partial f_1}{\partial y} \\ \frac{\partial f_2}{\partial x} & \frac{\partial f_2}{\partial y} \end{pmatrix} = \begin{pmatrix} 2x & -2y \\ 2y & 2x \end{pmatrix}.
\end{equation*}
Evaluando en $p_0 = (1,1)$:
\begin{equation*}
Df(1,1) = \begin{pmatrix} 2 & -2 \\ 2 & 2 \end{pmatrix}.
\end{equation*}
\end{enumerate}

\begin{enumerate}[label=\textbf{\arabic*.}]
\item \textbf{Determinante jacobiano:}
\begin{equation*}
\det Df(1,1) = (2)(2) - (-2)(2) = 4 + 4 = 8 \neq 0.
\end{equation*}
Como $f$ es de clase $C^\infty$ y $\det Df(1,1) \neq 0$, por el Teorema de la Función Inversa existe un entorno abierto $U$ de $(1,1)$ y un entorno abierto $V$ de $(0,2)$ tal que $f|_U: U \to V$ es un difeomorfismo.
\end{enumerate}

\begin{enumerate}[label=\textbf{\arabic*.}]
\item \textbf{Jacobiano de la inversa:}
\begin{equation*}
D(f^{-1})(0,2) = [Df(1,1)]^{-1} = \frac{1}{8} \begin{pmatrix} 2 & 2 \\ -2 & 2 \end{pmatrix} = \begin{pmatrix} 1/4 & 1/4 \\ -1/4 & 1/4 \end{pmatrix}.
\end{equation*}
\end{enumerate}

\begin{trapbox}
Confundir el punto donde se evalúa el inverso. Se debe evaluar $D(f^{-1})$ en la imagen $q_0 = (0,2)$, no en $(1,1)$.
\end{trapbox}

\vspace{0.5cm}
\hrule
\vspace{0.5cm}

\subsection*{Problema 1.2: Difeomorfismo afín en R$^2$ \hfill \color{primaryblue}\small[Nivel 1]}
\addcontentsline{toc}{subsection}{Problema 1.2: \texorpdfstring{Difeomorfismo afín en R$^2$}{Difeomorfismo afín en R^2} (Nivel 1)}

\begin{tcolorbox}[colback=white,colframe=gray!40,arc=1mm,boxrule=0.5pt,left=3mm,right=3mm,top=2mm,bottom=2mm]
\textbf{Enunciado:} Considere la transformación afín $T: \mathbb{R}^2 \to \mathbb{R}^2$ dada por $T(x,y) = (2x + 3y - 1, \, x - y + 4)$. Demuestre que $T$ es un difeomorfismo global y halle la expresión explícita de $T^{-1}(u,v)$.
\end{tcolorbox}

\begin{tcolorbox}[colback=yellow!4,colframe=accentamber!70,arc=1mm,boxrule=0.5pt,left=3mm,right=3mm,top=1.5mm,bottom=1.5mm]
\textbf{Pista Conceptual:} Escriba $T(X) = AX + B$. Verifique si la matriz asociada $A$ es invertible ($\det A \neq 0$). Luego resuelva el sistema lineal para $(x,y)$.
\end{tcolorbox}

\paragraph{Solución Paso a Paso:}\par
\begin{enumerate}[label=\textbf{\arabic*.}]
\item \textbf{Representación matricial:}
\begin{equation*}
\begin{pmatrix} u \\ v \end{pmatrix} = \begin{pmatrix} 2 & 3 \\ 1 & -1 \end{pmatrix} \begin{pmatrix} x \\ y \end{pmatrix} + \begin{pmatrix} -1 \\ 4 \end{pmatrix}.
\end{equation*}
La matriz lineal es $A = \begin{pmatrix} 2 & 3 \\ 1 & -1 \end{pmatrix}$.
\end{enumerate}

\begin{enumerate}[label=\textbf{\arabic*.}]
\item \textbf{Invertibilidad:}
\begin{equation*}
\det A = (2)(-1) - (3)(1) = -2 - 3 = -5 \neq 0.
\end{equation*}
Al ser una aplicación afín con matriz no singular, $T$ es una biyección diferenciable de $\mathbb{R}^2$ en $\mathbb{R}^2$ con inversa también diferenciable, es decir, un difeomorfismo global.
\end{enumerate}

\begin{enumerate}[label=\textbf{\arabic*.}]
\item \textbf{Inversa explícita:}
\begin{equation*}
A^{-1} = -\frac{1}{5} \begin{pmatrix} -1 & -3 \\ -1 & 2 \end{pmatrix} = \begin{pmatrix} 1/5 & 3/5 \\ 1/5 & -2/5 \end{pmatrix}.
\end{equation*}
Despejando $\begin{pmatrix} x \\ y \end{pmatrix} = A^{-1} \left( \begin{pmatrix} u \\ v \end{pmatrix} - \begin{pmatrix} -1 \\ 4 \end{pmatrix} \right)$:
\begin{equation*}
x = \frac{1}{5}(u+1) + \frac{3}{5}(v-4) = \frac{u + 3v - 11}{5},
\end{equation*}
\begin{equation*}
y = \frac{1}{5}(u+1) - \frac{2}{5}(v-4) = \frac{u - 2v + 9}{5}.
\end{equation*}
Por tanto, $T^{-1}(u,v) = \left( \frac{u + 3v - 11}{5}, \frac{u - 2v + 9}{5} \right)$.
\end{enumerate}

\begin{trapbox}
Olvidar restar el vector de traslación constante antes de multiplicar por la inversa de la matriz lineal.
\end{trapbox}

\vspace{0.5cm}
\hrule
\vspace{0.5cm}

\subsection*{Problema 1.3: Mapeo a coordenadas polares y su inversa \hfill \color{primaryblue}\small[Nivel 1]}
\addcontentsline{toc}{subsection}{Problema 1.3: \texorpdfstring{Mapeo a coordenadas polares y su inversa}{Mapeo a coordenadas polares y su inversa} (Nivel 1)}

\begin{tcolorbox}[colback=white,colframe=gray!40,arc=1mm,boxrule=0.5pt,left=3mm,right=3mm,top=2mm,bottom=2mm]
\textbf{Enunciado:} Sea $\Phi(r,\theta) = (r\cos\theta, r\sin\theta)$ definida en $U = \{(r,\theta) \in \mathbb{R}^2 : r > 0, \, 0 < \theta < \pi/2\}$. Verifique que $\Phi$ es localmente invertible en $(r_0, \theta_0) = (2, \pi/4)$ y calcule $D(\Phi^{-1})(\sqrt{2}, \sqrt{2})$.
\end{tcolorbox}

\begin{tcolorbox}[colback=yellow!4,colframe=accentamber!70,arc=1mm,boxrule=0.5pt,left=3mm,right=3mm,top=1.5mm,bottom=1.5mm]
\textbf{Pista Conceptual:} Calcule $J\Phi(r,\theta)$, halle su determinante que resulta ser $r$. Luego invierta la matriz $2\times 2$.
\end{tcolorbox}

\paragraph{Solución Paso a Paso:}\par
\begin{enumerate}[label=\textbf{\arabic*.}]
\item \textbf{Matriz jacobiana:}
\begin{equation*}
D\Phi(r,\theta) = \begin{pmatrix} \cos\theta & -r\sin\theta \\ \sin\theta & r\cos\theta \end{pmatrix}.
\end{equation*}
El determinante es $\det D\Phi(r,\theta) = r\cos^2\theta + r\sin^2\theta = r$.
\end{enumerate}

\begin{enumerate}[label=\textbf{\arabic*.}]
\item \textbf{Evaluación en $(2, \pi/4)$:}
Como $r_0 = 2 > 0$, $\det D\Phi(2, \pi/4) = 2 \neq 0$.
Por el Teorema de la Función Inversa, existe inversa local diferenciable en una vecindad de $q_0 = \Phi(2, \pi/4) = (\sqrt{2}, \sqrt{2})$.
\begin{equation*}
D\Phi(2, \pi/4) = \begin{pmatrix} \frac{\sqrt{2}}{2} & -\sqrt{2} \\ \frac{\sqrt{2}}{2} & \sqrt{2} \end{pmatrix}.
\end{equation*}
\end{enumerate}

\begin{enumerate}[label=\textbf{\arabic*.}]
\item \textbf{Cálculo de la inversa:}
\begin{equation*}
D(\Phi^{-1})(\sqrt{2}, \sqrt{2}) = [D\Phi(2, \pi/4)]^{-1} = \frac{1}{2} \begin{pmatrix} \sqrt{2} & \sqrt{2} \\ -\frac{\sqrt{2}}{2} & \frac{\sqrt{2}}{2} \end{pmatrix} = \begin{pmatrix} \frac{\sqrt{2}}{2} & \frac{\sqrt{2}}{2} \\ -\frac{\sqrt{2}}{4} & \frac{\sqrt{2}}{4} \end{pmatrix}.
\end{equation*}
\end{enumerate}

\begin{trapbox}
Evaluar $\det D\Phi$ en el origen $r=0$ donde no es invertible, o intercambiar filas y columnas en la matriz adjunta al invertir.
\end{trapbox}

\vspace{0.5cm}
\hrule
\vspace{0.5cm}

\subsection*{Problema 1.4: Invertibilidad local vs global (Exponencial compleja) \hfill \color{primaryblue}\small[Nivel 2]}
\addcontentsline{toc}{subsection}{Problema 1.4: \texorpdfstring{Invertibilidad local vs global (Exponencial compleja)}{Invertibilidad local vs global (Exponencial compleja)} (Nivel 2)}

\begin{tcolorbox}[colback=white,colframe=gray!40,arc=1mm,boxrule=0.5pt,left=3mm,right=3mm,top=2mm,bottom=2mm]
\textbf{Enunciado:} Sea $F(x,y) = (e^x\cos y, \, e^x\sin y)$ en $\mathbb{R}^2$. (a) Verifique que $F$ es localmente invertible en todo punto de $\mathbb{R}^2$ y halle $D(F^{-1})(1,0)$ correspondiente a $F(0,0)$. (b) Demuestre rigurosamente que $F$ NO es inyectiva globalmente en $\mathbb{R}^2$.
\end{tcolorbox}

\begin{tcolorbox}[colback=yellow!4,colframe=accentamber!70,arc=1mm,boxrule=0.5pt,left=3mm,right=3mm,top=1.5mm,bottom=1.5mm]
\textbf{Pista Conceptual:} Examine la periodicidad en la variable $y$. Considere $F(x, y + 2\pi)$.
\end{tcolorbox}

\paragraph{Solución Paso a Paso:}\par
\begin{enumerate}[label=\textbf{\arabic*.}]
\item \textbf{Jacobiano y determinante:}
\begin{equation*}
DF(x,y) = \begin{pmatrix} e^x\cos y & -e^x\sin y \\ e^x\sin y & e^x\cos y \end{pmatrix}.
\end{equation*}
\begin{equation*}
\det DF(x,y) = e^{2x}(\cos^2 y + \sin^2 y) = e^{2x} > 0, \quad \forall (x,y) \in \mathbb{R}^2.
\end{equation*}
Dado que el determinante es estrictamente positivo en todo $\mathbb{R}^2$ y $F \in C^\infty$, por el Teorema de la Función Inversa, $F$ es un difeomorfismo local en cualquier entorno de cualquier punto $(x,y)$.
\end{enumerate}

\begin{enumerate}[label=\textbf{\arabic*.}]
\item \textbf{Inversa local en $(1,0)$ para $p_0 = (0,0)$:}
\begin{equation*}
DF(0,0) = \begin{pmatrix} 1 & 0 \\ 0 & 1 \end{pmatrix} = I_2 \implies D(F^{-1})(1,0) = I_2^{-1} = \begin{pmatrix} 1 & 0 \\ 0 & 1 \end{pmatrix}.
\end{equation*}
\end{enumerate}

\begin{enumerate}[label=\textbf{\arabic*.}]
\item \textbf{No inyectividad global:}
Para cualquier $(x,y) \in \mathbb{R}^2$, tomemos el punto $(x, y + 2\pi) \neq (x,y)$:
\begin{equation*}
F(x, y + 2\pi) = (e^x\cos(y+2\pi), e^x\sin(y+2\pi)) = (e^x\cos y, e^x\sin y) = F(x,y).
\end{equation*}
Por tanto, existen infinitos pares distintos que tienen la misma imagen. $F$ no es inyectiva globalmente en $\mathbb{R}^2$ a pesar de que su determinante jacobiano no se anula en ningún punto.
\end{enumerate}

\begin{trapbox}
Asumir que $\det Df(x) \neq 0$ en todo el espacio implica inyectividad global. Este es uno de los errores más comunes en las evaluaciones UNI.
\end{trapbox}

\vspace{0.5cm}
\hrule
\vspace{0.5cm}

\subsection*{Problema 1.5: Difeomorfismo global por componentes monótonas \hfill \color{primaryblue}\small[Nivel 2]}
\addcontentsline{toc}{subsection}{Problema 1.5: \texorpdfstring{Difeomorfismo global por componentes monótonas}{Difeomorfismo global por componentes monótonas} (Nivel 2)}

\begin{tcolorbox}[colback=white,colframe=gray!40,arc=1mm,boxrule=0.5pt,left=3mm,right=3mm,top=2mm,bottom=2mm]
\textbf{Enunciado:} Sea $F: \mathbb{R}^2 \to \mathbb{R}^2$ dada por $F(x,y) = (x^3 + 3x, \, y^3 + y)$. Demuestre que $F$ es un difeomorfismo global de $\mathbb{R}^2$ sobre $\mathbb{R}^2$.
\end{tcolorbox}

\begin{tcolorbox}[colback=yellow!4,colframe=accentamber!70,arc=1mm,boxrule=0.5pt,left=3mm,right=3mm,top=1.5mm,bottom=1.5mm]
\textbf{Pista Conceptual:} Analice la derivada de cada componente por separado $g(t) = t^3 + ct$ y verifique biyectividad unidimensional y jacobiano.
\end{tcolorbox}

\paragraph{Solución Paso a Paso:}\par
\begin{enumerate}[label=\textbf{\arabic*.}]
\item \textbf{Clase de la función y matriz jacobiana:}
Las componentes $f_1(x) = x^3 + 3x$ y $f_2(y) = y^3 + y$ son polinomios, luego $F \in C^\infty(\mathbb{R}^2, \mathbb{R}^2)$.
\begin{equation*}
DF(x,y) = \begin{pmatrix} 3x^2 + 3 & 0 \\ 0 & 3y^2 + 1 \end{pmatrix}.
\end{equation*}
\end{enumerate}

\begin{enumerate}[label=\textbf{\arabic*.}]
\item \textbf{Determinante jacobiano:}
\begin{equation*}
\det DF(x,y) = (3x^2 + 3)(3y^2 + 1) \ge 3 \cdot 1 = 3 > 0, \quad \forall (x,y) \in \mathbb{R}^2.
\end{equation*}
Por tanto, $F$ es localmente invertible en todo punto.
\end{enumerate}

\begin{enumerate}[label=\textbf{\arabic*.}]
\item \textbf{Inyectividad y sobreyectividad global:}
Nótese que $F(x,y) = (f_1(x), f_2(y))$ es desacoplada.
- Para $f_1(x) = x^3 + 3x$: $f_1'(x) = 3x^2 + 3 \ge 3 > 0$, luego $f_1$ es estrictamente creciente (inyectiva). Además $\lim_{x \to \pm\infty} f_1(x) = \pm\infty$, y por el Teorema del Valor Intermedio, $f_1(\mathbb{R}) = \mathbb{R}$ (sobreyectiva).
- Análogamente para $f_2(y) = y^3 + y$: $f_2'(y) = 3y^2 + 1 \ge 1 > 0$, estrictamente creciente y sobreyectiva en $\mathbb{R}$.
Al ser cada componente biyectiva de $\mathbb{R}$ en $\mathbb{R}$, $F$ es una biyección global de $\mathbb{R}^2$ en $\mathbb{R}^2$.
\end{enumerate}

\begin{enumerate}[label=\textbf{\arabic*.}]
\item \textbf{Conclusión:} Al ser $F$ biyectiva, $C^\infty$ y con diferencial invertible en todo punto, $F$ es un difeomorfismo global.
\end{enumerate}

\begin{trapbox}
Olvidar demostrar la sobreyectividad; un difeomorfismo global exige biyección completa entre los espacios.
\end{trapbox}

\vspace{0.5cm}
\hrule
\vspace{0.5cm}

\subsection*{Problema 1.6: Cambio de coordenadas elementales simétricas \hfill \color{primaryblue}\small[Nivel 2]}
\addcontentsline{toc}{subsection}{Problema 1.6: \texorpdfstring{Cambio de coordenadas elementales simétricas}{Cambio de coordenadas elementales simétricas} (Nivel 2)}

\begin{tcolorbox}[colback=white,colframe=gray!40,arc=1mm,boxrule=0.5pt,left=3mm,right=3mm,top=2mm,bottom=2mm]
\textbf{Enunciado:} Sea el sistema $u = x + y$, $v = xy$. Considere el punto $p_0 = (2,1)$. Verifique si la transformación $F(x,y) = (u,v)$ es localmente invertible alrededor de $p_0$ y determine los puntos $(x,y) \in \mathbb{R}^2$ donde el teorema NO es aplicable.
\end{tcolorbox}

\begin{tcolorbox}[colback=yellow!4,colframe=accentamber!70,arc=1mm,boxrule=0.5pt,left=3mm,right=3mm,top=1.5mm,bottom=1.5mm]
\textbf{Pista Conceptual:} Calcule $\det DF(x,y)$ en términos de $x-y$.
\end{tcolorbox}

\paragraph{Solución Paso a Paso:}\par
\begin{enumerate}[label=\textbf{\arabic*.}]
\item \textbf{Matriz jacobiana:}
\begin{equation*}
DF(x,y) = \begin{pmatrix} 1 & 1 \\ y & x \end{pmatrix}.
\end{equation*}
El determinante es $\det DF(x,y) = x - y$.
\end{enumerate}

\begin{enumerate}[label=\textbf{\arabic*.}]
\item \textbf{Evaluación en $p_0 = (2,1)$:}
\begin{equation*}
\det DF(2,1) = 2 - 1 = 1 \neq 0.
\end{equation*}
Por el Teorema de la Función Inversa, $F$ es un difeomorfismo local alrededor de $(2,1)$.
\end{enumerate}

\begin{enumerate}[label=\textbf{\arabic*.}]
\item \textbf{Puntos singulares:}
El teorema no garantiza invertibilidad donde $\det DF(x,y) = 0$, es decir, cuando $x - y = 0 \iff y = x$.
En la diagonal $y = x$, $u = 2x$ y $v = x^2 \implies v = u^2 / 4$. Los dos puntos $(x,y)$ y $(y,x)$ producen el mismo valor $(u,v)$, por lo que la función se pliega sobre la diagonal y pierde la invertibilidad local.
\end{enumerate}

\begin{trapbox}
Afirmar que cuando $\det DF = 0$ la función nunca es invertible sin verificar si la hipótesis es de condición suficiente pero no necesaria en casos singulares degenerados.
\end{trapbox}

\vspace{0.5cm}
\hrule
\vspace{0.5cm}

\subsection*{Problema 1.7: Cinemática directa de robot planar (Examen UNI 2025-II) \hfill \color{primaryblue}\small[Nivel 3]}
\addcontentsline{toc}{subsection}{Problema 1.7: \texorpdfstring{Cinemática directa de robot planar (Examen UNI 2025-II)}{Cinemática directa de robot planar (Examen UNI 2025-II)} (Nivel 3)}

\begin{tcolorbox}[colback=white,colframe=gray!40,arc=1mm,boxrule=0.5pt,left=3mm,right=3mm,top=2mm,bottom=2mm]
\textbf{Enunciado:} Un manipulador robótico planar de 2 eslabones de longitudes $l_1, l_2 > 0$ tiene ángulos articulares $(\theta_1, \theta_2)$. Su cinemática viene dada por $F(\theta_1, \theta_2) = (l_1\cos\theta_1 + l_2\cos(\theta_1+\theta_2), \, l_1\sin\theta_1 + l_2\sin(\theta_1+\theta_2))$. (a) Calcule la matriz jacobiana $DF(\theta_1, \theta_2)$ y su determinante. (b) Determine qué condición geométrica debe cumplir $\theta_2$ para que $F$ sea localmente invertible. (c) Para $l_1 = l_2 = 1$ y $(\theta_1, \theta_2) = (\pi/4, \pi/6)$, verifique si es invertible y halle $\det D(F^{-1})$.
\end{tcolorbox}

\begin{tcolorbox}[colback=yellow!4,colframe=accentamber!70,arc=1mm,boxrule=0.5pt,left=3mm,right=3mm,top=1.5mm,bottom=1.5mm]
\textbf{Pista Conceptual:} Use identidades trigonométricas de adición: $\cos A \sin B - \sin A \cos B = \sin(B - A)$.
\end{tcolorbox}

\paragraph{Solución Paso a Paso:}\par
\begin{enumerate}[label=\textbf{\arabic*.}]
\item \textbf{Matriz Jacobiana:}
\begin{equation*}
DF(\theta_1, \theta_2) = \begin{pmatrix} -l_1\sin\theta_1 - l_2\sin(\theta_1+\theta_2) & -l_2\sin(\theta_1+\theta_2) \\ l_1\cos\theta_1 + l_2\cos(\theta_1+\theta_2) & l_2\cos(\theta_1+\theta_2) \end{pmatrix}.
\end{equation*}
\end{enumerate}

\begin{enumerate}[label=\textbf{\arabic*.}]
\item \textbf{Cálculo del determinante jacobiano:}
Calculamos el producto cruzado:
\begin{equation*}
\det DF = (-l_1\sin\theta_1 - l_2\sin(\theta_1+\theta_2))(l_2\cos(\theta_1+\theta_2)) - (-l_2\sin(\theta_1+\theta_2))(l_1\cos\theta_1 + l_2\cos(\theta_1+\theta_2))
\end{equation*}
\begin{equation*}
= -l_1 l_2 \sin\theta_1\cos(\theta_1+\theta_2) - l_2^2 \sin(\theta_1+\theta_2)\cos(\theta_1+\theta_2) + l_1 l_2 \sin(\theta_1+\theta_2)\cos\theta_1 + l_2^2 \sin(\theta_1+\theta_2)\cos(\theta_1+\theta_2)
\end{equation*}
\begin{equation*}
= l_1 l_2 [\sin(\theta_1+\theta_2)\cos\theta_1 - \cos(\theta_1+\theta_2)\sin\theta_1] = l_1 l_2 \sin((\theta_1+\theta_2) - \theta_1) = l_1 l_2 \sin\theta_2.
\end{equation*}
\end{enumerate}

\begin{enumerate}[label=\textbf{\arabic*.}]
\item \textbf{Condición de invertibilidad local:}
Para que sea localmente invertible se requiere $\det DF \neq 0$:
\begin{equation*}
l_1 l_2 \sin\theta_2 \neq 0 \iff \theta_2 \neq k\pi, \quad k \in \mathbb{Z}.
\end{equation*}
Físicamente, $\theta_2 = 0$ o $\theta_2 = \pi$ representan las singularidades cinemáticas del robot (brazo completamente estirado o completamente doblado sobre sí mismo).
\end{enumerate}

\begin{enumerate}[label=\textbf{\arabic*.}]
\item \textbf{Evaluación en $(\pi/4, \pi/6)$ con $l_1 = l_2 = 1$:}
\begin{equation*}
\det DF(\pi/4, \pi/6) = (1)(1)\sin(\pi/6) = \frac{1}{2} \neq 0.
\end{equation*}
Como $\det DF \neq 0$, $F$ es localmente invertible en dicho punto y:
\begin{equation*}
\det D(F^{-1}) = \frac{1}{\det DF} = \frac{1}{1/2} = 2.
\end{equation*}
\end{enumerate}

\begin{trapbox}
Desarrollar algebraicamente sin simplificar con la identidad del seno de la diferencia, lo que ocasiona errores masivos de cálculo en la PC.
\end{trapbox}

\vspace{0.5cm}
\hrule
\vspace{0.5cm}

\subsection*{Problema 1.8: Solución única de sistema no lineal mediante TFI (Examen UNI 2026-I) \hfill \color{primaryblue}\small[Nivel 3]}
\addcontentsline{toc}{subsection}{Problema 1.8: \texorpdfstring{Solución única de sistema no lineal mediante TFI (Examen UNI 2026-I)}{Solución única de sistema no lineal mediante TFI (Examen UNI 2026-I)} (Nivel 3)}

\begin{tcolorbox}[colback=white,colframe=gray!40,arc=1mm,boxrule=0.5pt,left=3mm,right=3mm,top=2mm,bottom=2mm]
\textbf{Enunciado:} Considere el sistema de ecuaciones no lineales:
\begin{equation*}
\begin{cases} x + y + \cos(xy) = 1 \\ x - y - \sin(xy) = 0 \end{cases}
\end{equation*}
Demuestre rigurosamente que dicho sistema admite solución única en alguna vecindad abierta del origen de coordenadas $(0,0)$.
\end{tcolorbox}

\begin{tcolorbox}[colback=yellow!4,colframe=accentamber!70,arc=1mm,boxrule=0.5pt,left=3mm,right=3mm,top=1.5mm,bottom=1.5mm]
\textbf{Pista Conceptual:} Defina $F(x,y) = (x+y+\cos(xy)-1, \, x-y-\sin(xy))$. Compruebe $F(0,0)=(0,0)$ y calcule $\det DF(0,0)$.
\end{tcolorbox}

\paragraph{Solución Paso a Paso:}\par
\begin{enumerate}[label=\textbf{\arabic*.}]
\item \textbf{Definición de la función vectorial:}
Definimos $F: \mathbb{R}^2 \to \mathbb{R}^2$ por
\begin{equation*}
F(x,y) = \begin{pmatrix} x + y + \cos(xy) \\ x - y - \sin(xy) \end{pmatrix}.
\end{equation*}
El sistema equivale a encontrar $(x,y)$ tal que $F(x,y) = (1,0)$.
\end{enumerate}

\begin{enumerate}[label=\textbf{\arabic*.}]
\item \textbf{Verificación del punto base:}
Evaluando en el origen $p_0 = (0,0)$:
\begin{equation*}
F(0,0) = (0 + 0 + \cos(0), \, 0 - 0 - \sin(0)) = (1, 0).
\end{equation*}
Por lo tanto, $(0,0)$ es en efecto una solución del sistema.
\end{enumerate}

\begin{enumerate}[label=\textbf{\arabic*.}]
\item \textbf{Cálculo de la matriz jacobiana:}
Derivando parcialmente:
\begin{equation*}
DF(x,y) = \begin{pmatrix} 1 - y\sin(xy) & 1 - x\sin(xy) \\ 1 - y\cos(xy) & -1 - x\cos(xy) \end{pmatrix}.
\end{equation*}
Evaluando en $(0,0)$:
\begin{equation*}
DF(0,0) = \begin{pmatrix} 1 - 0 & 1 - 0 \\ 1 - 0 & -1 - 0 \end{pmatrix} = \begin{pmatrix} 1 & 1 \\ 1 & -1 \end{pmatrix}.
\end{equation*}
\end{enumerate}

\begin{enumerate}[label=\textbf{\arabic*.}]
\item \textbf{Determinante jacobiano:}
\begin{equation*}
\det DF(0,0) = (1)(-1) - (1)(1) = -2 \neq 0.
\end{equation*}
Como las funciones componentes son infinitamente diferenciables ($C^\infty$) y $\det DF(0,0) = -2 \neq 0$, por el Teorema de la Función Inversa, existe una vecindad abierta $U$ de $(0,0)$ y una vecindad abierta $V$ de $(1,0)$ tal que $F|_U: U \to V$ es una biyección con inversa diferenciable.
\end{enumerate}

\begin{enumerate}[label=\textbf{\arabic*.}]
\item \textbf{Conclusión:} Para el valor $(1,0) \in V$, existe un único $(x,y) \in U$ tal que $F(x,y) = (1,0)$, lo que prueba la existencia y unicidad local.
\end{enumerate}

\begin{trapbox}
Intentar despejar analíticamente $x$ o $y$ igualando expresiones trigonométricas en vez de invocar el Teorema de la Función Inversa.
\end{trapbox}

\vspace{0.5cm}
\hrule
\vspace{0.5cm}

\subsection*{Problema 1.9: Mapeo no acotado con jacobiano estrictamente positivo \hfill \color{primaryblue}\small[Nivel 3]}
\addcontentsline{toc}{subsection}{Problema 1.9: \texorpdfstring{Mapeo no acotado con jacobiano estrictamente positivo}{Mapeo no acotado con jacobiano estrictamente positivo} (Nivel 3)}

\begin{tcolorbox}[colback=white,colframe=gray!40,arc=1mm,boxrule=0.5pt,left=3mm,right=3mm,top=2mm,bottom=2mm]
\textbf{Enunciado:} Sea $f(x,y) = (x + e^y, \, y - e^x)$. (a) Demuestre que $Df(x,y)$ es invertible en todo punto de $\mathbb{R}^2$. (b) Demuestre que $f$ es una aplicación abierta. (c) Determine si $f$ es inyectiva en todo $\mathbb{R}^2$.
\end{tcolorbox}

\begin{tcolorbox}[colback=yellow!4,colframe=accentamber!70,arc=1mm,boxrule=0.5pt,left=3mm,right=3mm,top=1.5mm,bottom=1.5mm]
\textbf{Pista Conceptual:} Analice $\det Df(x,y) = 1 + e^{x+y}$. Para inyectividad, suponga $f(x_1, y_1) = f(x_2, y_2)$ y use monotonía.
\end{tcolorbox}

\paragraph{Solución Paso a Paso:}\par
\begin{enumerate}[label=\textbf{\arabic*.}]
\item \textbf{Matriz jacobiana:}
\begin{equation*}
Df(x,y) = \begin{pmatrix} 1 & e^y \\ -e^x & 1 \end{pmatrix}.
\end{equation*}
\begin{equation*}
\det Df(x,y) = 1\cdot 1 - (e^y)(-e^x) = 1 + e^{x+y}.
\end{equation*}
Dado que $e^{x+y} > 0$ para todos $(x,y) \in \mathbb{R}^2$, se tiene $\det Df(x,y) \ge 1 > 0$. Por el Teorema de la Función Inversa, $f$ es un difeomorfismo local en todo punto.
\end{enumerate}

\begin{enumerate}[label=\textbf{\arabic*.}]
\item \textbf{Aplicación abierta:}
Por el teorema de la aplicación abierta derivado del TFI, como $f$ es localmente invertible en todo punto, la imagen de cualquier conjunto abierto es un conjunto abierto en $\mathbb{R}^2$.
\end{enumerate}

\begin{enumerate}[label=\textbf{\arabic*.}]
\item \textbf{Inyectividad global:}
Supongamos $f(x_1, y_1) = f(x_2, y_2)$:
\begin{equation*}
x_1 + e^{y_1} = x_2 + e^{y_2} \implies x_1 - x_2 = e^{y_2} - e^{y_1}
\end{equation*}
\begin{equation*}
y_1 - e^{x_1} = y_2 - e^{x_2} \implies y_1 - y_2 = e^{x_1} - e^{x_2}
\end{equation*}
Multiplicando miembro a miembro:
\begin{equation*}
(x_1 - x_2)(e^{x_1} - e^{x_2}) + (y_1 - y_2)(e^{y_1} - e^{y_2}) = 0.
\end{equation*}
Como la función exponencial $t \mapsto e^t$ es estrictamente creciente, $(a - b)(e^a - e^b) \ge 0$, anulándose únicamente si $a = b$.
Para que la suma de dos términos no negativos sea 0, ambos deben ser 0:
\begin{equation*}
x_1 = x_2 \quad \text{y} \quad y_1 = y_2.
\end{equation*}
Por lo tanto, $f$ es estrictamente inyectiva en todo $\mathbb{R}^2$.
\end{enumerate}

\begin{trapbox}
Concluir inyectividad global únicamente basándose en que $\det Df > 0$, sin realizar la demostración de inyectividad.
\end{trapbox}

\vspace{0.5cm}
\hrule
\vspace{0.5cm}

\subsection*{Problema 1.10: Mapeo acotado en R$^n$ y radio de inyectividad (Examen UNI 2025-II) \hfill \color{primaryblue}\small[Nivel 4]}
\addcontentsline{toc}{subsection}{Problema 1.10: \texorpdfstring{Mapeo acotado en R$^n$ y radio de inyectividad (Examen UNI 2025-II)}{Mapeo acotado en R^n y radio de inyectividad (Examen UNI 2025-II)} (Nivel 4)}

\begin{tcolorbox}[colback=white,colframe=gray!40,arc=1mm,boxrule=0.5pt,left=3mm,right=3mm,top=2mm,bottom=2mm]
\textbf{Enunciado:} Sea $f: \mathbb{R}^n \to \mathbb{R}^n$ definida por $f(x) = \frac{x}{1 + \|x\|^2}$. (a) Demuestre que $f$ es de clase $C^\infty$ y acotada en todo $\mathbb{R}^n$. (b) Calcule la matriz jacobiana $Jf(0)$ y verifique que $Jf(0) = I_n$. (c) Demuestre que $f$ es inyectiva en la bola abierta $B(0, 1)$ pero NO es inyectiva en todo $\mathbb{R}^n$.
\end{tcolorbox}

\begin{tcolorbox}[colback=yellow!4,colframe=accentamber!70,arc=1mm,boxrule=0.5pt,left=3mm,right=3mm,top=1.5mm,bottom=1.5mm]
\textbf{Pista Conceptual:} Observe que $\|f(x)\| = \frac{\|x\|}{1+\|x\|^2}$. Para la inyectividad, si $f(x) = f(y)$, analice sus normas.
\end{tcolorbox}

\paragraph{Solución Paso a Paso:}\par
\begin{enumerate}[label=\textbf{\arabic*.}]
\item \textbf{Clase $C^\infty$ y acotación:}
Como el denominador $1 + \|x\|^2 \ge 1 > 0$ nunca se anula y los polinomios son $C^\infty$, el cociente es de clase $C^\infty$.
Para la acotación, tomemos la norma euclidiana:
\begin{equation*}
\|f(x)\| = \frac{\|x\|}{1 + \|x\|^2}.
\end{equation*}
Por la desigualdad entre media aritmética y geométrica: $1 + \|x\|^2 \ge 2\|x\| \implies \frac{\|x\|}{1 + \|x\|^2} \le \frac{1}{2}$.
Luego $\|f(x)\| \le \frac{1}{2}$ para todo $x \in \mathbb{R}^n$, demostrando que $f$ es acotada.
\end{enumerate}

\begin{enumerate}[label=\textbf{\arabic*.}]
\item \textbf{Cálculo de la matriz jacobiana en el origen:}
La $i$-ésima componente es $f_i(x) = \frac{x_i}{1 + \sum_{k=1}^n x_k^2}$.
Derivando parcialmente:
\begin{equation*}
\frac{\partial f_i}{\partial x_j}(x) = \frac{\delta_{ij}(1 + \|x\|^2) - x_i (2x_j)}{(1 + \|x\|^2)^2}.
\end{equation*}
Evaluando en $x = 0$:
\begin{equation*}
\frac{\partial f_i}{\partial x_j}(0) = \frac{\delta_{ij}(1) - 0}{1} = \delta_{ij} \implies Jf(0) = I_n.
\end{equation*}
Al ser $Jf(0) = I_n$, $\det Jf(0) = 1 \neq 0$, por lo que $f$ es un difeomorfismo local alrededor del origen.
\end{enumerate}

\begin{enumerate}[label=\textbf{\arabic*.}]
\item \textbf{Análisis de inyectividad:}
Supongamos $f(x) = f(y)$. Tomando normas:
\begin{equation*}
\frac{\|x\|}{1 + \|x\|^2} = \frac{\|y\|}{1 + \|y\|^2} \implies g(\|x\|) = g(\|y\|), \quad \text{donde } g(r) = \frac{r}{1 + r^2}.
\end{equation*}
La función $g'(r) = \frac{1 - r^2}{(1 + r^2)^2} > 0$ para $r \in [0, 1)$.
Por ende, $g$ es estrictamente creciente en $[0, 1)$. Si $x, y \in B(0, 1)$, entonces $\|x\|, \|y\| < 1$, lo que obliga a $\|x\| = \|y\|$.
Sustituyendo en $f(x) = f(y)$: $\frac{x}{1+\|x\|^2} = \frac{y}{1+\|x\|^2} \implies x = y$.
Así, $f$ es inyectiva en $B(0, 1)$.
Sin embargo, en todo $\mathbb{R}^n$ NO es inyectiva porque para cualquier vector unitario $u$, $g(r) = g(1/r)$:
\begin{equation*}
f(2u) = \frac{2u}{1 + 4} = \frac{2}{5}u, \qquad f(1/2 u) = \frac{(1/2)u}{1 + 1/4} = \frac{1/2}{5/4}u = \frac{2}{5}u.
\end{equation*}
Como $2u \neq \frac{1}{2}u$, $f$ no es inyectiva en $\mathbb{R}^n$.
\end{enumerate}

\begin{trapbox}
Derivar usando notación escalar en vez de vectorial sin identificar la delta de Kronecker $\delta_{ij}$.
\end{trapbox}

\vspace{0.5cm}
\hrule
\vspace{0.5cm}

\subsection*{Problema 1.11: Teorema de Hadamard-Caccioppoli sobre difeomorfismo global \hfill \color{primaryblue}\small[Nivel 4]}
\addcontentsline{toc}{subsection}{Problema 1.11: \texorpdfstring{Teorema de Hadamard-Caccioppoli sobre difeomorfismo global}{Teorema de Hadamard-Caccioppoli sobre difeomorfismo global} (Nivel 4)}

\begin{tcolorbox}[colback=white,colframe=gray!40,arc=1mm,boxrule=0.5pt,left=3mm,right=3mm,top=2mm,bottom=2mm]
\textbf{Enunciado:} Sea $f: \mathbb{R}^n \to \mathbb{R}^n$ de clase $C^1$ tal que $\det Df(x) \neq 0$ para todo $x \in \mathbb{R}^n$. Suponga además que $f$ es propia (es decir, $\|f(x)\| \to \infty$ cuando $\|x\| \to \infty$). Demuestre que $f$ es un difeomorfismo global sobre $\mathbb{R}^n$.
\end{tcolorbox}

\begin{tcolorbox}[colback=yellow!4,colframe=accentamber!70,arc=1mm,boxrule=0.5pt,left=3mm,right=3mm,top=1.5mm,bottom=1.5mm]
\textbf{Pista Conceptual:} Demuestre que $f(\mathbb{R}^n)$ es abierto y cerrado a la vez en $\mathbb{R}^n$, y aplique conexidad.
\end{tcolorbox}

\paragraph{Solución Paso a Paso:}\par
\begin{enumerate}[label=\textbf{\arabic*.}]
\item \textbf{La imagen $f(\mathbb{R}^n)$ es un conjunto abierto:}
Como $\det Df(x) \neq 0$ para todo $x \in \mathbb{R}^n$, por el Teorema de la Función Inversa, $f$ es un difeomorfismo local. Toda aplicación con inversa local continua es una aplicación abierta. Dado que $\mathbb{R}^n$ es abierto, su imagen $f(\mathbb{R}^n)$ es un abierto no vacío de $\mathbb{R}^n$.
\end{enumerate}

\begin{enumerate}[label=\textbf{\arabic*.}]
\item \textbf{La imagen $f(\mathbb{R}^n)$ es un conjunto cerrado:}
Sea $y_k \in f(\mathbb{R}^n)$ una sucesión tal que $y_k \to y \in \mathbb{R}^n$.
Por definición, existen $x_k \in \mathbb{R}^n$ tales que $f(x_k) = y_k$.
Dado que $\{y_k\} \cup \{y\}$ es un conjunto compacto en $\mathbb{R}^n$ y $f$ es una aplicación propia, la preimagen $f^{-1}(\{y_k\} \cup \{y\})$ es compacta.
Por tanto, la sucesión $\{x_k\}$ está acotada. Por el Teorema de Bolzano-Weierstrass, posee una subsucesión convergente $x_{k_j} \to x^* \in \mathbb{R}^n$.
Por la continuidad de $f$:
\begin{equation*}
y = \lim y_{k_j} = \lim f(x_{k_j}) = f(x^*).
\end{equation*}
Por ende, $y \in f(\mathbb{R}^n)$, lo que demuestra que $f(\mathbb{R}^n)$ contiene todos sus puntos de acumulación, es decir, es un conjunto cerrado.
\end{enumerate}

\begin{enumerate}[label=\textbf{\arabic*.}]
\item \textbf{Sobreyectividad por conexidad:}
El espacio $\mathbb{R}^n$ es conexo. Los únicos subconjuntos de $\mathbb{R}^n$ que son simultáneamente abiertos y cerrados (clopen) son $\emptyset$ y $\mathbb{R}^n$.
Como $f(\mathbb{R}^n) \neq \emptyset$, concluimos que $f(\mathbb{R}^n) = \mathbb{R}^n$ (sobreyectiva).
\end{enumerate}

\begin{enumerate}[label=\textbf{\arabic*.}]
\item \textbf{Inyectividad por simple conexidad:}
Al ser $f$ un difeomorfismo local propio sobre $\mathbb{R}^n$, constituye una aplicación de recubrimiento (covering map). Dado que $\mathbb{R}^n$ es simplemente conexo para $n \ge 2$, todo recubrimiento conexo es un homeomorfismo trivial de 1 sola hoja. Por ende, $f$ es inyectiva y por tanto un difeomorfismo global.
\end{enumerate}

\begin{trapbox}
Olvidar probar que la imagen es cerrada al verificar solo el carácter abierto mediante el TFI.
\end{trapbox}

\vspace{0.5cm}
\hrule
\vspace{0.5cm}

\subsection*{Problema 1.12: Necesidad estricta de la hipótesis $C^1$ (Contraejemplo de Cauchy) \hfill \color{primaryblue}\small[Nivel 4]}
\addcontentsline{toc}{subsection}{Problema 1.12: \texorpdfstring{Necesidad estricta de la hipótesis $C^1$ (Contraejemplo de Cauchy)}{Necesidad estricta de la hipótesis C^1 (Contraejemplo de Cauchy)} (Nivel 4)}

\begin{tcolorbox}[colback=white,colframe=gray!40,arc=1mm,boxrule=0.5pt,left=3mm,right=3mm,top=2mm,bottom=2mm]
\textbf{Enunciado:} Considere la función $f: \mathbb{R} \to \mathbb{R}$ dada por:
\begin{equation*}
f(x) = \begin{cases} \frac{x}{2} + x^2 \sin\left(\frac{1}{x}\right), & x \neq 0 \\ 0, & x = 0 \end{cases}
\end{equation*}
(a) Demuestre que $f$ es diferenciable en todo $\mathbb{R}$ y que $f'(0) = 1/2 \neq 0$. (b) Demuestre que $f'$ no es continua en $0$. (c) Demuestre que $f$ NO es inyectiva en ningún intervalo abierto que contenga a $0$, evidenciando por qué la continuidad de la derivada es imprescindible en el TFI.
\end{tcolorbox}

\begin{tcolorbox}[colback=yellow!4,colframe=accentamber!70,arc=1mm,boxrule=0.5pt,left=3mm,right=3mm,top=1.5mm,bottom=1.5mm]
\textbf{Pista Conceptual:} Calcule $f'(0)$ por definición de límite y luego calcule $f'(x)$ para $x \neq 0$ mostrando que oscila con signos positivos y negativos cerca del origen.
\end{tcolorbox}

\paragraph{Solución Paso a Paso:}\par
\begin{enumerate}[label=\textbf{\arabic*.}]
\item \textbf{Diferenciabilidad en el origen:}
\begin{equation*}
f'(0) = \lim_{h \to 0} \frac{f(h) - f(0)}{h} = \lim_{h \to 0} \frac{\frac{h}{2} + h^2 \sin(1/h)}{h} = \lim_{h \to 0} \left(\frac{1}{2} + h \sin\left(\frac{1}{h}\right)\right).
\end{equation*}
Como $|h\sin(1/h)| \le |h| \to 0$, por el Teorema del Sándwich:
\begin{equation*}
f'(0) = \frac{1}{2} + 0 = \frac{1}{2} \neq 0.
\end{equation*}
\end{enumerate}

\begin{enumerate}[label=\textbf{\arabic*.}]
\item \textbf{Derivada fuera del origen y discontinuidad:}
Para $x \neq 0$, por las reglas usuales de derivación:
\begin{equation*}
f'(x) = \frac{1}{2} + 2x\sin\left(\frac{1}{x}\right) - \cos\left(\frac{1}{x}\right).
\end{equation*}
Cuando $x \to 0$, el término $\cos(1/x)$ oscila entre $-1$ y $1$.
Por tanto, $\lim_{x \to 0} f'(x)$ no existe. Esto demuestra que $f$ es diferenciable en todo $\mathbb{R}$, pero su derivada NO es continua en $x=0$, es decir, $f \notin C^1$.
\end{enumerate}

\begin{enumerate}[label=\textbf{\arabic*.}]
\item \textbf{Pérdida de la inyectividad local:}
En cualquier vecindad $(-\varepsilon, \varepsilon)$, existen puntos $x_k = \frac{1}{2k\pi}$ donde $\cos(1/x_k) = 1$:
\begin{equation*}
f'(x_k) = \frac{1}{2} + 0 - 1 = -\frac{1}{2} < 0.
\end{equation*}
Y existen puntos $y_k = \frac{1}{(2k+1)\pi}$ donde $\cos(1/y_k) = -1$:
\begin{equation*}
f'(y_k) = \frac{1}{2} + 0 - (-1) = \frac{3}{2} > 0.
\end{equation*}
Por el Teorema de Darboux para derivadas, $f'$ cambia de signo infinitas veces en cualquier entorno de $0$.
Por el Teorema de Rolle, al tener infinitos máximos y mínimos locales cerca de $0$, $f$ no es monótona y toma el mismo valor múltiples veces en cualquier intervalo que contenga al origen.
Por ende, $f$ no es invertible en ningún entorno de $0$, demostrando que la hipótesis $f \in C^1$ es estrictamente necesaria.
\end{enumerate}

\begin{trapbox}
Creer erróneamente que tener derivada no nula en un punto basta para asegurar que la función sea monótona en una vecindad.
\end{trapbox}

\vspace{0.5cm}
\hrule
\vspace{0.5cm}

\subsection*{Problema 1.13: Demostración constructiva del TFI por Contracción de Banach \hfill \color{primaryblue}\small[Nivel 5]}
\addcontentsline{toc}{subsection}{Problema 1.13: \texorpdfstring{Demostración constructiva del TFI por Contracción de Banach}{Demostración constructiva del TFI por Contracción de Banach} (Nivel 5)}

\begin{tcolorbox}[colback=white,colframe=gray!40,arc=1mm,boxrule=0.5pt,left=3mm,right=3mm,top=2mm,bottom=2mm]
\textbf{Enunciado:} Sea $U \subseteq \mathbb{R}^n$ abierto, $x_0 \in U$ y $f: U \to \mathbb{R}^n$ de clase $C^1$ con $A = Df(x_0)$ invertible. Defina la aplicación auxiliar $T_y(x) = x + A^{-1}(y - f(x))$. (a) Demuestre que $DT_y(x_0) = 0$. (b) Pruebe que existe una bola cerrada $\overline{B}(x_0, r)$ y un $\delta > 0$ tal que para todo $y \in B(f(x_0), \delta)$, $T_y$ es una contracción estricta en $\overline{B}(x_0, r)$. (c) Concluya la existencia de la inversa local.
\end{tcolorbox}

\begin{tcolorbox}[colback=yellow!4,colframe=accentamber!70,arc=1mm,boxrule=0.5pt,left=3mm,right=3mm,top=1.5mm,bottom=1.5mm]
\textbf{Pista Conceptual:} Aplique el Teorema del Valor Medio vectorial: $\|T_y(x_1) - T_y(x_2)\| \le \sup \|DT_y\| \|x_1 - x_2\|$, eligiendo el radio tal que $\|DT_y(x)\| \le 1/2$.
\end{tcolorbox}

\paragraph{Solución Paso a Paso:}\par
\begin{enumerate}[label=\textbf{\arabic*.}]
\item \textbf{Diferencial de $T_y$ en $x_0$:}
Nótese que $DT_y(x) = I_n - A^{-1} Df(x)$.
Evaluando en $x_0$:
\begin{equation*}
DT_y(x_0) = I_n - A^{-1} Df(x_0) = I_n - A^{-1}A = I_n - I_n = 0.
\end{equation*}
\end{enumerate}

\begin{enumerate}[label=\textbf{\arabic*.}]
\item \textbf{Construcción de la contracción:}
Dado que $f \in C^1$, la función $x \mapsto DT_y(x)$ es continua.
Por definición de continuidad, existe $r > 0$ tal que $\overline{B}(x_0, r) \subset U$ y:
\begin{equation*}
\|DT_y(x)\| \le \frac{1}{2}, \quad \forall x \in \overline{B}(x_0, r).
\end{equation*}
Por la desigualdad del valor medio:
\begin{equation*}
\|T_y(x_1) - T_y(x_2)\| \le \frac{1}{2} \|x_1 - x_2\|, \quad \forall x_1, x_2 \in \overline{B}(x_0, r).
\end{equation*}
Por tanto, $T_y$ es una aplicación lipschitziana con constante $k = 1/2 < 1$.
\end{enumerate}

\begin{enumerate}[label=\textbf{\arabic*.}]
\item \textbf{Invariancia de la bola $\overline{B}(x_0, r)$:}
Tomemos $y$ tal que $\|y - f(x_0)\| < \delta$ con $\delta = \frac{r}{2 \|A^{-1}\|}$.
Para cualquier $x \in \overline{B}(x_0, r)$:
\begin{equation*}
\|T_y(x) - x_0\| \le \|T_y(x) - T_y(x_0)\| + \|T_y(x_0) - x_0\|
\end{equation*}
\begin{equation*}
\le \frac{1}{2}\|x - x_0\| + \|A^{-1}(y - f(x_0))\| \le \frac{1}{2}r + \|A^{-1}\| \cdot \frac{r}{2\|A^{-1}\|} = \frac{r}{2} + \frac{r}{2} = r.
\end{equation*}
Así, $T_y(\overline{B}(x_0, r)) \subseteq \overline{B}(x_0, r)$.
\end{enumerate}

\begin{enumerate}[label=\textbf{\arabic*.}]
\item \textbf{Aplicación del Teorema del Punto Fijo de Banach:}
El espacio $\overline{B}(x_0, r) \subset \mathbb{R}^n$ es un espacio métrico completo. Como $T_y$ es una contracción de dicho espacio en sí mismo, por el Teorema del Punto Fijo de Banach existe un único punto fijo $x \in \overline{B}(x_0, r)$:
\begin{equation*}
T_y(x) = x \iff x + A^{-1}(y - f(x)) = x \iff A^{-1}(y - f(x)) = 0 \iff f(x) = y.
\end{equation*}
Esto demuestra la existencia y unicidad de la solución local $x = f^{-1}(y)$.
\end{enumerate}

\begin{trapbox}
Omitir verificar que $T_y$ aplica la bola cerrada en sí misma antes de invocar el teorema de Banach.
\end{trapbox}

\vspace{0.5cm}
\hrule
\vspace{0.5cm}

\subsection*{Problema 1.14: Teorema de Hadamard de la cota uniforme del jacobiano inverso \hfill \color{primaryblue}\small[Nivel 5]}
\addcontentsline{toc}{subsection}{Problema 1.14: \texorpdfstring{Teorema de Hadamard de la cota uniforme del jacobiano inverso}{Teorema de Hadamard de la cota uniforme del jacobiano inverso} (Nivel 5)}

\begin{tcolorbox}[colback=white,colframe=gray!40,arc=1mm,boxrule=0.5pt,left=3mm,right=3mm,top=2mm,bottom=2mm]
\textbf{Enunciado:} Sea $f: \mathbb{R}^n \to \mathbb{R}^n$ de clase $C^1$. Suponga que existe una constante $M > 0$ tal que para todo $x \in \mathbb{R}^n$, $Df(x)$ es invertible y $\|[Df(x)]^{-1}\| \le M$. Demuestre que $f$ es un difeomorfismo global de $\mathbb{R}^n$ sobre $\mathbb{R}^n$.
\end{tcolorbox}

\begin{tcolorbox}[colback=yellow!4,colframe=accentamber!70,arc=1mm,boxrule=0.5pt,left=3mm,right=3mm,top=1.5mm,bottom=1.5mm]
\textbf{Pista Conceptual:} Para mostrar que $f$ es sobreyectiva, use el método del levantamiento de caminos: para cualquier segmento de recta que una $y_0$ con $y_1$, integre la ecuación diferencial $x'(t) = [Df(x(t))]^{-1}(y_1 - y_0)$.
\end{tcolorbox}

\paragraph{Solución Paso a Paso:}\par
\begin{enumerate}[label=\textbf{\arabic*.}]
\item \textbf{Invertibilidad local:}
Como $\det Df(x) \neq 0$ para todo $x$, $f$ es un difeomorfismo local.
\end{enumerate}

\begin{enumerate}[label=\textbf{\arabic*.}]
\item \textbf{Crecimiento en el infinito (Aplicación propia):}
Para dos puntos $x_1, x_2$, consideremos el segmento $\gamma(t) = (1-t)x_1 + t x_2$.
Por el Teorema Fundamental del Cálculo aplicado a la inversa, o considerando la desigualdad:
\begin{equation*}
\|x_1 - x_2\| = \left\| \int_0^1 [Df(\gamma(t))]^{-1} \frac{d}{dt}f(\gamma(t)) dt \right\| \le \sup_t \|[Df(\gamma(t))]^{-1}\| \cdot \|f(x_1) - f(x_2)\|
\end{equation*}
\begin{equation*}
\implies \|x_1 - x_2\| \le M \|f(x_1) - f(x_2)\| \implies \|f(x_1) - f(x_2)\| \ge \frac{1}{M} \|x_1 - x_2\|.
\end{equation*}
Esta cota inferior uniforme demuestra inmediatamente dos propiedades fundamentales:
a) \textbf{Inyectividad global:} Si $f(x_1) = f(x_2)$, entonces $\|x_1 - x_2\| \le M(0) = 0 \implies x_1 = x_2$.
b) \textbf{Propia:} Si $\|x\| \to \infty$, entonces $\|f(x)\| \ge \|f(0)\| + \frac{1}{M}\|x\| \to \infty$.
\end{enumerate}

\begin{enumerate}[label=\textbf{\arabic*.}]
\item \textbf{Sobreyectividad y conclusión:}
Como vimos en el ejercicio 1.11, al ser $f$ localmente invertible y propia sobre el espacio conexo $\mathbb{R}^n$, la imagen $f(\mathbb{R}^n) = \mathbb{R}^n$.
Por tanto, $f$ es inyectiva, sobreyectiva, de clase $C^1$ y con diferencial no singular en todo punto, constituyendo un difeomorfismo global.
\end{enumerate}

\begin{trapbox}
Aplicar el teorema del valor medio vectorial como igualdad en lugar de la desigualdad integral con la cota de la norma de la matriz inversa.
\end{trapbox}

\vspace{0.5cm}
\hrule
\vspace{0.5cm}

\subsection*{Problema 1.15: Segunda derivada de la función inversa y forma bilineal \hfill \color{primaryblue}\small[Nivel 5]}
\addcontentsline{toc}{subsection}{Problema 1.15: \texorpdfstring{Segunda derivada de la función inversa y forma bilineal}{Segunda derivada de la función inversa y forma bilineal} (Nivel 5)}

\begin{tcolorbox}[colback=white,colframe=gray!40,arc=1mm,boxrule=0.5pt,left=3mm,right=3mm,top=2mm,bottom=2mm]
\textbf{Enunciado:} Sea $f: U \to V$ un difeomorfismo de clase $C^2$ entre abiertos de $\mathbb{R}^n$, y sea $g = f^{-1}$. A partir de la identidad $Dg(f(x)) \cdot Df(x) = I_n$, deduzca la fórmula bilineal intrínseca para la segunda diferencial $D^2 g(y)$ en términos de $Df(x)$ y $D^2 f(x)$, donde $y = f(x)$.
\end{tcolorbox}

\begin{tcolorbox}[colback=yellow!4,colframe=accentamber!70,arc=1mm,boxrule=0.5pt,left=3mm,right=3mm,top=1.5mm,bottom=1.5mm]
\textbf{Pista Conceptual:} Diferencie la función matricial $x \mapsto [Df(x)]^{-1}$ usando la regla $d(A^{-1}) = -A^{-1} (dA) A^{-1}$.
\end{tcolorbox}

\paragraph{Solución Paso a Paso:}\par
\begin{enumerate}[label=\textbf{\arabic*.}]
\item \textbf{Relación de primer orden:}
Para todo $x \in U$, siendo $y = f(x)$:
\begin{equation*}
Dg(y) = [Df(x)]^{-1}.
\end{equation*}
\end{enumerate}

\begin{enumerate}[label=\textbf{\arabic*.}]
\item \textbf{Diferenciación de la matriz inversa:}
Recordemos que si $A(x)$ es una función matricial invertible diferenciable, diferenciando la identidad $A(x) A(x)^{-1} = I$:
\begin{equation*}
dA \cdot A^{-1} + A \cdot d(A^{-1}) = 0 \implies d(A^{-1}) = -A^{-1} \cdot dA \cdot A^{-1}.
\end{equation*}
En nuestro caso, con $A(x) = Df(x)$:
\begin{equation*}
D([Df(x)]^{-1})(h) = -[Df(x)]^{-1} \cdot D^2 f(x)(h, \cdot) \cdot [Df(x)]^{-1}.
\end{equation*}
\end{enumerate}

\begin{enumerate}[label=\textbf{\arabic*.}]
\item \textbf{Aplicación de la regla de la cadena:}
Por otro lado, derivando $x \mapsto Dg(f(x))$ con respecto a $x$ en la dirección $h$:
\begin{equation*}
D(Dg \circ f)(x)(h) = D^2 g(f(x))(Df(x)h, \cdot).
\end{equation*}
Haciendo el cambio de vector dirección $k = Df(x)h \iff h = [Df(x)]^{-1} k$:
\begin{equation*}
D^2 g(y)(k, w) = -[Df(x)]^{-1} \cdot \left[ D^2 f(x)\left([Df(x)]^{-1}k, [Df(x)]^{-1}w\right) \right].
\end{equation*}
Esta es la fórmula tensorial/bilineal exacta para la segunda derivada de la transformación inversa.
\end{enumerate}

\begin{trapbox}
Derivar cada elemento escalarmente sin mantener el orden del producto no conmutativo de matrices y operadores bilineales.
\end{trapbox}

\vspace{0.5cm}
\hrule
\vspace{0.5cm}

\section{Teorema de la Función Implícita}

\subsection*{Problema 2.0A: Derivación implícita de la circunferencia en un punto regular \hfill \color{teal}\small[Nivel 0]}
\addcontentsline{toc}{subsection}{Problema 2.0A: Derivación implícita de la circunferencia en un punto regular (Nivel 0)}

\begin{tcolorbox}[colback=white,colframe=gray!40,arc=1mm,boxrule=0.5pt,left=3mm,right=3mm,top=2mm,bottom=2mm]
\textbf{Enunciado:} Sea la curva descrita implícitamente por $F(x,y) = x^2 + y^2 - 25 = 0$ y sea el punto $P_0 = (3,4)$.
\begin{enumerate}[label=\textbf{\alph*)}]
\item Verifique que $P_0$ pertenece a la curva.
\item Calcule $\frac{\partial F}{\partial x}(3,4)$ y $\frac{\partial F}{\partial y}(3,4)$.
\item Verifique que $\frac{\partial F}{\partial y}(3,4) \neq 0$, aplique la fórmula del TFI para hallar $\frac{dy}{dx}(3)$ y compruebe el resultado despejando $y = \sqrt{25 - x^2}$.
\end{enumerate}
\end{tcolorbox}

\begin{tcolorbox}[colback=yellow!4,colframe=accentamber!70,arc=1mm,boxrule=0.5pt,left=3mm,right=3mm,top=1.5mm,bottom=1.5mm]
\textbf{Pista Conceptual:} La fórmula del Teorema de la Función Implícita en el plano es $\frac{dy}{dx} = -\frac{\partial F / \partial x}{\partial F / \partial y}$.
\end{tcolorbox}

\paragraph{Solución Paso a Paso:}\par
\begin{enumerate}[label=\textbf{\arabic*.}]
\item \textbf{Pertenencia a la curva:}
\[ F(3,4) = 3^2 + 4^2 - 25 = 9 + 16 - 25 = 0. \]
El punto satisface la ecuación.
\item \textbf{Derivadas parciales:}
\[ \frac{\partial F}{\partial x} = 2x \implies \frac{\partial F}{\partial x}(3,4) = 2(3) = 6. \]
\[ \frac{\partial F}{\partial y} = 2y \implies \frac{\partial F}{\partial y}(3,4) = 2(4) = 8. \]
\item \textbf{Cálculo con la fórmula del TFI:}
Como $\frac{\partial F}{\partial y}(3,4) = 8 \neq 0$, por el TFI existe una función $y = \varphi(x)$ en una vecindad de $x=3$ tal que $\varphi(3) = 4$, y su derivada es:
\[ \frac{dy}{dx}(3) = - \frac{\frac{\partial F}{\partial x}(3,4)}{\frac{\partial F}{\partial y}(3,4)} = -\frac{6}{8} = -\frac{3}{4}. \]
\item \textbf{Comprobación por despeje:}
Cerca de $(3,4)$, $y > 0$, de modo que $y = \sqrt{25 - x^2} = (25 - x^2)^{1/2}$. Derivando:
\[ \frac{dy}{dx} = \frac{1}{2}(25 - x^2)^{-1/2}(-2x) = -\frac{x}{\sqrt{25 - x^2}} \implies \frac{dy}{dx}(3) = -\frac{3}{\sqrt{25-9}} = -\frac{3}{4}. \]
¡El resultado coincide con exactitud total!
\end{enumerate}

\begin{trapbox}
Olvidar el signo menos en la fórmula del TFI: $\frac{dy}{dx} = -\frac{F_x}{F_y}$. Recuerde que al diferenciar $F(x,y)=0$ se tiene $F_x dx + F_y dy = 0 \implies dy = -\frac{F_x}{F_y} dx$.
\end{trapbox}

\vspace{0.5cm}
\hrule
\vspace{0.5cm}
\subsection*{Problema 2.0B: Falla de despeje para y(x) pero éxito para x(y) (Punto con tangente vertical) \hfill \color{teal}\small[Nivel 0]}
\addcontentsline{toc}{subsection}{Problema 2.0B: Falla de despeje para y(x) pero éxito para x(y) (Punto con tangente vertical) (Nivel 0)}

\begin{tcolorbox}[colback=white,colframe=gray!40,arc=1mm,boxrule=0.5pt,left=3mm,right=3mm,top=2mm,bottom=2mm]
\textbf{Enunciado:} Considere la misma circunferencia $F(x,y) = x^2 + y^2 - 25 = 0$ en el punto extremo $P_1 = (5,0)$.
\begin{enumerate}[label=\textbf{\alph*)}]
\item Calcule $\frac{\partial F}{\partial y}(5,0)$ y explique por qué el TFI no garantiza que $y$ sea función diferenciable de $x$ alrededor de este punto.
\item Demuestre que, en cambio, la variable $x$ sí puede despejarse como función suave de $y$, es decir $x = \psi(y)$, calculando $\frac{dx}{dy}(0)$.
\end{enumerate}
\end{tcolorbox}

\begin{tcolorbox}[colback=yellow!4,colframe=accentamber!70,arc=1mm,boxrule=0.5pt,left=3mm,right=3mm,top=1.5mm,bottom=1.5mm]
\textbf{Pista Conceptual:} En $(5,0)$, la recta tangente a la circunferencia es la recta vertical $x = 5$. Una curva con recta tangente vertical no puede ser función de $x$.
\end{tcolorbox}

\paragraph{Solución Paso a Paso:}\par
\begin{enumerate}[label=\textbf{\arabic*.}]
\item \textbf{Evaluación de derivadas parciales en $(5,0)$:}
\[ \frac{\partial F}{\partial y}(5,0) = 2(0) = 0, \qquad \frac{\partial F}{\partial x}(5,0) = 2(5) = 10 \neq 0. \]
\item \textbf{Por qué falla $y = \varphi(x)$:}
Como $\frac{\partial F}{\partial y}(5,0) = 0$, la hipótesis del TFI para despejar $y(x)$ no se cumple. Geométricamente, en $(5,0)$ la recta tangente a la circunferencia es vertical ($x=5$). Para cualquier $x < 5$ cercano, existen dos valores correspondientes de $y$: $+\sqrt{25-x^2}$ y $-\sqrt{25-x^2}$, mientras que para $x > 5$ no existe ningún valor real. Por tanto, ninguna vecindad de $(5,0)$ permite definir $y$ como función unívoca de $x$.
\item \textbf{Éxito para $x = \psi(y)$:}
Como $\frac{\partial F}{\partial x}(5,0) = 10 \neq 0$, intercambiamos los roles de las variables en el TFI:
\[ \frac{dx}{dy}(0) = - \frac{\frac{\partial F}{\partial y}(5,0)}{\frac{\partial F}{\partial x}(5,0)} = -\frac{0}{10} = 0. \]
Efectivamente, localmente la curva se puede despejar como $x = \psi(y) = \sqrt{25 - y^2}$, la cual es suave para $y$ cercano a 0 y tiene tangente horizontal $\psi'(0) = 0$ respecto a $y$.
\end{enumerate}

\begin{trapbox}
Pensar que cuando $\frac{\partial F}{\partial y} = 0$ la curva no tiene sentido analítico. Solo significa que la recta tangente es vertical y que la variable que debe despejarse es $x$, no $y$.
\end{trapbox}

\vspace{0.5cm}
\hrule
\vspace{0.5cm}
\subsection*{Problema 2.0C: Función implícita en tres variables $z = \varphi(x,y)$ (Plano afín) \hfill \color{teal}\small[Nivel 0]}
\addcontentsline{toc}{subsection}{Problema 2.0C: Función implícita en tres variables $z = \varphi(x,y)$ (Plano afín) (Nivel 0)}

\begin{tcolorbox}[colback=white,colframe=gray!40,arc=1mm,boxrule=0.5pt,left=3mm,right=3mm,top=2mm,bottom=2mm]
\textbf{Enunciado:} Sea la ecuación implícita en $\mathbb{R}^3$ dada por $F(x,y,z) = 2x - 3y + 4z - 12 = 0$.
\begin{enumerate}[label=\textbf{\alph*)}]
\item Verifique que $\frac{\partial F}{\partial z} \neq 0$ en todo punto.
\item Aplique las fórmulas del TFI para calcular las derivadas parciales $\frac{\partial z}{\partial x}$ y $\frac{\partial z}{\partial y}$.
\item Despeje $z$ directamente y verifique las derivadas.
\end{enumerate}
\end{tcolorbox}

\begin{tcolorbox}[colback=yellow!4,colframe=accentamber!70,arc=1mm,boxrule=0.5pt,left=3mm,right=3mm,top=1.5mm,bottom=1.5mm]
\textbf{Pista Conceptual:} Para una superficie $F(x,y,z)=0$, las derivadas de $z(x,y)$ son $\frac{\partial z}{\partial x} = -\frac{F_x}{F_z}$ y $\frac{\partial z}{\partial y} = -\frac{F_y}{F_z}$.
\end{tcolorbox}

\paragraph{Solución Paso a Paso:}\par
\begin{enumerate}[label=\textbf{\arabic*.}]
\item \textbf{Cálculo de derivadas parciales de $F$:}
\[ \frac{\partial F}{\partial x} = 2, \qquad \frac{\partial F}{\partial y} = -3, \qquad \frac{\partial F}{\partial z} = 4. \]
Como $\frac{\partial F}{\partial z} = 4 \neq 0$, la variable dependiente $z$ puede expresarse como función suave $z = \varphi(x,y)$ en todo $\mathbb{R}^2$.
\item \textbf{Derivadas parciales con el TFI:}
\[ \frac{\partial z}{\partial x} = -\frac{\frac{\partial F}{\partial x}}{\frac{\partial F}{\partial z}} = -\frac{2}{4} = -\frac{1}{2}. \]
\[ \frac{\partial z}{\partial y} = -\frac{\frac{\partial F}{\partial y}}{\frac{\partial F}{\partial z}} = -\frac{-3}{4} = \frac{3}{4}. \]
\item \textbf{Verificación por despeje directo:}
\[ 2x - 3y + 4z - 12 = 0 \implies 4z = 12 - 2x + 3y \implies z = 3 - \frac{1}{2}x + \frac{3}{4}y. \]
Derivando directamente respecto a $x$ e $y$:
\[ \frac{\partial z}{\partial x} = -\frac{1}{2}, \qquad \frac{\partial z}{\partial y} = \frac{3}{4}, \]
coincidiendo plenamente con las fórmulas del teorema.
\end{enumerate}

\begin{trapbox}
Confundir el signo al dividir por un término negativo, como en $\frac{\partial z}{\partial y} = -\frac{-3}{4} = +\frac{3}{4}$.
\end{trapbox}

\vspace{0.5cm}
\hrule
\vspace{0.5cm}


\subsection*{Problema 2.1: Función implícita en curva plana (Folium de Descartes) \hfill \color{primaryblue}\small[Nivel 1]}
\addcontentsline{toc}{subsection}{Problema 2.1: \texorpdfstring{Función implícita en curva plana (Folium de Descartes)}{Función implícita en curva plana (Folium de Descartes)} (Nivel 1)}

\begin{tcolorbox}[colback=white,colframe=gray!40,arc=1mm,boxrule=0.5pt,left=3mm,right=3mm,top=2mm,bottom=2mm]
\textbf{Enunciado:} Sea la ecuación $F(x,y) = x^3 + y^3 - 3xy = 0$ en el punto $P_0 = (3/2, 3/2)$. (a) Verifique que se cumplen las hipótesis del Teorema de la Función Implícita para definir $y = \varphi(x)$ en un entorno de $x_0 = 3/2$. (b) Calcule $\varphi'(3/2)$ y halle la recta tangente a la curva en dicho punto.
\end{tcolorbox}

\begin{tcolorbox}[colback=yellow!4,colframe=accentamber!70,arc=1mm,boxrule=0.5pt,left=3mm,right=3mm,top=1.5mm,bottom=1.5mm]
\textbf{Pista Conceptual:} Compruebe que $F(3/2, 3/2) = 0$ y que $\frac{\partial F}{\partial y}(3/2, 3/2) \neq 0$. Aplique $\varphi'(x) = -\frac{F_x}{F_y}$.
\end{tcolorbox}

\paragraph{Solución Paso a Paso:}\par
\begin{enumerate}[label=\textbf{\arabic*.}]
\item \textbf{Verificación de pertenencia a la curva:}
\begin{equation*}
F(3/2, 3/2) = (3/2)^3 + (3/2)^3 - 3(3/2)(3/2) = \frac{27}{8} + \frac{27}{8} - \frac{27}{4} = \frac{54}{8} - \frac{54}{8} = 0.
\end{equation*}
\end{enumerate}

\begin{enumerate}[label=\textbf{\arabic*.}]
\item \textbf{Derivadas parciales y regularidad:}
Las funciones componentes son polinomios, luego $F \in C^\infty(\mathbb{R}^2)$.
\begin{equation*}
F_x(x,y) = 3x^2 - 3y, \qquad F_y(x,y) = 3y^2 - 3x.
\end{equation*}
Evaluando en $P_0 = (3/2, 3/2)$:
\begin{equation*}
F_x(3/2, 3/2) = 3(9/4) - 3(3/2) = \frac{27}{4} - \frac{18}{4} = \frac{9}{4},
\end{equation*}
\begin{equation*}
F_y(3/2, 3/2) = 3(9/4) - 3(3/2) = \frac{9}{4} \neq 0.
\end{equation*}
Como $F(3/2, 3/2) = 0$, $F \in C^1$ y $F_y(3/2, 3/2) \neq 0$, por el Teorema de la Función Implícita existe un intervalo abierto $I$ alrededor de $3/2$ y una función diferenciable única $y = \varphi(x)$ tal que $\varphi(3/2) = 3/2$ y $F(x, \varphi(x)) = 0$.
\end{enumerate}

\begin{enumerate}[label=\textbf{\arabic*.}]
\item \textbf{Cálculo de la derivada:}
\begin{equation*}
\varphi'(3/2) = -\frac{F_x(3/2, 3/2)}{F_y(3/2, 3/2)} = -\frac{9/4}{9/4} = -1.
\end{equation*}
\end{enumerate}

\begin{enumerate}[label=\textbf{\arabic*.}]
\item \textbf{Ecuación de la recta tangente:}
\begin{equation*}
y - \frac{3}{2} = -1\left(x - \frac{3}{2}\right) \implies x + y = 3.
\end{equation*}
\end{enumerate}

\begin{trapbox}
Olvidar el signo menos en la fórmula de la derivada implícita: $\varphi' = -F_x / F_y$.
\end{trapbox}

\vspace{0.5cm}
\hrule
\vspace{0.5cm}

\subsection*{Problema 2.2: Superficie implícita y plano tangente \hfill \color{primaryblue}\small[Nivel 1]}
\addcontentsline{toc}{subsection}{Problema 2.2: \texorpdfstring{Superficie implícita y plano tangente}{Superficie implícita y plano tangente} (Nivel 1)}

\begin{tcolorbox}[colback=white,colframe=gray!40,arc=1mm,boxrule=0.5pt,left=3mm,right=3mm,top=2mm,bottom=2mm]
\textbf{Enunciado:} Dada la ecuación $F(x,y,z) = x^2 + y^2 + z^2 - 14 = 0$ y el punto $P_0 = (1, 2, 3)$: (a) Demuestre que define implícitamente a $z = \varphi(x,y)$ en una vecindad de $(1,2)$. (b) Calcule el gradiente $\nabla\varphi(1,2)$ y la ecuación del plano tangente a la superficie en $P_0$.
\end{tcolorbox}

\begin{tcolorbox}[colback=yellow!4,colframe=accentamber!70,arc=1mm,boxrule=0.5pt,left=3mm,right=3mm,top=1.5mm,bottom=1.5mm]
\textbf{Pista Conceptual:} Calcule $F_z(1,2,3)$ y use $\frac{\partial z}{\partial x} = -\frac{F_x}{F_z}$, $\frac{\partial z}{\partial y} = -\frac{F_y}{F_z}$.
\end{tcolorbox}

\paragraph{Solución Paso a Paso:}\par
\begin{enumerate}[label=\textbf{\arabic*.}]
\item \textbf{Verificación de hipótesis:}
\begin{equation*}
F(1,2,3) = 1^2 + 2^2 + 3^2 - 14 = 1 + 4 + 9 - 14 = 0.
\end{equation*}
$F$ es un polinomio, por ende $C^\infty$.
Calculamos las derivadas parciales en $P_0(1,2,3)$:
\begin{equation*}
F_x = 2x \implies F_x(1,2,3) = 2,
\end{equation*}
\begin{equation*}
F_y = 2y \implies F_y(1,2,3) = 4,
\end{equation*}
\begin{equation*}
F_z = 2z \implies F_z(1,2,3) = 6 \neq 0.
\end{equation*}
Como $F_z(P_0) = 6 \neq 0$, por el Teorema de la Función Implícita existe una función única $z = \varphi(x,y)$ de clase $C^1$ en un entorno de $(1,2)$ con $\varphi(1,2) = 3$.
\end{enumerate}

\begin{enumerate}[label=\textbf{\arabic*.}]
\item \textbf{Cálculo del gradiente:}
\begin{equation*}
\frac{\partial \varphi}{\partial x}(1,2) = -\frac{F_x}{F_z} = -\frac{2}{6} = -\frac{1}{3},
\end{equation*}
\begin{equation*}
\frac{\partial \varphi}{\partial y}(1,2) = -\frac{F_y}{F_z} = -\frac{4}{6} = -\frac{2}{3}.
\end{equation*}
Por tanto, $\nabla\varphi(1,2) = \left(-\frac{1}{3}, -\frac{2}{3}\right)$.
\end{enumerate}

\begin{enumerate}[label=\textbf{\arabic*.}]
\item \textbf{Plano tangente:}
\begin{equation*}
z - 3 = -\frac{1}{3}(x - 1) - \frac{2}{3}(y - 2) \implies (x - 1) + 2(y - 2) + 3(z - 3) = 0 \iff x + 2y + 3z = 14.
\end{equation*}
\end{enumerate}

\begin{trapbox}
Confundir la variable dependiente con las independientes en el denominador de la fórmula del cociente de derivadas parciales.
\end{trapbox}

\vspace{0.5cm}
\hrule
\vspace{0.5cm}

\subsection*{Problema 2.3: Función implícita trascendente (Examen UNI 2025-II) \hfill \color{primaryblue}\small[Nivel 1]}
\addcontentsline{toc}{subsection}{Problema 2.3: \texorpdfstring{Función implícita trascendente (Examen UNI 2025-II)}{Función implícita trascendente (Examen UNI 2025-II)} (Nivel 1)}

\begin{tcolorbox}[colback=white,colframe=gray!40,arc=1mm,boxrule=0.5pt,left=3mm,right=3mm,top=2mm,bottom=2mm]
\textbf{Enunciado:} Sea la ecuación $F(x,y) = x e^{xy} + y - 1 = 0$. (a) Demuestre que existe una función $y = \varphi(x)$ en un entorno de $x = 0$ tal que $\varphi(0) = 1$. (b) Calcule $\varphi'(0)$ y $\varphi''(0)$.
\end{tcolorbox}

\begin{tcolorbox}[colback=yellow!4,colframe=accentamber!70,arc=1mm,boxrule=0.5pt,left=3mm,right=3mm,top=1.5mm,bottom=1.5mm]
\textbf{Pista Conceptual:} Verifique $F(0,1) = 0$, $F_y(0,1) \neq 0$, calcule la primera derivada y luego derive nuevamente respecto a $x$ considerando $y = \varphi(x)$.
\end{tcolorbox}

\paragraph{Solución Paso a Paso:}\par
\begin{enumerate}[label=\textbf{\arabic*.}]
\item \textbf{Verificación de hipótesis en $(0,1)$:}
\begin{equation*}
F(0,1) = 0\cdot e^0 + 1 - 1 = 0.
\end{equation*}
$F$ es composición de exponenciales y polinomios, luego $F \in C^\infty(\mathbb{R}^2)$.
Calculamos las derivadas parciales:
\begin{equation*}
F_x(x,y) = e^{xy} + x(y e^{xy}) = e^{xy}(1 + xy),
\end{equation*}
\begin{equation*}
F_y(x,y) = x(x e^{xy}) + 1 = x^2 e^{xy} + 1.
\end{equation*}
Evaluando en $(0,1)$:
\begin{equation*}
F_x(0,1) = e^0(1 + 0) = 1,
\end{equation*}
\begin{equation*}
F_y(0,1) = 0 + 1 = 1 \neq 0.
\end{equation*}
Al ser $F_y(0,1) \neq 0$, el Teorema de la Función Implícita garantiza la existencia de una función única $y = \varphi(x)$ diferenciable en una vecindad de $0$ con $\varphi(0) = 1$.
\end{enumerate}

\begin{enumerate}[label=\textbf{\arabic*.}]
\item \textbf{Cálculo de $\varphi'(0)$:}
\begin{equation*}
\varphi'(0) = -\frac{F_x(0,1)}{F_y(0,1)} = -\frac{1}{1} = -1.
\end{equation*}
\end{enumerate}

\begin{enumerate}[label=\textbf{\arabic*.}]
\item \textbf{Cálculo de $\varphi''(0)$:}
Diferenciamos la identidad $x e^{x\varphi(x)} + \varphi(x) - 1 = 0$ dos veces con respecto a $x$:
Primera derivada:
\begin{equation*}
e^{x\varphi} + x e^{x\varphi}(\varphi + x\varphi') + \varphi' = 0.
\end{equation*}
Segunda derivada:
\begin{equation*}
\frac{d}{dx}\left[e^{x\varphi}(1 + x\varphi + x^2\varphi') + \varphi'\right] = 0
\end{equation*}
Evaluando directamente en $x=0$, recordando que $\varphi(0) = 1$ y $\varphi'(0) = -1$:
\begin{equation*}
e^0 (\varphi(0) + 0) \cdot 1 + e^0 (0 + \varphi(0) + 0) + \varphi''(0) = 0
\end{equation*}
\begin{equation*}
\implies 1(1) + 1(1) + \varphi''(0) = 0 \implies 2 + \varphi''(0) = 0 \implies \varphi''(0) = -2.
\end{equation*}
\end{enumerate}

\begin{trapbox}
Olvidar aplicar la regla del producto y de la cadena a los términos donde aparece $y$ como función de $x$ al calcular la segunda derivada.
\end{trapbox}

\vspace{0.5cm}
\hrule
\vspace{0.5cm}

\subsection*{Problema 2.4: Sistema de 2 ecuaciones y 3 variables \hfill \color{primaryblue}\small[Nivel 2]}
\addcontentsline{toc}{subsection}{Problema 2.4: \texorpdfstring{Sistema de 2 ecuaciones y 3 variables}{Sistema de 2 ecuaciones y 3 variables} (Nivel 2)}

\begin{tcolorbox}[colback=white,colframe=gray!40,arc=1mm,boxrule=0.5pt,left=3mm,right=3mm,top=2mm,bottom=2mm]
\textbf{Enunciado:} Considere el sistema de ecuaciones:
\begin{equation*}
\begin{cases} x^2 + y^2 - z^2 = 0 \\ x + y + 2z = 4 \end{cases}
\end{equation*}
alrededor del punto $P_0 = (1, 1, \sqrt{2})$... espere, en $P_0(1,1,1)$: $1+1-1 \neq 0$. Tomemos $P_0 = (0, 2, 2)$: $0^2 + 2^2 - 2^2 = 0$ y $0 + 2 + 2(2) = 6$.
Ajustando la constante: $\begin{cases} F_1(x,y,z) = x^2 + y^2 - z^2 = 0 \\ F_2(x,y,z) = x + y + 2z = 6 \end{cases}$ en $P_0 = (0, 2, 2)$.
(a) Demuestre que el sistema define a $y$ y $z$ como funciones de $x$ cerca de $x=0$.
(b) Calcule $y'(0)$ y $z'(0)$.
\end{tcolorbox}

\begin{tcolorbox}[colback=yellow!4,colframe=accentamber!70,arc=1mm,boxrule=0.5pt,left=3mm,right=3mm,top=1.5mm,bottom=1.5mm]
\textbf{Pista Conceptual:} Calcule el determinante de la submatriz jacobiana $\frac{\partial(F_1, F_2)}{\partial(y,z)}$ en $P_0$. Si es no nulo, aplique la regla matricial.
\end{tcolorbox}

\paragraph{Solución Paso a Paso:}\par
\begin{enumerate}[label=\textbf{\arabic*.}]
\item \textbf{Evaluación del punto:}
$F_1(0,2,2) = 0^2 + 2^2 - 2^2 = 0$.
$F_2(0,2,2) = 0 + 2 + 2(2) = 6$. Ambas ecuaciones se satisfacen.
\end{enumerate}

\begin{enumerate}[label=\textbf{\arabic*.}]
\item \textbf{Submatriz jacobiana respecto a las variables dependientes $(y,z)$:}
\begin{equation*}
D_{(y,z)}F = \begin{pmatrix} \frac{\partial F_1}{\partial y} & \frac{\partial F_1}{\partial z} \\ \frac{\partial F_2}{\partial y} & \frac{\partial F_2}{\partial z} \end{pmatrix} = \begin{pmatrix} 2y & -2z \\ 1 & 2 \end{pmatrix}.
\end{equation*}
Evaluando en $P_0(0,2,2)$:
\begin{equation*}
D_{(y,z)}F(P_0) = \begin{pmatrix} 4 & -4 \\ 1 & 2 \end{pmatrix}.
\end{equation*}
Determinante:
\begin{equation*}
\det D_{(y,z)}F(P_0) = (4)(2) - (-4)(1) = 8 + 4 = 12 \neq 0.
\end{equation*}
Como $\det \neq 0$ y las funciones son $C^\infty$, el Teorema de la Función Implícita garantiza que $y = y(x)$ y $z = z(x)$ están unívocamente determinadas en un entorno de $x = 0$.
\end{enumerate}

\begin{enumerate}[label=\textbf{\arabic*.}]
\item \textbf{Cálculo de las derivadas:}
Submatriz respecto a la variable independiente $x$:
\begin{equation*}
D_x F = \begin{pmatrix} 2x \\ 1 \end{pmatrix} \implies D_x F(P_0) = \begin{pmatrix} 0 \\ 1 \end{pmatrix}.
\end{equation*}
Por la fórmula matricial del TFI:
\begin{equation*}
\begin{pmatrix} y'(0) \\ z'(0) \end{pmatrix} = - [D_{(y,z)}F(P_0)]^{-1} D_x F(P_0).
\end{equation*}
Invertimos la matriz $2\times 2$:
\begin{equation*}
[D_{(y,z)}F(P_0)]^{-1} = \frac{1}{12} \begin{pmatrix} 2 & 4 \\ -1 & 4 \end{pmatrix}.
\end{equation*}
Multiplicamos:
\begin{equation*}
\begin{pmatrix} y'(0) \\ z'(0) \end{pmatrix} = - \frac{1}{12} \begin{pmatrix} 2 & 4 \\ -1 & 4 \end{pmatrix} \begin{pmatrix} 0 \\ 1 \end{pmatrix} = - \frac{1}{12} \begin{pmatrix} 4 \\ 4 \end{pmatrix} = \begin{pmatrix} -1/3 \\ -1/3 \end{pmatrix}.
\end{equation*}
Por tanto, $y'(0) = -1/3$ y $z'(0) = -1/3$.
\end{enumerate}

\begin{trapbox}
Invertir incorrectamente las filas y columnas al calcular la matriz inversa o ignorar el signo negativo matricial global.
\end{trapbox}

\vspace{0.5cm}
\hrule
\vspace{0.5cm}

\subsection*{Problema 2.5: Vector tangente a curva intersección de superficies \hfill \color{primaryblue}\small[Nivel 2]}
\addcontentsline{toc}{subsection}{Problema 2.5: \texorpdfstring{Vector tangente a curva intersección de superficies}{Vector tangente a curva intersección de superficies} (Nivel 2)}

\begin{tcolorbox}[colback=white,colframe=gray!40,arc=1mm,boxrule=0.5pt,left=3mm,right=3mm,top=2mm,bottom=2mm]
\textbf{Enunciado:} Sea la curva $C$ en $\mathbb{R}^3$ definida por la intersección del cilindro $x^2 + y^2 = 2$ y el plano $x + y + z = 2$. En el punto $P_0 = (1, 1, 0)$: (a) Parametrice localmente la curva con respecto a $x$. (b) Halle un vector tangente a la curva en $P_0$.
\end{tcolorbox}

\begin{tcolorbox}[colback=yellow!4,colframe=accentamber!70,arc=1mm,boxrule=0.5pt,left=3mm,right=3mm,top=1.5mm,bottom=1.5mm]
\textbf{Pista Conceptual:} El vector tangente es $T = (1, y'(x_0), z'(x_0))$. Alternativamente verifique que $T$ es colineal a $\nabla F_1 \times \nabla F_2$.
\end{tcolorbox}

\paragraph{Solución Paso a Paso:}\par
\begin{enumerate}[label=\textbf{\arabic*.}]
\item \textbf{Verificación del sistema:}
$F_1(x,y,z) = x^2 + y^2 - 2 = 0 \implies F_1(1,1,0) = 1+1-2 = 0$.
$F_2(x,y,z) = x + y + z - 2 = 0 \implies F_2(1,1,0) = 1+1+0-2 = 0$.
\end{enumerate}

\begin{enumerate}[label=\textbf{\arabic*.}]
\item \textbf{Cálculo de derivadas implícitas con respecto a $x$:}
Jacobiano respecto a $(y,z)$:
\begin{equation*}
D_{(y,z)}F(1,1,0) = \begin{pmatrix} 2y & 0 \\ 1 & 1 \end{pmatrix}_{(1,1,0)} = \begin{pmatrix} 2 & 0 \\ 1 & 1 \end{pmatrix}.
\end{equation*}
Determinante: $\det = (2)(1) - 0 = 2 \neq 0$.
Derivada respecto a $x$:
\begin{equation*}
D_x F(1,1,0) = \begin{pmatrix} 2x \\ 1 \end{pmatrix}_{(1,1,0)} = \begin{pmatrix} 2 \\ 1 \end{pmatrix}.
\end{equation*}
Derivadas implícitas:
\begin{equation*}
\begin{pmatrix} y'(1) \\ z'(1) \end{pmatrix} = - \begin{pmatrix} 2 & 0 \\ 1 & 1 \end{pmatrix}^{-1} \begin{pmatrix} 2 \\ 1 \end{pmatrix} = - \begin{pmatrix} 1/2 & 0 \\ -1/2 & 1 \end{pmatrix} \begin{pmatrix} 2 \\ 1 \end{pmatrix} = - \begin{pmatrix} 1 \\ 0 \end{pmatrix} = \begin{pmatrix} -1 \\ 0 \end{pmatrix}.
\end{equation*}
\end{enumerate}

\begin{enumerate}[label=\textbf{\arabic*.}]
\item \textbf{Vector tangente:}
La parametrización local es $\gamma(x) = (x, y(x), z(x))$, por tanto el vector tangente es:
\begin{equation*}
T = \gamma'(1) = (1, y'(1), z'(1)) = (1, -1, 0).
\end{equation*}
\textit{Comprobación geométrica:}
$\nabla F_1(1,1,0) = (2, 2, 0)$, $\nabla F_2(1,1,0) = (1, 1, 1)$.
$\nabla F_1 \times \nabla F_2 = (2, -2, 0) = 2(1, -1, 0)$, lo que confirma la exactitud del resultado.
\end{enumerate}

\begin{trapbox}
No verificar que el gradiente de ambas superficies sea no paralelo en el punto de corte.
\end{trapbox}

\vspace{0.5cm}
\hrule
\vspace{0.5cm}

\subsection*{Problema 2.6: Derivadas parciales de segundo orden implícitas \hfill \color{primaryblue}\small[Nivel 2]}
\addcontentsline{toc}{subsection}{Problema 2.6: \texorpdfstring{Derivadas parciales de segundo orden implícitas}{Derivadas parciales de segundo orden implícitas} (Nivel 2)}

\begin{tcolorbox}[colback=white,colframe=gray!40,arc=1mm,boxrule=0.5pt,left=3mm,right=3mm,top=2mm,bottom=2mm]
\textbf{Enunciado:} Sea $z = z(x,y)$ la función implícita definida por $F(x,y,z) = z^3 - xz - y = 0$ cerca del punto $P_0(1, 0, 1)$. Calcule $\frac{\partial^2 z}{\partial x^2}(1,0)$ y $\frac{\partial^2 z}{\partial x \partial y}(1,0)$.
\end{tcolorbox}

\begin{tcolorbox}[colback=yellow!4,colframe=accentamber!70,arc=1mm,boxrule=0.5pt,left=3mm,right=3mm,top=1.5mm,bottom=1.5mm]
\textbf{Pista Conceptual:} Diferencie la identidad implícita respecto a $x$ y respecto a $y$, y vuelva a derivar aplicando la regla de la cadena.
\end{tcolorbox}

\paragraph{Solución Paso a Paso:}\par
\begin{enumerate}[label=\textbf{\arabic*.}]
\item \textbf{Evaluación de derivadas de primer orden:}
$F(1,0,1) = 1^3 - 1(1) - 0 = 0$.
$F_x = -z \implies F_x(P_0) = -1$.
$F_y = -1 \implies F_y(P_0) = -1$.
$F_z = 3z^2 - x \implies F_z(P_0) = 3(1)^2 - 1 = 2 \neq 0$.
Por ende:
\begin{equation*}
z_x(1,0) = -\frac{F_x}{F_z} = -\frac{-1}{2} = \frac{1}{2}, \qquad z_y(1,0) = -\frac{F_y}{F_z} = -\frac{-1}{2} = \frac{1}{2}.
\end{equation*}
\end{enumerate}

\begin{enumerate}[label=\textbf{\arabic*.}]
\item \textbf{Cálculo de $\frac{\partial^2 z}{\partial x^2}$:}
Derivamos la identidad $3z^2 z_x - z - x z_x = 0$ con respecto a $x$:
\begin{equation*}
3(2z z_x) z_x + 3z^2 z_{xx} - z_x - (z_x + x z_{xx}) = 0
\end{equation*}
\begin{equation*}
\implies 6z (z_x)^2 + (3z^2 - x) z_{xx} - 2z_x = 0.
\end{equation*}
Evaluando en $P_0(1,0,1)$ con $z=1, z_x = 1/2$:
\begin{equation*}
6(1)(1/4) + (3(1) - 1) z_{xx} - 2(1/2) = 0 \implies \frac{3}{2} + 2 z_{xx} - 1 = 0 \implies 2 z_{xx} = -\frac{1}{2} \implies z_{xx}(1,0) = -\frac{1}{4}.
\end{equation*}
\end{enumerate}

\begin{enumerate}[label=\textbf{\arabic*.}]
\item \textbf{Cálculo de $\frac{\partial^2 z}{\partial x \partial y}$:}
Derivamos la misma relación $3z^2 z_x - z - x z_x = 0$ respecto a $y$:
\begin{equation*}
6z z_y z_x + 3z^2 z_{xy} - z_y - x z_{xy} = 0 \implies 6z z_x z_y - z_y + (3z^2 - x) z_{xy} = 0.
\end{equation*}
Evaluando en $(1,0)$ con $z=1, z_x = 1/2, z_y = 1/2$:
\begin{equation*}
6(1)(1/2)(1/2) - 1/2 + 2 z_{xy} = 0 \implies \frac{3}{2} - \frac{1}{2} + 2 z_{xy} = 0 \implies 1 + 2 z_{xy} = 0 \implies z_{xy}(1,0) = -\frac{1}{2}.
\end{equation*}
\end{enumerate}

\begin{trapbox}
Olvidar que $z$ depende de $x$ e $y$ y tratar $z$ como constante al derivar el producto $z^2 z_x$.
\end{trapbox}

\vspace{0.5cm}
\hrule
\vspace{0.5cm}

\subsection*{Problema 2.7: Sistema no lineal 2x4 y derivadas parciales matriciales (Examen UNI 2026-I) \hfill \color{primaryblue}\small[Nivel 3]}
\addcontentsline{toc}{subsection}{Problema 2.7: \texorpdfstring{Sistema no lineal 2x4 y derivadas parciales matriciales (Examen UNI 2026-I)}{Sistema no lineal 2x4 y derivadas parciales matriciales (Examen UNI 2026-I)} (Nivel 3)}

\begin{tcolorbox}[colback=white,colframe=gray!40,arc=1mm,boxrule=0.5pt,left=3mm,right=3mm,top=2mm,bottom=2mm]
\textbf{Enunciado:} Considere la solución $P_0(x_0, y_0, u_0, v_0) = (0, 0, 0, 0)$ del sistema de ecuaciones:
\begin{equation*}
\begin{cases} x\cos(y) + e^{x-u} + u^2 v = 1 \\ x^2 + y + \sin(u) - v = 0 \end{cases}
\end{equation*}
(a) Explique rigurosamente por qué en un entorno del punto $P_0$, las variables $u, v$ pueden expresarse como funciones diferenciables de $x, y$, es decir, $u = \phi_1(x,y)$ y $v = \phi_2(x,y)$.
(b) Calcule las derivadas parciales $\frac{\partial u}{\partial x}, \frac{\partial u}{\partial y}, \frac{\partial v}{\partial x}, \frac{\partial v}{\partial y}$ en el punto $(x,y) = (0,0)$.
\end{tcolorbox}

\begin{tcolorbox}[colback=yellow!4,colframe=accentamber!70,arc=1mm,boxrule=0.5pt,left=3mm,right=3mm,top=1.5mm,bottom=1.5mm]
\textbf{Pista Conceptual:} Defina $F: \mathbb{R}^4 \to \mathbb{R}^2$. Calcule $D_{(u,v)}F(P_0)$ y compruebe que su determinante es no nulo. Luego use la fórmula $D_{(x,y)}\phi = -[D_{(u,v)}F]^{-1} D_{(x,y)}F$.
\end{tcolorbox}

\paragraph{Solución Paso a Paso:}\par
\begin{enumerate}[label=\textbf{\arabic*.}]
\item \textbf{Definición de la función vectorial y verificación del punto base:}
Definimos $F = (F_1, F_2): \mathbb{R}^4 \to \mathbb{R}^2$ por:
\begin{equation*}
F_1(x,y,u,v) = x\cos(y) + e^{x-u} + u^2 v - 1,
\end{equation*}
\begin{equation*}
F_2(x,y,u,v) = x^2 + y + \sin(u) - v.
\end{equation*}
Evaluamos en $P_0(0,0,0,0)$:
\begin{equation*}
F_1(0,0,0,0) = 0\cos(0) + e^0 + 0 - 1 = 0,
\end{equation*}
\begin{equation*}
F_2(0,0,0,0) = 0 + 0 + \sin(0) - 0 = 0.
\end{equation*}
Como las funciones están compuestas de funciones elementales suaves, $F \in C^\infty(\mathbb{R}^4, \mathbb{R}^2)$.
\end{enumerate}

\begin{enumerate}[label=\textbf{\arabic*.}]
\item \textbf{Submatriz jacobiana respecto a las variables dependientes $(u,v)$:}
\begin{equation*}
D_{(u,v)}F = \begin{pmatrix} \frac{\partial F_1}{\partial u} & \frac{\partial F_1}{\partial v} \\ \frac{\partial F_2}{\partial u} & \frac{\partial F_2}{\partial v} \end{pmatrix} = \begin{pmatrix} -e^{x-u} + 2uv & u^2 \\ \cos(u) & -1 \end{pmatrix}.
\end{equation*}
Evaluando en $P_0(0,0,0,0)$:
\begin{equation*}
D_{(u,v)}F(P_0) = \begin{pmatrix} -e^0 + 0 & 0 \\ \cos(0) & -1 \end{pmatrix} = \begin{pmatrix} -1 & 0 \\ 1 & -1 \end{pmatrix}.
\end{equation*}
Determinante:
\begin{equation*}
\det D_{(u,v)}F(P_0) = (-1)(-1) - (0)(1) = 1 \neq 0.
\end{equation*}
Dado que $F(P_0) = 0$, $F \in C^1$ y $\det D_{(u,v)}F(P_0) \neq 0$, por el Teorema de la Función Implícita existen una vecindad abierta $U \subset \mathbb{R}^2$ de $(0,0)$, una vecindad abierta $V \subset \mathbb{R}^2$ de $(0,0)$ y una única función diferenciable $\phi = (\phi_1, \phi_2): U \to V$ tal que $\phi(0,0) = (0,0)$ y $F(x, y, \phi_1(x,y), \phi_2(x,y)) = (0,0)$ para todo $(x,y) \in U$.
\end{enumerate}

\begin{enumerate}[label=\textbf{\arabic*.}]
\item \textbf{Cálculo de la submatriz jacobiana respecto a las variables independientes $(x,y)$:}
\begin{equation*}
D_{(x,y)}F = \begin{pmatrix} \cos(y) + e^{x-u} & -x\sin(y) \\ 2x & 1 \end{pmatrix} \implies D_{(x,y)}F(P_0) = \begin{pmatrix} 1 + 1 & 0 \\ 0 & 1 \end{pmatrix} = \begin{pmatrix} 2 & 0 \\ 0 & 1 \end{pmatrix}.
\end{equation*}
\end{enumerate}

\begin{enumerate}[label=\textbf{\arabic*.}]
\item \textbf{Inversa de la submatriz jacobiana:}
\begin{equation*}
[D_{(u,v)}F(P_0)]^{-1} = \begin{pmatrix} -1 & 0 \\ 1 & -1 \end{pmatrix}^{-1} = \frac{1}{1}\begin{pmatrix} -1 & 0 \\ -1 & -1 \end{pmatrix} = \begin{pmatrix} -1 & 0 \\ -1 & -1 \end{pmatrix}.
\end{equation*}
\end{enumerate}

\begin{enumerate}[label=\textbf{\arabic*.}]
\item \textbf{Multiplicación matricial y cálculo de derivadas parciales:}
\begin{equation*}
D_{(x,y)}\phi(0,0) = \begin{pmatrix} \frac{\partial u}{\partial x} & \frac{\partial u}{\partial y} \\ \frac{\partial v}{\partial x} & \frac{\partial v}{\partial y} \end{pmatrix} = - [D_{(u,v)}F(P_0)]^{-1} \cdot D_{(x,y)}F(P_0)
\end{equation*}
\begin{equation*}
= - \begin{pmatrix} -1 & 0 \\ -1 & -1 \end{pmatrix} \begin{pmatrix} 2 & 0 \\ 0 & 1 \end{pmatrix} = - \begin{pmatrix} -2 & 0 \\ -2 & -1 \end{pmatrix} = \begin{pmatrix} 2 & 0 \\ 2 & 1 \end{pmatrix}.
\end{equation*}
Por lo tanto:
\begin{equation*}
\frac{\partial u}{\partial x}(0,0) = 2, \quad \frac{\partial u}{\partial y}(0,0) = 0, \quad \frac{\partial v}{\partial x}(0,0) = 2, \quad \frac{\partial v}{\partial y}(0,0) = 1.
\end{equation*}
\end{enumerate}

\begin{trapbox}
Equivocarse en la fórmula de la inversa de matrices triangulares inferiores o no colocar el signo negativo frente al producto matricial.
\end{trapbox}

\vspace{0.5cm}
\hrule
\vspace{0.5cm}

\subsection*{Problema 2.8: Sistema no lineal 3x5 (Examen UNI 2024-II) \hfill \color{primaryblue}\small[Nivel 3]}
\addcontentsline{toc}{subsection}{Problema 2.8: \texorpdfstring{Sistema no lineal 3x5 (Examen UNI 2024-II)}{Sistema no lineal 3x5 (Examen UNI 2024-II)} (Nivel 3)}

\begin{tcolorbox}[colback=white,colframe=gray!40,arc=1mm,boxrule=0.5pt,left=3mm,right=3mm,top=2mm,bottom=2mm]
\textbf{Enunciado:} Considere el sistema de ecuaciones:
\begin{equation*}
\begin{cases} x^2 - y\cos(uv) + z^2 = 0 \\ x^2 + y^2 - \sin(uv) + 2z^2 = 2 \\ xy - \sin(u)\cos(v) + z = 0 \end{cases}
\end{equation*}
en un entorno del punto $P_0(x,y,z,u,v) = (1, 1, 0, \pi/2, 0)$.
(a) Demuestre que el sistema define a $x,y,z$ como funciones diferenciables de $(u,v)$ en una vecindad de $(\pi/2, 0)$.
(b) Calcule $\frac{\partial x}{\partial u}, \frac{\partial y}{\partial v}$ y $\frac{\partial z}{\partial v}$ en el punto $(\pi/2, 0)$.
\end{tcolorbox}

\begin{tcolorbox}[colback=yellow!4,colframe=accentamber!70,arc=1mm,boxrule=0.5pt,left=3mm,right=3mm,top=1.5mm,bottom=1.5mm]
\textbf{Pista Conceptual:} Calcule la matriz jacobiana $3\times 3$ respecto a $(x,y,z)$ en $P_0$. Luego diferencie implícitamente respecto a $u$ y $v$.
\end{tcolorbox}

\paragraph{Solución Paso a Paso:}\par
\begin{enumerate}[label=\textbf{\arabic*.}]
\item \textbf{Comprobación del punto:}
Para $(x,y,z,u,v) = (1, 1, 0, \pi/2, 0)$, $uv = 0$, $\cos(0)=1$, $\sin(0)=0$, $\sin(\pi/2)=1$, $\cos(0)=1$:
$F_1 = 1^2 - 1(1) + 0 = 0$.
$F_2 = 1^2 + 1^2 - 0 + 0 = 2$.
$F_3 = 1(1) - (1)(1) + 0 = 0$. Las 3 ecuaciones se cumplen.
\end{enumerate}

\begin{enumerate}[label=\textbf{\arabic*.}]
\item \textbf{Matriz jacobiana respecto a las variables dependientes $(x,y,z)$:}
\begin{equation*}
D_{(x,y,z)}F = \begin{pmatrix} 2x & -\cos(uv) & 2z \\ 2x & 2y & 4z \\ y & x & 1 \end{pmatrix}.
\end{equation*}
Evaluando en $P_0(1,1,0,\pi/2,0)$:
\begin{equation*}
D_{(x,y,z)}F(P_0) = \begin{pmatrix} 2 & -1 & 0 \\ 2 & 2 & 0 \\ 1 & 1 & 1 \end{pmatrix}.
\end{equation*}
Determinante (desarrollando por la tercera columna):
\begin{equation*}
\det = 1 \cdot \det\begin{pmatrix} 2 & -1 \\ 2 & 2 \end{pmatrix} = (2)(2) - (-1)(2) = 4 + 2 = 6 \neq 0.
\end{equation*}
Como $\det \neq 0$ y las funciones son $C^\infty$, el Teorema de la Función Implícita asegura que $x(u,v), y(u,v), z(u,v)$ existen y son diferenciables.
\end{enumerate}

\begin{enumerate}[label=\textbf{\arabic*.}]
\item \textbf{Derivadas respecto a $u$:}
Diferenciando el sistema con respecto a $u$ y evaluando en $P_0$ (con $v=0 \implies \frac{\partial(uv)}{\partial u} = v = 0$):
\begin{equation*}
\begin{pmatrix} 2 & -1 & 0 \\ 2 & 2 & 0 \\ 1 & 1 & 1 \end{pmatrix} \begin{pmatrix} x_u \\ y_u \\ z_u \end{pmatrix} = - \begin{pmatrix} y v \sin(uv) \\ -v\cos(uv) \\ -\cos(u)\cos(v) \end{pmatrix}_{P_0} = - \begin{pmatrix} 0 \\ 0 \\ -\cos(\pi/2)\cdot 1 \end{pmatrix} = \begin{pmatrix} 0 \\ 0 \\ 0 \end{pmatrix}.
\end{equation*}
Al ser un sistema lineal homogéneo con matriz invertible:
\begin{equation*}
x_u(\pi/2, 0) = 0, \quad y_u(\pi/2, 0) = 0, \quad z_u(\pi/2, 0) = 0.
\end{equation*}
\end{enumerate}

\begin{enumerate}[label=\textbf{\arabic*.}]
\item \textbf{Derivadas respecto a $v$:}
Diferenciando con respecto a $v$ (con $\frac{\partial(uv)}{\partial v} = u = \pi/2$):
\begin{equation*}
(F_1)_v = y u \sin(uv) \implies 1(\pi/2)\sin(0) = 0,
\end{equation*}
\begin{equation*}
(F_2)_v = -u \cos(uv) \implies -\pi/2 \cos(0) = -\pi/2,
\end{equation*}
\begin{equation*}
(F_3)_v = \sin(u)\sin(v) \implies \sin(\pi/2)\sin(0) = 0.
\end{equation*}
Por ende, el sistema para $(x_v, y_v, z_v)$ es:
\begin{equation*}
\begin{pmatrix} 2 & -1 & 0 \\ 2 & 2 & 0 \\ 1 & 1 & 1 \end{pmatrix} \begin{pmatrix} x_v \\ y_v \\ z_v \end{pmatrix} = - \begin{pmatrix} 0 \\ -\pi/2 \\ 0 \end{pmatrix} = \begin{pmatrix} 0 \\ \pi/2 \\ 0 \end{pmatrix}.
\end{equation*}
Resolviendo:
De la 1ra ecuación: $2x_v - y_v = 0 \implies y_v = 2x_v$.
En la 2da ecuación: $2x_v + 2(2x_v) = \pi/2 \implies 6x_v = \pi/2 \implies x_v = \frac{\pi}{12}$.
Entonces: $y_v = 2(\pi/12) = \frac{\pi}{6}$.
En la 3ra ecuación: $x_v + y_v + z_v = 0 \implies z_v = -(x_v + y_v) = -\left(\frac{\pi}{12} + \frac{2\pi}{12}\right) = -\frac{\pi}{4}$.
Conclusión: $\frac{\partial x}{\partial u} = 0, \quad \frac{\partial y}{\partial v} = \frac{\pi}{6}, \quad \frac{\partial z}{\partial v} = -\frac{\pi}{4}$.
\end{enumerate}

\begin{trapbox}
Olvidar que al derivar respecto a $v$, el argumento de las funciones trigonométricas contiene el producto $uv$, cuya derivada respecto a $v$ es $u = \pi/2$.
\end{trapbox}

\vspace{0.5cm}
\hrule
\vspace{0.5cm}

\subsection*{Problema 2.9: Identidad cíclica de las derivadas termodinámicas \hfill \color{primaryblue}\small[Nivel 3]}
\addcontentsline{toc}{subsection}{Problema 2.9: \texorpdfstring{Identidad cíclica de las derivadas termodinámicas}{Identidad cíclica de las derivadas termodinámicas} (Nivel 3)}

\begin{tcolorbox}[colback=white,colframe=gray!40,arc=1mm,boxrule=0.5pt,left=3mm,right=3mm,top=2mm,bottom=2mm]
\textbf{Enunciado:} Sea $F(x,y,z) = 0$ una relación implícita que define localmente a cada variable en términos de las otras dos (es decir, $x(y,z)$, $y(x,z)$ y $z(x,y)$), y suponga que las derivadas parciales $F_x, F_y, F_z$ son continuas y no nulas. Demuestre rigurosamente la regla de la cadena triple:
\begin{equation*}
\left(\frac{\partial x}{\partial y}\right)_z \left(\frac{\partial y}{\partial z}\right)_x \left(\frac{\partial z}{\partial x}\right)_y = -1.
\end{equation*}
\end{tcolorbox}

\begin{tcolorbox}[colback=yellow!4,colframe=accentamber!70,arc=1mm,boxrule=0.5pt,left=3mm,right=3mm,top=1.5mm,bottom=1.5mm]
\textbf{Pista Conceptual:} Aplique la fórmula del TFI para cada una de las tres derivadas parciales como cociente de derivadas de $F$.
\end{tcolorbox}

\paragraph{Solución Paso a Paso:}\par
\begin{enumerate}[label=\textbf{\arabic*.}]
\item \textbf{Fórmulas de las derivadas implícitas:}
Por el Teorema de la Función Implícita:
- Considerando $x = x(y,z)$ con $z$ constante:
\begin{equation*}
\left(\frac{\partial x}{\partial y}\right)_z = -\frac{F_y}{F_x}.
\end{equation*}
- Considerando $y = y(x,z)$ con $x$ constante:
\begin{equation*}
\left(\frac{\partial y}{\partial z}\right)_x = -\frac{F_z}{F_y}.
\end{equation*}
- Considerando $z = z(x,y)$ con $y$ constante:
\begin{equation*}
\left(\frac{\partial z}{\partial x}\right)_y = -\frac{F_x}{F_z}.
\end{equation*}
\end{enumerate}

\begin{enumerate}[label=\textbf{\arabic*.}]
\item \textbf{Producto de las tres derivadas:}
Multiplicando las tres expresiones:
\begin{equation*}
\left(\frac{\partial x}{\partial y}\right)_z \left(\frac{\partial y}{\partial z}\right)_x \left(\frac{\partial z}{\partial x}\right)_y = \left(-\frac{F_y}{F_x}\right) \left(-\frac{F_z}{F_y}\right) \left(-\frac{F_x}{F_z}\right)
\end{equation*}
\begin{equation*}
= (-1)^3 \cdot \frac{F_y \cdot F_z \cdot F_x}{F_x \cdot F_y \cdot F_z} = -1 \cdot 1 = -1.
\end{equation*}
Esta célebre identidad termodinámica (regla del producto triple cíclico) demuestra que las derivadas parciales no se comportan algebraicamente como fracciones ordinarias debido a los signos menos implícitos.
\end{enumerate}

\begin{trapbox}
Simplificar ingenuamente las fracciones cancelando numeradores y denominadores para obtener erroneamente $+1$.
\end{trapbox}

\vspace{0.5cm}
\hrule
\vspace{0.5cm}

\subsection*{Problema 2.10: Deducción de los Multiplicadores de Lagrange mediante el TFI \hfill \color{primaryblue}\small[Nivel 4]}
\addcontentsline{toc}{subsection}{Problema 2.10: \texorpdfstring{Deducción de los Multiplicadores de Lagrange mediante el TFI}{Deducción de los Multiplicadores de Lagrange mediante el TFI} (Nivel 4)}

\begin{tcolorbox}[colback=white,colframe=gray!40,arc=1mm,boxrule=0.5pt,left=3mm,right=3mm,top=2mm,bottom=2mm]
\textbf{Enunciado:} Sea $f: \mathbb{R}^n \to \mathbb{R}$ y $g: \mathbb{R}^n \to \mathbb{R}^k$ ($k < n$) funciones de clase $C^1$. Sea $x_0$ un extremo local de $f$ condicionado a la variedad $M = \{x \in \mathbb{R}^n : g(x) = 0\}$, y suponga que la matriz jacobiana $Dg(x_0)$ tiene rango pleno $k$. Demuestre rigurosamente mediante el Teorema de la Función Implícita que existen escalares $\lambda_1, \dots, \lambda_k \in \mathbb{R}$ tales que:
\begin{equation*}
\nabla f(x_0) = \sum_{i=1}^k \lambda_i \nabla g_i(x_0).
\end{equation*}
\end{tcolorbox}

\begin{tcolorbox}[colback=yellow!4,colframe=accentamber!70,arc=1mm,boxrule=0.5pt,left=3mm,right=3mm,top=1.5mm,bottom=1.5mm]
\textbf{Pista Conceptual:} Reordene las variables para escribir $x = (u,v)$ con $u\in\mathbb{R}^{n-k}, v\in\mathbb{R}^k$ tal que $D_v g(x_0)$ sea invertible. Aplique TFI para escribir $v = \varphi(u)$ y use que el gradiente de la función libre compuesta en un extremo libre debe anularse.
\end{tcolorbox}

\paragraph{Solución Paso a Paso:}\par
\begin{enumerate}[label=\textbf{\arabic*.}]
\item \textbf{Partición de variables e invocación del TFI:}
Dado que $\text{rango } Dg(x_0) = k$, existen $k$ columnas linealmente independientes. Reordenando coordenadas si es necesario, denotemos $x = (u, v)$ donde $u \in \mathbb{R}^{n-k}$ y $v \in \mathbb{R}^k$ con $x_0 = (u_0, v_0)$, tal que la submatriz cuadrada:
\begin{equation*}
D_v g(u_0, v_0) \in \mathbb{R}^{k \times k} \quad \text{es invertible}.
\end{equation*}
Por el Teorema de la Función Implícita, existe una vecindad $U$ de $u_0$ y una función de clase $C^1$, $\varphi: U \to \mathbb{R}^k$, tal que $\varphi(u_0) = v_0$ y:
\begin{equation*}
g(u, \varphi(u)) = 0, \quad \forall u \in U.
\end{equation*}
Además, diferenciando esta identidad:
\begin{equation*}
D_u g(u_0, v_0) + D_v g(u_0, v_0) \cdot D\varphi(u_0) = 0 \implies D\varphi(u_0) = - [D_v g(u_0, v_0)]^{-1} D_u g(u_0, v_0).
\end{equation*}
\end{enumerate}

\begin{enumerate}[label=\textbf{\arabic*.}]
\item \textbf{Condición de extremo libre en las variables $u$:}
Definimos la función compuesta $\psi(u) = f(u, \varphi(u))$ en el abierto $U$.
Como $x_0$ es un extremo local condicionado a $M$, $u_0$ es un extremo local libre de $\psi$ en el abierto $U$.
Por el Teorema de Fermat multivariable:
\begin{equation*}
D\psi(u_0) = 0.
\end{equation*}
Por la regla de la cadena:
\begin{equation*}
D\psi(u_0) = D_u f(u_0, v_0) + D_v f(u_0, v_0) \cdot D\varphi(u_0) = 0.
\end{equation*}
Sustituyendo la expresión de $D\varphi(u_0)$:
\begin{equation*}
D_u f(u_0, v_0) - D_v f(u_0, v_0) [D_v g(u_0, v_0)]^{-1} D_u g(u_0, v_0) = 0.
\end{equation*}
\end{enumerate}

\begin{enumerate}[label=\textbf{\arabic*.}]
\item \textbf{Construcción de los multiplicadores:}
Definamos el vector fila de multiplicadores $\lambda^T \in \mathbb{R}^{1 \times k}$ como:
\begin{equation*}
\lambda^T = D_v f(u_0, v_0) [D_v g(u_0, v_0)]^{-1} \iff \lambda^T D_v g(u_0, v_0) = D_v f(u_0, v_0).
\end{equation*}
Sustituyendo $\lambda^T$ en la ecuación para $u$:
\begin{equation*}
D_u f(u_0, v_0) = \lambda^T D_u g(u_0, v_0).
\end{equation*}
Uniendo ambas identidades:
\begin{equation*}
Df(x_0) = (D_u f(x_0), D_v f(x_0)) = \lambda^T (D_u g(x_0), D_v g(x_0)) = \lambda^T Dg(x_0).
\end{equation*}
Tomando transpuesta se obtiene la clásica ecuación vectorial: $\nabla f(x_0) = \sum_{i=1}^k \lambda_i \nabla g_i(x_0)$.
\end{enumerate}

\begin{trapbox}
Omitir la condición de rango pleno de $Dg(x_0)$ (cualificación de restricciones); sin ella, el gradiente de la restricción puede anularse y el teorema falla.
\end{trapbox}

\vspace{0.5cm}
\hrule
\vspace{0.5cm}

\subsection*{Problema 2.11: Superficies cónicas y convergencia al origen de planos tangentes \hfill \color{primaryblue}\small[Nivel 4]}
\addcontentsline{toc}{subsection}{Problema 2.11: \texorpdfstring{Superficies cónicas y convergencia al origen de planos tangentes}{Superficies cónicas y convergencia al origen de planos tangentes} (Nivel 4)}

\begin{tcolorbox}[colback=white,colframe=gray!40,arc=1mm,boxrule=0.5pt,left=3mm,right=3mm,top=2mm,bottom=2mm]
\textbf{Enunciado:} Sea $F: \mathbb{R}^2 \to \mathbb{R}$ diferenciable y sea la superficie $S$ dada implícitamente por $H(x,y,z) = F\left(\frac{x}{z}, \frac{y}{z}\right) = 0$ con $z \neq 0$. Demuestre que el plano tangente a $S$ en cualquier punto regular $P_0(x_0, y_0, z_0)$ pasa necesariamente por el origen $(0,0,0)$.
\end{tcolorbox}

\begin{tcolorbox}[colback=yellow!4,colframe=accentamber!70,arc=1mm,boxrule=0.5pt,left=3mm,right=3mm,top=1.5mm,bottom=1.5mm]
\textbf{Pista Conceptual:} Calcule $\nabla H$ usando la regla de la cadena y evalúe la ecuación del plano tangente $\nabla H(P_0) \cdot (x - x_0, y - y_0, z - z_0) = 0$ en $(0,0,0)$. Compruebe que $\nabla H(P_0) \cdot P_0 = 0$.
\end{tcolorbox}

\paragraph{Solución Paso a Paso:}\par
\begin{enumerate}[label=\textbf{\arabic*.}]
\item \textbf{Cálculo del gradiente de $H$:}
Sean $u = x/z$ y $v = y/z$. Por la regla de la cadena:
\begin{equation*}
H_x = \frac{\partial F}{\partial u} \frac{\partial u}{\partial x} + \frac{\partial F}{\partial v} \frac{\partial v}{\partial x} = F_u \cdot \frac{1}{z},
\end{equation*}
\begin{equation*}
H_y = \frac{\partial F}{\partial u} \frac{\partial u}{\partial y} + \frac{\partial F}{\partial v} \frac{\partial v}{\partial y} = F_v \cdot \frac{1}{z},
\end{equation*}
\begin{equation*}
H_z = \frac{\partial F}{\partial u} \frac{\partial u}{\partial z} + \frac{\partial F}{\partial v} \frac{\partial v}{\partial z} = F_u \cdot \left(-\frac{x}{z^2}\right) + F_v \cdot \left(-\frac{y}{z^2}\right) = -\frac{1}{z} \left( u F_u + v F_v \right).
\end{equation*}
\end{enumerate}

\begin{enumerate}[label=\textbf{\arabic*.}]
\item \textbf{Evaluación del producto escalar $\nabla H(P_0) \cdot P_0$:}
Realizamos el producto escalar entre el vector gradiente y el vector posición $P_0 = (x_0, y_0, z_0)$:
\begin{equation*}
\nabla H(P_0) \cdot P_0 = x_0 H_x + y_0 H_y + z_0 H_z
\end{equation*}
\begin{equation*}
= x_0 \left(\frac{F_u}{z_0}\right) + y_0 \left(\frac{F_v}{z_0}\right) + z_0 \left(-\frac{x_0 F_u + y_0 F_v}{z_0^2}\right)
\end{equation*}
\begin{equation*}
= \frac{x_0 F_u + y_0 F_v}{z_0} - \frac{x_0 F_u + y_0 F_v}{z_0} = 0.
\end{equation*}
\end{enumerate}

\begin{enumerate}[label=\textbf{\arabic*.}]
\item \textbf{Ecuación del plano tangente:}
El plano tangente en $P_0$ tiene por ecuación:
\begin{equation*}
\nabla H(P_0) \cdot (X - P_0) = 0 \iff \nabla H(P_0) \cdot X - \nabla H(P_0) \cdot P_0 = 0.
\end{equation*}
Como acabamos de demostrar que $\nabla H(P_0) \cdot P_0 = 0$, la ecuación del plano tangente se reduce a:
\begin{equation*}
\nabla H(P_0) \cdot X = 0.
\end{equation*}
Evaluando en el origen $X = (0,0,0)$: $\nabla H(P_0) \cdot (0,0,0) = 0$, lo que satisface idénticamente la ecuación.
Por tanto, el plano tangente pasa por el origen para todo punto regular.
\end{enumerate}

\begin{trapbox}
Derivar $1/z$ incorrectamente o no relacionar la condición de paso por el origen con el producto escalar $\nabla H \cdot P_0 = 0$.
\end{trapbox}

\vspace{0.5cm}
\hrule
\vspace{0.5cm}

\subsection*{Problema 2.12: Bifurcación en puntos singulares donde det DyF = 0 \hfill \color{primaryblue}\small[Nivel 4]}
\addcontentsline{toc}{subsection}{Problema 2.12: \texorpdfstring{Bifurcación en puntos singulares donde det DyF = 0}{Bifurcación en puntos singulares donde det DyF = 0} (Nivel 4)}

\begin{tcolorbox}[colback=white,colframe=gray!40,arc=1mm,boxrule=0.5pt,left=3mm,right=3mm,top=2mm,bottom=2mm]
\textbf{Enunciado:} Considere la ecuación $F(x,y) = y^2 - x(x - 1)^2 = 0$. (a) Determine todos los puntos donde $F(x,y) = 0$ y $\nabla F(x,y) = (0,0)$. (b) Analice si es posible despejar $y = \varphi(x)$ en un entorno del punto singular $(1,0)$ y clasifique geométricamente la curva cerca de dicho punto.
\end{tcolorbox}

\begin{tcolorbox}[colback=yellow!4,colframe=accentamber!70,arc=1mm,boxrule=0.5pt,left=3mm,right=3mm,top=1.5mm,bottom=1.5mm]
\textbf{Pista Conceptual:} Halle las soluciones del sistema $F=0, F_x=0, F_y=0$. Luego use desarrollo de Taylor local para aproximar la curva.
\end{tcolorbox}

\paragraph{Solución Paso a Paso:}\par
\begin{enumerate}[label=\textbf{\arabic*.}]
\item \textbf{Búsqueda de puntos singulares:}
Gradiente de $F$:
\begin{equation*}
F_x = -(x - 1)^2 - 2x(x - 1) = -(x - 1)[(x - 1) + 2x] = -(x - 1)(3x - 1),
\end{equation*}
\begin{equation*}
F_y = 2y.
\end{equation*}
Para que $\nabla F = (0,0)$, requerimos $y = 0$ y $(x=1$ o $x=1/3)$.
Evaluando en $F(x,y) = 0$:
- Si $x=1$: $F(1,0) = 0^2 - 1(0)^2 = 0 \implies (1,0)$ es un punto singular.
- Si $x=1/3$: $F(1/3, 0) = 0 - (1/3)(-2/3)^2 = -4/27 \neq 0$, no pertenece a la curva.
Por tanto, el único punto singular de la curva es $P_0 = (1,0)$.
\end{enumerate}

\begin{enumerate}[label=\textbf{\arabic*.}]
\item \textbf{Falla de las hipótesis del TFI:}
En $P_0(1,0)$, $F_y(1,0) = 0$. El Teorema de la Función Implícita no es aplicable.
\end{enumerate}

\begin{enumerate}[label=\textbf{\arabic*.}]
\item \textbf{Análisis de Taylor y comportamiento local:}
Hagamos el cambio de coordenadas centrado en $(1,0)$: sea $u = x - 1$.
Entonces $x = 1 + u$, y la ecuación queda:
\begin{equation*}
y^2 = (1 + u) u^2 = u^2 + u^3.
\end{equation*}
Cerca de $u = 0$, $u^3$ es despreciable respecto a $u^2$:
\begin{equation*}
y^2 \approx u^2 \implies y \approx \pm u \implies y \approx \pm (x - 1).
\end{equation*}
La curva posee un \textbf{punto nodal (cruce de dos ramas tangentes a las rectas $y = x - 1$ y $y = -(x - 1)$)}.
Al haber dos valores de $y$ para cada $x$ en una vecindad de $1$, es imposible expresar $y$ como una función univaluada de $x$. Esto ilustra la bifurcación geométrica generada por la nulidad del jacobiano.
\end{enumerate}

\begin{trapbox}
Intentar aplicar la fórmula del cociente cuando $F_y = 0$, lo que conduce a una división por cero.
\end{trapbox}

\vspace{0.5cm}
\hrule
\vspace{0.5cm}

\subsection*{Problema 2.13: Demostración rigurosa del TFI mediante el Teorema de la Función Inversa \hfill \color{primaryblue}\small[Nivel 5]}
\addcontentsline{toc}{subsection}{Problema 2.13: \texorpdfstring{Demostración rigurosa del TFI mediante el Teorema de la Función Inversa}{Demostración rigurosa del TFI mediante el Teorema de la Función Inversa} (Nivel 5)}

\begin{tcolorbox}[colback=white,colframe=gray!40,arc=1mm,boxrule=0.5pt,left=3mm,right=3mm,top=2mm,bottom=2mm]
\textbf{Enunciado:} Sea $U \subseteq \mathbb{R}^n \times \mathbb{R}^m$ un abierto, $(x_0, y_0) \in U$ y $F: U \to \mathbb{R}^m$ de clase $C^1$ con $F(x_0, y_0) = 0$ y $D_y F(x_0, y_0)$ invertible. Defina la aplicación auxiliar:
\begin{equation*}
G(x,y) = (x, F(x,y)) \in \mathbb{R}^n \times \mathbb{R}^m.
\end{equation*}
Demuestre el Teorema de la Función Implícita a partir del Teorema de la Función Inversa aplicado a $G$.
\end{tcolorbox}

\begin{tcolorbox}[colback=yellow!4,colframe=accentamber!70,arc=1mm,boxrule=0.5pt,left=3mm,right=3mm,top=1.5mm,bottom=1.5mm]
\textbf{Pista Conceptual:} Calcule la matriz jacobiana por bloques de $G$, demuestre que su determinante es $\det D_y F(x_0, y_0) \neq 0$ y construya la inversa $G^{-1}(x, 0)$.
\end{tcolorbox}

\paragraph{Solución Paso a Paso:}\par
\begin{enumerate}[label=\textbf{\arabic*.}]
\item \textbf{Matriz jacobiana por bloques de $G$:}
La función $G: U \subset \mathbb{R}^{n+m} \to \mathbb{R}^{n+m}$ está dada por $G(x,y) = (x, F(x,y))$.
Su matriz diferencial en $(x_0, y_0)$ expresada en bloques es:
\begin{equation*}
DG(x_0, y_0) = \begin{pmatrix} I_n & 0 \\ D_x F(x_0, y_0) & D_y F(x_0, y_0) \end{pmatrix}.
\end{equation*}
El determinante de una matriz triangular por bloques es el producto de los determinantes de los bloques diagonales:
\begin{equation*}
\det DG(x_0, y_0) = \det(I_n) \cdot \det(D_y F(x_0, y_0)) = 1 \cdot \det(D_y F(x_0, y_0)) \neq 0.
\end{equation*}
\end{enumerate}

\begin{enumerate}[label=\textbf{\arabic*.}]
\item \textbf{Aplicación del Teorema de la Función Inversa:}
Como $G$ es de clase $C^1$ y $\det DG(x_0, y_0) \neq 0$, por el Teorema de la Función Inversa existen un entorno abierto $W$ de $(x_0, y_0)$ y un entorno abierto $V$ de $G(x_0, y_0) = (x_0, 0)$ tal que $G|_W: W \to V$ es un difeomorfismo de clase $C^1$.
\end{enumerate}

\begin{enumerate}[label=\textbf{\arabic*.}]
\item \textbf{Estructura de la función inversa $G^{-1}$:}
Sea $H = G^{-1}: V \to W$. Dado que $G(x,y) = (x, F(x,y))$, la primera componente de $G$ es la proyección canónica $x$.
Por tanto, la inversa debe respetar la primera coordenada:
\begin{equation*}
H(x, z) = (x, \psi(x, z)),
\end{equation*}
para alguna función diferenciable $\psi: V \to \mathbb{R}^m$.
\end{enumerate}

\begin{enumerate}[label=\textbf{\arabic*.}]
\item \textbf{Construcción de la función implícita $\varphi$:}
Para los puntos donde $z = 0$, el conjunto $U_0 = \{x \in \mathbb{R}^n : (x, 0) \in V\}$ es un abierto de $\mathbb{R}^n$ que contiene a $x_0$.
Definimos $\varphi(x) = \psi(x, 0)$ para $x \in U_0$.
Entonces $(x, \varphi(x)) = H(x, 0) \in W$, y aplicando $G$ a ambos lados:
\begin{equation*}
G(x, \varphi(x)) = G(H(x, 0)) = (x, 0).
\end{equation*}
Por la definición de $G$:
\begin{equation*}
(x, F(x, \varphi(x))) = (x, 0) \implies F(x, \varphi(x)) = 0, \quad \forall x \in U_0.
\end{equation*}
Además $\varphi$ es de clase $C^1$ por ser composición de la aplicación $C^1$ $H$ con la inclusión continua $x \mapsto (x, 0)$. Esto completa la demostración.
\end{enumerate}

\begin{trapbox}
No justificar rigurosamente por qué la primera componente de $G^{-1}$ es idénticamente la coordenada $x$.
\end{trapbox}

\vspace{0.5cm}
\hrule
\vspace{0.5cm}

\subsection*{Problema 2.14: Regularidad óptima: Herencia de clase $C^k$ de la función implícita \hfill \color{primaryblue}\small[Nivel 5]}
\addcontentsline{toc}{subsection}{Problema 2.14: \texorpdfstring{Regularidad óptima: Herencia de clase $C^k$ de la función implícita}{Regularidad óptima: Herencia de clase C^k de la función implícita} (Nivel 5)}

\begin{tcolorbox}[colback=white,colframe=gray!40,arc=1mm,boxrule=0.5pt,left=3mm,right=3mm,top=2mm,bottom=2mm]
\textbf{Enunciado:} Sea $F: U \subseteq \mathbb{R}^n \times \mathbb{R}^m \to \mathbb{R}^m$ de clase $C^k$ con $k \ge 1$, y sea $\varphi: V \to \mathbb{R}^m$ la función continua dada por el TFI tal que $F(x, \varphi(x)) = 0$. Demuestre por inducción sobre $k$ que $\varphi$ es exactamente de clase $C^k$.
\end{tcolorbox}

\begin{tcolorbox}[colback=yellow!4,colframe=accentamber!70,arc=1mm,boxrule=0.5pt,left=3mm,right=3mm,top=1.5mm,bottom=1.5mm]
\textbf{Pista Conceptual:} Use la relación $D\varphi(x) = -[D_y F(x, \varphi(x))]^{-1} D_x F(x, \varphi(x))$ y el hecho de que la inversión matricial $A \mapsto A^{-1}$ es de clase $C^\infty$.
\end{tcolorbox}

\paragraph{Solución Paso a Paso:}\par
\begin{enumerate}[label=\textbf{\arabic*.}]
\item \textbf{Caso base ($k = 1$):}
El Teorema de la Función Implícita estándar asegura que si $F \in C^1$, entonces $\varphi$ es diferenciable y satisface:
\begin{equation*}
D\varphi(x) = - [D_y F(x, \varphi(x))]^{-1} D_x F(x, \varphi(x)).
\end{equation*}
Como $F \in C^1$, las funciones $(x,y) \mapsto D_x F(x,y)$ y $(x,y) \mapsto D_y F(x,y)$ son continuas.
La aplicación de inversión matricial $\text{Inv}: GL(m, \mathbb{R}) \to GL(m, \mathbb{R})$ dada por $A \mapsto A^{-1}$ es racional en los coeficientes (por la regla de Cramer), por ende es $C^\infty$.
Al ser $\varphi$ continua, la composición $x \mapsto [D_y F(x, \varphi(x))]^{-1}$ es continua, y por tanto $D\varphi(x)$ es producto de funciones continuas, concluyendo que $D\varphi$ es continua, es decir, $\varphi \in C^1$.
\end{enumerate}

\begin{enumerate}[label=\textbf{\arabic*.}]
\item \textbf{Paso inductivo:}
Supongamos como hipótesis de inducción que la propiedad es cierta para $k-1 \ge 1$: toda función $F \in C^{k-1}$ genera una función implícita $\varphi \in C^{k-1}$.
Sea ahora $F \in C^k$. Por la hipótesis inductiva, como $F \in C^{k-1}$, ya sabemos que $\varphi \in C^{k-1}$.
Analicemos la expresión de su diferencial:
\begin{equation*}
D\varphi(x) = - \text{Inv}(D_y F(x, \varphi(x))) \cdot D_x F(x, \varphi(x)).
\end{equation*}
En esta expresión:
- $D_x F$ y $D_y F$ son de clase $C^{k-1}$ (al ser derivadas de una función $C^k$).
- $\varphi$ es de clase $C^{k-1}$ (por hipótesis inductiva).
- La inversión de matrices es $C^\infty$.
Por el teorema de composición y producto de funciones diferenciables, la composición de funciones $C^{k-1}$ es de clase $C^{k-1}$.
Por consiguiente, la función $x \mapsto D\varphi(x)$ es de clase $C^{k-1}$.
Tener una derivada de clase $C^{k-1}$ implica por definición que $\varphi$ es de clase $C^k$.
\end{enumerate}

\begin{enumerate}[label=\textbf{\arabic*.}]
\item \textbf{Conclusión:} Por inducción matemática, para todo $k \ge 1$ (incluyendo $k = \infty$), si $F \in C^k$, entonces $\varphi \in C^k$.
\end{enumerate}

\begin{trapbox}
Olvidar justificar la suavidad de la función inversa de matrices $A \mapsto A^{-1}$.
\end{trapbox}

\vspace{0.5cm}
\hrule
\vspace{0.5cm}

\subsection*{Problema 2.15: Teorema de la submersión y estructura de subvariedades \hfill \color{primaryblue}\small[Nivel 5]}
\addcontentsline{toc}{subsection}{Problema 2.15: \texorpdfstring{Teorema de la submersión y estructura de subvariedades}{Teorema de la submersión y estructura de subvariedades} (Nivel 5)}

\begin{tcolorbox}[colback=white,colframe=gray!40,arc=1mm,boxrule=0.5pt,left=3mm,right=3mm,top=2mm,bottom=2mm]
\textbf{Enunciado:} Sea $U \subseteq \mathbb{R}^{n+m}$ un abierto y $F: U \to \mathbb{R}^m$ una submersión de clase $C^k$ ($k \ge 1$) (es decir, $\text{rango } DF(p) = m$ para todo $p \in U$). Demuestre que para cualquier valor regular $c \in \mathbb{R}^m$, el conjunto de nivel $M = F^{-1}(c)$ es una subvariedad diferenciable de dimensión $n$ de $\mathbb{R}^{n+m}$.
\end{tcolorbox}

\begin{tcolorbox}[colback=yellow!4,colframe=accentamber!70,arc=1mm,boxrule=0.5pt,left=3mm,right=3mm,top=1.5mm,bottom=1.5mm]
\textbf{Pista Conceptual:} Muestre que cada punto de $M$ tiene una vecindad que es difeomorfa a un abierto de $\mathbb{R}^n \times \{0\}^m$ mediante una carta local construida con el TFI.
\end{tcolorbox}

\paragraph{Solución Paso a Paso:}\par
\begin{enumerate}[label=\textbf{\arabic*.}]
\item \textbf{Definición de carta de subvariedad:}
Un conjunto $M \subset \mathbb{R}^{n+m}$ es una subvariedad de dimensión $n$ y clase $C^k$ si para cada punto $p_0 \in M$ existe un abierto $V \subset \mathbb{R}^{n+m}$ conteniendo a $p_0$ y un difeomorfismo $\Phi: V \to W \subset \mathbb{R}^{n+m}$ tal que:
\begin{equation*}
\Phi(V \cap M) = W \cap (\mathbb{R}^n \times \{0\}^m).
\end{equation*}
\end{enumerate}

\begin{enumerate}[label=\textbf{\arabic*.}]
\item \textbf{Construcción de la carta local:}
Sea $p_0 = (x_0, y_0) \in M$, de modo que $F(x_0, y_0) = c$.
Como $\text{rango } DF(p_0) = m$, existen $m$ columnas linealmente independientes. Reordenando coordenadas si es necesario, denotemos los puntos como $(x,y) \in \mathbb{R}^n \times \mathbb{R}^m$ tal que:
\begin{equation*}
\det D_y F(x_0, y_0) \neq 0.
\end{equation*}
Definamos la aplicación $\Phi: U \to \mathbb{R}^{n+m}$ por:
\begin{equation*}
\Phi(x,y) = (x, F(x,y) - c).
\end{equation*}
La matriz diferencial en $p_0$ es:
\begin{equation*}
D\Phi(p_0) = \begin{pmatrix} I_n & 0 \\ D_x F(p_0) & D_y F(p_0) \end{pmatrix}.
\end{equation*}
Su determinante es $\det D\Phi(p_0) = \det(I_n) \det(D_y F(p_0)) \neq 0$.
Por el Teorema de la Función Inversa, existe un abierto $V \ni p_0$ tal que $\Phi|_V$ es un difeomorfismo $C^k$ sobre un abierto $W = \Phi(V) \subset \mathbb{R}^{n+m}$.
\end{enumerate}

\begin{enumerate}[label=\textbf{\arabic*.}]
\item \textbf{Verificación de la condición de variedad:}
Para cualquier $(x,y) \in V$:
\begin{equation*}
(x,y) \in M \iff F(x,y) = c \iff F(x,y) - c = 0 \iff \Phi(x,y) = (x, 0) \in \mathbb{R}^n \times \{0\}^m.
\end{equation*}
Por ende:
\begin{equation*}
\Phi(V \cap M) = W \cap (\mathbb{R}^n \times \{0\}^m).
\end{equation*}
Esto demuestra que el difeomorfismo $\Phi$ ``endereza'' localmente el conjunto de nivel $M$ transformándolo en el subespacio euclidiano estándar $\mathbb{R}^n$.
Por tanto, $M = F^{-1}(c)$ es una subvariedad diferenciable de dimensión $n$ en $\mathbb{R}^{n+m}$.
\end{enumerate}

\begin{trapbox}
Confundir la dimensión de la variedad ($n$) con la codimensión impuesta por las restricciones ($m$).
\end{trapbox}

\vspace{0.5cm}
\hrule
\vspace{0.5cm}

\section{Particiones y Sumas de Darboux/Riemann}

\subsection*{Problema 3.1: Partición regular y geometría de subrectángulos \hfill \color{primaryblue}\small[Nivel 1]}
\addcontentsline{toc}{subsection}{Problema 3.1: \texorpdfstring{Partición regular y geometría de subrectángulos}{Partición regular y geometría de subrectángulos} (Nivel 1)}

\begin{tcolorbox}[colback=white,colframe=gray!40,arc=1mm,boxrule=0.5pt,left=3mm,right=3mm,top=2mm,bottom=2mm]
\textbf{Enunciado:} Sea el rectángulo $R = [0, 2] \times [0, 4] \subset \mathbb{R}^2$. Considere la partición regular $P = P_1 \times P_2$ con $m = 2$ subdivisiones en el eje $x$ y $n = 2$ subdivisiones en el eje $y$. (a) Determine explícitamente los subintervalos de $P_1$ y $P_2$. (b) Liste los 4 subrectángulos $R_{ij}$ y verifique que la suma de sus áreas es igual al área de $R$.
\end{tcolorbox}

\begin{tcolorbox}[colback=yellow!4,colframe=accentamber!70,arc=1mm,boxrule=0.5pt,left=3mm,right=3mm,top=1.5mm,bottom=1.5mm]
\textbf{Pista Conceptual:} El paso uniforme es $\Delta x = (b-a)/m$ y $\Delta y = (d-c)/n$. Cada subrectángulo es $R_{ij} = [x_{i-1}, x_i] \times [y_{j-1}, y_j]$.
\end{tcolorbox}

\paragraph{Solución Paso a Paso:}\par
\begin{enumerate}[label=\textbf{\arabic*.}]
\item \textbf{Subdivisiones regulares:}
- Para el intervalo $[0, 2]$ con $m = 2$:
\begin{equation*}
\Delta x = \frac{2 - 0}{2} = 1 \implies x_0 = 0, \ x_1 = 1, \ x_2 = 2 \implies P_1 = \{0, 1, 2\}.
\end{equation*}
- Para el intervalo $[0, 4]$ con $n = 2$:
\begin{equation*}
\Delta y = \frac{4 - 0}{2} = 2 \implies y_0 = 0, \ y_1 = 2, \ y_2 = 4 \implies P_2 = \{0, 2, 4\}.
\end{equation*}
\end{enumerate}

\begin{enumerate}[label=\textbf{\arabic*.}]
\item \textbf{Subrectángulos:}
La partición producto $P = P_1 \times P_2$ define $2 \times 2 = 4$ subrectángulos:
- $R_{11} = [0, 1] \times [0, 2]$,
- $R_{12} = [0, 1] \times [2, 4]$,
- $R_{21} = [1, 2] \times [0, 2]$,
- $R_{22} = [1, 2] \times [2, 4]$.
\end{enumerate}

\begin{enumerate}[label=\textbf{\arabic*.}]
\item \textbf{Cálculo de áreas:}
Cada subrectángulo tiene área:
\begin{equation*}
\text{Área}(R_{ij}) = \Delta x \cdot \Delta y = 1 \cdot 2 = 2.
\end{equation*}
La suma total de áreas es:
\begin{equation*}
\sum_{i=1}^2 \sum_{j=1}^2 \text{Área}(R_{ij}) = 4 \times 2 = 8.
\end{equation*}
Por otra parte, el área del rectángulo original es $\text{Área}(R) = (2 - 0)(4 - 0) = 8$. Se comprueba la partición exacta.
\end{enumerate}

\begin{trapbox}
Confundir los índices de los puntos extremos de la partición con el número de subintervalos generados.
\end{trapbox}

\vspace{0.5cm}
\hrule
\vspace{0.5cm}

\subsection*{Problema 3.2: Cálculo explícito de sumas de Darboux inferior y superior \hfill \color{primaryblue}\small[Nivel 1]}
\addcontentsline{toc}{subsection}{Problema 3.2: \texorpdfstring{Cálculo explícito de sumas de Darboux inferior y superior}{Cálculo explícito de sumas de Darboux inferior y superior} (Nivel 1)}

\begin{tcolorbox}[colback=white,colframe=gray!40,arc=1mm,boxrule=0.5pt,left=3mm,right=3mm,top=2mm,bottom=2mm]
\textbf{Enunciado:} Sea $f(x,y) = 2x + 3y$ definida sobre el cuadrado $R = [0, 1] \times [0, 1]$. Considere la partición uniforme $P_1 = P_2 = \{0, 1/2, 1\}$. Calcule la suma inferior $L(f,P)$, la suma superior $U(f,P)$ y la diferencia $U(f,P) - L(f,P)$.
\end{tcolorbox}

\begin{tcolorbox}[colback=yellow!4,colframe=accentamber!70,arc=1mm,boxrule=0.5pt,left=3mm,right=3mm,top=1.5mm,bottom=1.5mm]
\textbf{Pista Conceptual:} Como $f$ es creciente en ambas variables, el mínimo en cada subrectángulo se alcanza en la esquina inferior izquierda y el máximo en la superior derecha.
\end{tcolorbox}

\paragraph{Solución Paso a Paso:}\par
\begin{enumerate}[label=\textbf{\arabic*.}]
\item \textbf{Geometría de la partición:}
Hay $2 \times 2 = 4$ subrectángulos, cada uno con área $\Delta A = \frac{1}{2} \times \frac{1}{2} = \frac{1}{4}$.
Los subrectángulos son $R_{ij} = [\frac{i-1}{2}, \frac{i}{2}] \times [\frac{j-1}{2}, \frac{j}{2}]$ para $i, j \in \{1, 2\}$.
\end{enumerate}

\begin{enumerate}[label=\textbf{\arabic*.}]
\item \textbf{Ínfimos y sumas inferiores $L(f,P)$:}
Al ser $f(x,y) = 2x + 3y$ monótona creciente en $x$ e $y$, el ínfimo en $R_{ij}$ es:
\begin{equation*}
m_{ij} = \inf_{R_{ij}} f = f\left(\frac{i-1}{2}, \frac{j-1}{2}\right) = 2\left(\frac{i-1}{2}\right) + 3\left(\frac{j-1}{2}\right) = (i-1) + \frac{3}{2}(j-1).
\end{equation*}
- $m_{11} = f(0, 0) = 0$,
- $m_{12} = f(0, 1/2) = 3/2$,
- $m_{21} = f(1/2, 0) = 1$,
- $m_{22} = f(1/2, 1/2) = 1 + 3/2 = 5/2$.
Suma inferior:
\begin{equation*}
L(f,P) = \sum_{i,j} m_{ij} \Delta A = \left(0 + \frac{3}{2} + 1 + \frac{5}{2}\right) \frac{1}{4} = 5 \cdot \frac{1}{4} = \frac{5}{4} = 1.25.
\end{equation*}
\end{enumerate}

\begin{enumerate}[label=\textbf{\arabic*.}]
\item \textbf{Supremos y sumas superiores $U(f,P)$:}
El supremo en $R_{ij}$ se alcanza en la esquina superior derecha:
\begin{equation*}
M_{ij} = \sup_{R_{ij}} f = f\left(\frac{i}{2}, \frac{j}{2}\right) = i + \frac{3}{2}j.
\end{equation*}
- $M_{11} = f(1/2, 1/2) = 5/2$,
- $M_{12} = f(1/2, 1) = 1 + 3 = 4$,
- $M_{21} = f(1, 1/2) = 2 + 3/2 = 7/2$,
- $M_{22} = f(1, 1) = 2 + 3 = 5$.
Suma superior:
\begin{equation*}
U(f,P) = \sum_{i,j} M_{ij} \Delta A = \left(\frac{5}{2} + 4 + \frac{7}{2} + 5\right) \frac{1}{4} = 15 \cdot \frac{1}{4} = \frac{15}{4} = 3.75.
\end{equation*}
\end{enumerate}

\begin{enumerate}[label=\textbf{\arabic*.}]
\item \textbf{Diferencia de Darboux:}
\begin{equation*}
U(f,P) - L(f,P) = \frac{15}{4} - \frac{5}{4} = \frac{10}{4} = \frac{5}{2} = 2.5.
\end{equation*}
Nótese que la integral exacta es $\iint_R (2x+3y)dA = 1 + \frac{3}{2} = 2.5$, y se cumple estrictamente $1.25 < 2.5 < 3.75$.
\end{enumerate}

\begin{trapbox}
Evaluar el ínfimo o supremo en puntos interiores en vez de los vértices extremos para funciones afines.
\end{trapbox}

\vspace{0.5cm}
\hrule
\vspace{0.5cm}

\subsection*{Problema 3.3: Suma de Riemann con puntos de muestra en el centro \hfill \color{primaryblue}\small[Nivel 1]}
\addcontentsline{toc}{subsection}{Problema 3.3: \texorpdfstring{Suma de Riemann con puntos de muestra en el centro}{Suma de Riemann con puntos de muestra en el centro} (Nivel 1)}

\begin{tcolorbox}[colback=white,colframe=gray!40,arc=1mm,boxrule=0.5pt,left=3mm,right=3mm,top=2mm,bottom=2mm]
\textbf{Enunciado:} Sea $f(x,y) = xy$ en el rectángulo $R = [0, 2] \times [0, 2]$. Considere la partición regular con $m = n = 2$. Elija como puntos de muestra $\xi_{ij} = (x_i^*, y_j^*)$ los centros geométricos de cada subrectángulo. Calcule la suma de Riemann $S(f, P, \xi)$ y compárela con el valor exacto de la integral.
\end{tcolorbox}

\begin{tcolorbox}[colback=yellow!4,colframe=accentamber!70,arc=1mm,boxrule=0.5pt,left=3mm,right=3mm,top=1.5mm,bottom=1.5mm]
\textbf{Pista Conceptual:} El centro de cada subrectángulo $[x_{i-1}, x_i] \times [y_{j-1}, y_j]$ es $(\frac{x_{i-1}+x_i}{2}, \frac{y_{j-1}+y_j}{2})$.
\end{tcolorbox}

\paragraph{Solución Paso a Paso:}\par
\begin{enumerate}[label=\textbf{\arabic*.}]
\item \textbf{Subrectángulos y áreas:}
Con $m=n=2$ en $[0,2]\times[0,2]$, $\Delta x = 1$ y $\Delta y = 1$, con $\Delta A = 1$.
Los subintervalos son $[0,1]$ y $[1,2]$ en cada eje.
\end{enumerate}

\begin{enumerate}[label=\textbf{\arabic*.}]
\item \textbf{Puntos medios:}
- $R_{11} = [0,1]\times[0,1] \implies \xi_{11} = (1/2, 1/2)$,
- $R_{12} = [0,1]\times[1,2] \implies \xi_{12} = (1/2, 3/2)$,
- $R_{21} = [1,2]\times[0,1] \implies \xi_{21} = (3/2, 1/2)$,
- $R_{22} = [1,2]\times[1,2] \implies \xi_{22} = (3/2, 3/2)$.
\end{enumerate}

\begin{enumerate}[label=\textbf{\arabic*.}]
\item \textbf{Cálculo de la suma de Riemann:}
\begin{equation*}
S(f,P,\xi) = \sum_{i=1}^2 \sum_{j=1}^2 f(\xi_{ij}) \Delta A
\end{equation*}
\begin{equation*}
= f(1/2, 1/2)(1) + f(1/2, 3/2)(1) + f(3/2, 1/2)(1) + f(3/2, 3/2)(1)
\end{equation*}
\begin{equation*}
= \frac{1}{4} + \frac{3}{4} + \frac{3}{4} + \frac{9}{4} = \frac{1 + 3 + 3 + 9}{4} = \frac{16}{4} = 4.
\end{equation*}
\end{enumerate}

\begin{enumerate}[label=\textbf{\arabic*.}]
\item \textbf{Comparación con el valor exacto:}
\begin{equation*}
\iint_R xy \, dA = \left(\int_0^2 x \, dx\right) \left(\int_0^2 y \, dy\right) = \left[\frac{x^2}{2}\right]_0^2 \left[\frac{y^2}{2}\right]_0^2 = 2 \times 2 = 4.
\end{equation*}
La regla del punto medio coincide exactamente con el valor real de la integral debido a la linealidad simétrica de la función bilineal $xy$ respecto al punto medio.
\end{enumerate}

\begin{trapbox}
Tomar puntos muestrales fuera de los respectivos subrectángulos.
\end{trapbox}

\vspace{0.5cm}
\hrule
\vspace{0.5cm}

\subsection*{Problema 3.4: Suma de Riemann finita en esquina superior derecha (Examen UNI 2024-II) \hfill \color{primaryblue}\small[Nivel 2]}
\addcontentsline{toc}{subsection}{Problema 3.4: \texorpdfstring{Suma de Riemann finita en esquina superior derecha (Examen UNI 2024-II)}{Suma de Riemann finita en esquina superior derecha (Examen UNI 2024-II)} (Nivel 2)}

\begin{tcolorbox}[colback=white,colframe=gray!40,arc=1mm,boxrule=0.5pt,left=3mm,right=3mm,top=2mm,bottom=2mm]
\textbf{Enunciado:} Sea el rectángulo $R = [0, 1] \times [0, 2]$. Usamos las particiones $\{x_i = \frac{i}{m} : i=0,\dots,m\}$ de $[0,1]$ y $\{y_j = \frac{j}{n} : j=0,\dots,2n\}$ de $[0,2]$ para calcular la suma de Riemann:
\begin{equation*}
S_{m,n} = \sum_{i=1}^m \sum_{j=1}^{2n} f(x_i^*, y_j^*) \text{Área}(R_{ij}),
\end{equation*}
donde $(x_i^*, y_j^*)$ es la esquina superior derecha de cada subrectángulo y $f(x,y) = x + y$. Determine el valor numérico exacto de $S_{2,2}$.
\end{tcolorbox}

\begin{tcolorbox}[colback=yellow!4,colframe=accentamber!70,arc=1mm,boxrule=0.5pt,left=3mm,right=3mm,top=1.5mm,bottom=1.5mm]
\textbf{Pista Conceptual:} Para $m=2$ y $n=2$, el número de subdivisiones en $y$ es $2n = 4$. Identifique $\Delta x = 1/2$, $\Delta y = 2/4 = 1/2$. Los puntos de muestra son $(i/2, j/2)$ con $i=1,2$ y $j=1,2,3,4$.
\end{tcolorbox}

\paragraph{Solución Paso a Paso:}\par
\begin{enumerate}[label=\textbf{\arabic*.}]
\item \textbf{Parámetros de la partición:}
Con $m = 2$ y $n = 2$:
- En $[0, 1]$, $m = 2$ subintervalos: $x_i = i/2$ para $i=0,1,2 \implies \Delta x = 1/2$.
- En $[0, 2]$, el número total de subintervalos es $2n = 4$: $y_j = j/2$ para $j=0,1,2,3,4 \implies \Delta y = 1/2$.
El área de cada subrectángulo es:
\begin{equation*}
\text{Área}(R_{ij}) = \Delta x \cdot \Delta y = \frac{1}{2} \cdot \frac{1}{2} = \frac{1}{4}.
\end{equation*}
Hay en total $2 \times 4 = 8$ subrectángulos.
\end{enumerate}

\begin{enumerate}[label=\textbf{\arabic*.}]
\item \textbf{Puntos muestrales de la esquina superior derecha:}
\begin{equation*}
(x_i^*, y_j^*) = (x_i, y_j) = \left(\frac{i}{2}, \frac{j}{2}\right), \quad i \in \{1, 2\}, \ j \in \{1, 2, 3, 4\}.
\end{equation*}
\end{enumerate}

\begin{enumerate}[label=\textbf{\arabic*.}]
\item \textbf{Cálculo de la doble sumatoria:}
\begin{equation*}
S_{2,2} = \sum_{i=1}^2 \sum_{j=1}^4 \left(\frac{i}{2} + \frac{j}{2}\right) \frac{1}{4} = \frac{1}{8} \sum_{i=1}^2 \sum_{j=1}^4 (i + j).
\end{equation*}
Desarrollamos la suma interna para cada $i$:
\begin{equation*}
\sum_{j=1}^4 (i + j) = 4i + (1 + 2 + 3 + 4) = 4i + 10.
\end{equation*}
Sumando sobre $i = 1, 2$:
\begin{equation*}
\sum_{i=1}^2 (4i + 10) = (4(1) + 10) + (4(2) + 10) = 14 + 18 = 32.
\end{equation*}
\end{enumerate}

\begin{enumerate}[label=\textbf{\arabic*.}]
\item \textbf{Multiplicación por el factor de área:}
\begin{equation*}
S_{2,2} = \frac{1}{8} \times 32 = 4.
\end{equation*}
\textit{Nota:} El límite analítico de la integral es $3$. La suma superior $S_{2,2} = 4$ la sobrestima debido a que las esquinas superiores derechas toman el valor máximo en cada celda.
\end{enumerate}

\begin{trapbox}
Tomar solo $n=2$ subintervalos en el eje $y$ en vez de $2n=4$ como establece explícitamente el enunciado del examen de la UNI.
\end{trapbox}

\vspace{0.5cm}
\hrule
\vspace{0.5cm}

\subsection*{Problema 3.5: Lema fundamental del refinamiento de particiones \hfill \color{primaryblue}\small[Nivel 2]}
\addcontentsline{toc}{subsection}{Problema 3.5: \texorpdfstring{Lema fundamental del refinamiento de particiones}{Lema fundamental del refinamiento de particiones} (Nivel 2)}

\begin{tcolorbox}[colback=white,colframe=gray!40,arc=1mm,boxrule=0.5pt,left=3mm,right=3mm,top=2mm,bottom=2mm]
\textbf{Enunciado:} Sea $f: R \to \mathbb{R}$ una función acotada sobre un rectángulo cerrado $R \subset \mathbb{R}^2$. Demuestre que si $P'$ es un refinamiento de la partición $P$ (es decir, $P \subset P'$), entonces:
\begin{equation*}
L(f, P) \le L(f, P') \quad \text{y} \quad U(f, P') \le U(f, P).
\end{equation*}
\end{tcolorbox}

\begin{tcolorbox}[colback=yellow!4,colframe=accentamber!70,arc=1mm,boxrule=0.5pt,left=3mm,right=3mm,top=1.5mm,bottom=1.5mm]
\textbf{Pista Conceptual:} Demuéstrelo primero añadiendo un solo punto a la partición en una de las coordenadas y use la propiedad del ínfimo sobre subconjuntos: si $A \subset B$, entonces $\inf B \le \inf A$.
\end{tcolorbox}

\paragraph{Solución Paso a Paso:}\par
\begin{enumerate}[label=\textbf{\arabic*.}]
\item \textbf{Reducción a la adición de un solo punto:}
Toda partición refinada $P'$ se obtiene de $P$ agregando un número finito de puntos de corte adicionales. Por inducción finita, basta demostrar la desigualdad cuando agregamos un único punto adicional $c$ a la partición $P_1$ del intervalo $[a,b]$.
\end{enumerate}

\begin{enumerate}[label=\textbf{\arabic*.}]
\item \textbf{Demostración para la suma inferior:}
Supongamos que el subintervalo $[x_{k-1}, x_k]$ se divide en dos: $[x_{k-1}, c]$ y $[c, x_k]$.
Esto divide cada subrectángulo $R_{kj} = [x_{k-1}, x_k] \times [y_{j-1}, y_j]$ en dos subrectángulos disjuntos $R'_{kj} = [x_{k-1}, c] \times [y_{j-1}, y_j]$ y $R''_{kj} = [c, x_k] \times [y_{j-1}, y_j]$.
Nótese que:
\begin{equation*}
\text{Área}(R_{kj}) = \text{Área}(R'_{kj}) + \text{Área}(R''_{kj}).
\end{equation*}
Definamos los ínfimos:
\begin{equation*}
m_{kj} = \inf_{R_{kj}} f, \qquad m'_{kj} = \inf_{R'_{kj}} f, \qquad m''_{kj} = \inf_{R''_{kj}} f.
\end{equation*}
Como $R'_{kj} \subset R_{kj}$ y $R''_{kj} \subset R_{kj}$, el ínfimo sobre un subconjunto es mayor o igual que el ínfimo sobre el conjunto completo:
\begin{equation*}
m'_{kj} \ge m_{kj} \quad \text{y} \quad m''_{kj} \ge m_{kj}.
\end{equation*}
Por consiguiente:
\begin{equation*}
m'_{kj} \text{Área}(R'_{kj}) + m''_{kj} \text{Área}(R''_{kj}) \ge m_{kj} \text{Área}(R'_{kj}) + m_{kj} \text{Área}(R''_{kj}) = m_{kj} \text{Área}(R_{kj}).
\end{equation*}
Sumando sobre todos los índices $j$, todos los demás rectángulos no afectados permanecen idénticos, con lo cual:
\begin{equation*}
L(f, P') \ge L(f, P).
\end{equation*}
\end{enumerate}

\begin{enumerate}[label=\textbf{\arabic*.}]
\item \textbf{Demostración para la suma superior:}
Análogamente, para los supremos se tiene que si $A \subset B$, entonces $\sup A \le \sup B$.
Por tanto $M'_{kj} \le M_{kj}$ y $M''_{kj} \le M_{kj}$, lo que produce:
\begin{equation*}
M'_{kj} \text{Área}(R'_{kj}) + M''_{kj} \text{Área}(R''_{kj}) \le M_{kj} \text{Área}(R_{kj}) \implies U(f, P') \le U(f, P).
\end{equation*}
Esto completa la demostración del lema.
\end{enumerate}

\begin{trapbox}
Invertir el sentido de la desigualdad del ínfimo sobre subconjuntos: recordar que achicar el conjunto agranda (o mantiene) el ínfimo.
\end{trapbox}

\vspace{0.5cm}
\hrule
\vspace{0.5cm}

\subsection*{Problema 3.6: Comparación universal de sumas de Darboux para particiones distintas \hfill \color{primaryblue}\small[Nivel 2]}
\addcontentsline{toc}{subsection}{Problema 3.6: \texorpdfstring{Comparación universal de sumas de Darboux para particiones distintas}{Comparación universal de sumas de Darboux para particiones distintas} (Nivel 2)}

\begin{tcolorbox}[colback=white,colframe=gray!40,arc=1mm,boxrule=0.5pt,left=3mm,right=3mm,top=2mm,bottom=2mm]
\textbf{Enunciado:} Demuestre que si $f: R \to \mathbb{R}$ es una función acotada y $P_1, P_2$ son dos particiones arbitrarias cualesquiera del rectángulo $R$, entonces:
\begin{equation*}
L(f, P_1) \le U(f, P_2).
\end{equation*}
Deduzca de ello que $\sup_P L(f, P) \le \inf_P U(f, P)$.
\end{tcolorbox}

\begin{tcolorbox}[colback=yellow!4,colframe=accentamber!70,arc=1mm,boxrule=0.5pt,left=3mm,right=3mm,top=1.5mm,bottom=1.5mm]
\textbf{Pista Conceptual:} Considere la partición común refinada $P^* = P_1 \cup P_2$ y aplique el lema del refinamiento.
\end{tcolorbox}

\paragraph{Solución Paso a Paso:}\par
\begin{enumerate}[label=\textbf{\arabic*.}]
\item \textbf{Construcción de la partición unión común:}
Dadas dos particiones $P_1 = P_{1,x} \times P_{1,y}$ y $P_2 = P_{2,x} \times P_{2,y}$, definimos la partición común (refinamiento común):
\begin{equation*}
P^* = (P_{1,x} \cup P_{2,x}) \times (P_{1,y} \cup P_{2,y}).
\end{equation*}
Por construcción, $P^*$ refina tanto a $P_1$ como a $P_2$, es decir, $P_1 \subset P^*$ y $P_2 \subset P^*$.
\end{enumerate}

\begin{enumerate}[label=\textbf{\arabic*.}]
\item \textbf{Cadena de desigualdades:}
Aplicando el lema del refinamiento (ejercicio 3.5) a $P_1$ y $P^*$:
\begin{equation*}
L(f, P_1) \le L(f, P^*).
\end{equation*}
Para cualquier partición fija (en particular para $P^*$), es evidente por definición que toda suma inferior es menor o igual que la suma superior correspondiente:
\begin{equation*}
L(f, P^*) \le U(f, P^*).
\end{equation*}
Aplicando nuevamente el lema del refinamiento a $P_2$ y $P^*$:
\begin{equation*}
U(f, P^*) \le U(f, P_2).
\end{equation*}
Encadenando los tres resultados:
\begin{equation*}
L(f, P_1) \le L(f, P^*) \le U(f, P^*) \le U(f, P_2) \implies L(f, P_1) \le U(f, P_2).
\end{equation*}
\end{enumerate}

\begin{enumerate}[label=\textbf{\arabic*.}]
\item \textbf{Paso al supremo e ínfimo:}
Fijemos $P_2$. La desigualdad $L(f, P) \le U(f, P_2)$ se cumple para toda partición $P$. Por consiguiente, $U(f, P_2)$ es una cota superior del conjunto $\{L(f, P) : P \text{ partición}\}$.
Por definición de supremo (la menor cota superior):
\begin{equation*}
\sup_P L(f, P) \le U(f, P_2).
\end{equation*}
Dado que esta última desigualdad es válida para toda partición $P_2$, el número $\sup_P L(f, P)$ es una cota inferior del conjunto $\{U(f, P) : P \text{ partición}\}$.
Por definición de ínfimo (la mayor cota inferior):
\begin{equation*}
\sup_P L(f, P) \le \inf_P U(f, P).
\end{equation*}
Esta desigualdad universal es la base de toda la teoría de integración de Darboux.
\end{enumerate}

\begin{trapbox}
Intentar comparar $L(f,P_1)$ con $U(f,P_2)$ celda a celda sin recurrir a la partición común $P^*$.
\end{trapbox}

\vspace{0.5cm}
\hrule
\vspace{0.5cm}

\subsection*{Problema 3.7: Cálculo analítico del volumen por límite de Riemann (Examen UNI 2024-II) \hfill \color{primaryblue}\small[Nivel 3]}
\addcontentsline{toc}{subsection}{Problema 3.7: \texorpdfstring{Cálculo analítico del volumen por límite de Riemann (Examen UNI 2024-II)}{Cálculo analítico del volumen por límite de Riemann (Examen UNI 2024-II)} (Nivel 3)}

\begin{tcolorbox}[colback=white,colframe=gray!40,arc=1mm,boxrule=0.5pt,left=3mm,right=3mm,top=2mm,bottom=2mm]
\textbf{Enunciado:} Sea el rectángulo $R = [0, 1] \times [0, 2]$. Usamos las particiones uniformes $P_m = \{x_i = \frac{i}{m} : i=0,\dots,m\}$ y $P_{2n} = \{y_j = \frac{j}{n} : j=0,\dots,2n\}$ para particionar $R$. Evaluando la función $f(x,y) = x + y$ en los puntos de la esquina superior derecha $(x_i^*, y_j^*) = (i/m, j/n)$ de cada subrectángulo:
(a) Deduzca la expresión algebraica cerrada de la suma de Riemann $S_{m,n}$ en términos de $m$ y $n$.
(b) Calcule el límite $\lim_{m,n\to\infty} S_{m,n}$ y demuestre que el volumen del sólido debajo de la gráfica es exactamente $3$ sin calcular ninguna integral por primitivas.
\end{tcolorbox}

\begin{tcolorbox}[colback=yellow!4,colframe=accentamber!70,arc=1mm,boxrule=0.5pt,left=3mm,right=3mm,top=1.5mm,bottom=1.5mm]
\textbf{Pista Conceptual:} Exprese el área de cada celda como $\frac{1}{mn}$. Separe las sumatorias dobles y aplique la fórmula de Gauss $\sum_{k=1}^N k = \frac{N(N+1)}{2}$.
\end{tcolorbox}

\paragraph{Solución Paso a Paso:}\par
\begin{enumerate}[label=\textbf{\arabic*.}]
\item \textbf{Geometría de la celda y área:}
- Subintervalos en $x$: $[x_{i-1}, x_i] = [\frac{i-1}{m}, \frac{i}{m}] \implies \Delta x = \frac{1}{m}$.
- Subintervalos en $y$: $[y_{j-1}, y_j] = [\frac{j-1}{n}, \frac{j}{n}]$ para $j = 1, \dots, 2n \implies \Delta y = \frac{j}{n} - \frac{j-1}{n} = \frac{1}{n}$.
El área de cada uno de los $2mn$ subrectángulos es:
\begin{equation*}
\text{Área}(R_{ij}) = \Delta x \cdot \Delta y = \frac{1}{m} \cdot \frac{1}{n} = \frac{1}{mn}.
\end{equation*}
\end{enumerate}

\begin{enumerate}[label=\textbf{\arabic*.}]
\item \textbf{Deducción algebraica de la doble sumatoria:}
Evaluamos en la esquina superior derecha $(x_i^*, y_j^*) = \left(\frac{i}{m}, \frac{j}{n}\right)$:
\begin{equation*}
S_{m,n} = \sum_{i=1}^m \sum_{j=1}^{2n} \left(\frac{i}{m} + \frac{j}{n}\right) \frac{1}{mn}
\end{equation*}
\begin{equation*}
= \frac{1}{m^2 n} \sum_{i=1}^m \sum_{j=1}^{2n} i + \frac{1}{m n^2} \sum_{i=1}^m \sum_{j=1}^{2n} j.
\end{equation*}
\end{enumerate}

\begin{enumerate}[label=\textbf{\arabic*.}]
\item \textbf{Cálculo de cada sumatoria usando la fórmula de Gauss:}
- Para el primer término:
\begin{equation*}
\sum_{i=1}^m \sum_{j=1}^{2n} i = \left(\sum_{j=1}^{2n} 1\right) \left(\sum_{i=1}^m i\right) = (2n) \cdot \frac{m(m+1)}{2} = n m (m + 1).
\end{equation*}
Multiplicando por el factor $\frac{1}{m^2 n}$:
\begin{equation*}
\frac{1}{m^2 n} \cdot n m (m + 1) = \frac{m + 1}{m} = 1 + \frac{1}{m}.
\end{equation*}
- Para el segundo término:
\begin{equation*}
\sum_{i=1}^m \sum_{j=1}^{2n} j = \left(\sum_{i=1}^m 1\right) \left(\sum_{j=1}^{2n} j\right) = m \cdot \frac{(2n)(2n + 1)}{2} = m n (2n + 1).
\end{equation*}
Multiplicando por el factor $\frac{1}{m n^2}$:
\begin{equation*}
\frac{1}{m n^2} \cdot m n (2n + 1) = \frac{2n + 1}{n} = 2 + \frac{1}{n}.
\end{equation*}
\end{enumerate}

\begin{enumerate}[label=\textbf{\arabic*.}]
\item \textbf{Expresión cerrada de $S_{m,n}$:}
\begin{equation*}
S_{m,n} = \left(1 + \frac{1}{m}\right) + \left(2 + \frac{1}{n}\right) = 3 + \frac{1}{m} + \frac{1}{n}.
\end{equation*}
\end{enumerate}

\begin{enumerate}[label=\textbf{\arabic*.}]
\item \textbf{Paso al límite y volumen del sólido:}
Tomando el límite cuando las normas de las particiones tienden a cero ($m \to \infty, n \to \infty$):
\begin{equation*}
\lim_{m,n \to \infty} S_{m,n} = \lim_{m \to \infty} \left(1 + \frac{1}{m}\right) + \lim_{n \to \infty} \left(2 + \frac{1}{n}\right) = 1 + 2 = 3.
\end{equation*}
Por la definición de integral de Riemann en bloques, el volumen debajo de la gráfica de $f$ sobre el rectángulo $R$ es rigurosamente:
\begin{equation*}
\iint_R (x + y) \, dA = \lim_{m,n \to \infty} S_{m,n} = 3 \, u^3.
\end{equation*}
\end{enumerate}

\begin{trapbox}
Olvidar que el límite superior del índice $j$ es $2n$ y escribir $n(n+1)/2$ en vez de $(2n)(2n+1)/2$.
\end{trapbox}

\vspace{0.5cm}
\hrule
\vspace{0.5cm}

\subsection*{Problema 3.8: Sumas de Riemann para función cuadrática bidimensional \hfill \color{primaryblue}\small[Nivel 3]}
\addcontentsline{toc}{subsection}{Problema 3.8: \texorpdfstring{Sumas de Riemann para función cuadrática bidimensional}{Sumas de Riemann para función cuadrática bidimensional} (Nivel 3)}

\begin{tcolorbox}[colback=white,colframe=gray!40,arc=1mm,boxrule=0.5pt,left=3mm,right=3mm,top=2mm,bottom=2mm]
\textbf{Enunciado:} Sea $f(x,y) = x^2 + y$ en el cuadrado $R = [0, 1] \times [0, 1]$. Dividiendo el intervalo $[0,1]$ en $n$ partes iguales en ambos ejes y eligiendo las esquinas superiores derechas $(i/n, j/n)$: (a) Halle la suma de Riemann $S_n$. (b) Calcule $\lim_{n\to\infty} S_n$ y verifique el resultado con la integral doble.
\end{tcolorbox}

\begin{tcolorbox}[colback=yellow!4,colframe=accentamber!70,arc=1mm,boxrule=0.5pt,left=3mm,right=3mm,top=1.5mm,bottom=1.5mm]
\textbf{Pista Conceptual:} Utilice la suma de cuadrados $\sum_{i=1}^n i^2 = \frac{n(n+1)(2n+1)}{6}$ y la suma lineal $\sum_{j=1}^n j = \frac{n(n+1)}{2}$.
\end{tcolorbox}

\paragraph{Solución Paso a Paso:}\par
\begin{enumerate}[label=\textbf{\arabic*.}]
\item \textbf{Geometría:}
$\Delta x = 1/n$, $\Delta y = 1/n$, $\Delta A = 1/n^2$.
Puntos muestrales: $(x_i^*, y_j^*) = (i/n, j/n)$.
\end{enumerate}

\begin{enumerate}[label=\textbf{\arabic*.}]
\item \textbf{Doble sumatoria:}
\begin{equation*}
S_n = \sum_{i=1}^n \sum_{j=1}^n \left(\frac{i^2}{n^2} + \frac{j}{n}\right) \frac{1}{n^2} = \frac{1}{n^4} \sum_{i=1}^n \sum_{j=1}^n i^2 + \frac{1}{n^3} \sum_{i=1}^n \sum_{j=1}^n j
\end{equation*}
\begin{equation*}
= \frac{1}{n^4} \cdot n \cdot \left(\sum_{i=1}^n i^2\right) + \frac{1}{n^3} \cdot n \cdot \left(\sum_{j=1}^n j\right) = \frac{1}{n^3} \cdot \frac{n(n+1)(2n+1)}{6} + \frac{1}{n^2} \cdot \frac{n(n+1)}{2}.
\end{equation*}
\end{enumerate}

\begin{enumerate}[label=\textbf{\arabic*.}]
\item \textbf{Simplificación:}
\begin{equation*}
S_n = \frac{(n+1)(2n+1)}{6n^2} + \frac{n+1}{2n} = \frac{2n^2 + 3n + 1}{6n^2} + \frac{n+1}{2n} = \left(\frac{1}{3} + \frac{1}{2n} + \frac{1}{6n^2}\right) + \left(\frac{1}{2} + \frac{1}{2n}\right).
\end{equation*}
\end{enumerate}

\begin{enumerate}[label=\textbf{\arabic*.}]
\item \textbf{Límite:}
\begin{equation*}
\lim_{n\to\infty} S_n = \frac{1}{3} + \frac{1}{2} = \frac{5}{6}.
\end{equation*}
Comprobación por cálculo integral:
\begin{equation*}
\iint_{[0,1]^2} (x^2 + y) dA = \int_0^1 x^2 dx \int_0^1 1 dy + \int_0^1 1 dx \int_0^1 y dy = \frac{1}{3}(1) + 1\left(\frac{1}{2}\right) = \frac{5}{6}.
\end{equation*}
\end{enumerate}

\begin{trapbox}
Omitir el factor $n$ que proviene de sumar una constante respecto al otro índice mudo de la sumatoria.
\end{trapbox}

\vspace{0.5cm}
\hrule
\vspace{0.5cm}

\subsection*{Problema 3.9: Sumas de Darboux sobre un paralelepípedo en R$^3$ \hfill \color{primaryblue}\small[Nivel 3]}
\addcontentsline{toc}{subsection}{Problema 3.9: \texorpdfstring{Sumas de Darboux sobre un paralelepípedo en R$^3$}{Sumas de Darboux sobre un paralelepípedo en R^3} (Nivel 3)}

\begin{tcolorbox}[colback=white,colframe=gray!40,arc=1mm,boxrule=0.5pt,left=3mm,right=3mm,top=2mm,bottom=2mm]
\textbf{Enunciado:} Sea $B = [0, 1] \times [0, 1] \times [0, 1]$ el cubo unitario y $f(x,y,z) = x + y + z$. Para la partición uniforme con $m = n = p = 2$ en cada eje, calcule la suma inferior $L(f, P)$ y la suma superior $U(f, P)$ y compruebe que $L(f,P) \le \iiint_B f dV \le U(f,P)$.
\end{tcolorbox}

\begin{tcolorbox}[colback=yellow!4,colframe=accentamber!70,arc=1mm,boxrule=0.5pt,left=3mm,right=3mm,top=1.5mm,bottom=1.5mm]
\textbf{Pista Conceptual:} El cubo se divide en $2^3 = 8$ subcubos iguales de volumen $\Delta V = 1/8$. Evalúe el mínimo y máximo de la función creciente en los vértices correspondientes.
\end{tcolorbox}

\paragraph{Solución Paso a Paso:}\par
\begin{enumerate}[label=\textbf{\arabic*.}]
\item \textbf{Geometría de la partición tridimensional:}
El cubo unitario se particiona mediante $P_1 = P_2 = P_3 = \{0, 1/2, 1\}$.
Se forman $2 \times 2 \times 2 = 8$ subcubos:
\begin{equation*}
B_{ijk} = \left[\frac{i-1}{2}, \frac{i}{2}\right] \times \left[\frac{j-1}{2}, \frac{j}{2}\right] \times \left[\frac{k-1}{2}, \frac{k}{2}\right], \quad i,j,k \in \{1, 2\}.
\end{equation*}
El volumen de cada subcubo es $\Delta V = (1/2)^3 = 1/8$.
\end{enumerate}

\begin{enumerate}[label=\textbf{\arabic*.}]
\item \textbf{Suma inferior $L(f,P)$:}
Al ser $f(x,y,z) = x+y+z$ creciente en todas las variables:
\begin{equation*}
m_{ijk} = \inf_{B_{ijk}} f = f\left(\frac{i-1}{2}, \frac{j-1}{2}, \frac{k-1}{2}\right) = \frac{(i-1) + (j-1) + (k-1)}{2}.
\end{equation*}
Sumando sobre los 8 cubos:
\begin{equation*}
L(f,P) = \sum_{i,j,k} m_{ijk} \Delta V = \frac{1}{8} \sum_{i=1}^2 \sum_{j=1}^2 \sum_{k=1}^2 \frac{(i-1) + (j-1) + (k-1)}{2}.
\end{equation*}
Cada variable $(i-1)$ toma el valor 0 (4 veces) y el valor 1 (4 veces).
Por simetría:
\begin{equation*}
\sum (i-1) = 4, \quad \sum (j-1) = 4, \quad \sum (k-1) = 4 \implies \text{Suma total} = \frac{12}{2} = 6.
\end{equation*}
\begin{equation*}
L(f,P) = \frac{1}{8} \times 6 = \frac{3}{4} = 0.75.
\end{equation*}
\end{enumerate}

\begin{enumerate}[label=\textbf{\arabic*.}]
\item \textbf{Suma superior $U(f,P)$:}
\begin{equation*}
M_{ijk} = \sup_{B_{ijk}} f = f\left(\frac{i}{2}, \frac{j}{2}, \frac{k}{2}\right) = \frac{i + j + k}{2}.
\end{equation*}
Cada variable $i$ toma el valor 1 (4 veces) y el valor 2 (4 veces), sumando $4(1) + 4(2) = 12$.
La suma total es $\frac{12 + 12 + 12}{2} = \frac{36}{2} = 18$.
\begin{equation*}
U(f,P) = \frac{1}{8} \times 18 = \frac{9}{4} = 2.25.
\end{equation*}
\end{enumerate}

\begin{enumerate}[label=\textbf{\arabic*.}]
\item \textbf{Comparación con la integral exacta:}
\begin{equation*}
\iiint_B (x+y+z) dV = \int_0^1 x dx + \int_0^1 y dy + \int_0^1 z dz = \frac{1}{2} + \frac{1}{2} + \frac{1}{2} = \frac{3}{2} = 1.5.
\end{equation*}
Efectivamente se cumple: $0.75 \le 1.5 \le 2.25$.
\end{enumerate}

\begin{trapbox}
Contar erróneamente las permutaciones de los índices al calcular la suma tridimensional.
\end{trapbox}

\vspace{0.5cm}
\hrule
\vspace{0.5cm}

\subsection*{Problema 3.10: Función de Dirichlet bidimensional y no integrabilidad \hfill \color{primaryblue}\small[Nivel 4]}
\addcontentsline{toc}{subsection}{Problema 3.10: \texorpdfstring{Función de Dirichlet bidimensional y no integrabilidad}{Función de Dirichlet bidimensional y no integrabilidad} (Nivel 4)}

\begin{tcolorbox}[colback=white,colframe=gray!40,arc=1mm,boxrule=0.5pt,left=3mm,right=3mm,top=2mm,bottom=2mm]
\textbf{Enunciado:} Sea $R = [0, 1] \times [0, 1]$ y defina la función de Dirichlet en 2D por:
\begin{equation*}
f(x,y) = \begin{cases} 1, & \text{si } x \in \mathbb{Q} \text{ e } y \in \mathbb{Q} \\ 0, & \text{en otro caso} \end{cases}
\end{equation*}
(a) Demuestre que para cualquier partición $P$ de $R$, se tiene $L(f, P) = 0$ y $U(f, P) = 1$.
(b) Concluya formalmente a partir de la definición de Darboux que $f$ NO es integrable sobre $R$.
\end{tcolorbox}

\begin{tcolorbox}[colback=yellow!4,colframe=accentamber!70,arc=1mm,boxrule=0.5pt,left=3mm,right=3mm,top=1.5mm,bottom=1.5mm]
\textbf{Pista Conceptual:} Utilice la densidad de los números racionales $\mathbb{Q}$ y de los números irracionales $\mathbb{R} \setminus \mathbb{Q}$ en cualquier intervalo abierto no vacío.
\end{tcolorbox}

\paragraph{Solución Paso a Paso:}\par
\begin{enumerate}[label=\textbf{\arabic*.}]
\item \textbf{Densidad en subrectángulos:}
Sea $P = P_1 \times P_2$ una partición arbitraria de $R$. Cada subrectángulo es de la forma:
\begin{equation*}
R_{ij} = [x_{i-1}, x_i] \times [y_{j-1}, y_j], \quad \text{con } x_{i-1} < x_i \text{ y } y_{j-1} < y_j.
\end{equation*}
- Por la densidad de $\mathbb{Q}$ en $\mathbb{R}$, existen $q_x \in (x_{i-1}, x_i) \cap \mathbb{Q}$ y $q_y \in (y_{j-1}, y_j) \cap \mathbb{Q}$.
Por ende, $(q_x, q_y) \in R_{ij}$ y $f(q_x, q_y) = 1$.
Esto implica que:
\begin{equation*}
M_{ij} = \sup_{R_{ij}} f = 1, \quad \forall i, j.
\end{equation*}
- Por la densidad de los irracionales $\mathbb{R} \setminus \mathbb{Q}$ en $\mathbb{R}$, existe $r_x \in (x_{i-1}, x_i) \setminus \mathbb{Q}$ y $q_y \in (y_{j-1}, y_j)$.
Entonces $(r_x, q_y) \in R_{ij}$ no tiene ambas coordenadas racionales, luego $f(r_x, q_y) = 0$.
Como $f \ge 0$, esto implica que:
\begin{equation*}
m_{ij} = \inf_{R_{ij}} f = 0, \quad \forall i, j.
\end{equation*}
\end{enumerate}

\begin{enumerate}[label=\textbf{\arabic*.}]
\item \textbf{Cálculo de las sumas de Darboux:}
Para cualquier partición $P$:
\begin{equation*}
L(f, P) = \sum_{i,j} m_{ij} \text{Área}(R_{ij}) = \sum_{i,j} 0 \cdot \text{Área}(R_{ij}) = 0,
\end{equation*}
\begin{equation*}
U(f, P) = \sum_{i,j} M_{ij} \text{Área}(R_{ij}) = \sum_{i,j} 1 \cdot \text{Área}(R_{ij}) = \text{Área}(R) = 1 \cdot 1 = 1.
\end{equation*}
\end{enumerate}

\begin{enumerate}[label=\textbf{\arabic*.}]
\item \textbf{Cálculo de integrales superior e inferior de Darboux:}
\begin{equation*}
\underline{\iint_R} f \, dA = \sup_P L(f, P) = \sup \{0\} = 0,
\end{equation*}
\begin{equation*}
\overline{\iint_R} f \, dA = \inf_P U(f, P) = \inf \{1\} = 1.
\end{equation*}
\end{enumerate}

\begin{enumerate}[label=\textbf{\arabic*.}]
\item \textbf{Conclusión:}
Una función es integrable en el sentido de Riemann/Darboux si y sólo si $\underline{\iint_R} f = \overline{\iint_R} f$.
Como $0 \neq 1$, las integrales de Darboux no coinciden, demostrando que la función de Dirichlet no es integrable Riemann.
\end{enumerate}

\begin{trapbox}
Confundir integrabilidad de Riemann con integrabilidad de Lebesgue (en Lebesgue la integral sí existe y vale 0).
\end{trapbox}

\vspace{0.5cm}
\hrule
\vspace{0.5cm}

\subsection*{Problema 3.11: Función modificada con discontinuidades de contenido cero \hfill \color{primaryblue}\small[Nivel 4]}
\addcontentsline{toc}{subsection}{Problema 3.11: \texorpdfstring{Función modificada con discontinuidades de contenido cero}{Función modificada con discontinuidades de contenido cero} (Nivel 4)}

\begin{tcolorbox}[colback=white,colframe=gray!40,arc=1mm,boxrule=0.5pt,left=3mm,right=3mm,top=2mm,bottom=2mm]
\textbf{Enunciado:} Sea $R = [0, 1] \times [0, 1]$. Defina $f(x,y) = 1$ si $y = x$ y $f(x,y) = 0$ si $y \neq x$. Demuestre rigurosamente que $f$ es integrable en $R$ calculando que para todo $\varepsilon > 0$ existe una partición $P$ tal que $U(f, P) - L(f, P) < \varepsilon$, y halle el valor de $\iint_R f dA$.
\end{tcolorbox}

\begin{tcolorbox}[colback=yellow!4,colframe=accentamber!70,arc=1mm,boxrule=0.5pt,left=3mm,right=3mm,top=1.5mm,bottom=1.5mm]
\textbf{Pista Conceptual:} Observe que $L(f,P) = 0$ siempre. Para $U(f,P)$, cubra la diagonal $y=x$ con subrectángulos de una partición uniforme con $n$ subdivisiones.
\end{tcolorbox}

\paragraph{Solución Paso a Paso:}\par
\begin{enumerate}[label=\textbf{\arabic*.}]
\item \textbf{Suma inferior nula:}
En cualquier subrectángulo $R_{ij}$, existen infinitos puntos fuera de la recta diagonal $y = x$. En esos puntos $f(x,y) = 0$.
Por ende, $m_{ij} = \inf_{R_{ij}} f = 0$ para todo $i,j$.
Luego $L(f, P) = 0$ para toda partición $P$.
\end{enumerate}

\begin{enumerate}[label=\textbf{\arabic*.}]
\item \textbf{Control de la suma superior con partición regular de tamaño $n$:}
Tomemos una partición uniforme $P_n$ dividiendo ambos ejes en $n$ subintervalos de longitud $1/n$.
Hay $n^2$ subrectángulos de área $1/n^2$.
La función $f$ vale 1 únicamente en la diagonal $y = x$.
¿Cuántos subrectángulos intersecan a la diagonal $y = x$?
Un subrectángulo $R_{ij} = [\frac{i-1}{n}, \frac{i}{n}] \times [\frac{j-1}{n}, \frac{j}{n}]$ interseca a la diagonal $y = x$ si y sólo si los intervalos $[\frac{i-1}{n}, \frac{i}{n}]$ y $[\frac{j-1}{n}, \frac{j}{n}]$ se solapan, lo cual ocurre exactamente cuando $i = j$.
Por tanto, exactamente $n$ subrectángulos tocan la diagonal (los que están sobre la diagonal principal).
En estos $n$ subrectángulos, $M_{ii} = 1$.
En los restantes $n^2 - n$ subrectángulos, $f$ es idénticamente 0, luego $M_{ij} = 0$.
\end{enumerate}

\begin{enumerate}[label=\textbf{\arabic*.}]
\item \textbf{Cálculo de la suma superior:}
\begin{equation*}
U(f, P_n) = \sum_{i=1}^n M_{ii} \text{Área}(R_{ii}) = n \cdot \left(1 \cdot \frac{1}{n^2}\right) = \frac{n}{n^2} = \frac{1}{n}.
\end{equation*}
\end{enumerate}

\begin{enumerate}[label=\textbf{\arabic*.}]
\item \textbf{Verificación del criterio de integrabilidad:}
Dado cualquier $\varepsilon > 0$, elijamos $n \in \mathbb{N}$ tal que $n > 1/\varepsilon \implies \frac{1}{n} < \varepsilon$.
Entonces para la partición $P_n$:
\begin{equation*}
U(f, P_n) - L(f, P_n) = \frac{1}{n} - 0 = \frac{1}{n} < \varepsilon.
\end{equation*}
Por el criterio de integrabilidad de Riemann, $f$ es integrable sobre $R$.
\end{enumerate}

\begin{enumerate}[label=\textbf{\arabic*.}]
\item \textbf{Valor de la integral:}
Como $0 = L(f, P_n) \le \iint_R f dA \le U(f, P_n) = 1/n$ para todo $n$:
\begin{equation*}
\iint_R f \, dA = 0.
\end{equation*}
\end{enumerate}

\begin{trapbox}
Creer que si una función tiene infinitos puntos de discontinuidad no puede ser integrable en el sentido de Riemann.
\end{trapbox}

\vspace{0.5cm}
\hrule
\vspace{0.5cm}

\subsection*{Problema 3.12: Caracterización de Darboux por la oscilación total \hfill \color{primaryblue}\small[Nivel 4]}
\addcontentsline{toc}{subsection}{Problema 3.12: \texorpdfstring{Caracterización de Darboux por la oscilación total}{Caracterización de Darboux por la oscilación total} (Nivel 4)}

\begin{tcolorbox}[colback=white,colframe=gray!40,arc=1mm,boxrule=0.5pt,left=3mm,right=3mm,top=2mm,bottom=2mm]
\textbf{Enunciado:} Sea $f: R \to \mathbb{R}$ acotada en el rectángulo $R$. Se define la oscilación de $f$ en un conjunto $E \subset R$ por $\omega(f, E) = \sup_{x,y \in E} |f(x) - f(y)| = \sup_E f - \inf_E f$. Demuestre que para cualquier partición $P = \{R_1, \dots, R_k\}$ de $R$:
\begin{equation*}
U(f, P) - L(f, P) = \sum_{j=1}^k \omega(f, R_j) \text{Área}(R_j).
\end{equation*}
\end{tcolorbox}

\begin{tcolorbox}[colback=yellow!4,colframe=accentamber!70,arc=1mm,boxrule=0.5pt,left=3mm,right=3mm,top=1.5mm,bottom=1.5mm]
\textbf{Pista Conceptual:} Substituya la definición de $U(f,P) = \sum M_j \text{Área}(R_j)$ y $L(f,P) = \sum m_j \text{Área}(R_j)$ y agrupe los sumandos.
\end{tcolorbox}

\paragraph{Solución Paso a Paso:}\par
\begin{enumerate}[label=\textbf{\arabic*.}]
\item \textbf{Definiciones básicas:}
Para cada subrectángulo $R_j$, definimos:
\begin{equation*}
M_j = \sup_{x \in R_j} f(x), \qquad m_j = \inf_{x \in R_j} f(x).
\end{equation*}
Por las propiedades del supremo e ínfimo, la oscilación satisface:
\begin{equation*}
\omega(f, R_j) = \sup_{u, v \in R_j} |f(u) - f(v)| = M_j - m_j.
\end{equation*}
En efecto, para cualesquiera $u, v \in R_j$: $m_j \le f(u) \le M_j$ y $m_j \le f(v) \le M_j$, lo que implica $f(u) - f(v) \le M_j - m_j$ y $f(v) - f(u) \le M_j - m_j$, luego $|f(u) - f(v)| \le M_j - m_j$.
Además, por definición de supremo e ínfimo, para todo $\eta > 0$ existen $u, v \in R_j$ con $f(u) > M_j - \eta/2$ y $f(v) < m_j + \eta/2$, de modo que $|f(u) - f(v)| > M_j - m_j - \eta$.
\end{enumerate}

\begin{enumerate}[label=\textbf{\arabic*.}]
\item \textbf{Diferencia de sumas:}
Por definición de las sumas de Darboux:
\begin{equation*}
U(f, P) = \sum_{j=1}^k M_j \text{Área}(R_j), \qquad L(f, P) = \sum_{j=1}^k m_j \text{Área}(R_j).
\end{equation*}
Restando ambas expresiones y factorizando $\text{Área}(R_j)$:
\begin{equation*}
U(f, P) - L(f, P) = \sum_{j=1}^k (M_j - m_j) \text{Área}(R_j) = \sum_{j=1}^k \omega(f, R_j) \text{Área}(R_j).
\end{equation*}
Esta formulación en términos de oscilación muestra intuitivamente que la integrabilidad depende de que la oscilación ponderada promedio de la función se pueda hacer arbitrariamente pequeña refinando la partición.
\end{enumerate}

\begin{trapbox}
No demostrar la igualdad $\sup |f(u)-f(v)| = M_j - m_j$ antes de usarla.
\end{trapbox}

\vspace{0.5cm}
\hrule
\vspace{0.5cm}

\subsection*{Problema 3.13: Equivalencia entre integrabilidad de Darboux y de Riemann \hfill \color{primaryblue}\small[Nivel 5]}
\addcontentsline{toc}{subsection}{Problema 3.13: \texorpdfstring{Equivalencia entre integrabilidad de Darboux y de Riemann}{Equivalencia entre integrabilidad de Darboux y de Riemann} (Nivel 5)}

\begin{tcolorbox}[colback=white,colframe=gray!40,arc=1mm,boxrule=0.5pt,left=3mm,right=3mm,top=2mm,bottom=2mm]
\textbf{Enunciado:} Sea $f: R \to \mathbb{R}$ acotada en el rectángulo $R$. Demuestre rigurosamente que $f$ es integrable en el sentido de Darboux (es decir, $\sup_P L(f,P) = \inf_P U(f,P) = I$) si y sólo si $f$ es integrable en el sentido de Riemann (existe $I \in \mathbb{R}$ tal que para todo $\varepsilon > 0$ existe $\delta > 0$ tal que para toda partición con norma $\|P\| < \delta$ y cualquier elección de puntos de muestra $\xi$, $|S(f, P, \xi) - I| < \varepsilon$).
\end{tcolorbox}

\begin{tcolorbox}[colback=yellow!4,colframe=accentamber!70,arc=1mm,boxrule=0.5pt,left=3mm,right=3mm,top=1.5mm,bottom=1.5mm]
\textbf{Pista Conceptual:} Observe que para cualquier elección de puntos de muestra $\xi$, $L(f,P) \le S(f,P,\xi) \le U(f,P)$. Utilice el Lema de Darboux sobre la aproximación uniforme de las sumas al tender la norma de la partición a cero.
\end{tcolorbox}

\paragraph{Solución Paso a Paso:}\par
\begin{enumerate}[label=\textbf{\arabic*.}]
\item \textbf{Acotación fundamental de las sumas de Riemann:}
Para cualquier partición $P = \{R_1, \dots, R_k\}$ y cualquier conjunto de puntos muestrales $\xi = \{\xi_1, \dots, \xi_k\}$ con $\xi_j \in R_j$:
Como $m_j \le f(\xi_j) \le M_j$, multiplicando por $\text{Área}(R_j) > 0$ y sumando:
\begin{equation*}
L(f, P) \le S(f, P, \xi) \le U(f, P).
\end{equation*}
\end{enumerate}

\begin{enumerate}[label=\textbf{\arabic*.}]
\item \textbf{Necesidad (Riemann $\implies$ Darboux):}
Supongamos que $f$ es integrable Riemann con valor $I$.
Dado $\varepsilon > 0$, existe $\delta > 0$ tal que si $\|P\| < \delta$, entonces $|S(f, P, \xi) - I| < \varepsilon/3$ para todo $\xi$.
Por propiedades del supremo e ínfimo, para dicha partición $P$ existen elecciones de puntos muestrales $\xi'$ y $\xi''$ tales que:
\begin{equation*}
U(f, P) < S(f, P, \xi') + \frac{\varepsilon}{3} < (I + \frac{\varepsilon}{3}) + \frac{\varepsilon}{3} = I + \frac{2\varepsilon}{3},
\end{equation*}
\begin{equation*}
L(f, P) > S(f, P, \xi'') - \frac{\varepsilon}{3} > (I - \frac{\varepsilon}{3}) - \frac{\varepsilon}{3} = I - \frac{2\varepsilon}{3}.
\end{equation*}
Restando ambas desigualdades:
\begin{equation*}
U(f, P) - L(f, P) < \frac{4\varepsilon}{3} < 2\varepsilon.
\end{equation*}
Como esto se cumple para particiones con norma pequeña, $\inf_P U(f,P) - \sup_P L(f,P) = 0$, lo que prueba la integrabilidad de Darboux.
\end{enumerate}

\begin{enumerate}[label=\textbf{\arabic*.}]
\item \textbf{Suficiencia (Darboux $\implies$ Riemann):}
Supongamos que $\sup L(f,P) = \inf U(f,P) = I$.
Por el Lema de Darboux: para todo $\varepsilon > 0$, existe $\delta > 0$ tal que para cualquier partición $P$ con malla $\|P\| < \delta$:
\begin{equation*}
U(f, P) < I + \varepsilon \quad \text{y} \quad L(f, P) > I - \varepsilon.
\end{equation*}
Dado que para cualquier selección $\xi$ se cumple $L(f, P) \le S(f, P, \xi) \le U(f, P)$, tenemos:
\begin{equation*}
I - \varepsilon < L(f, P) \le S(f, P, \xi) \le U(f, P) < I + \varepsilon \implies |S(f, P, \xi) - I| < \varepsilon.
\end{equation*}
Esto demuestra que el límite de Riemann existe y coincide con el valor común de Darboux $I$.
\end{enumerate}

\begin{trapbox}
Omitir el lema de Darboux que conecta el supremo/ínfimo abstracto con particiones cuya norma es geométricamente pequeña.
\end{trapbox}

\vspace{0.5cm}
\hrule
\vspace{0.5cm}

\subsection*{Problema 3.14: Aproximación por funciones escalonadas y densidad $L^1$ \hfill \color{primaryblue}\small[Nivel 5]}
\addcontentsline{toc}{subsection}{Problema 3.14: \texorpdfstring{Aproximación por funciones escalonadas y densidad $L^1$}{Aproximación por funciones escalonadas y densidad L^1} (Nivel 5)}

\begin{tcolorbox}[colback=white,colframe=gray!40,arc=1mm,boxrule=0.5pt,left=3mm,right=3mm,top=2mm,bottom=2mm]
\textbf{Enunciado:} Sea $f: R \to \mathbb{R}$ integrable Riemann en el rectángulo compacto $R$. Demuestre que para todo $\varepsilon > 0$, existen dos funciones escalonadas $\phi, \psi: R \to \mathbb{R}$ tales que:
\begin{equation*}
\phi(x,y) \le f(x,y) \le \psi(x,y), \quad \forall (x,y) \in R, \qquad \text{y} \qquad \iint_R (\psi - \phi) \, dA < \varepsilon.
\end{equation*}
\end{tcolorbox}

\begin{tcolorbox}[colback=yellow!4,colframe=accentamber!70,arc=1mm,boxrule=0.5pt,left=3mm,right=3mm,top=1.5mm,bottom=1.5mm]
\textbf{Pista Conceptual:} Construya $\phi$ y $\psi$ constantes en el interior de cada subrectángulo de una partición de Darboux que satisfaga $U(f,P) - L(f,P) < \varepsilon$.
\end{tcolorbox}

\paragraph{Solución Paso a Paso:}\par
\begin{enumerate}[label=\textbf{\arabic*.}]
\item \textbf{Definición de función escalonada:}
Una función $s: R \to \mathbb{R}$ es escalonada si existe una partición $P = \{R_1, \dots, R_k\}$ de $R$ tal que $s$ es constante en el interior de cada subrectángulo $R_j$.
La integral de una función escalonada es $\iint_R s \, dA = \sum_{j=1}^k c_j \text{Área}(R_j)$ donde $c_j = s|_{R_j^\circ}$.
\end{enumerate}

\begin{enumerate}[label=\textbf{\arabic*.}]
\item \textbf{Construcción a partir de una partición de Darboux:}
Por el criterio de integrabilidad de Riemann, como $f$ es integrable, dado $\varepsilon > 0$ existe una partición $P = \{R_1, \dots, R_k\}$ tal que:
\begin{equation*}
U(f, P) - L(f, P) < \varepsilon.
\end{equation*}
Para cada $j \in \{1, \dots, k\}$, definamos:
\begin{equation*}
m_j = \inf_{R_j} f, \qquad M_j = \sup_{R_j} f.
\end{equation*}
Definimos las funciones escalonadas $\phi, \psi: R \to \mathbb{R}$ por:
\begin{equation*}
\phi(x,y) = m_j \quad \text{si } (x,y) \in R_j^\circ, \qquad \psi(x,y) = M_j \quad \text{si } (x,y) \in R_j^\circ.
\end{equation*}
(En las fronteras de los subrectángulos, asignamos el valor mínimo global de $f$ para $\phi$ y el máximo global para $\psi$).
\end{enumerate}

\begin{enumerate}[label=\textbf{\arabic*.}]
\item \textbf{Verificación de las desigualdades:}
Por definición de ínfimo y supremo, para todo $(x,y) \in R$:
\begin{equation*}
\phi(x,y) \le f(x,y) \le \psi(x,y).
\end{equation*}
Además, integrando la diferencia:
\begin{equation*}
\iint_R (\psi - \phi) \, dA = \sum_{j=1}^k (M_j - m_j) \text{Área}(R_j) = U(f, P) - L(f, P) < \varepsilon.
\end{equation*}
Esto demuestra que el espacio de funciones escalonadas es denso respecto a la métrica $L^1$ en el conjunto de funciones integrables de Riemann.
\end{enumerate}

\begin{trapbox}
Olvidar tratar la definición en las fronteras de los subrectángulos (conjuntos de medida nula).
\end{trapbox}

\vspace{0.5cm}
\hrule
\vspace{0.5cm}

\subsection*{Problema 3.15: Volumen de tetraedro tridimensional por sumas de Darboux \hfill \color{primaryblue}\small[Nivel 5]}
\addcontentsline{toc}{subsection}{Problema 3.15: \texorpdfstring{Volumen de tetraedro tridimensional por sumas de Darboux}{Volumen de tetraedro tridimensional por sumas de Darboux} (Nivel 5)}

\begin{tcolorbox}[colback=white,colframe=gray!40,arc=1mm,boxrule=0.5pt,left=3mm,right=3mm,top=2mm,bottom=2mm]
\textbf{Enunciado:} Sea el tetraedro $T \subset \mathbb{R}^3$ acotado por los planos coordenados $x=0, y=0, z=0$ y el plano $x + y + z = 1$. Considere el cubo contenedor $B = [0,1]^3$ y la función característica $\chi_T$. Mediante particiones regulares de $B$ en $n^3$ cubos elementales, demuestre rigurosamente que el volumen de $T$ es $\lim_{n\to\infty} L(\chi_T, P_n) = 1/6$.
\end{tcolorbox}

\begin{tcolorbox}[colback=yellow!4,colframe=accentamber!70,arc=1mm,boxrule=0.5pt,left=3mm,right=3mm,top=1.5mm,bottom=1.5mm]
\textbf{Pista Conceptual:} Identifique qué cubos elementales $[\frac{i-1}{n}, \frac{i}{n}] \times [\frac{j-1}{n}, \frac{j}{n}] \times [\frac{k-1}{n}, \frac{k}{n}]$ están totalmente contenidos dentro de $x+y+z \le 1$. La condición es $i + j + k \le n$. Use la fórmula de la suma de combinatorios $\sum_{s=3}^n \binom{s-1}{2} = \binom{n}{3}$.
\end{tcolorbox}

\paragraph{Solución Paso a Paso:}\par
\begin{enumerate}[label=\textbf{\arabic*.}]
\item \textbf{Geometría y condición de contención estricta:}
Un subcubo $B_{ijk} = [\frac{i-1}{n}, \frac{i}{n}] \times [\frac{j-1}{n}, \frac{j}{n}] \times [\frac{k-1}{n}, \frac{k}{n}]$ está completamente contenido en $T$ si y sólo si su vértice más alejado del origen está en $T$:
\begin{equation*}
\frac{i}{n} + \frac{j}{n} + \frac{k}{n} \le 1 \iff i + j + k \le n.
\end{equation*}
Para estos cubos, $\inf_{B_{ijk}} \chi_T = 1$.
Para cualquier otro subcubo, contiene al menos un punto donde $x+y+z > 1$ (donde $\chi_T = 0$), luego $\inf_{B_{ijk}} \chi_T = 0$.
\end{enumerate}

\begin{enumerate}[label=\textbf{\arabic*.}]
\item \textbf{Cálculo combinatorio de cubos contenidos:}
El número de tripletas de enteros positivos $(i, j, k) \ge 1$ tales que $i + j + k \le n$ es un problema clásico de combinatoria (separadores y barras):
Sea $s = i + j + k$. Para un $s$ fijo ($3 \le s \le n$), el número de soluciones enteras positivas es $\binom{s-1}{2}$.
Sumando sobre $s$:
\begin{equation*}
N(n) = \sum_{s=3}^n \binom{s-1}{2} = \binom{2}{2} + \binom{3}{2} + \dots + \binom{n-1}{2} = \binom{n}{3} = \frac{n(n-1)(n-2)}{6}.
\end{equation*}
\end{enumerate}

\begin{enumerate}[label=\textbf{\arabic*.}]
\item \textbf{Cálculo de la suma inferior:}
Cada subcubo tiene volumen $\Delta V = \frac{1}{n^3}$. Por tanto:
\begin{equation*}
L(\chi_T, P_n) = N(n) \cdot \frac{1}{n^3} = \frac{n(n-1)(n-2)}{6 n^3} = \frac{1}{6} \left(1 - \frac{1}{n}\right)\left(1 - \frac{2}{n}\right).
\end{equation*}
\end{enumerate}

\begin{enumerate}[label=\textbf{\arabic*.}]
\item \textbf{Paso al límite:}
\begin{equation*}
\text{Vol}(T) = \lim_{n \to \infty} L(\chi_T, P_n) = \lim_{n \to \infty} \frac{1}{6} \left(1 - \frac{1}{n}\right)\left(1 - \frac{2}{n}\right) = \frac{1}{6} \cdot 1 \cdot 1 = \frac{1}{6}.
\end{equation*}
Este cálculo elemental verifica directamente la fórmula de la pirámide $\frac{1}{3} \times \text{área base} \times \text{altura} = \frac{1}{3} \times \frac{1}{2} \times 1 = \frac{1}{6}$ sin recurrir a integrales iteradas.
\end{enumerate}

\begin{trapbox}
Incluir en la suma inferior cubos que cortan la frontera oblicua del tetraedro.
\end{trapbox}

\vspace{0.5cm}
\hrule
\vspace{0.5cm}

\section{Criterios y Propiedades de Integrabilidad}

\subsection*{Problema 4.0A: Criterio de Riemann para una función constante \hfill \color{teal}\small[Nivel 0]}
\addcontentsline{toc}{subsection}{Problema 4.0A: Criterio de Riemann para una función constante (Nivel 0)}

\begin{tcolorbox}[colback=white,colframe=gray!40,arc=1mm,boxrule=0.5pt,left=3mm,right=3mm,top=2mm,bottom=2mm]
\textbf{Enunciado:} Sea $f(x,y) = c$ una función constante sobre un rectángulo $R = [a,b] \times [c,d]$, con $c \in \mathbb{R}$.
\begin{enumerate}[label=\textbf{\alph*)}]
\item Demuestre que para cualquier partición $P$ de $R$, la suma inferior $L(f,P)$ y la suma superior $U(f,P)$ son exactamente iguales a $c \cdot \text{Área}(R)$.
\item Concluya mediante el Criterio de Riemann que $f$ es integrable y calcule el valor de $\iint_R c \, dA$.
\end{enumerate}
\end{tcolorbox}

\begin{tcolorbox}[colback=yellow!4,colframe=accentamber!70,arc=1mm,boxrule=0.5pt,left=3mm,right=3mm,top=1.5mm,bottom=1.5mm]
\textbf{Pista Conceptual:} En cualquier subconjunto no vacío, el ínfimo y el supremo de una constante son la propia constante: $\inf c = \sup c = c$.
\end{tcolorbox}

\paragraph{Solución Paso a Paso:}\par
\begin{enumerate}[label=\textbf{\arabic*.}]
\item \textbf{Extremos en cada subrectángulo:}
Para cualquier partición $P$ y para todo subrectángulo $R_{ij}$:
\[ m_{ij} = \inf_{R_{ij}} c = c, \qquad M_{ij} = \sup_{R_{ij}} c = c. \]
\item \textbf{Cálculo de sumas de Darboux:}
\[ L(f,P) = \sum_{i,j} c \cdot \text{Área}(R_{ij}) = c \sum_{i,j} \text{Área}(R_{ij}) = c \cdot \text{Área}(R). \]
\[ U(f,P) = \sum_{i,j} c \cdot \text{Área}(R_{ij}) = c \sum_{i,j} \text{Área}(R_{ij}) = c \cdot \text{Área}(R). \]
\item \textbf{Aplicación del Criterio de Riemann:}
Para cualquier $\varepsilon > 0$, tomamos cualquier partición $P$ y se tiene:
\[ U(f,P) - L(f,P) = c \cdot \text{Área}(R) - c \cdot \text{Área}(R) = 0 < \varepsilon. \]
Como la diferencia es 0 (menor que todo $\varepsilon > 0$), $f$ es integrable Riemann y:
\[ \iint_R c \, dA = L(f,P) = U(f,P) = c \cdot \text{Área}(R). \]
\end{enumerate}

\begin{trapbox}
Pensar que se necesita una partición muy fina con $\Delta x \to 0$ para una función constante. Para una constante, cualquier partición satisface $U-L=0$.
\end{trapbox}

\vspace{0.5cm}
\hrule
\vspace{0.5cm}
\subsection*{Problema 4.0B: Linealidad y monotonía de la integral doble \hfill \color{teal}\small[Nivel 0]}
\addcontentsline{toc}{subsection}{Problema 4.0B: Linealidad y monotonía de la integral doble (Nivel 0)}

\begin{tcolorbox}[colback=white,colframe=gray!40,arc=1mm,boxrule=0.5pt,left=3mm,right=3mm,top=2mm,bottom=2mm]
\textbf{Enunciado:} Sean $f(x,y) = x$ y $g(x,y) = x + y^2$ dos funciones continuas sobre el cuadrado unitario $R = [0,1] \times [0,1]$. Se sabe que $\iint_R x \, dA = \frac{1}{2}$ y $\iint_R y^2 \, dA = \frac{1}{3}$.
\begin{enumerate}[label=\textbf{\alph*)}]
\item Utilice la propiedad de linealidad para calcular $\iint_R (x + y^2) \, dA$.
\item Demuestre que $f(x,y) \le g(x,y)$ en todo $(x,y) \in R$ y verifique que $\iint_R f \, dA \le \iint_R g \, dA$.
\end{enumerate}
\end{tcolorbox}

\begin{tcolorbox}[colback=yellow!4,colframe=accentamber!70,arc=1mm,boxrule=0.5pt,left=3mm,right=3mm,top=1.5mm,bottom=1.5mm]
\textbf{Pista Conceptual:} La linealidad establece que $\iint (\alpha f + \beta g) dA = \alpha \iint f dA + \beta \iint g dA$. La monotonía garantiza que si $f \le g$, entonces sus integrales conservan el orden.
\end{tcolorbox}

\paragraph{Solución Paso a Paso:}\par
\begin{enumerate}[label=\textbf{\arabic*.}]
\item \textbf{Cálculo por linealidad:}
\[ \iint_R (x + y^2) \, dA = \iint_R x \, dA + \iint_R y^2 \, dA = \frac{1}{2} + \frac{1}{3} = \frac{5}{6}. \]
\item \textbf{Comprobación de la monotonía puntual:}
Para todo $(x,y) \in [0,1] \times [0,1]$, se tiene $y^2 \ge 0$. Por lo tanto:
\[ g(x,y) - f(x,y) = (x + y^2) - x = y^2 \ge 0 \implies f(x,y) \le g(x,y). \]
\item \textbf{Monotonía de las integrales:}
\[ \iint_R f \, dA = \frac{1}{2} = \frac{3}{6} \le \frac{5}{6} = \iint_R g \, dA. \]
Se confirma experimentalmente la propiedad de monotonía.
\end{enumerate}

\begin{trapbox}
Separar productos de funciones dentro de la integral creyendo erróneamente que $\iint (f \cdot g) dA = (\iint f dA)(\iint g dA)$. La linealidad solo aplica a sumas y productos por escalares.
\end{trapbox}

\vspace{0.5cm}
\hrule
\vspace{0.5cm}
\subsection*{Problema 4.0C: Integrabilidad con una discontinuidad aislada en el origen \hfill \color{teal}\small[Nivel 0]}
\addcontentsline{toc}{subsection}{Problema 4.0C: Integrabilidad con una discontinuidad aislada en el origen (Nivel 0)}

\begin{tcolorbox}[colback=white,colframe=gray!40,arc=1mm,boxrule=0.5pt,left=3mm,right=3mm,top=2mm,bottom=2mm]
\textbf{Enunciado:} Considere la función $f: [0,1] \times [0,1] \to \mathbb{R}$ definida por $f(0,0) = 5$ y $f(x,y) = 0$ para todo $(x,y) \neq (0,0)$.
\begin{enumerate}[label=\textbf{\alph*)}]
\item Para $\delta \in (0,1)$, construya una partición que aísle al origen en un subcuadrado $[0,\delta] \times [0,\delta]$.
\item Calcule $L(f,P)$ y $U(f,P)$, y demuestre que $U(f,P) - L(f,P) = 5\delta^2$.
\item Concluya que $f$ es integrable Riemann y que su integral es 0.
\end{enumerate}
\end{tcolorbox}

\begin{tcolorbox}[colback=yellow!4,colframe=accentamber!70,arc=1mm,boxrule=0.5pt,left=3mm,right=3mm,top=1.5mm,bottom=1.5mm]
\textbf{Pista Conceptual:} En el subcuadrado que contiene al origen, el supremo es 5 y el ínfimo es 0. En todos los demás subrectángulos la función es constantemente 0.
\end{tcolorbox}

\paragraph{Solución Paso a Paso:}\par
\begin{enumerate}[label=\textbf{\arabic*.}]
\item \textbf{Construcción de la partición:}
Tomamos $P_1 = \{0, \delta, 1\}$ y $P_2 = \{0, \delta, 1\}$. Esto divide al cuadrado en 4 subrectángulos: $R_{11} = [0, \delta] \times [0, \delta]$ de área $\delta^2$, y tres subrectángulos que no contienen al origen.
\item \textbf{Cálculo de extremos:}
En $R_{11}$, $f(0,0) = 5$ y $f(x,y) = 0$ en el resto: $m_{11} = 0$, $M_{11} = 5$. En los demás subrectángulos: $m_{ij} = M_{ij} = 0$.
\item \textbf{Sumas de Darboux:}
\[ L(f,P) = 0 \cdot \delta^2 + 0 = 0, \qquad U(f,P) = 5 \cdot \delta^2 + 0 = 5\delta^2. \]
\item \textbf{Criterio de Riemann:}
Dado cualquier $\varepsilon > 0$, basta elegir $\delta < \sqrt{\frac{\varepsilon}{5}}$. Entonces:
\[ U(f,P) - L(f,P) = 5\delta^2 < 5\left(\frac{\varepsilon}{5}\right) = \varepsilon. \]
Como $U(f,P) - L(f,P) < \varepsilon$, $f$ es integrable Riemann, y al ser $L(f,P)=0$, $\iint_R f \, dA = 0$.
\end{enumerate}

\begin{trapbox}
Creer que la presencia de una discontinuidad anula automáticamente la integrabilidad. Un punto aislado tiene área cero (medida nula de Jordan) y no altera el valor de la integral.
\end{trapbox}

\vspace{0.5cm}
\hrule
\vspace{0.5cm}


\subsection*{Problema 4.1: Propiedad de acotación básica en rectángulos \hfill \color{primaryblue}\small[Nivel 1]}
\addcontentsline{toc}{subsection}{Problema 4.1: \texorpdfstring{Propiedad de acotación básica en rectángulos}{Propiedad de acotación básica en rectángulos} (Nivel 1)}

\begin{tcolorbox}[colback=white,colframe=gray!40,arc=1mm,boxrule=0.5pt,left=3mm,right=3mm,top=2mm,bottom=2mm]
\textbf{Enunciado:} Sea $f: R \to \mathbb{R}$ integrable en el rectángulo $R = [a,b] \times [c,d]$, y suponga que existen constantes reales $m, M$ tales que $m \le f(x,y) \le M$ para todo $(x,y) \in R$. Demuestre que:
\begin{equation*}
m \cdot \text{Área}(R) \le \iint_R f(x,y) \, dA \le M \cdot \text{Área}(R).
\end{equation*}
\end{tcolorbox}

\begin{tcolorbox}[colback=yellow!4,colframe=accentamber!70,arc=1mm,boxrule=0.5pt,left=3mm,right=3mm,top=1.5mm,bottom=1.5mm]
\textbf{Pista Conceptual:} Aplique la definición de sumas de Darboux observando que para cualquier partición $P$, $m \le m_j \le M_j \le M$.
\end{tcolorbox}

\paragraph{Solución Paso a Paso:}\par
\begin{enumerate}[label=\textbf{\arabic*.}]
\item \textbf{Acotación en subrectángulos:}
Sea $P = \{R_1, \dots, R_k\}$ cualquier partición de $R$. Para cada subrectángulo $R_j$:
Como $m \le f(x,y) \le M$ para todo $(x,y) \in R_j$, se tiene inmediatamente:
\begin{equation*}
m \le m_j = \inf_{R_j} f \quad \text{y} \quad M_j = \sup_{R_j} f \le M.
\end{equation*}
\end{enumerate}

\begin{enumerate}[label=\textbf{\arabic*.}]
\item \textbf{Acotación de sumas de Darboux:}
Multiplicando por el área $\text{Área}(R_j) > 0$:
\begin{equation*}
m \sum_{j=1}^k \text{Área}(R_j) \le \sum_{j=1}^k m_j \text{Área}(R_j) = L(f, P),
\end{equation*}
\begin{equation*}
U(f, P) = \sum_{j=1}^k M_j \text{Área}(R_j) \le M \sum_{j=1}^k \text{Área}(R_j).
\end{equation*}
Como $\sum_{j=1}^k \text{Área}(R_j) = \text{Área}(R) = (b-a)(d-c)$, obtenemos:
\begin{equation*}
m \cdot \text{Área}(R) \le L(f, P) \le U(f, P) \le M \cdot \text{Área}(R).
\end{equation*}
\end{enumerate}

\begin{enumerate}[label=\textbf{\arabic*.}]
\item \textbf{Paso a la integral:}
Como $f$ es integrable, $L(f, P) \le \iint_R f \, dA \le U(f, P)$ para toda partición $P$.
Por consiguiente:
\begin{equation*}
m \cdot \text{Área}(R) \le \iint_R f(x,y) \, dA \le M \cdot \text{Área}(R).
\end{equation*}
\end{enumerate}

\begin{trapbox}
Olvidar multiplicar las cotas $m$ y $M$ por el área total del rectángulo.
\end{trapbox}

\vspace{0.5cm}
\hrule
\vspace{0.5cm}

\subsection*{Problema 4.2: Desigualdad integral con exponencial cuadrática (Examen UNI 2026-I) \hfill \color{primaryblue}\small[Nivel 1]}
\addcontentsline{toc}{subsection}{Problema 4.2: \texorpdfstring{Desigualdad integral con exponencial cuadrática (Examen UNI 2026-I)}{Desigualdad integral con exponencial cuadrática (Examen UNI 2026-I)} (Nivel 1)}

\begin{tcolorbox}[colback=white,colframe=gray!40,arc=1mm,boxrule=0.5pt,left=3mm,right=3mm,top=2mm,bottom=2mm]
\textbf{Enunciado:} Demuestre analíticamente que se cumplen las siguientes desigualdades:
\begin{equation*}
e^{-1} \le \iint_{[0,1] \times [0,1]} e^{x^2 - y^2} \, dA \le e.
\end{equation*}
Justifique detalladamente el cálculo de los valores extremos del integrando.
\end{tcolorbox}

\begin{tcolorbox}[colback=yellow!4,colframe=accentamber!70,arc=1mm,boxrule=0.5pt,left=3mm,right=3mm,top=1.5mm,bottom=1.5mm]
\textbf{Pista Conceptual:} Halle el mínimo absoluto y el máximo absoluto de $g(x,y) = x^2 - y^2$ en el cuadrado unitario $[0,1]^2$. Como la función exponencial es estrictamente creciente, preserve el orden.
\end{tcolorbox}

\paragraph{Solución Paso a Paso:}\par
\begin{enumerate}[label=\textbf{\arabic*.}]
\item \textbf{Análisis de los extremos del exponente en el cuadrado unitario:}
Sea $g(x,y) = x^2 - y^2$ para $(x,y) \in [0,1] \times [0,1]$.
Dado que $x \in [0,1]$ e $y \in [0,1]$:
\begin{equation*}
0 \le x^2 \le 1, \qquad 0 \le y^2 \le 1 \implies -1 \le -y^2 \le 0.
\end{equation*}
Sumando ambas desigualdades:
\begin{equation*}
0 + (-1) \le x^2 - y^2 \le 1 + 0 \implies -1 \le x^2 - y^2 \le 1.
\end{equation*}
- El valor mínimo $-1$ se alcanza en el punto $(0, 1) \in [0,1]^2$: $g(0,1) = 0^2 - 1^2 = -1$.
- El valor máximo $1$ se alcanza en el punto $(1, 0) \in [0,1]^2$: $g(1,0) = 1^2 - 0^2 = 1$.
Por ende, los extremos absolutos en el compacto $[0,1]^2$ son exactamente $m = -1$ y $M = 1$.
\end{enumerate}

\begin{enumerate}[label=\textbf{\arabic*.}]
\item \textbf{Comportamiento monótono de la exponencial:}
Como la función exponencial $t \mapsto e^t$ es estrictamente creciente en todo $\mathbb{R}$:
\begin{equation*}
-1 \le x^2 - y^2 \le 1 \implies e^{-1} \le e^{x^2 - y^2} \le e^1 = e, \quad \forall (x,y) \in [0,1]^2.
\end{equation*}
\end{enumerate}

\begin{enumerate}[label=\textbf{\arabic*.}]
\item \textbf{Integración de la desigualdad de funciones continuas:}
Por la propiedad de monotonía y acotación de la integral de Riemann:
\begin{equation*}
\iint_{[0,1]^2} e^{-1} \, dA \le \iint_{[0,1]^2} e^{x^2 - y^2} \, dA \le \iint_{[0,1]^2} e \, dA.
\end{equation*}
Como el área del cuadrado es $\text{Área}([0,1]^2) = (1 - 0)(1 - 0) = 1$:
\begin{equation*}
\iint_{[0,1]^2} e^{-1} \, dA = e^{-1} \cdot 1 = e^{-1}, \qquad \iint_{[0,1]^2} e \, dA = e \cdot 1 = e.
\end{equation*}
Concluimos rigurosamente:
\begin{equation*}
e^{-1} \le \iint_{[0,1] \times [0,1]} e^{x^2 - y^2} \, dA \le e.
\end{equation*}
\end{enumerate}

\begin{trapbox}
Buscar puntos críticos en el interior mediante gradiente nulo sin percatarse de que el gradiente solo se anula en $(0,0)$ donde vale 0, y que los extremos ocurren en la frontera.
\end{trapbox}

\vspace{0.5cm}
\hrule
\vspace{0.5cm}

\subsection*{Problema 4.3: Linealidad de la integral sobre rectángulos \hfill \color{primaryblue}\small[Nivel 1]}
\addcontentsline{toc}{subsection}{Problema 4.3: \texorpdfstring{Linealidad de la integral sobre rectángulos}{Linealidad de la integral sobre rectángulos} (Nivel 1)}

\begin{tcolorbox}[colback=white,colframe=gray!40,arc=1mm,boxrule=0.5pt,left=3mm,right=3mm,top=2mm,bottom=2mm]
\textbf{Enunciado:} Demuestre que si $f, g: R \to \mathbb{R}$ son integrables en el rectángulo $R$ y $\alpha \in \mathbb{R}$, entonces $\alpha f$ es integrable y $\iint_R \alpha f \, dA = \alpha \iint_R f \, dA$.
\end{tcolorbox}

\begin{tcolorbox}[colback=yellow!4,colframe=accentamber!70,arc=1mm,boxrule=0.5pt,left=3mm,right=3mm,top=1.5mm,bottom=1.5mm]
\textbf{Pista Conceptual:} Separe los casos $\alpha > 0$, $\alpha = 0$ y $\alpha < 0$. Para $\alpha < 0$, observe que $\sup(\alpha f) = \alpha \inf f$.
\end{tcolorbox}

\paragraph{Solución Paso a Paso:}\par
\begin{enumerate}[label=\textbf{\arabic*.}]
\item \textbf{Caso $\alpha = 0$:}
Si $\alpha = 0$, $\alpha f = 0$. La función idénticamente nula tiene sumas de Darboux $L = U = 0$, por lo que su integral es $0 = 0 \cdot \iint_R f dA$.
\end{enumerate}

\begin{enumerate}[label=\textbf{\arabic*.}]
\item \textbf{Caso $\alpha > 0$:}
Para cualquier subconjunto $E \subset R$:
\begin{equation*}
\sup_E (\alpha f) = \alpha \sup_E f, \qquad \inf_E (\alpha f) = \alpha \inf_E f.
\end{equation*}
Por tanto, para cualquier partición $P$:
\begin{equation*}
U(\alpha f, P) = \alpha U(f, P), \qquad L(\alpha f, P) = \alpha L(f, P).
\end{equation*}
Restando:
\begin{equation*}
U(\alpha f, P) - L(\alpha f, P) = \alpha [U(f, P) - L(f, P)].
\end{equation*}
Dado $\varepsilon > 0$, como $f$ es integrable, existe $P$ tal que $U(f, P) - L(f, P) < \varepsilon / \alpha$.
Entonces $U(\alpha f, P) - L(\alpha f, P) < \varepsilon$, lo que prueba que $\alpha f$ es integrable.
Además $\inf_P U(\alpha f, P) = \alpha \inf_P U(f, P) = \alpha \iint_R f dA$.
\end{enumerate}

\begin{enumerate}[label=\textbf{\arabic*.}]
\item \textbf{Caso $\alpha = -1$:}
Para cualquier conjunto $E$: $\sup_E (-f) = -\inf_E f$ e $\inf_E (-f) = -\sup_E f$.
Entonces $U(-f, P) = -L(f, P)$ y $L(-f, P) = -U(f, P)$.
Luego $U(-f, P) - L(-f, P) = U(f, P) - L(f, P) < \varepsilon$, demostrando integrabilidad de $-f$.
Además $\iint_R (-f) dA = \inf_P U(-f, P) = \inf_P (-L(f, P)) = -\sup_P L(f, P) = -\iint_R f dA$.
\end{enumerate}

\begin{enumerate}[label=\textbf{\arabic*.}]
\item \textbf{Caso general $\alpha < 0$:}
Escribiendo $\alpha = -|\alpha|$ con $|\alpha| > 0$, combinando los casos 2 y 3 se concluye que $\iint_R \alpha f dA = \alpha \iint_R f dA$.
\end{enumerate}

\begin{trapbox}
No tratar por separado el caso $\alpha < 0$ donde se intercambian el supremo y el ínfimo.
\end{trapbox}

\vspace{0.5cm}
\hrule
\vspace{0.5cm}

\subsection*{Problema 4.4: Integrabilidad no implica continuidad (Examen UNI 2026-I) \hfill \color{primaryblue}\small[Nivel 2]}
\addcontentsline{toc}{subsection}{Problema 4.4: \texorpdfstring{Integrabilidad no implica continuidad (Examen UNI 2026-I)}{Integrabilidad no implica continuidad (Examen UNI 2026-I)} (Nivel 2)}

\begin{tcolorbox}[colback=white,colframe=gray!40,arc=1mm,boxrule=0.5pt,left=3mm,right=3mm,top=2mm,bottom=2mm]
\textbf{Enunciado:} Determine la veracidad o falsedad de la siguiente afirmación, justificando rigurosamente su respuesta mediante un contraejemplo constructivo:
``Si $f: [0, 1] \times [0, 1] \to \mathbb{R}$ es integrable, entonces $f$ es continua''.
\end{tcolorbox}

\begin{tcolorbox}[colback=yellow!4,colframe=accentamber!70,arc=1mm,boxrule=0.5pt,left=3mm,right=3mm,top=1.5mm,bottom=1.5mm]
\textbf{Pista Conceptual:} Construya una función escalonada simple con un salto a lo largo de una recta (por ejemplo $x=1/2$) y demuestre que es integrable a pesar de no ser continua.
\end{tcolorbox}

\paragraph{Solución Paso a Paso:}\par
\begin{enumerate}[label=\textbf{\arabic*.}]
\item \textbf{Calificación:} La afirmación es \textbf{FALSA}.
\end{enumerate}

\begin{enumerate}[label=\textbf{\arabic*.}]
\item \textbf{Construcción del contraejemplo:}
Definamos la función escalón $f: [0, 1] \times [0, 1] \to \mathbb{R}$ dada por:
\begin{equation*}
f(x,y) = \begin{cases} 0, & \text{si } 0 \le x < 1/2 \\ 1, & \text{si } 1/2 \le x \le 1 \end{cases}
\end{equation*}
\end{enumerate}

\begin{enumerate}[label=\textbf{\arabic*.}]
\item \textbf{Demostración de que NO es continua:}
Tomemos cualquier punto sobre el segmento de discontinuidad, por ejemplo $p_0 = (1/2, 1/2)$.
Nótese que $f(1/2, 1/2) = 1$.
Consideremos la sucesión $p_k = (1/2 - 1/k, 1/2)$ para $k \ge 3$.
Es claro que $p_k \to p_0$ en $[0,1]^2$ cuando $k \to \infty$.
Sin embargo:
\begin{equation*}
f(p_k) = f(1/2 - 1/k, 1/2) = 0, \quad \forall k \implies \lim_{k\to\infty} f(p_k) = 0 \neq f(p_0) = 1.
\end{equation*}
Por lo tanto, $f$ es discontinua en todo el segmento $\{1/2\} \times [0, 1]$.
\end{enumerate}

\begin{enumerate}[label=\textbf{\arabic*.}]
\item \textbf{Demostración de que SÍ es integrable:}
Dado cualquier $\varepsilon > 0$, elijamos $0 < \delta < \min(1/4, \varepsilon/2)$.
Consideremos la partición de $[0,1]$ en $x$: $P_1 = \{0, 1/2 - \delta, 1/2 + \delta, 1\}$ y la partición trivial de $[0,1]$ en $y$: $P_2 = \{0, 1\}$.
La partición producto $P = P_1 \times P_2$ divide el cuadrado en 3 subrectángulos:
- $R_1 = [0, 1/2 - \delta] \times [0, 1]$: $f(x,y) = 0$ idénticamente $\implies M_1 = m_1 = 0 \implies \omega_1 = 0$.
- $R_2 = [1/2 - \delta, 1/2 + \delta] \times [0, 1]$: contiene puntos de ambos lados, luego $m_2 = 0$ y $M_2 = 1 \implies \omega_2 = 1$.
  El área de $R_2$ es $\text{Área}(R_2) = (1/2 + \delta - (1/2 - \delta)) \times 1 = 2\delta$.
- $R_3 = [1/2 + \delta, 1] \times [0, 1]$: $f(x,y) = 1$ idénticamente $\implies M_3 = m_3 = 1 \implies \omega_3 = 0$.
Calculando la diferencia de Darboux:
\begin{equation*}
U(f, P) - L(f, P) = \sum_{j=1}^3 (M_j - m_j)\text{Área}(R_j) = 0 + 1 \cdot (2\delta) + 0 = 2\delta < \varepsilon.
\end{equation*}
Por el criterio de Riemann, $f$ es integrable sobre $[0,1]^2$, completando la refutación analítica.
\end{enumerate}

\begin{trapbox}
Dar un contraejemplo sin probar formalmente la integrabilidad con sumas de Darboux.
\end{trapbox}

\vspace{0.5cm}
\hrule
\vspace{0.5cm}

\subsection*{Problema 4.5: Monotonía y conservación de desigualdades \hfill \color{primaryblue}\small[Nivel 2]}
\addcontentsline{toc}{subsection}{Problema 4.5: \texorpdfstring{Monotonía y conservación de desigualdades}{Monotonía y conservación de desigualdades} (Nivel 2)}

\begin{tcolorbox}[colback=white,colframe=gray!40,arc=1mm,boxrule=0.5pt,left=3mm,right=3mm,top=2mm,bottom=2mm]
\textbf{Enunciado:} Sean $f, g: R \to \mathbb{R}$ funciones integrables en el rectángulo $R$ tales que $f(x,y) \le g(x,y)$ para todo $(x,y) \in R$. Demuestre que:
\begin{equation*}
\iint_R f(x,y) \, dA \le \iint_R g(x,y) \, dA.
\end{equation*}
\end{tcolorbox}

\begin{tcolorbox}[colback=yellow!4,colframe=accentamber!70,arc=1mm,boxrule=0.5pt,left=3mm,right=3mm,top=1.5mm,bottom=1.5mm]
\textbf{Pista Conceptual:} Defina $h(x,y) = g(x,y) - f(x,y) \ge 0$, aplique linealidad y demuestre que la integral de una función no negativa es no negativa.
\end{tcolorbox}

\paragraph{Solución Paso a Paso:}\par
\begin{enumerate}[label=\textbf{\arabic*.}]
\item \textbf{Integral de una función no negativa:}
Sea $h = g - f$. Como $f(x,y) \le g(x,y)$, se tiene $h(x,y) \ge 0$ para todo $(x,y) \in R$.
Por la linealidad de la integral, $h$ es integrable en $R$.
Para cualquier partición $P = \{R_1, \dots, R_k\}$ de $R$:
Como $h(x,y) \ge 0$ en cada $R_j$, su ínfimo satisface:
\begin{equation*}
m_j(h) = \inf_{R_j} h \ge 0, \quad \forall j.
\end{equation*}
Por ende, la suma inferior es no negativa:
\begin{equation*}
L(h, P) = \sum_{j=1}^k m_j(h) \text{Área}(R_j) \ge 0.
\end{equation*}
Dado que la integral de Riemann es el supremo de las sumas inferiores:
\begin{equation*}
\iint_R h(x,y) \, dA = \sup_P L(h, P) \ge L(h, P) \ge 0.
\end{equation*}
\end{enumerate}

\begin{enumerate}[label=\textbf{\arabic*.}]
\item \textbf{Aplicación a $f$ y $g$:}
Usando la linealidad de la integral:
\begin{equation*}
\iint_R g(x,y) \, dA - \iint_R f(x,y) \, dA = \iint_R (g(x,y) - f(x,y)) \, dA = \iint_R h(x,y) \, dA \ge 0.
\end{equation*}
Sumando $\iint_R f \, dA$ a ambos lados:
\begin{equation*}
\iint_R f(x,y) \, dA \le \iint_R g(x,y) \, dA.
\end{equation*}
\end{enumerate}

\begin{trapbox}
Asumir que si $f < g$ estrictamente en casi todo punto entonces la integral es estrictamente menor, sin considerar conjuntos de contenido cero donde podrían diferir.
\end{trapbox}

\vspace{0.5cm}
\hrule
\vspace{0.5cm}

\subsection*{Problema 4.6: Desigualdad triangular e integrabilidad del valor absoluto \hfill \color{primaryblue}\small[Nivel 2]}
\addcontentsline{toc}{subsection}{Problema 4.6: \texorpdfstring{Desigualdad triangular e integrabilidad del valor absoluto}{Desigualdad triangular e integrabilidad del valor absoluto} (Nivel 2)}

\begin{tcolorbox}[colback=white,colframe=gray!40,arc=1mm,boxrule=0.5pt,left=3mm,right=3mm,top=2mm,bottom=2mm]
\textbf{Enunciado:} Sea $f: R \to \mathbb{R}$ integrable en el rectángulo $R$. (a) Demuestre que $|f|$ es integrable en $R$. (b) Demuestre la desigualdad triangular para integrales:
\begin{equation*}
\left| \iint_R f(x,y) \, dA \right| \le \iint_R |f(x,y)| \, dA.
\end{equation*}
\end{tcolorbox}

\begin{tcolorbox}[colback=yellow!4,colframe=accentamber!70,arc=1mm,boxrule=0.5pt,left=3mm,right=3mm,top=1.5mm,bottom=1.5mm]
\textbf{Pista Conceptual:} Use la desigualdad inversa del valor absoluto: $||f(u)| - |f(v)|| \le |f(u) - f(v)|$, lo que implica $\omega(|f|, R_j) \le \omega(f, R_j)$. Luego use $-|f| \le f \le |f|$.
\end{tcolorbox}

\paragraph{Solución Paso a Paso:}\par
\begin{enumerate}[label=\textbf{\arabic*.}]
\item \textbf{Integrabilidad de $|f|$ mediante oscilación:}
Por la desigualdad triangular inversa para números reales:
\begin{equation*}
\big| |f(u)| - |f(v)| \big| \le |f(u) - f(v)|, \quad \forall u, v \in R_j.
\end{equation*}
Tomando el supremo sobre todos los pares $u, v \in R_j$:
\begin{equation*}
\omega(|f|, R_j) = \sup_{u,v \in R_j} \big| |f(u)| - |f(v)| \big| \le \sup_{u,v \in R_j} |f(u) - f(v)| = \omega(f, R_j).
\end{equation*}
Multiplicando por $\text{Área}(R_j)$ y sumando sobre toda partición $P$:
\begin{equation*}
U(|f|, P) - L(|f|, P) = \sum_{j} \omega(|f|, R_j)\text{Área}(R_j) \le \sum_{j} \omega(f, R_j)\text{Área}(R_j) = U(f, P) - L(f, P).
\end{equation*}
Como $f$ es integrable, para todo $\varepsilon > 0$ existe $P$ tal que $U(f, P) - L(f, P) < \varepsilon$.
Por tanto $U(|f|, P) - L(|f|, P) < \varepsilon$, lo que demuestra por el criterio de Riemann que $|f|$ es integrable.
\end{enumerate}

\begin{enumerate}[label=\textbf{\arabic*.}]
\item \textbf{Demostración de la desigualdad triangular:}
Para todo $(x,y) \in R$, por propiedades del valor absoluto:
\begin{equation*}
-|f(x,y)| \le f(x,y) \le |f(x,y)|.
\end{equation*}
Aplicando la propiedad de monotonía de la integral:
\begin{equation*}
-\iint_R |f(x,y)| \, dA \le \iint_R f(x,y) \, dA \le \iint_R |f(x,y)| \, dA.
\end{equation*}
Esta doble desigualdad equivale exactamente a:
\begin{equation*}
\left| \iint_R f(x,y) \, dA \right| \le \iint_R |f(x,y)| \, dA.
\end{equation*}
\end{enumerate}

\begin{trapbox}
Creer que el recíproco es cierto: si $|f|$ es integrable, $f$ NO es necesariamente integrable (ej. $f = 1$ en racionales y $-1$ en irracionales).
\end{trapbox}

\vspace{0.5cm}
\hrule
\vspace{0.5cm}

\subsection*{Problema 4.7: Criterio de integrabilidad de Riemann-Darboux \hfill \color{primaryblue}\small[Nivel 3]}
\addcontentsline{toc}{subsection}{Problema 4.7: \texorpdfstring{Criterio de integrabilidad de Riemann-Darboux}{Criterio de integrabilidad de Riemann-Darboux} (Nivel 3)}

\begin{tcolorbox}[colback=white,colframe=gray!40,arc=1mm,boxrule=0.5pt,left=3mm,right=3mm,top=2mm,bottom=2mm]
\textbf{Enunciado:} Sea $f: R \to \mathbb{R}$ acotada en el rectángulo $R$. Demuestre el Criterio de Riemann:
$f$ es integrable en $R$ si y sólo si para todo $\varepsilon > 0$, existe una partición $P$ de $R$ tal que:
\begin{equation*}
U(f, P) - L(f, P) < \varepsilon.
\end{equation*}
\end{tcolorbox}

\begin{tcolorbox}[colback=yellow!4,colframe=accentamber!70,arc=1mm,boxrule=0.5pt,left=3mm,right=3mm,top=1.5mm,bottom=1.5mm]
\textbf{Pista Conceptual:} Recuerde que $\underline{\int} f = \sup_P L(f,P)$ y $\overline{\int} f = \inf_P U(f,P)$. La integrabilidad exige que $\overline{\int} f - \underline{\int} f = 0$.
\end{tcolorbox}

\paragraph{Solución Paso a Paso:}\par
\begin{enumerate}[label=\textbf{\arabic*.}]
\item \textbf{Demostración de la necesidad ($\implies$):}
Supongamos que $f$ es integrable en $R$, es decir:
\begin{equation*}
I = \sup_P L(f, P) = \inf_P U(f, P).
\end{equation*}
Sea $\varepsilon > 0$. Por la definición de ínfimo como la mayor cota inferior, existe una partición $P_1$ tal que:
\begin{equation*}
U(f, P_1) < I + \frac{\varepsilon}{2}.
\end{equation*}
Por la definición de supremo como la menor cota superior, existe una partición $P_2$ tal que:
\begin{equation*}
L(f, P_2) > I - \frac{\varepsilon}{2}.
\end{equation*}
Sea $P = P_1 \cup P_2$ el refinamiento común. Por el lema del refinamiento:
\begin{equation*}
U(f, P) \le U(f, P_1) < I + \frac{\varepsilon}{2},
\end{equation*}
\begin{equation*}
L(f, P) \ge L(f, P_2) > I - \frac{\varepsilon}{2} \iff -L(f, P) < -I + \frac{\varepsilon}{2}.
\end{equation*}
Sumando ambas desigualdades:
\begin{equation*}
U(f, P) - L(f, P) < \left(I + \frac{\varepsilon}{2}\right) + \left(-I + \frac{\varepsilon}{2}\right) = \varepsilon.
\end{equation*}
\end{enumerate}

\begin{enumerate}[label=\textbf{\arabic*.}]
\item \textbf{Demostración de la suficiencia ($\impliedby$):}
Supongamos que para todo $\varepsilon > 0$ existe una partición $P$ tal que $U(f, P) - L(f, P) < \varepsilon$.
Para dicha partición $P$, como $L(f, P) \le \underline{\int} f \le \overline{\int} f \le U(f, P)$, tenemos:
\begin{equation*}
0 \le \overline{\iint_R} f \, dA - \underline{\iint_R} f \, dA \le U(f, P) - L(f, P) < \varepsilon.
\end{equation*}
Como la desigualdad $0 \le \overline{\iint_R} f - \underline{\iint_R} f < \varepsilon$ es válida para todo número positivo $\varepsilon > 0$, la única posibilidad numérica es que:
\begin{equation*}
\overline{\iint_R} f \, dA - \underline{\iint_R} f \, dA = 0 \implies \overline{\iint_R} f \, dA = \underline{\iint_R} f \, dA.
\end{equation*}
Por definición de Darboux, esto demuestra que $f$ es integrable en $R$.
\end{enumerate}

\begin{trapbox}
Tomar particiones distintas $P_1$ y $P_2$ sin unirlas en una partición común $P = P_1 \cup P_2$.
\end{trapbox}

\vspace{0.5cm}
\hrule
\vspace{0.5cm}

\subsection*{Problema 4.8: Toda función continua sobre un rectángulo compacto es integrable \hfill \color{primaryblue}\small[Nivel 3]}
\addcontentsline{toc}{subsection}{Problema 4.8: \texorpdfstring{Toda función continua sobre un rectángulo compacto es integrable}{Toda función continua sobre un rectángulo compacto es integrable} (Nivel 3)}

\begin{tcolorbox}[colback=white,colframe=gray!40,arc=1mm,boxrule=0.5pt,left=3mm,right=3mm,top=2mm,bottom=2mm]
\textbf{Enunciado:} Sea $R = [a,b] \times [c,d] \subset \mathbb{R}^2$ un rectángulo cerrado y acotado, y sea $f: R \to \mathbb{R}$ continua en $R$. Demuestre rigurosamente, empleando el Teorema de Heine-Cantor y el Criterio de Riemann, que $f$ es integrable sobre $R$.
\end{tcolorbox}

\begin{tcolorbox}[colback=yellow!4,colframe=accentamber!70,arc=1mm,boxrule=0.5pt,left=3mm,right=3mm,top=1.5mm,bottom=1.5mm]
\textbf{Pista Conceptual:} Como $R$ es compacto, $f$ es uniformemente continua: $\forall \varepsilon' > 0, \exists \delta > 0$ tal que si $\|p - q\| < \delta$, entonces $|f(p) - f(q)| < \varepsilon'$. Elija la norma de la partición menor que $\delta$.
\end{tcolorbox}

\paragraph{Solución Paso a Paso:}\par
\begin{enumerate}[label=\textbf{\arabic*.}]
\item \textbf{Continuidad uniforme (Teorema de Heine-Cantor):}
El rectángulo $R = [a,b] \times [c,d]$ es un conjunto cerrado y acotado en $\mathbb{R}^2$, luego es compacto por el Teorema de Heine-Borel.
Por el Teorema de Heine-Cantor, toda función continua definida sobre un conjunto compacto es uniformemente continua.
Esto significa que para todo $\varepsilon > 0$, existe $\delta > 0$ tal que:
\begin{equation*}
\forall p, q \in R, \quad \|p - q\| < \delta \implies |f(p) - f(q)| < \frac{\varepsilon}{\text{Área}(R)}.
\end{equation*}
\end{enumerate}

\begin{enumerate}[label=\textbf{\arabic*.}]
\item \textbf{Elección de la partición con norma menor que $\delta$:}
Elijamos una partición regular $P = \{R_1, \dots, R_k\}$ con $m$ y $n$ subdivisiones suficientemente grandes de modo que el diámetro de cada subrectángulo $R_j$ sea menor que $\delta$:
\begin{equation*}
\text{diám}(R_j) = \sqrt{(\Delta x)^2 + (\Delta y)^2} < \delta, \quad \forall j \in \{1, \dots, k\}.
\end{equation*}
\end{enumerate}

\begin{enumerate}[label=\textbf{\arabic*.}]
\item \textbf{Control de la oscilación en cada subrectángulo:}
Dado que $R_j$ es compacto y $f$ es continua, por el Teorema de Weierstrass $f$ alcanza su valor mínimo $m_j$ en algún punto $p_j \in R_j$ y su valor máximo $M_j$ en algún punto $q_j \in R_j$:
\begin{equation*}
m_j = f(p_j), \qquad M_j = f(q_j).
\end{equation*}
Dado que $p_j, q_j \in R_j$, la distancia entre ellos cumple $\|p_j - q_j\| \le \text{diám}(R_j) < \delta$.
Por la continuidad uniforme:
\begin{equation*}
M_j - m_j = f(q_j) - f(p_j) = |f(q_j) - f(p_j)| < \frac{\varepsilon}{\text{Área}(R)}.
\end{equation*}
\end{enumerate}

\begin{enumerate}[label=\textbf{\arabic*.}]
\item \textbf{Verificación del Criterio de Riemann:}
Calculamos la diferencia de sumas de Darboux:
\begin{equation*}
U(f, P) - L(f, P) = \sum_{j=1}^k (M_j - m_j) \text{Área}(R_j) < \sum_{j=1}^k \frac{\varepsilon}{\text{Área}(R)} \text{Área}(R_j)
\end{equation*}
\begin{equation*}
= \frac{\varepsilon}{\text{Área}(R)} \sum_{j=1}^k \text{Área}(R_j) = \frac{\varepsilon}{\text{Área}(R)} \cdot \text{Área}(R) = \varepsilon.
\end{equation*}
Por el Criterio de Riemann (ejercicio 4.7), $f$ es integrable sobre $R$.
\end{enumerate}

\begin{trapbox}
Usar continuidad puntual en vez de continuidad uniforme, lo que haría que $\delta$ dependa de cada punto impidiendo acotar todos los subrectángulos simultáneamente.
\end{trapbox}

\vspace{0.5cm}
\hrule
\vspace{0.5cm}

\subsection*{Problema 4.9: Integrabilidad con discontinuidades sobre curvas suaves \hfill \color{primaryblue}\small[Nivel 3]}
\addcontentsline{toc}{subsection}{Problema 4.9: \texorpdfstring{Integrabilidad con discontinuidades sobre curvas suaves}{Integrabilidad con discontinuidades sobre curvas suaves} (Nivel 3)}

\begin{tcolorbox}[colback=white,colframe=gray!40,arc=1mm,boxrule=0.5pt,left=3mm,right=3mm,top=2mm,bottom=2mm]
\textbf{Enunciado:} Sea $R = [0,1] \times [0,1]$ y sea $\gamma: [0,1] \to [0,1]$ una función continua. Considere la gráfica $C = \{(x, \gamma(x)) : x \in [0,1]\}$. Demuestre que $C$ tiene contenido bidimensional nulo (contenido cero en $\mathbb{R}^2$), y deduzca que toda función acotada $f: R \to \mathbb{R}$ cuyas discontinuidades estén contenidas en $C$ es integrable sobre $R$.
\end{tcolorbox}

\begin{tcolorbox}[colback=yellow!4,colframe=accentamber!70,arc=1mm,boxrule=0.5pt,left=3mm,right=3mm,top=1.5mm,bottom=1.5mm]
\textbf{Pista Conceptual:} Como $\gamma$ es continua en $[0,1]$, es uniformemente continua. Cubra la curva con una cantidad finita de rectángulos de base $\Delta x$ y altura menor a $\varepsilon$.
\end{tcolorbox}

\paragraph{Solución Paso a Paso:}\par
\begin{enumerate}[label=\textbf{\arabic*.}]
\item \textbf{Definición de contenido cero:}
Un conjunto $E \subset \mathbb{R}^2$ tiene contenido cero si para todo $\varepsilon > 0$, existe una cantidad \textit{finita} de rectángulos abiertos o cerrados $\{J_1, \dots, J_N\}$ tales que $E \subset \bigcup_{k=1}^N J_k$ y $\sum_{k=1}^N \text{Área}(J_k) < \varepsilon$.
\end{enumerate}

\begin{enumerate}[label=\textbf{\arabic*.}]
\item \textbf{Demostración de que la gráfica continua $C$ tiene contenido cero:}
Como $\gamma: [0,1] \to [0,1]$ es continua sobre un compacto, es uniformemente continua:
Dado $\varepsilon > 0$, existe $\delta > 0$ tal que si $|x - x'| < \delta$, entonces $|\gamma(x) - \gamma(x')| < \varepsilon / 2$.
Particionemos $[0,1]$ en $N$ subintervalos $[x_{k-1}, x_k]$ de longitud igual $\Delta x = 1/N < \delta$.
En cada subintervalo $[x_{k-1}, x_k]$, definamos:
\begin{equation*}
c_k = \min_{x \in [x_{k-1}, x_k]} \gamma(x), \qquad d_k = \max_{x \in [x_{k-1}, x_k]} \gamma(x).
\end{equation*}
Por la continuidad uniforme, $d_k - c_k < \varepsilon / 2$.
Consideremos los rectángulos:
\begin{equation*}
J_k = [x_{k-1}, x_k] \times [c_k, d_k], \quad k = 1, \dots, N.
\end{equation*}
Por construcción, la gráfica de $\gamma$ sobre $[x_{k-1}, x_k]$ está completamente contenida en $J_k$, luego $C \subset \bigcup_{k=1}^N J_k$.
El área total cubierta es:
\begin{equation*}
\sum_{k=1}^N \text{Área}(J_k) = \sum_{k=1}^N (x_k - x_{k-1})(d_k - c_k) < \sum_{k=1}^N \Delta x \cdot \frac{\varepsilon}{2} = \frac{\varepsilon}{2} \sum_{k=1}^N \Delta x = \frac{\varepsilon}{2} \cdot 1 = \frac{\varepsilon}{2} < \varepsilon.
\end{equation*}
Por tanto, la gráfica $C$ tiene contenido cero.
\end{enumerate}

\begin{enumerate}[label=\textbf{\arabic*.}]
\item \textbf{Integrabilidad de $f$:}
Dado que $f$ es acotada ($\|f\|_\infty \le M$) y sus discontinuidades están en $C$, el conjunto de subrectángulos que cortan a $C$ se puede encerrar en una región de área total menor que $\varepsilon / (4M)$.
En el resto del rectángulo $R$, $f$ es continua, por lo que su oscilación se puede hacer menor que $\varepsilon / (2\text{Área}(R))$.
Sumando ambas contribuciones, $U(f,P) - L(f,P) < \varepsilon$, demostrando la integrabilidad de $f$.
\end{enumerate}

\begin{trapbox}
Confundir contenido cero (cubrimiento finito) con medida cero de Lebesgue (cubrimiento numerable).
\end{trapbox}

\vspace{0.5cm}
\hrule
\vspace{0.5cm}

\subsection*{Problema 4.10: Cotas de integrales dobles para sumas continuas (Examen UNI 2024-II) \hfill \color{primaryblue}\small[Nivel 4]}
\addcontentsline{toc}{subsection}{Problema 4.10: \texorpdfstring{Cotas de integrales dobles para sumas continuas (Examen UNI 2024-II)}{Cotas de integrales dobles para sumas continuas (Examen UNI 2024-II)} (Nivel 4)}

\begin{tcolorbox}[colback=white,colframe=gray!40,arc=1mm,boxrule=0.5pt,left=3mm,right=3mm,top=2mm,bottom=2mm]
\textbf{Enunciado:} Sean $f: [a,b] \to \mathbb{R}$ y $g: [c,d] \to \mathbb{R}$ dos funciones continuas de modo que:
\begin{equation*}
0 \le n_1 \le f(x) \le N_1, \quad \forall x \in [a,b], \qquad 0 \le n_2 \le g(y) \le N_2, \quad \forall y \in [c,d].
\end{equation*}
Demuestre rigurosamente que se cumple la siguiente desigualdad de promedios integrales:
\begin{equation*}
(n_1 + n_2) \le \frac{\int_a^b \int_c^d (f(x) + g(y)) \, dy \, dx}{(b - a)(d - c)} \le (N_1 + N_2).
\end{equation*}
\end{tcolorbox}

\begin{tcolorbox}[colback=yellow!4,colframe=accentamber!70,arc=1mm,boxrule=0.5pt,left=3mm,right=3mm,top=1.5mm,bottom=1.5mm]
\textbf{Pista Conceptual:} Exprese la integral doble como suma de dos integrales simples separables: $\int_a^b \int_c^d (f(x)+g(y)) dy dx = (d-c)\int_a^b f(x)dx + (b-a)\int_c^d g(y)dy$. Luego acote cada integral unidimensional.
\end{tcolorbox}

\paragraph{Solución Paso a Paso:}\par
\begin{enumerate}[label=\textbf{\arabic*.}]
\item \textbf{Descomposición lineal de la integral doble:}
Dado que $f(x)$ y $g(y)$ son continuas, la suma $F(x,y) = f(x) + g(y)$ es continua sobre el rectángulo $R = [a,b] \times [c,d]$.
Por la linealidad de la integral y el Teorema de Fubini:
\begin{equation*}
\int_a^b \int_c^d (f(x) + g(y)) \, dy \, dx = \int_a^b \left[ \int_c^d f(x) \, dy + \int_c^d g(y) \, dy \right] dx
\end{equation*}
Como $f(x)$ no depende de $y$: $\int_c^d f(x) dy = f(x) \int_c^d 1 dy = (d - c) f(x)$.
Por tanto:
\begin{equation*}
\int_a^b \int_c^d (f(x) + g(y)) \, dy \, dx = (d - c) \int_a^b f(x) \, dx + (b - a) \int_c^d g(y) \, dy.
\end{equation*}
\end{enumerate}

\begin{enumerate}[label=\textbf{\arabic*.}]
\item \textbf{Acotación de las integrales unidimensionales:}
Por hipótesis, $n_1 \le f(x) \le N_1$ en $[a,b]$. Integrando en $[a,b]$:
\begin{equation*}
n_1(b - a) \le \int_a^b f(x) \, dx \le N_1(b - a).
\end{equation*}
Multiplicando por la longitud positiva $(d - c) > 0$:
\begin{equation*}
n_1(b - a)(d - c) \le (d - c)\int_a^b f(x) \, dx \le N_1(b - a)(d - c).
\end{equation*}
Análogamente, por hipótesis $n_2 \le g(y) \le N_2$ en $[c,d]$. Integrando en $[c,d]$:
\begin{equation*}
n_2(d - c) \le \int_c^d g(y) \, dy \le N_2(d - c).
\end{equation*}
Multiplicando por la longitud positiva $(b - a) > 0$:
\begin{equation*}
n_2(b - a)(d - c) \le (b - a)\int_c^d g(y) \, dy \le N_2(b - a)(d - c).
\end{equation*}
\end{enumerate}

\begin{enumerate}[label=\textbf{\arabic*.}]
\item \textbf{Suma de desigualdades:}
Sumando miembro a miembro ambas desigualdades:
\begin{equation*}
(n_1 + n_2)(b - a)(d - c) \le (d - c)\int_a^b f(x) \, dx + (b - a)\int_c^d g(y) \, dy \le (N_1 + N_2)(b - a)(d - c)
\end{equation*}
\begin{equation*}
\implies (n_1 + n_2)(b - a)(d - c) \le \int_a^b \int_c^d (f(x) + g(y)) \, dy \, dx \le (N_1 + N_2)(b - a)(d - c).
\end{equation*}
\end{enumerate}

\begin{enumerate}[label=\textbf{\arabic*.}]
\item \textbf{Conclusión dividiendo entre el área $(b-a)(d-c) > 0$:}
\begin{equation*}
(n_1 + n_2) \le \frac{\int_a^b \int_c^d (f(x) + g(y)) \, dy \, dx}{(b - a)(d - c)} \le (N_1 + N_2).
\end{equation*}
Esto completa con exactitud la demostración requerida en el examen de la UNI.
\end{enumerate}

\begin{trapbox}
Intentar acotar directamente $f(x)+g(y)$ sin descomponer la integral, lo cual es matemáticamente correcto pero no demuestra la factorización analítica en términos de las áreas unidimensionales.
\end{trapbox}

\vspace{0.5cm}
\hrule
\vspace{0.5cm}

\subsection*{Problema 4.11: Teorema de Lebesgue-Vitali para integrabilidad en bloques \hfill \color{primaryblue}\small[Nivel 4]}
\addcontentsline{toc}{subsection}{Problema 4.11: \texorpdfstring{Teorema de Lebesgue-Vitali para integrabilidad en bloques}{Teorema de Lebesgue-Vitali para integrabilidad en bloques} (Nivel 4)}

\begin{tcolorbox}[colback=white,colframe=gray!40,arc=1mm,boxrule=0.5pt,left=3mm,right=3mm,top=2mm,bottom=2mm]
\textbf{Enunciado:} Enuncie con precisión el Teorema de Lebesgue-Vitali sobre la caracterización de la integrabilidad de Riemann en bloques de $\mathbb{R}^n$. Aplique este teorema para demostrar formalmente que la función característica del disco unitario $\chi_{D}$ (donde $D = \{(x,y) : x^2 + y^2 \le 1\}$) es integrable Riemann sobre el cuadrado $[-1,1] \times [-1,1]$.
\end{tcolorbox}

\begin{tcolorbox}[colback=yellow!4,colframe=accentamber!70,arc=1mm,boxrule=0.5pt,left=3mm,right=3mm,top=1.5mm,bottom=1.5mm]
\textbf{Pista Conceptual:} Identifique los puntos de discontinuidad de $\chi_D$. Demuestre que el conjunto de discontinuidades es exactamente la circunferencia frontera $\partial D$, y pruebe que una circunferencia tiene medida cero de Lebesgue en $\mathbb{R}^2$.
\end{tcolorbox}

\paragraph{Solución Paso a Paso:}\par
\begin{enumerate}[label=\textbf{\arabic*.}]
\item \textbf{Enunciado del Teorema de Lebesgue-Vitali:}
Una función acotada $f: R \subset \mathbb{R}^n \to \mathbb{R}$ definida sobre un bloque compacto $R$ es integrable en el sentido de Riemann si y sólo si el conjunto de sus puntos de discontinuidad:
\begin{equation*}
D_f = \{x \in R : f \text{ es discontinua en } x\}
\end{equation*}
es un conjunto de medida cero de Lebesgue en $\mathbb{R}^n$ (es decir, $f$ es continua casi en todas partes según la medida de Lebesgue).
\end{enumerate}

\begin{enumerate}[label=\textbf{\arabic*.}]
\item \textbf{Identificación de los puntos de discontinuidad de $\chi_D$:}
Sea $R = [-1,1]^2$.
- Si $(x,y) \in D^\circ$ (interior del disco, $x^2 + y^2 < 1$), existe una bola abierta $B((x,y), r) \subset D^\circ$ donde $\chi_D \equiv 1$ constante, luego $\chi_D$ es continua.
- Si $(x,y) \in R \setminus D$ (exterior del disco, $x^2 + y^2 > 1$), existe una bola abierta donde $\chi_D \equiv 0$ constante, luego $\chi_D$ es continua.
- Si $(x,y) \in \partial D = \{(x,y) : x^2 + y^2 = 1\}$, cualquier bola abierta centrada en $(x,y)$ contiene tanto puntos interiores a $D$ (donde $\chi_D = 1$) como exteriores a $D$ (donde $\chi_D = 0$).
Por tanto, $\chi_D$ es discontinua exactamente sobre la frontera $\partial D = S^1$.
\end{enumerate}

\begin{enumerate}[label=\textbf{\arabic*.}]
\item \textbf{Medida cero de la circunferencia:}
La circunferencia $S^1$ es la imagen de la curva diferenciable $\gamma(t) = (\cos t, \sin t)$ para $t \in [0, 2\pi]$.
Como $\gamma$ es de clase $C^1$ (suave), la imagen de un intervalo compacto bajo una aplicación $C^1$ en un espacio de dimensión mayor ($1 < 2$) tiene medida cero de Lebesgue en $\mathbb{R}^2$.
En efecto, $S^1$ se puede descomponer en dos gráficas de funciones continuas $y = \pm \sqrt{1 - x^2}$ para $x \in [-1,1]$, y como probamos en el ejercicio 4.9, la gráfica de toda función continua sobre un compacto tiene contenido nulo y por ende medida cero de Lebesgue.
\end{enumerate}

\begin{enumerate}[label=\textbf{\arabic*.}]
\item \textbf{Conclusión:}
Como el conjunto de discontinuidades $D_{\chi_D} = S^1$ tiene medida cero de Lebesgue y $\chi_D$ es acotada ($0 \le \chi_D \le 1$), por el Teorema de Lebesgue-Vitali $\chi_D$ es integrable Riemann sobre $R$.
\end{enumerate}

\begin{trapbox}
Afirmar que la frontera de cualquier conjunto tiene medida cero; existen conjuntos fractales patológicos donde la frontera tiene medida positiva (conjuntos no medibles Jordan).
\end{trapbox}

\vspace{0.5cm}
\hrule
\vspace{0.5cm}

\subsection*{Problema 4.12: Integrabilidad del producto mediante identidad de polarización \hfill \color{primaryblue}\small[Nivel 4]}
\addcontentsline{toc}{subsection}{Problema 4.12: \texorpdfstring{Integrabilidad del producto mediante identidad de polarización}{Integrabilidad del producto mediante identidad de polarización} (Nivel 4)}

\begin{tcolorbox}[colback=white,colframe=gray!40,arc=1mm,boxrule=0.5pt,left=3mm,right=3mm,top=2mm,bottom=2mm]
\textbf{Enunciado:} Demuestre que si $f, g: R \to \mathbb{R}$ son integrables en el rectángulo $R$, entonces el producto $f \cdot g$ es integrable en $R$.
\end{tcolorbox}

\begin{tcolorbox}[colback=yellow!4,colframe=accentamber!70,arc=1mm,boxrule=0.5pt,left=3mm,right=3mm,top=1.5mm,bottom=1.5mm]
\textbf{Pista Conceptual:} Demuestre primero que si $h$ es integrable, $h^2$ es integrable usando la cota $|h(u)^2 - h(v)^2| \le 2M |h(u) - h(v)|$. Luego utilice la identidad de polarización $fg = \frac{1}{4}((f+g)^2 - (f-g)^2)$.
\end{tcolorbox}

\paragraph{Solución Paso a Paso:}\par
\begin{enumerate}[label=\textbf{\arabic*.}]
\item \textbf{Integrabilidad del cuadrado $h^2$:}
Sea $h$ integrable en $R$. Como $h$ es integrable, es acotada: existe $M > 0$ tal que $|h(x,y)| \le M$ para todo $(x,y) \in R$.
Para cualesquiera $u, v \in R_j$:
\begin{equation*}
|h(u)^2 - h(v)^2| = |h(u) + h(v)| \cdot |h(u) - h(v)| \le (|h(u)| + |h(v)|) |h(u) - h(v)| \le 2M |h(u) - h(v)|.
\end{equation*}
Tomando el supremo sobre todos los pares $u, v \in R_j$:
\begin{equation*}
\omega(h^2, R_j) \le 2M \omega(h, R_j).
\end{equation*}
Multiplicando por el área $\text{Área}(R_j)$ y sumando sobre una partición $P$:
\begin{equation*}
U(h^2, P) - L(h^2, P) \le 2M [U(h, P) - L(h, P)].
\end{equation*}
Dado $\varepsilon > 0$, por la integrabilidad de $h$ existe $P$ tal que $U(h, P) - L(h, P) < \frac{\varepsilon}{2M}$.
Entonces $U(h^2, P) - L(h^2, P) < \varepsilon$, lo que demuestra por el Criterio de Riemann que $h^2$ es integrable.
\end{enumerate}

\begin{enumerate}[label=\textbf{\arabic*.}]
\item \textbf{Uso de la identidad de polarización algebraica:}
Recordemos la identidad algebraica elemental:
\begin{equation*}
f(x,y) \cdot g(x,y) = \frac{1}{4} \left[ (f(x,y) + g(x,y))^2 - (f(x,y) - g(x,y))^2 \right].
\end{equation*}
\end{enumerate}

\begin{enumerate}[label=\textbf{\arabic*.}]
\item \textbf{Conclusión por el álgebra de funciones integrables:}
- Como $f$ y $g$ son integrables, por linealidad $f + g$ y $f - g$ son integrables.
- Por el paso 1, los cuadrados $(f+g)^2$ y $(f-g)^2$ son integrables.
- Nuevamente por linealidad, la resta $(f+g)^2 - (f-g)^2$ y el producto por el escalar $1/4$ son integrables.
Por lo tanto, $f \cdot g$ es integrable sobre $R$.
\end{enumerate}

\begin{trapbox}
Intentar acotar directamente $\omega(fg) \le \omega(f)\omega(g)$, lo cual es falso si las funciones toman valores negativos o mayores que 1.
\end{trapbox}

\vspace{0.5cm}
\hrule
\vspace{0.5cm}

\subsection*{Problema 4.13: Teorema de composición de funciones integrables \hfill \color{primaryblue}\small[Nivel 5]}
\addcontentsline{toc}{subsection}{Problema 4.13: \texorpdfstring{Teorema de composición de funciones integrables}{Teorema de composición de funciones integrables} (Nivel 5)}

\begin{tcolorbox}[colback=white,colframe=gray!40,arc=1mm,boxrule=0.5pt,left=3mm,right=3mm,top=2mm,bottom=2mm]
\textbf{Enunciado:} Sea $f: R \to [a,b]$ una función integrable Riemann en el rectángulo $R$, y sea $\phi: [a,b] \to \mathbb{R}$ una función continua. Demuestre rigurosamente que la función compuesta $\phi \circ f: R \to \mathbb{R}$ es integrable Riemann sobre $R$.
\end{tcolorbox}

\begin{tcolorbox}[colback=yellow!4,colframe=accentamber!70,arc=1mm,boxrule=0.5pt,left=3mm,right=3mm,top=1.5mm,bottom=1.5mm]
\textbf{Pista Conceptual:} Use la continuidad uniforme de $\phi$ en $[a,b]$. Divida los subrectángulos de una partición en dos clases: los de oscilación pequeña de $f$, y los de oscilación grande (cuya área total es pequeña).
\end{tcolorbox}

\paragraph{Solución Paso a Paso:}\par
\begin{enumerate}[label=\textbf{\arabic*.}]
\item \textbf{Continuidad uniforme de $\phi$:}
Como $[a,b]$ es compacto y $\phi$ es continua, $\phi$ es acotada ($\|\phi\|_\infty \le K$) y uniformemente continua en $[a,b]$:
Dado $\varepsilon > 0$, existe $\eta > 0$ (pudiendo asumir $\eta < \varepsilon$) tal que:
\begin{equation*}
\forall s, t \in [a,b], \quad |s - t| < \eta \implies |\phi(s) - \phi(t)| < \frac{\varepsilon}{2\text{Área}(R)}.
\end{equation*}
\end{enumerate}

\begin{enumerate}[label=\textbf{\arabic*.}]
\item \textbf{Partición de subrectángulos según la oscilación de $f$:}
Como $f$ es integrable, existe una partición $P = \{R_1, \dots, R_m\}$ tal que:
\begin{equation*}
U(f, P) - L(f, P) = \sum_{j=1}^m \omega(f, R_j) \text{Área}(R_j) < \eta^2.
\end{equation*}
Dividimos los índices $\{1, \dots, m\}$ en dos conjuntos disjuntos $A$ y $B$:
- $A = \{j : \omega(f, R_j) < \eta\}$ (subrectángulos donde la oscilación de $f$ es pequeña).
- $B = \{j : \omega(f, R_j) \ge \eta\}$ (subrectángulos donde la oscilación de $f$ es grande).
\end{enumerate}

\begin{enumerate}[label=\textbf{\arabic*.}]
\item \textbf{Cota del área del conjunto $B$:}
Nótese que:
\begin{equation*}
\eta \sum_{j \in B} \text{Área}(R_j) \le \sum_{j \in B} \omega(f, R_j) \text{Área}(R_j) \le U(f, P) - L(f, P) < \eta^2.
\end{equation*}
Dividiendo entre $\eta > 0$:
\begin{equation*}
\sum_{j \in B} \text{Área}(R_j) < \eta < \frac{\varepsilon}{4K}.
\end{equation*}
\end{enumerate}

\begin{enumerate}[label=\textbf{\arabic*.}]
\item \textbf{Acotación de la diferencia de Darboux para $\phi \circ f$:}
Para $j \in A$: si $u, v \in R_j$, entonces $|f(u) - f(v)| \le \omega(f, R_j) < \eta$.
Por la continuidad uniforme de $\phi$: $|\phi(f(u)) - \phi(f(v))| < \frac{\varepsilon}{2\text{Área}(R)}$, luego $\omega(\phi \circ f, R_j) \le \frac{\varepsilon}{2\text{Área}(R)}$.
Para $j \in B$: $\omega(\phi \circ f, R_j) \le 2\|\phi\|_\infty \le 2K$.
Sumando sobre todos los subrectángulos:
\begin{equation*}
U(\phi \circ f, P) - L(\phi \circ f, P) = \sum_{j \in A} \omega(\phi \circ f, R_j)\text{Área}(R_j) + \sum_{j \in B} \omega(\phi \circ f, R_j)\text{Área}(R_j)
\end{equation*}
\begin{equation*}
< \frac{\varepsilon}{2\text{Área}(R)} \sum_{j \in A} \text{Área}(R_j) + 2K \sum_{j \in B} \text{Área}(R_j) < \frac{\varepsilon}{2\text{Área}(R)} \text{Área}(R) + 2K \left(\frac{\varepsilon}{4K}\right) = \frac{\varepsilon}{2} + \frac{\varepsilon}{2} = \varepsilon.
\end{equation*}
Por el Criterio de Riemann, $\phi \circ f$ es integrable.
\end{enumerate}

\begin{trapbox}
Creer que la composición en orden inverso $f \circ \phi$ (donde la función externa es integrable y la interna es continua) también es siempre integrable Riemann (es falso, existen contraejemplos célebres).
\end{trapbox}

\vspace{0.5cm}
\hrule
\vspace{0.5cm}

\subsection*{Problema 4.14: Conjuntos medibles Jordan y condición de frontera con contenido nulo \hfill \color{primaryblue}\small[Nivel 5]}
\addcontentsline{toc}{subsection}{Problema 4.14: \texorpdfstring{Conjuntos medibles Jordan y condición de frontera con contenido nulo}{Conjuntos medibles Jordan y condición de frontera con contenido nulo} (Nivel 5)}

\begin{tcolorbox}[colback=white,colframe=gray!40,arc=1mm,boxrule=0.5pt,left=3mm,right=3mm,top=2mm,bottom=2mm]
\textbf{Enunciado:} Sea $\Omega \subset \mathbb{R}^2$ un conjunto acotado. Demuestre que la función característica $\chi_\Omega$ es integrable Riemann sobre un rectángulo contenedor $R \supset \Omega$ si y sólo si la frontera topológica $\partial \Omega$ tiene contenido cero (es decir, $\Omega$ es medible según Jordan).
\end{tcolorbox}

\begin{tcolorbox}[colback=yellow!4,colframe=accentamber!70,arc=1mm,boxrule=0.5pt,left=3mm,right=3mm,top=1.5mm,bottom=1.5mm]
\textbf{Pista Conceptual:} Demuestre que el conjunto de discontinuidades de $\chi_\Omega$ es exactamente $\partial \Omega$. Luego use que una función acotada es integrable Riemann con discontinuidades en un conjunto cerrado si y sólo si dicho conjunto tiene contenido cero.
\end{tcolorbox}

\paragraph{Solución Paso a Paso:}\par
\begin{enumerate}[label=\textbf{\arabic*.}]
\item \textbf{Identificación topológica de discontinuidades:}
Para cualquier conjunto $\Omega \subset \mathbb{R}^2$, el espacio se particiona en tres conjuntos disjuntos abiertos y su frontera común:
\begin{equation*}
\mathbb{R}^2 = \Omega^\circ \cup (\mathbb{R}^2 \setminus \overline{\Omega}) \cup \partial \Omega.
\end{equation*}
- Para $p \in \Omega^\circ$: existe una bola $B(p, r) \subset \Omega$. En dicha bola, $\chi_\Omega(x) = 1$ para todo $x$, por lo que $\chi_\Omega$ es constante y continua en $p$.
- Para $p \in \mathbb{R}^2 \setminus \overline{\Omega}$: existe una bola donde $\chi_\Omega(x) = 0$ para todo $x$, luego es continua en $p$.
- Para $p \in \partial \Omega$: toda bola abierta $B(p, r)$ contiene puntos de $\Omega$ (donde $\chi_\Omega = 1$) y puntos de $\Omega^c$ (donde $\chi_\Omega = 0$).
Por tanto, $\chi_\Omega$ oscila entre 0 y 1 en cualquier vecindad de $p$, lo que implica que la oscilación es $\omega(\chi_\Omega, p) = 1 > 0$.
En consecuencia, el conjunto de puntos de discontinuidad de $\chi_\Omega$ es exactamente:
\begin{equation*}
D_{\chi_\Omega} = \partial \Omega.
\end{equation*}
\end{enumerate}

\begin{enumerate}[label=\textbf{\arabic*.}]
\item \textbf{Equivalencia de integrabilidad y contenido cero:}
- \textbf{Necesidad ($\implies$):}
Si $\chi_\Omega$ es integrable Riemann, dado $\varepsilon > 0$ existe una partición $P$ de $R$ tal que:
\begin{equation*}
U(\chi_\Omega, P) - L(\chi_\Omega, P) = \sum_{j} (M_j - m_j)\text{Área}(R_j) < \varepsilon.
\end{equation*}
Nótese que $M_j - m_j = 1$ si $R_j \cap \partial \Omega \neq \emptyset$, y $M_j - m_j = 0$ en caso contrario.
Por tanto, la familia finita de subrectángulos $\{R_j : R_j \cap \partial \Omega \neq \emptyset\}$ recubre a $\partial \Omega$ y su área total suma exactamente $U - L < \varepsilon$.
Esto prueba por definición que $\partial \Omega$ tiene contenido cero.
- \textbf{Suficiencia ($\impliedby$):}
Si $\partial \Omega$ tiene contenido cero, para todo $\varepsilon > 0$ existe una familia finita de rectángulos $\{J_1, \dots, J_m\}$ que cubre a $\partial \Omega$ con área total menor que $\varepsilon$.
Extendiendo los lados de estos rectángulos formamos una partición $P$ de $R$ tal que los subrectángulos que tocan $\partial \Omega$ tienen área total menor que $\varepsilon$.
Entonces $U(\chi_\Omega, P) - L(\chi_\Omega, P) < \varepsilon$, probando la integrabilidad de $\chi_\Omega$.
Por definición, un conjunto $\Omega$ con $\chi_\Omega$ integrable se denomina \textbf{medible según Jordan}.
\end{enumerate}

\begin{trapbox}
Confundir la frontera $\partial \Omega$ con el conjunto $\Omega$ en sí mismo.
\end{trapbox}

\vspace{0.5cm}
\hrule
\vspace{0.5cm}

\subsection*{Problema 4.15: Desigualdad de Cauchy-Schwarz para integrales dobles \hfill \color{primaryblue}\small[Nivel 5]}
\addcontentsline{toc}{subsection}{Problema 4.15: \texorpdfstring{Desigualdad de Cauchy-Schwarz para integrales dobles}{Desigualdad de Cauchy-Schwarz para integrales dobles} (Nivel 5)}

\begin{tcolorbox}[colback=white,colframe=gray!40,arc=1mm,boxrule=0.5pt,left=3mm,right=3mm,top=2mm,bottom=2mm]
\textbf{Enunciado:} Sean $f, g: R \to \mathbb{R}$ funciones integrables en el rectángulo $R$. (a) Demuestre rigurosamente la desigualdad de Cauchy-Schwarz:
\begin{equation*}
\left( \iint_R f(x,y) g(x,y) \, dA \right)^2 \le \left( \iint_R f(x,y)^2 \, dA \right) \left( \iint_R g(x,y)^2 \, dA \right).
\end{equation*}
(b) Para funciones continuas, determine la condición necesaria y suficiente para que se dé la igualdad estricta.
\end{tcolorbox}

\begin{tcolorbox}[colback=yellow!4,colframe=accentamber!70,arc=1mm,boxrule=0.5pt,left=3mm,right=3mm,top=1.5mm,bottom=1.5mm]
\textbf{Pista Conceptual:} Considere el polinomio cuadrático en $\lambda \in \mathbb{R}$: $P(\lambda) = \iint_R (f(x,y) + \lambda g(x,y))^2 dA \ge 0$. Analice el discriminante $\Delta \le 0$.
\end{tcolorbox}

\paragraph{Solución Paso a Paso:}\par
\begin{enumerate}[label=\textbf{\arabic*.}]
\item \textbf{Construcción de la forma cuadrática no negativa:}
Para cualquier parámetro real $\lambda \in \mathbb{R}$, consideremos la función:
\begin{equation*}
h_\lambda(x,y) = (f(x,y) + \lambda g(x,y))^2.
\end{equation*}
Al ser un cuadrado de números reales, $h_\lambda(x,y) \ge 0$ para todo $(x,y) \in R$.
Por la linealidad de la integral y la propiedad de monotonía, su integral sobre $R$ es no negativa:
\begin{equation*}
P(\lambda) = \iint_R (f(x,y) + \lambda g(x,y))^2 \, dA \ge 0, \quad \forall \lambda \in \mathbb{R}.
\end{equation*}
\end{enumerate}

\begin{enumerate}[label=\textbf{\arabic*.}]
\item \textbf{Desarrollo del polinomio cuadrático:}
Expandiendo el binomio al cuadrado:
\begin{equation*}
P(\lambda) = \iint_R \left[ f(x,y)^2 + 2\lambda f(x,y)g(x,y) + \lambda^2 g(x,y)^2 \right] dA
\end{equation*}
\begin{equation*}
= \lambda^2 \left(\iint_R g^2 \, dA\right) + 2\lambda \left(\iint_R fg \, dA\right) + \left(\iint_R f^2 \, dA\right).
\end{equation*}
Sean los coeficientes:
\begin{equation*}
A = \iint_R g^2 \, dA, \qquad B = \iint_R fg \, dA, \qquad C = \iint_R f^2 \, dA.
\end{equation*}
Entonces $P(\lambda) = A \lambda^2 + 2B \lambda + C \ge 0$ para todo $\lambda \in \mathbb{R}$.
\end{enumerate}

\begin{enumerate}[label=\textbf{\arabic*.}]
\item \textbf{Análisis del discriminante:}
- Si $A = 0$, como $g^2 \ge 0$, la integral es cero. Por la desigualdad de Cauchy-Schwarz elemental para sumas de Riemann, $B = 0$, y la desigualdad $0 \le 0$ se cumple trivialmente.
- Si $A > 0$: Un polinomio cuadrático con coeficiente principal positivo satisface $P(\lambda) \ge 0$ para todo $\lambda \in \mathbb{R}$ si y sólo si no tiene dos raíces reales distintas, es decir, su discriminante reducido debe ser menor o igual que cero:
\begin{equation*}
\Delta' = B^2 - AC \le 0 \iff B^2 \le AC.
\end{equation*}
Sustituyendo $A, B, C$:
\begin{equation*}
\left( \iint_R f(x,y) g(x,y) \, dA \right)^2 \le \left( \iint_R f(x,y)^2 \, dA \right) \left( \iint_R g(x,y)^2 \, dA \right).
\end{equation*}
\end{enumerate}

\begin{enumerate}[label=\textbf{\arabic*.}]
\item \textbf{Condición de igualdad para funciones continuas:}
La igualdad se cumple si y sólo si $\Delta' = 0$, lo que ocurre si y sólo si existe un $\lambda_0 \in \mathbb{R}$ tal que $P(\lambda_0) = 0$:
\begin{equation*}
\iint_R (f(x,y) + \lambda_0 g(x,y))^2 \, dA = 0.
\end{equation*}
Como $(f + \lambda_0 g)^2$ es una función continua no negativa cuya integral sobre $R$ es cero, debe anularse idénticamente en todo punto de $R$:
\begin{equation*}
(f(x,y) + \lambda_0 g(x,y))^2 = 0 \iff f(x,y) = -\lambda_0 g(x,y), \quad \forall (x,y) \in R.
\end{equation*}
Es decir, la igualdad se da si y sólo si $f$ y $g$ son linealmente dependientes (una es múltiplo escalar de la otra).
\end{enumerate}

\begin{trapbox}
Omitir el caso en que $\iint g^2 dA = 0$ al dividir entre el coeficiente principal.
\end{trapbox}

\vspace{0.5cm}
\hrule
\vspace{0.5cm}

\section{Teorema de Fubini y Teorema del Valor Medio}

\subsection*{Problema 5.0A: Integración iterada básica en dos órdenes distintos \hfill \color{teal}\small[Nivel 0]}
\addcontentsline{toc}{subsection}{Problema 5.0A: Integración iterada básica en dos órdenes distintos (Nivel 0)}

\begin{tcolorbox}[colback=white,colframe=gray!40,arc=1mm,boxrule=0.5pt,left=3mm,right=3mm,top=2mm,bottom=2mm]
\textbf{Enunciado:} Sea $f(x,y) = x \cdot y$ sobre el rectángulo $R = [0,2] \times [0,3]$.
\begin{enumerate}[label=\textbf{\alph*)}]
\item Calcule la integral doble integrando primero respecto a $y$ y luego respecto a $x$: $I_1 = \int_0^2 \left( \int_0^3 x y \, dy \right) dx$.
\item Calcule la misma integral invirtiendo el orden: $I_2 = \int_0^3 \left( \int_0^2 x y \, dx \right) dy$.
\item Verifique que $I_1 = I_2$.
\end{enumerate}
\end{tcolorbox}

\begin{tcolorbox}[colback=yellow!4,colframe=accentamber!70,arc=1mm,boxrule=0.5pt,left=3mm,right=3mm,top=1.5mm,bottom=1.5mm]
\textbf{Pista Conceptual:} Al integrar respecto a $y$, la variable $x$ se comporta como una constante que puede salir fuera de la integral interna.
\end{tcolorbox}

\paragraph{Solución Paso a Paso:}\par
\begin{enumerate}[label=\textbf{\arabic*.}]
\item \textbf{Primer orden (respecto a $y$ primero):}
\[ \int_0^3 x y \, dy = x \left[ \frac{y^2}{2} \right]_0^3 = \frac{9}{2}x. \]
Integrando respecto a $x$ de 0 a 2:
\[ I_1 = \int_0^2 \frac{9}{2}x \, dx = \frac{9}{2} \left[ \frac{x^2}{2} \right]_0^2 = \frac{9}{2}(2) = 9. \]
\item \textbf{Segundo orden (respecto a $x$ primero):}
\[ \int_0^2 x y \, dx = y \left[ \frac{x^2}{2} \right]_0^2 = 2y. \]
Integrando respecto a $y$ de 0 a 3:
\[ I_2 = \int_0^3 2y \, dy = \left[ y^2 \right]_0^3 = 9. \]
\item \textbf{Conclusión:}
\[ I_1 = 9 = I_2. \]
Ambos órdenes coinciden plenamente, ilustrando el Teorema de Fubini.
\end{enumerate}

\begin{trapbox}
Tratar a $x$ como variable al integrar respecto a $y$, o confundir los límites de integración de cada variable ($x$ va de 0 a 2; $y$ va de 0 a 3).
\end{trapbox}

\vspace{0.5cm}
\hrule
\vspace{0.5cm}
\subsection*{Problema 5.0B: Factorización de integrales con variables desacopladas \hfill \color{teal}\small[Nivel 0]}
\addcontentsline{toc}{subsection}{Problema 5.0B: Factorización de integrales con variables desacopladas (Nivel 0)}

\begin{tcolorbox}[colback=white,colframe=gray!40,arc=1mm,boxrule=0.5pt,left=3mm,right=3mm,top=2mm,bottom=2mm]
\textbf{Enunciado:} Demuestre que si $f(x,y) = g(x) \cdot h(y)$ sobre un rectángulo $R = [a,b] \times [c,d]$, entonces:
\[ \iint_R g(x)h(y) \, dA = \left( \int_a^b g(x) \, dx \right) \cdot \left( \int_c^d h(y) \, dy \right). \]
Utilice este resultado para calcular $\iint_R \sin(x) e^y \, dA$ sobre $R = [0, \pi] \times [0, 1]$.
\end{tcolorbox}

\begin{tcolorbox}[colback=yellow!4,colframe=accentamber!70,arc=1mm,boxrule=0.5pt,left=3mm,right=3mm,top=1.5mm,bottom=1.5mm]
\textbf{Pista Conceptual:} Extraiga $g(x)$ fuera de la integral de $y$, y luego observe que la integral de $h(y)$ es un número constante que puede salir fuera de la integral de $x$.
\end{tcolorbox}

\paragraph{Solución Paso a Paso:}\par
\begin{enumerate}[label=\textbf{\arabic*.}]
\item \textbf{Demostración de la propiedad de factorización:}
Por el Teorema de Fubini:
\[ \iint_R g(x)h(y) \, dA = \int_a^b \left( \int_c^d g(x)h(y) \, dy \right) dx = \int_a^b g(x) \left( \int_c^d h(y) \, dy \right) dx. \]
La integral interior $K = \int_c^d h(y) dy$ es constante respecto a $x$, por lo que sale de la integral exterior:
\[ = \left( \int_c^d h(y) \, dy \right) \cdot \int_a^b g(x) \, dx = \left( \int_a^b g(x) \, dx \right) \left( \int_c^d h(y) \, dy \right). \]
\item \textbf{Aplicación al ejemplo:}
\[ \iint_R \sin(x) e^y \, dA = \left( \int_0^\pi \sin x \, dx \right) \cdot \left( \int_0^1 e^y \, dy \right) = [-\cos x]_0^\pi \cdot [e^y]_0^1 = 2(e - 1) = 2e - 2. \]
\end{enumerate}

\begin{trapbox}
Intentar aplicar esta fórmula de factorización cuando el dominio no es un rectángulo (por ejemplo si los límites de $y$ dependen de $x$), o cuando la función no es un producto factorizable (como $\sin(x+y)$).
\end{trapbox}

\vspace{0.5cm}
\hrule
\vspace{0.5cm}
\subsection*{Problema 5.0C: Verificación explícita del Teorema del Valor Medio \hfill \color{teal}\small[Nivel 0]}
\addcontentsline{toc}{subsection}{Problema 5.0C: Verificación explícita del Teorema del Valor Medio (Nivel 0)}

\begin{tcolorbox}[colback=white,colframe=gray!40,arc=1mm,boxrule=0.5pt,left=3mm,right=3mm,top=2mm,bottom=2mm]
\textbf{Enunciado:} Sea $f(x,y) = x$ sobre el rectángulo $R = [0,2] \times [0,4]$.
\begin{enumerate}[label=\textbf{\alph*)}]
\item Calcule el área de $R$ y la integral $\iint_R x \, dA$.
\item Calcule el valor promedio $\bar{f} = \frac{1}{\text{Área}(R)} \iint_R f(x,y) \, dA$.
\item Halle explícitamente al menos un punto $(x_0, y_0) \in R$ tal que $f(x_0, y_0) = \bar{f}$.
\end{enumerate}
\end{tcolorbox}

\begin{tcolorbox}[colback=yellow!4,colframe=accentamber!70,arc=1mm,boxrule=0.5pt,left=3mm,right=3mm,top=1.5mm,bottom=1.5mm]
\textbf{Pista Conceptual:} El Teorema del Valor Medio garantiza que si $f$ es continua y $R$ es conexo, el valor promedio se alcanza en algún punto interior o de borde del dominio.
\end{tcolorbox}

\paragraph{Solución Paso a Paso:}\par
\begin{enumerate}[label=\textbf{\arabic*.}]
\item \textbf{Cálculo del área:}
\[ \text{Área}(R) = (2 - 0)(4 - 0) = 8. \]
\item \textbf{Cálculo de la integral doble:}
\[ \iint_R x \, dA = \int_0^4 \left( \int_0^2 x \, dx \right) dy = \int_0^4 2 \, dy = 8. \]
\item \textbf{Valor promedio:}
\[ \bar{f} = \frac{1}{\text{Área}(R)} \iint_R x \, dA = \frac{1}{8}(8) = 1. \]
\item \textbf{Búsqueda del punto $(x_0, y_0) \in R$:}
Queremos $f(x_0, y_0) = 1$. Como $f(x,y) = x$, basta tomar $x_0 = 1$ con cualquier $y_0 \in [0,4]$. Por ejemplo, $(x_0, y_0) = (1, 2) \in R \implies f(1,2) = 1 = \bar{f}$, comprobando el Teorema del Valor Medio.
\end{enumerate}

\begin{trapbox}
Olvidar dividir entre el área del dominio al calcular el valor promedio $\bar{f}$, o buscar un punto fuera del rectángulo $R$.
\end{trapbox}

\vspace{0.5cm}
\hrule
\vspace{0.5cm}


\subsection*{Problema 5.1: Inversión del orden de integración con Fubini (Examen UNI 2026-I) \hfill \color{primaryblue}\small[Nivel 1]}
\addcontentsline{toc}{subsection}{Problema 5.1: \texorpdfstring{Inversión del orden de integración con Fubini (Examen UNI 2026-I)}{Inversión del orden de integración con Fubini (Examen UNI 2026-I)} (Nivel 1)}

\begin{tcolorbox}[colback=white,colframe=gray!40,arc=1mm,boxrule=0.5pt,left=3mm,right=3mm,top=2mm,bottom=2mm]
\textbf{Enunciado:} Considere la integral iterada:
\begin{equation*}
I = \int_1^3 \int_0^1 y e^{xy} \, dy \, dx.
\end{equation*}
(a) Invierta el orden de integración aplicando el Teorema de Fubini. (b) Calcule el valor exacto de la integral y verifique que $I = \frac{e^3}{3} - e + \frac{2}{3}$.
\end{tcolorbox}

\begin{tcolorbox}[colback=yellow!4,colframe=accentamber!70,arc=1mm,boxrule=0.5pt,left=3mm,right=3mm,top=1.5mm,bottom=1.5mm]
\textbf{Pista Conceptual:} La función $f(x,y) = y e^{xy}$ es continua en el rectángulo $[1,3] \times [0,1]$. Integre primero respecto a $x$, donde $y$ actúa como constante, facilitando la primitiva inmediata de $y e^{xy}$.
\end{tcolorbox}

\paragraph{Solución Paso a Paso:}\par
\begin{enumerate}[label=\textbf{\arabic*.}]
\item \textbf{Aplicabilidad del Teorema de Fubini:}
El dominio de integración es el rectángulo $R = [1, 3] \times [0, 1]$.
La función integrando $f(x,y) = y e^{xy}$ es continua en todo $\mathbb{R}^2$, luego es continua en el compacto $R$.
Por el Teorema de Fubini en rectángulos, la integral doble existe y coincide con las dos integrales iteradas:
\begin{equation*}
\int_1^3 \int_0^1 y e^{xy} \, dy \, dx = \int_0^1 \int_1^3 y e^{xy} \, dx \, dy.
\end{equation*}
\end{enumerate}

\begin{enumerate}[label=\textbf{\arabic*.}]
\item \textbf{Cálculo de la integral interna respecto a $x$:}
Para $y$ fija, la derivada de $e^{xy}$ respecto a $x$ es precisamente $\frac{\partial}{\partial x}(e^{xy}) = y e^{xy}$.
Por tanto, la antiderivada inmediata es $e^{xy}$:
\begin{equation*}
\int_1^3 y e^{xy} \, dx = \left[ e^{xy} \right]_{x=1}^{x=3} = e^{3y} - e^y.
\end{equation*}
\end{enumerate}

\begin{enumerate}[label=\textbf{\arabic*.}]
\item \textbf{Cálculo de la integral externa respecto a $y$:}
Integrando de $0$ a $1$:
\begin{equation*}
I = \int_0^1 (e^{3y} - e^y) \, dy = \left[ \frac{e^{3y}}{3} - e^y \right]_0^1
\end{equation*}
Evaluando en los límites superior e inferior:
\begin{equation*}
= \left( \frac{e^3}{3} - e \right) - \left( \frac{e^0}{3} - e^0 \right) = \frac{e^3}{3} - e - \left(\frac{1}{3} - 1\right)
\end{equation*}
\begin{equation*}
= \frac{e^3}{3} - e - \left(-\frac{2}{3}\right) = \frac{e^3}{3} - e + \frac{2}{3}.
\end{equation*}
Se demuestra con total precisión la igualdad propuesta en el examen oficial de la UNI.
\end{enumerate}

\begin{trapbox}
Intentar integrar primero por partes respecto a $y$ en lugar de invertir el orden respecto a $x$, lo que complica innecesariamente los cálculos.
\end{trapbox}

\vspace{0.5cm}
\hrule
\vspace{0.5cm}

\subsection*{Problema 5.2: Integral triple en un paralelepípedo tridimensional \hfill \color{primaryblue}\small[Nivel 1]}
\addcontentsline{toc}{subsection}{Problema 5.2: \texorpdfstring{Integral triple en un paralelepípedo tridimensional}{Integral triple en un paralelepípedo tridimensional} (Nivel 1)}

\begin{tcolorbox}[colback=white,colframe=gray!40,arc=1mm,boxrule=0.5pt,left=3mm,right=3mm,top=2mm,bottom=2mm]
\textbf{Enunciado:} Sea el bloque $B = [0, 1] \times [0, 2] \times [0, 3] \subset \mathbb{R}^3$. Calcule la integral triple:
\begin{equation*}
I = \iiint_B (x + 2y + 3z) \, dV
\end{equation*}
mediante integrales iteradas aplicando el Teorema de Fubini.
\end{tcolorbox}

\begin{tcolorbox}[colback=yellow!4,colframe=accentamber!70,arc=1mm,boxrule=0.5pt,left=3mm,right=3mm,top=1.5mm,bottom=1.5mm]
\textbf{Pista Conceptual:} Use la linealidad para separar en tres integrales simples, multiplicando cada término por el área de las dos dimensiones restantes.
\end{tcolorbox}

\paragraph{Solución Paso a Paso:}\par
\begin{enumerate}[label=\textbf{\arabic*.}]
\item \textbf{Descomposición por linealidad:}
Por el Teorema de Fubini para paralelepípedos en $\mathbb{R}^3$:
\begin{equation*}
\iiint_B (x + 2y + 3z) \, dV = \iiint_B x \, dV + 2 \iiint_B y \, dV + 3 \iiint_B z \, dV.
\end{equation*}
\end{enumerate}

\begin{enumerate}[label=\textbf{\arabic*.}]
\item \textbf{Cálculo de cada término independiente:}
- Para $\iiint_B x \, dV$:
\begin{equation*}
\left(\int_0^1 x \, dx\right) \left(\int_0^2 1 \, dy\right) \left(\int_0^3 1 \, dz\right) = \left[\frac{x^2}{2}\right]_0^1 \times 2 \times 3 = \frac{1}{2} \times 6 = 3.
\end{equation*}
- Para $2 \iiint_B y \, dV$:
\begin{equation*}
2 \left(\int_0^1 1 \, dx\right) \left(\int_0^2 y \, dy\right) \left(\int_0^3 1 \, dz\right) = 2 \times 1 \times \left[\frac{y^2}{2}\right]_0^2 \times 3 = 2 \times 1 \times 2 \times 3 = 12.
\end{equation*}
- Para $3 \iiint_B z \, dV$:
\begin{equation*}
3 \left(\int_0^1 1 \, dx\right) \left(\int_0^2 1 \, dy\right) \left(\int_0^3 z \, dz\right) = 3 \times 1 \times 2 \times \left[\frac{z^2}{2}\right]_0^3 = 6 \times \frac{9}{2} = 27.
\end{equation*}
\end{enumerate}

\begin{enumerate}[label=\textbf{\arabic*.}]
\item \textbf{Suma total:}
\begin{equation*}
I = 3 + 12 + 27 = 42.
\end{equation*}
\end{enumerate}

\begin{trapbox}
Olvidar los factores multiplicativos constantes de las longitudes de los otros intervalos al integrar funciones independientes de esas variables.
\end{trapbox}

\vspace{0.5cm}
\hrule
\vspace{0.5cm}

\subsection*{Problema 5.3: Integral doble trigonométrica en rectángulo \hfill \color{primaryblue}\small[Nivel 1]}
\addcontentsline{toc}{subsection}{Problema 5.3: \texorpdfstring{Integral doble trigonométrica en rectángulo}{Integral doble trigonométrica en rectángulo} (Nivel 1)}

\begin{tcolorbox}[colback=white,colframe=gray!40,arc=1mm,boxrule=0.5pt,left=3mm,right=3mm,top=2mm,bottom=2mm]
\textbf{Enunciado:} Calcule la integral doble:
\begin{equation*}
\iint_R \sin(x + y) \, dA, \qquad \text{donde } R = [0, \pi] \times [0, \pi/2].
\end{equation*}
\end{tcolorbox}

\begin{tcolorbox}[colback=yellow!4,colframe=accentamber!70,arc=1mm,boxrule=0.5pt,left=3mm,right=3mm,top=1.5mm,bottom=1.5mm]
\textbf{Pista Conceptual:} Aplique Fubini directo o use la identidad $\sin(x+y) = \sin x \cos y + \cos x \sin y$.
\end{tcolorbox}

\paragraph{Solución Paso a Paso:}\par
\begin{enumerate}[label=\textbf{\arabic*.}]
\item \textbf{Planteamiento de la integral iterada:}
Por Fubini:
\begin{equation*}
I = \int_0^\pi \left( \int_0^{\pi/2} \sin(x + y) \, dy \right) dx.
\end{equation*}
\end{enumerate}

\begin{enumerate}[label=\textbf{\arabic*.}]
\item \textbf{Integral interna respecto a $y$:}
\begin{equation*}
\int_0^{\pi/2} \sin(x + y) \, dy = [-\cos(x + y)]_{y=0}^{y=\pi/2} = -\cos(x + \pi/2) - (-\cos(x)) = \sin(x) + \cos(x),
\end{equation*}
pues $\cos(x + \pi/2) = -\sin(x)$.
\end{enumerate}

\begin{enumerate}[label=\textbf{\arabic*.}]
\item \textbf{Integral externa respecto a $x$:}
\begin{equation*}
\int_0^\pi (\sin x + \cos x) \, dx = [-\cos x + \sin x]_0^\pi
\end{equation*}
\begin{equation*}
= (-\cos\pi + \sin\pi) - (-\cos 0 + \sin 0) = (-(-1) + 0) - (-1 + 0) = 1 - (-1) = 2.
\end{equation*}
Por tanto, $\iint_R \sin(x + y) \, dA = 2$.
\end{enumerate}

\begin{trapbox}
Confundir los signos de las antiderivadas trigonométricas: $\int \sin u du = -\cos u$.
\end{trapbox}

\vspace{0.5cm}
\hrule
\vspace{0.5cm}

\subsection*{Problema 5.4: Teorema de factorización para funciones separables \hfill \color{primaryblue}\small[Nivel 2]}
\addcontentsline{toc}{subsection}{Problema 5.4: \texorpdfstring{Teorema de factorización para funciones separables}{Teorema de factorización para funciones separables} (Nivel 2)}

\begin{tcolorbox}[colback=white,colframe=gray!40,arc=1mm,boxrule=0.5pt,left=3mm,right=3mm,top=2mm,bottom=2mm]
\textbf{Enunciado:} Sea $R = [a,b] \times [c,d]$. Demuestre que si $f(x,y) = g(x) h(y)$, donde $g: [a,b] \to \mathbb{R}$ y $h: [c,d] \to \mathbb{R}$ son continuas, entonces:
\begin{equation*}
\iint_R f(x,y) \, dA = \left( \int_a^b g(x) \, dx \right) \left( \int_c^d h(y) \, dy \right).
\end{equation*}
\end{tcolorbox}

\begin{tcolorbox}[colback=yellow!4,colframe=accentamber!70,arc=1mm,boxrule=0.5pt,left=3mm,right=3mm,top=1.5mm,bottom=1.5mm]
\textbf{Pista Conceptual:} Aplique Fubini sacando $g(x)$ fuera de la integral en $y$, y observe que la integral en $y$ es un número real constante que luego sale de la integral en $x$.
\end{tcolorbox}

\paragraph{Solución Paso a Paso:}\par
\begin{enumerate}[label=\textbf{\arabic*.}]
\item \textbf{Continuidad del producto:}
Como $g$ y $h$ son continuas, su producto $f(x,y) = g(x)h(y)$ es continuo en el rectángulo compacto $R$.
Por el Teorema de Fubini:
\begin{equation*}
\iint_R f(x,y) \, dA = \int_a^b \left( \int_c^d g(x) h(y) \, dy \right) dx.
\end{equation*}
\end{enumerate}

\begin{enumerate}[label=\textbf{\arabic*.}]
\item \textbf{Factorización de la variable independiente:}
En la integral interna con respecto a $y$, la variable $x$ se mantiene constante, por lo que $g(x)$ actúa como factor constante frente a la integración en $y$:
\begin{equation*}
\int_c^d g(x) h(y) \, dy = g(x) \int_c^d h(y) \, dy.
\end{equation*}
\end{enumerate}

\begin{enumerate}[label=\textbf{\arabic*.}]
\item \textbf{Factorización de la constante escalar:}
Definamos el número real $K = \int_c^d h(y) dy \in \mathbb{R}$.
Sustituyendo en la integral externa:
\begin{equation*}
\int_a^b \left( g(x) \cdot K \right) dx = K \int_a^b g(x) \, dx = \left( \int_c^d h(y) \, dy \right) \left( \int_a^b g(x) \, dx \right).
\end{equation*}
Por conmutatividad del producto de números reales:
\begin{equation*}
\iint_R f(x,y) \, dA = \left( \int_a^b g(x) \, dx \right) \left( \int_c^d h(y) \, dy \right).
\end{equation*}
\end{enumerate}

\begin{trapbox}
Aplicar esta propiedad a regiones no rectangulares donde los límites de integración de una variable dependen de la otra.
\end{trapbox}

\vspace{0.5cm}
\hrule
\vspace{0.5cm}

\subsection*{Problema 5.5: Evaluación por Fubini de función racional en cuadrado \hfill \color{primaryblue}\small[Nivel 2]}
\addcontentsline{toc}{subsection}{Problema 5.5: \texorpdfstring{Evaluación por Fubini de función racional en cuadrado}{Evaluación por Fubini de función racional en cuadrado} (Nivel 2)}

\begin{tcolorbox}[colback=white,colframe=gray!40,arc=1mm,boxrule=0.5pt,left=3mm,right=3mm,top=2mm,bottom=2mm]
\textbf{Enunciado:} Calcule la integral doble:
\begin{equation*}
\iint_{[0,1] \times [0,1]} \frac{1}{(x + y + 1)^2} \, dA.
\end{equation*}
\end{tcolorbox}

\begin{tcolorbox}[colback=yellow!4,colframe=accentamber!70,arc=1mm,boxrule=0.5pt,left=3mm,right=3mm,top=1.5mm,bottom=1.5mm]
\textbf{Pista Conceptual:} Integre primero con respecto a $x$ usando la primitiva inmediata $-\frac{1}{x+y+1}$ y luego integre con respecto a $y$ usando logaritmos.
\end{tcolorbox}

\paragraph{Solución Paso a Paso:}\par
\begin{enumerate}[label=\textbf{\arabic*.}]
\item \textbf{Planteamiento de Fubini:}
\begin{equation*}
I = \int_0^1 \int_0^1 \frac{1}{(x + y + 1)^2} \, dx \, dy.
\end{equation*}
\end{enumerate}

\begin{enumerate}[label=\textbf{\arabic*.}]
\item \textbf{Integral interna en $x$:}
\begin{equation*}
\int_0^1 \frac{dx}{(x + y + 1)^2} = \left[ -\frac{1}{x + y + 1} \right]_{x=0}^{x=1} = -\frac{1}{1 + y + 1} - \left(-\frac{1}{0 + y + 1}\right) = \frac{1}{y + 1} - \frac{1}{y + 2}.
\end{equation*}
\end{enumerate}

\begin{enumerate}[label=\textbf{\arabic*.}]
\item \textbf{Integral externa en $y$:}
\begin{equation*}
\int_0^1 \left( \frac{1}{y + 1} - \frac{1}{y + 2} \right) dy = [\ln(y + 1) - \ln(y + 2)]_0^1 = \left[ \ln\left(\frac{y+1}{y+2}\right) \right]_0^1
\end{equation*}
\begin{equation*}
= \ln\left(\frac{2}{3}\right) - \ln\left(\frac{1}{2}\right) = \ln\left(\frac{2/3}{1/2}\right) = \ln\left(\frac{4}{3}\right).
\end{equation*}
Por tanto, la integral vale $\ln(4/3) \approx 0.2877$.
\end{enumerate}

\begin{trapbox}
Cometer errores de signos en la integración de potencias negativas: $\int u^{-2} du = -u^{-1}$.
\end{trapbox}

\vspace{0.5cm}
\hrule
\vspace{0.5cm}

\subsection*{Problema 5.6: Integral iterada tridimensional con exponencial \hfill \color{primaryblue}\small[Nivel 2]}
\addcontentsline{toc}{subsection}{Problema 5.6: \texorpdfstring{Integral iterada tridimensional con exponencial}{Integral iterada tridimensional con exponencial} (Nivel 2)}

\begin{tcolorbox}[colback=white,colframe=gray!40,arc=1mm,boxrule=0.5pt,left=3mm,right=3mm,top=2mm,bottom=2mm]
\textbf{Enunciado:} Calcule la integral triple en el cubo unitario $B = [0, 1]^3$:
\begin{equation*}
\iiint_B x y z \, e^{x^2 + y^2 + z^2} \, dx \, dy \, dz.
\end{equation*}
\end{tcolorbox}

\begin{tcolorbox}[colback=yellow!4,colframe=accentamber!70,arc=1mm,boxrule=0.5pt,left=3mm,right=3mm,top=1.5mm,bottom=1.5mm]
\textbf{Pista Conceptual:} Observe que la función es separable: $x y z e^{x^2+y^2+z^2} = (x e^{x^2})(y e^{y^2})(z e^{z^2})$.
\end{tcolorbox}

\paragraph{Solución Paso a Paso:}\par
\begin{enumerate}[label=\textbf{\arabic*.}]
\item \textbf{Separación de variables en tres factores:}
Por las leyes de los exponentes:
\begin{equation*}
f(x,y,z) = (x e^{x^2}) \cdot (y e^{y^2}) \cdot (z e^{z^2}).
\end{equation*}
Por la extensión de la propiedad separable a paralelepípedos en $\mathbb{R}^3$:
\begin{equation*}
\iiint_B f(x,y,z) \, dV = \left( \int_0^1 x e^{x^2} dx \right) \left( \int_0^1 y e^{y^2} dy \right) \left( \int_0^1 z e^{z^2} dz \right).
\end{equation*}
\end{enumerate}

\begin{enumerate}[label=\textbf{\arabic*.}]
\item \textbf{Cálculo de la integral unidimensional común:}
Sea $J = \int_0^1 t e^{t^2} dt$. Con el cambio de variable $u = t^2 \implies du = 2t dt$:
\begin{equation*}
J = \frac{1}{2} \int_0^1 e^u du = \frac{1}{2} [e^u]_0^1 = \frac{e - 1}{2}.
\end{equation*}
\end{enumerate}

\begin{enumerate}[label=\textbf{\arabic*.}]
\item \textbf{Cálculo de la integral triple:}
\begin{equation*}
I = J \cdot J \cdot J = J^3 = \left(\frac{e - 1}{2}\right)^3 = \frac{(e - 1)^3}{8}.
\end{equation*}
\end{enumerate}

\begin{trapbox}
Olvidar el factor $1/2$ que surge del diferencial $du = 2t dt$ al integrar $t e^{t^2}$.
\end{trapbox}

\vspace{0.5cm}
\hrule
\vspace{0.5cm}

\subsection*{Problema 5.7: Demostración constructiva del Teorema del Valor Medio (Examen UNI 2026-I) \hfill \color{primaryblue}\small[Nivel 3]}
\addcontentsline{toc}{subsection}{Problema 5.7: \texorpdfstring{Demostración constructiva del Teorema del Valor Medio (Examen UNI 2026-I)}{Demostración constructiva del Teorema del Valor Medio (Examen UNI 2026-I)} (Nivel 3)}

\begin{tcolorbox}[colback=white,colframe=gray!40,arc=1mm,boxrule=0.5pt,left=3mm,right=3mm,top=2mm,bottom=2mm]
\textbf{Enunciado:} Sea $f: [2, 5] \times [3, 7] \to \mathbb{R}$ una función continua. Demuestre rigurosamente que existe un punto $(x_0, y_0) \in [2, 5] \times [3, 7]$ tal que:
\begin{equation*}
\int_2^5 \int_3^7 f(x,y) \, dy \, dx = 12 \, f(x_0, y_0).
\end{equation*}
\textit{Nota oficial del examen UNI:} Lo que se pide probar en este problema es conocido como el Teorema del Valor Medio para integrales. Por ende, NO es válido usar el teorema mencionado para demostrar el mismo teorema; la demostración debe ser autocontenida a partir de propiedades fundamentales de compacidad y conexidad.
\end{tcolorbox}

\begin{tcolorbox}[colback=yellow!4,colframe=accentamber!70,arc=1mm,boxrule=0.5pt,left=3mm,right=3mm,top=1.5mm,bottom=1.5mm]
\textbf{Pista Conceptual:} Como el rectángulo $R$ es compacto y conexo, su imagen $f(R)$ es un intervalo cerrado $[m, M]$. Use la acotación de la integral para probar que el promedio $\mu = \frac{1}{12}\iint_R f dA$ satisface $m \le \mu \le M$, y aplique el Teorema del Valor Intermedio.
\end{tcolorbox}

\paragraph{Solución Paso a Paso:}\par
\begin{enumerate}[label=\textbf{\arabic*.}]
\item \textbf{Propiedades topológicas del dominio:}
El dominio $R = [2, 5] \times [3, 7] \subset \mathbb{R}^2$ es el producto cartesiano de dos intervalos cerrados y acotados.
Por el Teorema de Heine-Borel, $R$ es un conjunto compacto.
Además, como $[2,5]$ y $[3,7]$ son conexos, el producto cartesiano $R$ es un conjunto conexo (incluso conexo por caminos).
\end{enumerate}

\begin{enumerate}[label=\textbf{\arabic*.}]
\item \textbf{Imagen continua de un compacto y conexo:}
Dado que $f: R \to \mathbb{R}$ es continua:
- Por el Teorema de Weierstrass (compacidad), $f$ alcanza su mínimo absoluto $m$ y su máximo absoluto $M$ en puntos de $R$:
\begin{equation*}
m = \min_{(x,y)\in R} f(x,y) = f(p_{\min}), \qquad M = \max_{(x,y)\in R} f(x,y) = f(p_{\max}).
\end{equation*}
- Por el Teorema del Valor Intermedio (conexidad), la imagen continua de un conjunto conexo es conexa en $\mathbb{R}$, es decir, es un intervalo:
\begin{equation*}
\text{Rang}(f) = f(R) = [m, M].
\end{equation*}
Esto significa que para cualquier valor $\mu \in [m, M]$, existe algún $(x_0, y_0) \in R$ tal que $f(x_0, y_0) = \mu$.
\end{enumerate}

\begin{enumerate}[label=\textbf{\arabic*.}]
\item \textbf{Cálculo del área del rectángulo:}
\begin{equation*}
\text{Área}(R) = (5 - 2) \times (7 - 3) = 3 \times 4 = 12.
\end{equation*}
\end{enumerate}

\begin{enumerate}[label=\textbf{\arabic*.}]
\item \textbf{Acotación fundamental de la integral:}
Como $m \le f(x,y) \le M$ para todo $(x,y) \in R$, por la propiedad de monotonía de la integral de Riemann:
\begin{equation*}
\iint_R m \, dA \le \iint_R f(x,y) \, dA \le \iint_R M \, dA
\end{equation*}
\begin{equation*}
\implies m \iint_R 1 \, dA \le \iint_R f(x,y) \, dA \le M \iint_R 1 \, dA
\end{equation*}
\begin{equation*}
\implies m \cdot \text{Área}(R) \le \iint_R f(x,y) \, dA \le M \cdot \text{Área}(R)
\end{equation*}
\begin{equation*}
\implies 12 m \le \int_2^5 \int_3^7 f(x,y) \, dy \, dx \le 12 M.
\end{equation*}
\end{enumerate}

\begin{enumerate}[label=\textbf{\arabic*.}]
\item \textbf{Paso al promedio y existencia del punto intermedio:}
Dividiendo toda la desigualdad entre el área positiva $12 > 0$:
\begin{equation*}
m \le \frac{1}{12} \int_2^5 \int_3^7 f(x,y) \, dy \, dx \le M.
\end{equation*}
Sea el número real:
\begin{equation*}
\mu = \frac{1}{12} \int_2^5 \int_3^7 f(x,y) \, dy \, dx.
\end{equation*}
Como hemos probado que $m \le \mu \le M$ y que $\text{Rang}(f) = [m, M]$, el valor $\mu$ pertenece necesariamente a la imagen de $f$.
Por el Teorema del Valor Intermedio, existe al menos un punto $(x_0, y_0) \in [2, 5] \times [3, 7]$ tal que:
\begin{equation*}
f(x_0, y_0) = \mu = \frac{1}{12} \int_2^5 \int_3^7 f(x,y) \, dy \, dx.
\end{equation*}
Multiplicando por 12 a ambos miembros:
\begin{equation*}
\int_2^5 \int_3^7 f(x,y) \, dy \, dx = 12 \, f(x_0, y_0).
\end{equation*}
Esto concluye la demostración rigurosa exigida en el examen oficial de la UNI.
\end{enumerate}

\begin{trapbox}
Invocar directamente el Teorema del Valor Medio para justificar el problema, lo que anula automáticamente la calificación según la advertencia expresa del enunciado.
\end{trapbox}

\vspace{0.5cm}
\hrule
\vspace{0.5cm}

\subsection*{Problema 5.8: Integral con valor absoluto y regiones de signos (Examen UNI 2026-I) \hfill \color{primaryblue}\small[Nivel 3]}
\addcontentsline{toc}{subsection}{Problema 5.8: \texorpdfstring{Integral con valor absoluto y regiones de signos (Examen UNI 2026-I)}{Integral con valor absoluto y regiones de signos (Examen UNI 2026-I)} (Nivel 3)}

\begin{tcolorbox}[colback=white,colframe=gray!40,arc=1mm,boxrule=0.5pt,left=3mm,right=3mm,top=2mm,bottom=2mm]
\textbf{Enunciado:} Considere la integral doble:
\begin{equation*}
I = \iint_R \frac{|x + my| + x}{|mx - y| - y} \, dA,
\end{equation*}
donde $m > 0$ es una constante real fija y el dominio es el rectángulo:
\begin{equation*}
R = \{(x,y) \in \mathbb{R}^2 : -2m \le x \le -m, \ -m^2 \le y \le 1\}.
\end{equation*}
(a) Esboce geométricamente el rectángulo $R$ y las rectas $L_1: x + my = 0$ y $L_2: mx - y = 0$.
(b) Determine el signo de las expresiones $x + my$ y $mx - y$ para todos los puntos $(x,y) \in R$.
(c) Justifique rigurosamente la existencia de la integral doble sin calcularla.
(d) Calcule el valor exacto de la integral y demuestre que $I = (m^4 - 1)\ln\sqrt{2}$.
\end{tcolorbox}

\begin{tcolorbox}[colback=yellow!4,colframe=accentamber!70,arc=1mm,boxrule=0.5pt,left=3mm,right=3mm,top=1.5mm,bottom=1.5mm]
\textbf{Pista Conceptual:} Analice los signos de $x+my$ y $mx-y$ en $R$. Como $x < 0$ y $y \le 1$, demuestre que ambas cantidades son no positivas. Simplifique el integrando cancelando términos y aplique Fubini.
\end{tcolorbox}

\paragraph{Solución Paso a Paso:}\par
\begin{enumerate}[label=\textbf{\arabic*.}]
\item \textbf{Análisis de signos en el rectángulo $R$:}
El rectángulo $R$ está delimitado por:
\begin{equation*}
-2m \le x \le -m < 0, \qquad -m^2 \le y \le 1.
\end{equation*}
- Para la expresión $x + my$:
Como $x \le -m$ y $y \le 1$ con $m > 0$:
\begin{equation*}
x + my \le -m + m(1) = -m + m = 0 \implies x + my \le 0, \quad \forall (x,y) \in R.
\end{equation*}
Por tanto: $|x + my| = -(x + my) = -x - my$.
- Para la expresión $mx - y$:
Como $x \le -m < 0 \implies mx \le -m^2$, y dado que $y \ge -m^2 \implies -y \le m^2$:
\begin{equation*}
mx - y \le -m^2 + m^2 = 0 \implies mx - y \le 0, \quad \forall (x,y) \in R.
\end{equation*}
Por tanto: $|mx - y| = -(mx - y) = y - mx$.
\end{enumerate}

\begin{enumerate}[label=\textbf{\arabic*.}]
\item \textbf{Simplificación del integrando:}
Sustituyendo los valores absolutos en la regla de correspondencia de $f(x,y)$:
\begin{equation*}
f(x,y) = \frac{(-x - my) + x}{(y - mx) - y} = \frac{-my}{-mx} = \frac{y}{x}.
\end{equation*}
Nótese que en el rectángulo $R$, $x \in [-2m, -m]$, por lo que $x$ nunca se anula ($x \neq 0$).
\end{enumerate}

\begin{enumerate}[label=\textbf{\arabic*.}]
\item \textbf{Justificación de la existencia:}
La función $f(x,y) = \frac{y}{x}$ es continua en todo el semiplano $x < 0$.
Como el rectángulo $R$ es compacto y está completamente contenido en el semiplano izquierdo $x \le -m < 0$, $f$ es continua en $R$.
Al ser continua sobre un rectángulo compacto, por el teorema fundamental de integrabilidad, $f$ es integrable sobre $R$, lo que garantiza la existencia de la integral sin necesidad de calcularla previamente.
\end{enumerate}

\begin{enumerate}[label=\textbf{\arabic*.}]
\item \textbf{Cálculo de la integral mediante Fubini:}
Por el Teorema de Fubini para funciones separables:
\begin{equation*}
I = \iint_R \frac{y}{x} \, dA = \left( \int_{-2m}^{-m} \frac{1}{x} \, dx \right) \left( \int_{-m^2}^1 y \, dy \right).
\end{equation*}
- Cálculo de la integral en $x$:
\begin{equation*}
\int_{-2m}^{-m} \frac{1}{x} \, dx = [\ln|x|]_{-2m}^{-m} = \ln|-m| - \ln|-2m| = \ln(m) - \ln(2m) = \ln\left(\frac{m}{2m}\right) = \ln\left(\frac{1}{2}\right) = -\ln(2).
\end{equation*}
Nótese que $-\ln 2 = \ln(2^{-1}) = \ln(1/2) = \ln(1/\sqrt{2}^2) = -2\ln\sqrt{2}$.
- Cálculo de la integral en $y$:
\begin{equation*}
\int_{-m^2}^1 y \, dy = \left[ \frac{y^2}{2} \right]_{-m^2}^1 = \frac{1^2}{2} - \frac{(-m^2)^2}{2} = \frac{1 - m^4}{2}.
\end{equation*}
- Multiplicación de ambos resultados:
\begin{equation*}
I = (-\ln 2) \left(\frac{1 - m^4}{2}\right) = \ln(2) \left(\frac{m^4 - 1}{2}\right) = (m^4 - 1) \frac{\ln 2}{2} = (m^4 - 1) \ln(2^{1/2}) = (m^4 - 1) \ln\sqrt{2}.
\end{equation*}
Queda demostrada rigurosamente la fórmula solicitada en el examen oficial de la UNI.
\end{enumerate}

\begin{trapbox}
Olvidar que $x < 0$ al calcular la primitiva de $1/x$, lo que requiere usar $\ln|x|$ en lugar de $\ln(x)$.
\end{trapbox}

\vspace{0.5cm}
\hrule
\vspace{0.5cm}

\subsection*{Problema 5.9: Demostración de Fubini para funciones continuas en rectángulos \hfill \color{primaryblue}\small[Nivel 3]}
\addcontentsline{toc}{subsection}{Problema 5.9: \texorpdfstring{Demostración de Fubini para funciones continuas en rectángulos}{Demostración de Fubini para funciones continuas en rectángulos} (Nivel 3)}

\begin{tcolorbox}[colback=white,colframe=gray!40,arc=1mm,boxrule=0.5pt,left=3mm,right=3mm,top=2mm,bottom=2mm]
\textbf{Enunciado:} Sea $R = [a,b] \times [c,d]$ y $f: R \to \mathbb{R}$ continua. Defina la función de área de la sección transversal $A(x) = \int_c^d f(x,y) dy$. Demuestre que $A(x)$ es continua en $[a,b]$ y que:
\begin{equation*}
\iint_R f(x,y) \, dA = \int_a^b A(x) \, dx = \int_a^b \left( \int_c^d f(x,y) \, dy \right) dx.
\end{equation*}
\end{tcolorbox}

\begin{tcolorbox}[colback=yellow!4,colframe=accentamber!70,arc=1mm,boxrule=0.5pt,left=3mm,right=3mm,top=1.5mm,bottom=1.5mm]
\textbf{Pista Conceptual:} Use la continuidad uniforme de $f$ en $R$ para demostrar que $A(x)$ es continua. Luego muestre que las sumas de Riemann de $A(x)$ en $[a,b]$ son sumas de Riemann de $f$ en $R$.
\end{tcolorbox}

\paragraph{Solución Paso a Paso:}\par
\begin{enumerate}[label=\textbf{\arabic*.}]
\item \textbf{Continuidad de la función sección transversal $A(x)$:}
Dado que $R$ es compacto y $f$ es continua, $f$ es uniformemente continua:
Para todo $\varepsilon > 0$, existe $\delta > 0$ tal que si $|x_1 - x_2| < \delta$ e $|y_1 - y_2| < \delta$, entonces $|f(x_1, y) - f(x_2, y)| < \frac{\varepsilon}{d - c}$.
Sean $x_1, x_2 \in [a,b]$ con $|x_1 - x_2| < \delta$. Entonces:
\begin{equation*}
|A(x_1) - A(x_2)| = \left| \int_c^d (f(x_1, y) - f(x_2, y)) \, dy \right| \le \int_c^d |f(x_1, y) - f(x_2, y)| \, dy
\end{equation*}
\begin{equation*}
< \int_c^d \frac{\varepsilon}{d - c} \, dy = \frac{\varepsilon}{d - c}(d - c) = \varepsilon.
\end{equation*}
Esto demuestra que $A(x)$ es continua en $[a,b]$, luego es integrable Riemann.
\end{enumerate}

\begin{enumerate}[label=\textbf{\arabic*.}]
\item \textbf{Aproximación por sumas de Darboux:}
Sea $P = P_1 \times P_2$ una partición de $R$, donde $P_1 = \{x_0, \dots, x_m\}$ y $P_2 = \{y_0, \dots, y_n\}$.
Para cualquier $x \in [x_{i-1}, x_i]$ y cualquier $y \in [y_{j-1}, y_j]$, denotemos $m_{ij} = \inf_{R_{ij}} f$ y $M_{ij} = \sup_{R_{ij}} f$.
Integrando en $y$ sobre $[c,d]$:
\begin{equation*}
\sum_{j=1}^n m_{ij} \Delta y_j \le \int_c^d f(x,y) \, dy = A(x) \le \sum_{j=1}^n M_{ij} \Delta y_j.
\end{equation*}
Integrando ahora en $x$ sobre $[x_{i-1}, x_i]$:
\begin{equation*}
\sum_{j=1}^n m_{ij} \Delta y_j \Delta x_i \le \int_{x_{i-1}}^{x_i} A(x) \, dx \le \sum_{j=1}^n M_{ij} \Delta y_j \Delta x_i.
\end{equation*}
Sumando sobre todos los subintervalos $i = 1, \dots, m$:
\begin{equation*}
L(f, P) \le \int_a^b A(x) \, dx \le U(f, P).
\end{equation*}
\end{enumerate}

\begin{enumerate}[label=\textbf{\arabic*.}]
\item \textbf{Conclusión:}
Como $f$ es continua, es integrable, por lo que $\sup_P L(f, P) = \inf_P U(f, P) = \iint_R f dA$.
Dado que el número fijo $\int_a^b A(x) dx$ está atrapado entre $L(f,P)$ y $U(f,P)$ para toda partición $P$:
\begin{equation*}
\int_a^b A(x) \, dx = \iint_R f(x,y) \, dA.
\end{equation*}
Esto demuestra rigurosamente el Teorema de Fubini para funciones continuas.
\end{enumerate}

\begin{trapbox}
Asumir la integrabilidad de $A(x)$ sin demostrar previamente su continuidad.
\end{trapbox}

\vspace{0.5cm}
\hrule
\vspace{0.5cm}

\subsection*{Problema 5.10: Falla de Fubini en integrando no acotado (Contraejemplo de Lichtenstein) \hfill \color{primaryblue}\small[Nivel 4]}
\addcontentsline{toc}{subsection}{Problema 5.10: \texorpdfstring{Falla de Fubini en integrando no acotado (Contraejemplo de Lichtenstein)}{Falla de Fubini en integrando no acotado (Contraejemplo de Lichtenstein)} (Nivel 4)}

\begin{tcolorbox}[colback=white,colframe=gray!40,arc=1mm,boxrule=0.5pt,left=3mm,right=3mm,top=2mm,bottom=2mm]
\textbf{Enunciado:} Considere la función $f: [0, 1] \times [0, 1] \to \mathbb{R}$ definida por:
\begin{equation*}
f(x,y) = \begin{cases} \frac{x^2 - y^2}{(x^2 + y^2)^2}, & (x,y) \neq (0,0) \\ 0, & (x,y) = (0,0) \end{cases}
\end{equation*}
(a) Calcule la integral iterada $I_1 = \int_0^1 \left(\int_0^1 f(x,y) dy\right) dx$.
(b) Calcule la integral iterada $I_2 = \int_0^1 \left(\int_0^1 f(x,y) dx\right) dy$.
(c) Verifique que $I_1 = \pi/4$ y $I_2 = -\pi/4$. Explique por qué este resultado no contradice el Teorema de Fubini.
\end{tcolorbox}

\begin{tcolorbox}[colback=yellow!4,colframe=accentamber!70,arc=1mm,boxrule=0.5pt,left=3mm,right=3mm,top=1.5mm,bottom=1.5mm]
\textbf{Pista Conceptual:} Observe que $\frac{\partial}{\partial y}\left(\frac{y}{x^2+y^2}\right) = \frac{x^2 - y^2}{(x^2+y^2)^2}$. Evalúe con cuidado en el límite inferior.
\end{tcolorbox}

\paragraph{Solución Paso a Paso:}\par
\begin{enumerate}[label=\textbf{\arabic*.}]
\item \textbf{Identificación de la antiderivada:}
Derivando el cociente respecto a $y$:
\begin{equation*}
\frac{\partial}{\partial y} \left( \frac{y}{x^2 + y^2} \right) = \frac{1(x^2 + y^2) - y(2y)}{(x^2 + y^2)^2} = \frac{x^2 - y^2}{(x^2 + y^2)^2}.
\end{equation*}
Por simetría, respecto a $x$:
\begin{equation*}
\frac{\partial}{\partial x} \left( \frac{-x}{x^2 + y^2} \right) = \frac{-(x^2 + y^2) - (-x)(2x)}{(x^2 + y^2)^2} = \frac{x^2 - y^2}{(x^2 + y^2)^2}.
\end{equation*}
\end{enumerate}

\begin{enumerate}[label=\textbf{\arabic*.}]
\item \textbf{Cálculo de la primera integral iterada $I_1$:}
Para cada $x \in (0, 1]$ fijo:
\begin{equation*}
\int_0^1 f(x,y) \, dy = \left[ \frac{y}{x^2 + y^2} \right]_{y=0}^{y=1} = \frac{1}{x^2 + 1} - 0 = \frac{1}{1 + x^2}.
\end{equation*}
Integrando ahora respecto a $x$:
\begin{equation*}
I_1 = \int_0^1 \frac{1}{1 + x^2} \, dx = [\arctan x]_0^1 = \arctan(1) - \arctan(0) = \frac{\pi}{4}.
\end{equation*}
\end{enumerate}

\begin{enumerate}[label=\textbf{\arabic*.}]
\item \textbf{Cálculo de la segunda integral iterada $I_2$:}
Para cada $y \in (0, 1]$ fijo:
\begin{equation*}
\int_0^1 f(x,y) \, dx = \left[ -\frac{x}{x^2 + y^2} \right]_{x=0}^{x=1} = -\frac{1}{1 + y^2} - 0 = -\frac{1}{1 + y^2}.
\end{equation*}
Integrando ahora respecto a $y$:
\begin{equation*}
I_2 = \int_0^1 \left( -\frac{1}{1 + y^2} \right) dy = -[\arctan y]_0^1 = -\frac{\pi}{4}.
\end{equation*}
\end{enumerate}

\begin{enumerate}[label=\textbf{\arabic*.}]
\item \textbf{Explicación teórica de la no contradicción:}
Obtenemos $I_1 = \pi/4 \neq I_2 = -\pi/4$.
Esto \textbf{NO contradice} el Teorema de Fubini porque la hipótesis fundamental del Teorema de Fubini es que la función sea \textbf{integrable Riemann en el bloque}, lo cual exige que $f$ sea una función \textbf{acotada}.
En este caso, cuando $(x,y) \to (0,0)$ a lo largo de la recta $y = 0$, $f(x,0) = 1/x^2 \to \infty$.
Al no ser acotada en $[0,1]^2$, $f$ no es integrable según Riemann. Por ende, las hipótesis de Fubini no se satisfacen.
\end{enumerate}

\begin{trapbox}
Creer que las integrales iteradas siempre coinciden independientemente de si la función es acotada o no.
\end{trapbox}

\vspace{0.5cm}
\hrule
\vspace{0.5cm}

\subsection*{Problema 5.11: Teorema del Valor Medio Ponderado para Integrales Dobles \hfill \color{primaryblue}\small[Nivel 4]}
\addcontentsline{toc}{subsection}{Problema 5.11: \texorpdfstring{Teorema del Valor Medio Ponderado para Integrales Dobles}{Teorema del Valor Medio Ponderado para Integrales Dobles} (Nivel 4)}

\begin{tcolorbox}[colback=white,colframe=gray!40,arc=1mm,boxrule=0.5pt,left=3mm,right=3mm,top=2mm,bottom=2mm]
\textbf{Enunciado:} Sean $f, g: R \to \mathbb{R}$ continuas en el rectángulo compacto y conexo $R$, con $g(x,y) \ge 0$ para todo $(x,y) \in R$. Demuestre que existe $(x_0, y_0) \in R$ tal que:
\begin{equation*}
\iint_R f(x,y) g(x,y) \, dA = f(x_0, y_0) \iint_R g(x,y) \, dA.
\end{equation*}
\end{tcolorbox}

\begin{tcolorbox}[colback=yellow!4,colframe=accentamber!70,arc=1mm,boxrule=0.5pt,left=3mm,right=3mm,top=1.5mm,bottom=1.5mm]
\textbf{Pista Conceptual:} Si $\iint_R g dA = 0$, $g \equiv 0$ trivialmente. Si no, defina $\mu = \frac{\iint_R fg dA}{\iint_R g dA}$ y acote usando $m \le f \le M$.
\end{tcolorbox}

\paragraph{Solución Paso a Paso:}\par
\begin{enumerate}[label=\textbf{\arabic*.}]
\item \textbf{Caso trivial $\iint_R g \, dA = 0$:}
Como $g$ es continua y $g \ge 0$, $\iint_R g dA = 0 \iff g(x,y) = 0$ para todo $(x,y) \in R$.
En ese caso, $fg \equiv 0$, y ambos lados de la igualdad valen 0 para cualquier punto $(x_0, y_0) \in R$.
\end{enumerate}

\begin{enumerate}[label=\textbf{\arabic*.}]
\item \textbf{Caso no trivial $\iint_R g \, dA > 0$:}
Dado que $R$ es compacto y $f$ es continua, por el Teorema de Weierstrass existen puntos $p_{\min}, p_{\max} \in R$ tales que:
\begin{equation*}
m = f(p_{\min}) \le f(x,y) \le f(p_{\max}) = M, \quad \forall (x,y) \in R.
\end{equation*}
Como $g(x,y) \ge 0$:
\begin{equation*}
m \cdot g(x,y) \le f(x,y) g(x,y) \le M \cdot g(x,y), \quad \forall (x,y) \in R.
\end{equation*}
Integrando sobre $R$:
\begin{equation*}
m \iint_R g(x,y) \, dA \le \iint_R f(x,y) g(x,y) \, dA \le M \iint_R g(x,y) \, dA.
\end{equation*}
Dividiendo entre el número estrictamente positivo $K = \iint_R g(x,y) dA > 0$:
\begin{equation*}
m \le \frac{\iint_R f(x,y) g(x,y) \, dA}{\iint_R g(x,y) \, dA} \le M.
\end{equation*}
\end{enumerate}

\begin{enumerate}[label=\textbf{\arabic*.}]
\item \textbf{Invocación del Teorema del Valor Intermedio:}
Sea el promedio ponderado $\mu = \frac{\iint_R fg dA}{\iint_R g dA} \in [m, M]$.
Como $R$ es conexo y $f$ es continua, la imagen $f(R)$ es conexa en $\mathbb{R}$, es decir, contiene todo el intervalo $[m, M]$.
Por ende, existe algún punto $(x_0, y_0) \in R$ tal que:
\begin{equation*}
f(x_0, y_0) = \mu = \frac{\iint_R f(x,y) g(x,y) \, dA}{\iint_R g(x,y) \, dA}.
\end{equation*}
Multiplicando por $\iint_R g dA$ se concluye la tesis.
\end{enumerate}

\begin{trapbox}
Omitir la hipótesis $g(x,y) \ge 0$; si $g$ cambia de signo, la desigualdad de funciones no se preserva al multiplicar por $g$.
\end{trapbox}

\vspace{0.5cm}
\hrule
\vspace{0.5cm}

\subsection*{Problema 5.12: Regla de Leibniz para derivación bajo el signo integral mediante Fubini \hfill \color{primaryblue}\small[Nivel 4]}
\addcontentsline{toc}{subsection}{Problema 5.12: \texorpdfstring{Regla de Leibniz para derivación bajo el signo integral mediante Fubini}{Regla de Leibniz para derivación bajo el signo integral mediante Fubini} (Nivel 4)}

\begin{tcolorbox}[colback=white,colframe=gray!40,arc=1mm,boxrule=0.5pt,left=3mm,right=3mm,top=2mm,bottom=2mm]
\textbf{Enunciado:} Sea $f(x,y)$ continua en $R = [a,b] \times [c,d]$ tal que la derivada parcial $\frac{\partial f}{\partial x}(x,y)$ existe y es continua en $R$. Defina $F(x) = \int_c^d f(x,y) dy$. Demuestre rigurosamente, empleando el Teorema de Fubini y el Teorema Fundamental del Cálculo, que $F$ es diferenciable en $[a,b]$ y:
\begin{equation*}
F'(x) = \int_c^d \frac{\partial f}{\partial x}(x,y) \, dy.
\end{equation*}
\end{tcolorbox}

\begin{tcolorbox}[colback=yellow!4,colframe=accentamber!70,arc=1mm,boxrule=0.5pt,left=3mm,right=3mm,top=1.5mm,bottom=1.5mm]
\textbf{Pista Conceptual:} Sea $g(x,y) = \frac{\partial f}{\partial x}(x,y)$. Integre $g$ en $[a, x] \times [c, d]$ aplicando Fubini y luego derive usando el TFC.
\end{tcolorbox}

\paragraph{Solución Paso a Paso:}\par
\begin{enumerate}[label=\textbf{\arabic*.}]
\item \textbf{Definición de la función derivada auxiliar:}
Sea $g(x,y) = \frac{\partial f}{\partial x}(x,y)$.
Por hipótesis, $g$ es continua en el rectángulo compacto $R = [a,b] \times [c,d]$.
Por el Teorema Fundamental del Cálculo en una variable:
\begin{equation*}
f(x, y) - f(a, y) = \int_a^x g(t, y) \, dt, \quad \forall (x,y) \in R.
\end{equation*}
\end{enumerate}

\begin{enumerate}[label=\textbf{\arabic*.}]
\item \textbf{Integración con respecto a $y$ e intercambio por Fubini:}
Integrando en $y \in [c, d]$:
\begin{equation*}
\int_c^d f(x, y) \, dy - \int_c^d f(a, y) \, dy = \int_c^d \left( \int_a^x g(t, y) \, dt \right) dy.
\end{equation*}
Por definición, el lado izquierdo es $F(x) - F(a)$.
Para el lado derecho, como $g$ es continua en el rectángulo compacto $[a,x] \times [c,d]$, por el Teorema de Fubini podemos intercambiar el orden de integración:
\begin{equation*}
\int_c^d \left( \int_a^x g(t, y) \, dt \right) dy = \int_a^x \left( \int_c^d g(t, y) \, dy \right) dt.
\end{equation*}
Definamos la función $H(t) = \int_c^d g(t, y) \, dy$.
Como probamos en el ejercicio 5.9, al ser $g$ continua, $H(t)$ es continua en $[a,b]$.
Entonces tenemos la identidad:
\begin{equation*}
F(x) = F(a) + \int_a^x H(t) \, dt.
\end{equation*}
\end{enumerate}

\begin{enumerate}[label=\textbf{\arabic*.}]
\item \textbf{Diferenciación mediante el Teorema Fundamental del Cálculo:}
Como $H$ es continua, por el Teorema Fundamental del Cálculo, la función integral $x \mapsto \int_a^x H(t) dt$ es diferenciable y su derivada es exactamente $H(x)$:
\begin{equation*}
F'(x) = \frac{d}{dx}\left[ F(a) + \int_a^x H(t) \, dt \right] = 0 + H(x) = \int_c^d g(x, y) \, dy = \int_c^d \frac{\partial f}{\partial x}(x,y) \, dy.
\end{equation*}
Esto demuestra la Regla de Leibniz mediante Fubini con total rigor analítico.
\end{enumerate}

\begin{trapbox}
Tomar límites del cociente de Newton dentro de la integral sin justificar el paso al límite bajo el signo de la integral con la continuidad uniforme.
\end{trapbox}

\vspace{0.5cm}
\hrule
\vspace{0.5cm}

\subsection*{Problema 5.13: Teorema de Fubini para funciones integrables arbitrarias \hfill \color{primaryblue}\small[Nivel 5]}
\addcontentsline{toc}{subsection}{Problema 5.13: \texorpdfstring{Teorema de Fubini para funciones integrables arbitrarias}{Teorema de Fubini para funciones integrables arbitrarias} (Nivel 5)}

\begin{tcolorbox}[colback=white,colframe=gray!40,arc=1mm,boxrule=0.5pt,left=3mm,right=3mm,top=2mm,bottom=2mm]
\textbf{Enunciado:} Sea $f: R = [a,b] \times [c,d] \to \mathbb{R}$ una función integrable Riemann (no necesariamente continua). Para cada $x \in [a,b]$, defina las integrales de Darboux unidimensionales:
\begin{equation*}
\underline{I}(x) = \underline{\int_c^d} f(x,y) \, dy, \qquad \overline{I}(x) = \overline{\int_c^d} f(x,y) \, dy.
\end{equation*}
Demuestre que $\underline{I}$ y $\overline{I}$ son integrables Riemann en $[a,b]$ y que:
\begin{equation*}
\iint_R f(x,y) \, dA = \int_a^b \underline{I}(x) \, dx = \int_a^b \overline{I}(x) \, dx.
\end{equation*}
\end{tcolorbox}

\begin{tcolorbox}[colback=yellow!4,colframe=accentamber!70,arc=1mm,boxrule=0.5pt,left=3mm,right=3mm,top=1.5mm,bottom=1.5mm]
\textbf{Pista Conceptual:} Para cualquier partición producto $P = P_1 \times P_2$, demuestre que $L(f, P) \le L(\underline{I}, P_1) \le U(\overline{I}, P_1) \le U(f, P)$.
\end{tcolorbox}

\paragraph{Solución Paso a Paso:}\par
\begin{enumerate}[label=\textbf{\arabic*.}]
\item \textbf{Acotación con particiones de Darboux:}
Sea $P = P_1 \times P_2$ una partición arbitraria de $R$, donde $P_1 = \{x_0, \dots, x_m\}$ y $P_2 = \{y_0, \dots, y_n\}$.
Para cada subrectángulo $R_{ij} = [x_{i-1}, x_i] \times [y_{j-1}, y_j]$, denotemos $m_{ij} = \inf_{R_{ij}} f$ y $M_{ij} = \sup_{R_{ij}} f$.
Para cualquier $x \in [x_{i-1}, x_i]$ fijo y cualquier $y \in [y_{j-1}, y_j]$, se tiene:
\begin{equation*}
m_{ij} \le f(x,y) \le M_{ij}.
\end{equation*}
Considerando la función $y \mapsto f(x,y)$ en $[c,d]$ con la partición $P_2$:
\begin{equation*}
\sum_{j=1}^n m_{ij} \Delta y_j \le \underline{I}(x) \le \overline{I}(x) \le \sum_{j=1}^n M_{ij} \Delta y_j, \quad \forall x \in [x_{i-1}, x_i].
\end{equation*}
\end{enumerate}

\begin{enumerate}[label=\textbf{\arabic*.}]
\item \textbf{Supremos e ínfimos sobre $[x_{i-1}, x_i]$:}
Como la cota inferior $\sum_{j=1}^n m_{ij} \Delta y_j$ es constante respecto a $x$, es menor o igual al ínfimo de $\underline{I}$ en $[x_{i-1}, x_i]$:
\begin{equation*}
\sum_{j=1}^n m_{ij} \Delta y_j \le \inf_{x \in [x_{i-1}, x_i]} \underline{I}(x).
\end{equation*}
Multiplicando por $\Delta x_i > 0$ y sumando sobre $i = 1, \dots, m$:
\begin{equation*}
L(f, P) = \sum_{i=1}^m \sum_{j=1}^n m_{ij} \Delta x_i \Delta y_j \le L(\underline{I}, P_1).
\end{equation*}
De forma análoga para la cota superior:
\begin{equation*}
U(\overline{I}, P_1) \le U(f, P).
\end{equation*}
Uniendo ambas cadenas de desigualdades con el hecho evidente de que $L(\underline{I}, P_1) \le U(\underline{I}, P_1) \le U(\overline{I}, P_1)$:
\begin{equation*}
L(f, P) \le L(\underline{I}, P_1) \le U(\underline{I}, P_1) \le U(f, P),
\end{equation*}
\begin{equation*}
L(f, P) \le L(\overline{I}, P_1) \le U(\overline{I}, P_1) \le U(f, P).
\end{equation*}
\end{enumerate}

\begin{enumerate}[label=\textbf{\arabic*.}]
\item \textbf{Criterio de Riemann y conclusión:}
Como $f$ es integrable en $R$, dado $\varepsilon > 0$ existe una partición $P$ tal que $U(f, P) - L(f, P) < \varepsilon$.
Por tanto:
\begin{equation*}
U(\underline{I}, P_1) - L(\underline{I}, P_1) \le U(f, P) - L(f, P) < \varepsilon.
\end{equation*}
Esto prueba que $\underline{I}(x)$ es integrable en $[a,b]$. Del mismo modo, $\overline{I}(x)$ es integrable en $[a,b]$.
Además, como $\iint_R f dA$ y las integrales de $\underline{I}$ y $\overline{I}$ están todas atrapadas entre $L(f,P)$ y $U(f,P)$ para toda partición, coinciden exactamente:
\begin{equation*}
\iint_R f(x,y) \, dA = \int_a^b \left( \underline{\int_c^d} f(x,y) \, dy \right) dx = \int_a^b \left( \overline{\int_c^d} f(x,y) \, dy \right) dx.
\end{equation*}
\end{enumerate}

\begin{trapbox}
Asumir que para todo $x$, la sección transversal $y \mapsto f(x,y)$ es integrable Riemann; existen funciones integrables en 2D donde para algunos $x$ la sección no es integrable, requiriendo usar integrales superior e inferior.
\end{trapbox}

\vspace{0.5cm}
\hrule
\vspace{0.5cm}

\subsection*{Problema 5.14: Cálculo de la Integral Gaussiana por Fubini y comparación en cuadrados \hfill \color{primaryblue}\small[Nivel 5]}
\addcontentsline{toc}{subsection}{Problema 5.14: \texorpdfstring{Cálculo de la Integral Gaussiana por Fubini y comparación en cuadrados}{Cálculo de la Integral Gaussiana por Fubini y comparación en cuadrados} (Nivel 5)}

\begin{tcolorbox}[colback=white,colframe=gray!40,arc=1mm,boxrule=0.5pt,left=3mm,right=3mm,top=2mm,bottom=2mm]
\textbf{Enunciado:} Sea $I_M = \int_{-M}^M e^{-x^2} dx$. Considere el cuadrado $R_M = [-M, M] \times [-M, M]$. (a) Exprese $\iint_{R_M} e^{-(x^2 + y^2)} dA$ en términos de $I_M$ usando Fubini. (b) Utilice los discos inscritos $D_M = \{x^2+y^2 \le M^2\}$ y circunscritos $D_{M\sqrt{2}} = \{x^2+y^2 \le 2M^2\}$ para acotar la integral doble. (c) Deduzca rigurosamente el valor de la integral de Poisson-Gauss $\int_{-\infty}^\infty e^{-x^2} dx = \sqrt{\pi}$.
\end{tcolorbox}

\begin{tcolorbox}[colback=yellow!4,colframe=accentamber!70,arc=1mm,boxrule=0.5pt,left=3mm,right=3mm,top=1.5mm,bottom=1.5mm]
\textbf{Pista Conceptual:} En coordenadas polares, la integral en un disco $D_r$ es $\int_0^{2\pi}\int_0^r e^{-\rho^2}\rho d\rho d\theta = \pi(1 - e^{-r^2})$. Acote $\iint_{D_M} \le \iint_{R_M} \le \iint_{D_{M\sqrt{2}}}$ y tome límite $M \to \infty$.
\end{tcolorbox}

\paragraph{Solución Paso a Paso:}\par
\begin{enumerate}[label=\textbf{\arabic*.}]
\item \textbf{Factorización mediante Fubini en el cuadrado $R_M$:}
La función $f(x,y) = e^{-(x^2 + y^2)} = e^{-x^2} e^{-y^2}$ es continua y separable.
Por el Teorema de Fubini:
\begin{equation*}
\iint_{R_M} e^{-(x^2 + y^2)} \, dA = \left(\int_{-M}^M e^{-x^2} dx\right) \left(\int_{-M}^M e^{-y^2} dy\right) = (I_M)^2.
\end{equation*}
\end{enumerate}

\begin{enumerate}[label=\textbf{\arabic*.}]
\item \textbf{Inclusión geométrica de discos y cuadrados:}
Sea $D_r = \{(x,y) \in \mathbb{R}^2 : x^2 + y^2 \le r^2\}$.
Geométricamente:
- El disco inscrito en el cuadrado $[-M, M]^2$ tiene radio $M$: $D_M \subset R_M$.
- El disco circunscrito al cuadrado $[-M, M]^2$ pasa por los vértices $(\pm M, \pm M)$ y tiene radio $\sqrt{M^2 + M^2} = M\sqrt{2}$: $R_M \subset D_{M\sqrt{2}}$.
Por consiguiente:
\begin{equation*}
D_M \subset R_M \subset D_{M\sqrt{2}}.
\end{equation*}
\end{enumerate}

\begin{enumerate}[label=\textbf{\arabic*.}]
\item \textbf{Cálculo de integrales en los discos en coordenadas polares:}
Como el integrando es no negativo $e^{-(x^2+y^2)} > 0$, la integral es monótona con respecto al dominio:
\begin{equation*}
\iint_{D_M} e^{-(x^2 + y^2)} \, dA \le \iint_{R_M} e^{-(x^2 + y^2)} \, dA \le \iint_{D_{M\sqrt{2}}} e^{-(x^2 + y^2)} \, dA.
\end{equation*}
Calculamos la integral en un disco de radio $R$:
\begin{equation*}
\iint_{D_R} e^{-(x^2+y^2)} dA = \int_0^{2\pi} \int_0^R e^{-\rho^2} \rho \, d\rho \, d\theta = 2\pi \left[ -\frac{1}{2} e^{-\rho^2} \right]_0^R = \pi (1 - e^{-R^2}).
\end{equation*}
Aplicando a los dos discos:
\begin{equation*}
\pi (1 - e^{-M^2}) \le (I_M)^2 \le \pi (1 - e^{-2M^2}).
\end{equation*}
\end{enumerate}

\begin{enumerate}[label=\textbf{\arabic*.}]
\item \textbf{Paso al límite cuando $M \to \infty$:}
Tomando el límite cuando $M \to \infty$:
\begin{equation*}
\lim_{M\to\infty} \pi (1 - e^{-M^2}) = \pi (1 - 0) = \pi,
\end{equation*}
\begin{equation*}
\lim_{M\to\infty} \pi (1 - e^{-2M^2}) = \pi (1 - 0) = \pi.
\end{equation*}
Por el Teorema del Sándwich:
\begin{equation*}
\lim_{M\to\infty} (I_M)^2 = \pi \implies \left( \int_{-\infty}^\infty e^{-x^2} \, dx \right)^2 = \pi.
\end{equation*}
Como $e^{-x^2} > 0$, la integral es positiva:
\begin{equation*}
\int_{-\infty}^\infty e^{-x^2} \, dx = \sqrt{\pi}.
\end{equation*}
\end{enumerate}

\begin{trapbox}
Calcular la integral gaussiana usando coordenadas polares directamente sobre todo el plano sin justificar el encaje de dominios cuadrados y discos.
\end{trapbox}

\vspace{0.5cm}
\hrule
\vspace{0.5cm}

\subsection*{Problema 5.15: Teorema de Fubini n-dimensional y reducción secuencial \hfill \color{primaryblue}\small[Nivel 5]}
\addcontentsline{toc}{subsection}{Problema 5.15: \texorpdfstring{Teorema de Fubini n-dimensional y reducción secuencial}{Teorema de Fubini n-dimensional y reducción secuencial} (Nivel 5)}

\begin{tcolorbox}[colback=white,colframe=gray!40,arc=1mm,boxrule=0.5pt,left=3mm,right=3mm,top=2mm,bottom=2mm]
\textbf{Enunciado:} Sea $B = \prod_{i=1}^n [a_i, b_i]$ un bloque cerrado en $\mathbb{R}^n$ y $f: B \to \mathbb{R}$ continua. Demuestre formalmente por inducción sobre la dimensión $n$ que la integral múltiple $\int_B f dV_n$ es igual a la integral iterada unidimensional secuencial:
\begin{equation*}
\int_B f \, dV_n = \int_{a_1}^{b_1} \dots \left( \int_{a_n}^{b_n} f(x_1, \dots, x_n) \, dx_n \right) \dots dx_1.
\end{equation*}
\end{tcolorbox}

\begin{tcolorbox}[colback=yellow!4,colframe=accentamber!70,arc=1mm,boxrule=0.5pt,left=3mm,right=3mm,top=1.5mm,bottom=1.5mm]
\textbf{Pista Conceptual:} Considere la descomposición $B = B_{n-1} \times [a_n, b_n]$ y aplique el caso base bidimensional $n=2$ entre un bloque $(n-1)$-dimensional y un intervalo.
\end{tcolorbox}

\paragraph{Solución Paso a Paso:}\par
\begin{enumerate}[label=\textbf{\arabic*.}]
\item \textbf{Caso base ($n = 2$):}
Para $n = 2$, el bloque es $B = [a_1, b_1] \times [a_2, b_2]$. El Teorema de Fubini bidimensional (demostrado en el ejercicio 5.9) asegura que para $f$ continua:
\begin{equation*}
\iint_B f \, dA = \int_{a_1}^{b_1} \left( \int_{a_2}^{b_2} f(x_1, x_2) \, dx_2 \right) dx_1.
\end{equation*}
\end{enumerate}

\begin{enumerate}[label=\textbf{\arabic*.}]
\item \textbf{Paso inductivo:}
Supongamos como hipótesis de inducción que el teorema es válido para cualquier bloque en dimensión $n-1$:
Para cualquier bloque $B' = \prod_{i=1}^{n-1} [a_i, b_i] \subset \mathbb{R}^{n-1}$ y cualquier función continua $g: B' \to \mathbb{R}$:
\begin{equation*}
\int_{B'} g \, dV_{n-1} = \int_{a_1}^{b_1} \dots \left( \int_{a_{n-1}}^{b_{n-1}} g(x_1, \dots, x_{n-1}) \, dx_{n-1} \right) \dots dx_1.
\end{equation*}
\end{enumerate}

\begin{enumerate}[label=\textbf{\arabic*.}]
\item \textbf{Descomposición del bloque $n$-dimensional:}
Sea $B = \prod_{i=1}^n [a_i, b_i] \subset \mathbb{R}^n$.
Podemos identificar canónicamente $\mathbb{R}^n = \mathbb{R}^{n-1} \times \mathbb{R}$, de modo que:
\begin{equation*}
B = B' \times [a_n, b_n], \quad \text{con } B' = \prod_{i=1}^{n-1} [a_i, b_i].
\end{equation*}
Aplicando el Teorema de Fubini para el producto de dos bloques (uno de dimensión $n-1$ y otro de dimensión 1):
\begin{equation*}
\int_B f \, dV_n = \int_{B'} \left( \int_{a_n}^{b_n} f(x_1, \dots, x_n) \, dx_n \right) dV_{n-1}(x_1, \dots, x_{n-1}).
\end{equation*}
Definamos la función $g: B' \to \mathbb{R}$ dada por:
\begin{equation*}
g(x_1, \dots, x_{n-1}) = \int_{a_n}^{b_n} f(x_1, \dots, x_n) \, dx_n.
\end{equation*}
Como $f$ es continua en el compacto $B$, $g$ es continua en $B'$.
\end{enumerate}

\begin{enumerate}[label=\textbf{\arabic*.}]
\item \textbf{Aplicación de la hipótesis inductiva a $g$:}
Aplicando la hipótesis de inducción a la integral sobre $B'$ de la función continua $g$:
\begin{equation*}
\int_{B'} g \, dV_{n-1} = \int_{a_1}^{b_1} \dots \left( \int_{a_{n-1}}^{b_{n-1}} g(x_1, \dots, x_{n-1}) \, dx_{n-1} \right) \dots dx_1.
\end{equation*}
Sustituyendo la definición de $g$:
\begin{equation*}
\int_B f \, dV_n = \int_{a_1}^{b_1} \dots \int_{a_{n-1}}^{b_{n-1}} \left( \int_{a_n}^{b_n} f(x_1, \dots, x_n) \, dx_n \right) dx_{n-1} \dots dx_1.
\end{equation*}
Por inducción matemática, el resultado es válido para todo $n \ge 2$.
\end{enumerate}

\begin{trapbox}
No justificar la continuidad de la función reducida $g$ en $\mathbb{R}^{n-1}$ antes de aplicar la hipótesis inductiva.
\end{trapbox}

\vspace{0.5cm}
\hrule
\vspace{0.5cm}

\section{Bonus: Regiones Generales Tipo 1 y 2}

\subsection*{Problema B.1: Área de región plana comprendida entre parábola y recta (Tipo 1) \hfill \color{primaryblue}\small[Nivel 2]}
\addcontentsline{toc}{subsection}{Problema B.1: \texorpdfstring{Área de región plana comprendida entre parábola y recta (Tipo 1)}{Área de región plana comprendida entre parábola y recta (Tipo 1)} (Nivel 2)}

\begin{tcolorbox}[colback=white,colframe=gray!40,arc=1mm,boxrule=0.5pt,left=3mm,right=3mm,top=2mm,bottom=2mm]
\textbf{Enunciado:} Sea la región plana $D$ acotada por la parábola $y = x^2$ y la recta $y = x + 2$. (a) Determine los puntos de intersección y clasifique $D$ como región de Tipo 1. (b) Calcule el área de $D$ mediante la integral doble $\iint_D 1 \, dA$.
\end{tcolorbox}

\begin{tcolorbox}[colback=yellow!4,colframe=accentamber!70,arc=1mm,boxrule=0.5pt,left=3mm,right=3mm,top=1.5mm,bottom=1.5mm]
\textbf{Pista Conceptual:} Resuelva $x^2 = x + 2$ para hallar el intervalo en $x$. En cada $x$, $y$ varía desde la curva inferior $y = x^2$ hasta la superior $y = x + 2$.
\end{tcolorbox}

\paragraph{Solución Paso a Paso:}\par
\begin{enumerate}[label=\textbf{\arabic*.}]
\item \textbf{Puntos de corte:}
\begin{equation*}
x^2 = x + 2 \implies x^2 - x - 2 = 0 \implies (x - 2)(x + 1) = 0 \implies x = -1, \quad x = 2.
\end{equation*}
Para $x \in [-1, 2]$, la recta está por encima de la parábola: $x + 2 \ge x^2$.
Por tanto, $D$ es una región de Tipo 1:
\begin{equation*}
D = \{(x,y) \in \mathbb{R}^2 : -1 \le x \le 2, \ x^2 \le y \le x + 2\}.
\end{equation*}
\end{enumerate}

\begin{enumerate}[label=\textbf{\arabic*.}]
\item \textbf{Cálculo del área por integral doble:}
\begin{equation*}
\text{Área}(D) = \iint_D 1 \, dA = \int_{-1}^2 \left( \int_{x^2}^{x+2} 1 \, dy \right) dx = \int_{-1}^2 ((x + 2) - x^2) \, dx
\end{equation*}
\begin{equation*}
= \left[ \frac{x^2}{2} + 2x - \frac{x^3}{3} \right]_{-1}^2
\end{equation*}
\begin{equation*}
= \left( \frac{4}{2} + 4 - \frac{8}{3} \right) - \left( \frac{1}{2} - 2 + \frac{1}{3} \right) = \left(6 - \frac{8}{3}\right) - \left(-\frac{3}{2} + \frac{1}{3}\right) = \frac{10}{3} - \left(-\frac{7}{6}\right) = \frac{20 + 7}{6} = \frac{27}{6} = \frac{9}{2} = 4.5 \, u^2.
\end{equation*}
\end{enumerate}

\begin{trapbox}
Invertir el orden de las funciones techo y piso: la fórmula es $\int (y_{\text{sup}} - y_{\text{inf}}) dx$.
\end{trapbox}

\vspace{0.5cm}
\hrule
\vspace{0.5cm}

\subsection*{Problema B.2: Inversión del orden de integración en región no rectangular \hfill \color{primaryblue}\small[Nivel 3]}
\addcontentsline{toc}{subsection}{Problema B.2: \texorpdfstring{Inversión del orden de integración en región no rectangular}{Inversión del orden de integración en región no rectangular} (Nivel 3)}

\begin{tcolorbox}[colback=white,colframe=gray!40,arc=1mm,boxrule=0.5pt,left=3mm,right=3mm,top=2mm,bottom=2mm]
\textbf{Enunciado:} Considere la integral:
\begin{equation*}
I = \int_0^1 \int_y^1 \sin(x^2) \, dx \, dy.
\end{equation*}
(a) Esboce la región de integración $D$ y clasifíquela como región de Tipo 2. (b) Invierta el orden de integración expresando $D$ como región de Tipo 1 y evalúe la integral exactamente.
\end{tcolorbox}

\begin{tcolorbox}[colback=yellow!4,colframe=accentamber!70,arc=1mm,boxrule=0.5pt,left=3mm,right=3mm,top=1.5mm,bottom=1.5mm]
\textbf{Pista Conceptual:} La región original está acotada por $y \le x \le 1$ con $0 \le y \le 1$. Al invertir, $x \in [0,1]$ y para cada $x$, $0 \le y \le x$.
\end{tcolorbox}

\paragraph{Solución Paso a Paso:}\par
\begin{enumerate}[label=\textbf{\arabic*.}]
\item \textbf{Descripción original (Tipo 2):}
\begin{equation*}
D = \{(x,y) \in \mathbb{R}^2 : 0 \le y \le 1, \ y \le x \le 1\}.
\end{equation*}
Geométricamente es el triángulo delimitado por las rectas $y = 0$, $x = 1$ y la diagonal $y = x$.
\end{enumerate}

\begin{enumerate}[label=\textbf{\arabic*.}]
\item \textbf{Descripción invertida (Tipo 1):}
Proyectando sobre el eje $x$, los valores de $x$ varían de $0$ a $1$.
Para cada $x$ fijo, la variable $y$ va desde la recta inferior $y = 0$ hasta la recta superior $y = x$:
\begin{equation*}
D = \{(x,y) \in \mathbb{R}^2 : 0 \le x \le 1, \ 0 \le y \le x\}.
\end{equation*}
\end{enumerate}

\begin{enumerate}[label=\textbf{\arabic*.}]
\item \textbf{Cálculo de la integral:}
\begin{equation*}
I = \int_0^1 \left( \int_0^x \sin(x^2) \, dy \right) dx = \int_0^1 \sin(x^2) [y]_0^x \, dx = \int_0^1 x \sin(x^2) \, dx.
\end{equation*}
Haciendo el cambio de variable $u = x^2 \implies du = 2x dx$:
\begin{equation*}
I = \frac{1}{2} \int_0^1 \sin(u) \, du = \frac{1}{2} [-\cos u]_0^1 = \frac{1 - \cos(1)}{2} \approx 0.2298.
\end{equation*}
\end{enumerate}

\begin{trapbox}
Intentar integrar $\sin(x^2)$ con respecto a $x$ sin invertir el orden; no posee antiderivada elemental.
\end{trapbox}

\vspace{0.5cm}
\hrule
\vspace{0.5cm}

\subsection*{Problema B.3: Volumen bajo paraboloide sobre triángulo (Sólido Tipo 1) \hfill \color{primaryblue}\small[Nivel 3]}
\addcontentsline{toc}{subsection}{Problema B.3: \texorpdfstring{Volumen bajo paraboloide sobre triángulo (Sólido Tipo 1)}{Volumen bajo paraboloide sobre triángulo (Sólido Tipo 1)} (Nivel 3)}

\begin{tcolorbox}[colback=white,colframe=gray!40,arc=1mm,boxrule=0.5pt,left=3mm,right=3mm,top=2mm,bottom=2mm]
\textbf{Enunciado:} Calcule el volumen del sólido $S$ acotado superiormente por el paraboloide $z = x^2 + y^2$, inferiormente por el plano $z = 0$, y lateralmente por los planos $x = 0$, $y = 0$ y $x + y = 1$.
\end{tcolorbox}

\begin{tcolorbox}[colback=yellow!4,colframe=accentamber!70,arc=1mm,boxrule=0.5pt,left=3mm,right=3mm,top=1.5mm,bottom=1.5mm]
\textbf{Pista Conceptual:} La proyección en el plano $xy$ es el triángulo rectángulo $D = \{0 \le x \le 1, \, 0 \le y \le 1-x\}$. Integre $\iint_D (x^2+y^2) dA$.
\end{tcolorbox}

\paragraph{Solución Paso a Paso:}\par
\begin{enumerate}[label=\textbf{\arabic*.}]
\item \textbf{Dominio base en el plano $xy$:}
El triángulo delimitado por los ejes y $x+y=1$ es de Tipo 1:
\begin{equation*}
D = \{(x,y) : 0 \le x \le 1, \ 0 \le y \le 1 - x\}.
\end{equation*}
\end{enumerate}

\begin{enumerate}[label=\textbf{\arabic*.}]
\item \textbf{Planteamiento de la integral doble:}
\begin{equation*}
V = \iint_D (x^2 + y^2) \, dA = \int_0^1 \left( \int_0^{1-x} (x^2 + y^2) \, dy \right) dx.
\end{equation*}
\end{enumerate}

\begin{enumerate}[label=\textbf{\arabic*.}]
\item \textbf{Integral interna en $y$:}
\begin{equation*}
\int_0^{1-x} (x^2 + y^2) \, dy = \left[ x^2 y + \frac{y^3}{3} \right]_0^{1-x} = x^2(1 - x) + \frac{(1 - x)^3}{3} = x^2 - x^3 + \frac{1 - 3x + 3x^2 - x^3}{3}
\end{equation*}
\begin{equation*}
= x^2 - x^3 + \frac{1}{3} - x + x^2 - \frac{x^3}{3} = 2x^2 - \frac{4}{3}x^3 - x + \frac{1}{3}.
\end{equation*}
\end{enumerate}

\begin{enumerate}[label=\textbf{\arabic*.}]
\item \textbf{Integral externa en $x$:}
\begin{equation*}
V = \int_0^1 \left( 2x^2 - \frac{4}{3}x^3 - x + \frac{1}{3} \right) dx = \left[ \frac{2x^3}{3} - \frac{x^4}{3} - \frac{x^2}{2} + \frac{x}{3} \right]_0^1
\end{equation*}
\begin{equation*}
= \frac{2}{3} - \frac{1}{3} - \frac{1}{2} + \frac{1}{3} = \frac{2}{3} - \frac{1}{2} = \frac{1}{6} \, u^3.
\end{equation*}
\end{enumerate}

\begin{trapbox}
Olvidar elevar el binomio $(1-x)^3$ o cometer errores aritméticos al sumar fracciones.
\end{trapbox}

\vspace{0.5cm}
\hrule
\vspace{0.5cm}

\subsection*{Problema B.4: Masa y centroide de una lámina con densidad lineal \hfill \color{primaryblue}\small[Nivel 4]}
\addcontentsline{toc}{subsection}{Problema B.4: \texorpdfstring{Masa y centroide de una lámina con densidad lineal}{Masa y centroide de una lámina con densidad lineal} (Nivel 4)}

\begin{tcolorbox}[colback=white,colframe=gray!40,arc=1mm,boxrule=0.5pt,left=3mm,right=3mm,top=2mm,bottom=2mm]
\textbf{Enunciado:} Sea una lámina delgada que ocupa la región semicircular $D = \{(x,y) \in \mathbb{R}^2 : x^2 + y^2 \le 1, \ y \ge 0\}$ con función de densidad $\rho(x,y) = y$. Calcule la masa total $M$ y la ordenada del centro de masa $\bar{y} = \frac{1}{M}\iint_D y \rho(x,y) dA$.
\end{tcolorbox}

\begin{tcolorbox}[colback=yellow!4,colframe=accentamber!70,arc=1mm,boxrule=0.5pt,left=3mm,right=3mm,top=1.5mm,bottom=1.5mm]
\textbf{Pista Conceptual:} Puede usar la descripción cartesiana de Tipo 1: $-1 \le x \le 1, 0 \le y \le \sqrt{1-x^2}$ o coordenadas polares para simplificar.
\end{tcolorbox}

\paragraph{Solución Paso a Paso:}\par
\begin{enumerate}[label=\textbf{\arabic*.}]
\item \textbf{Cálculo de la masa total:}
\begin{equation*}
\text{Masa } M = \iint_D y \, dA = \int_{-1}^1 \left( \int_0^{\sqrt{1-x^2}} y \, dy \right) dx = \int_{-1}^1 \left[ \frac{y^2}{2} \right]_0^{\sqrt{1-x^2}} dx
\end{equation*}
\begin{equation*}
= \int_{-1}^1 \frac{1 - x^2}{2} \, dx = \left[ \frac{x}{2} - \frac{x^3}{6} \right]_{-1}^1 = \left(\frac{1}{2} - \frac{1}{6}\right) - \left(-\frac{1}{2} + \frac{1}{6}\right) = \frac{1}{3} - \left(-\frac{1}{3}\right) = \frac{2}{3}.
\end{equation*}
\end{enumerate}

\begin{enumerate}[label=\textbf{\arabic*.}]
\item \textbf{Cálculo del momento $M_{xz}$:}
\begin{equation*}
M_{xz} = \iint_D y \cdot \rho(x,y) \, dA = \iint_D y^2 \, dA.
\end{equation*}
En coordenadas polares $x = r\cos\theta, y = r\sin\theta$ con $r \in [0,1], \theta \in [0,\pi]$:
\begin{equation*}
M_{xz} = \int_0^\pi \int_0^1 (r\sin\theta)^2 r \, dr \, d\theta = \left(\int_0^\pi \sin^2\theta \, d\theta\right) \left(\int_0^1 r^3 \, dr\right)
\end{equation*}
\begin{equation*}
= \frac{\pi}{2} \times \left[\frac{r^4}{4}\right]_0^1 = \frac{\pi}{2} \times \frac{1}{4} = \frac{\pi}{8}.
\end{equation*}
\end{enumerate}

\begin{enumerate}[label=\textbf{\arabic*.}]
\item \textbf{Ordenada del centro de masa:}
\begin{equation*}
\bar{y} = \frac{M_{xz}}{M} = \frac{\pi/8}{2/3} = \frac{3\pi}{16} \approx 0.589.
\end{equation*}
Por simetría, $\bar{x} = 0$. Por ende el centro de masa es $(0, \frac{3\pi}{16})$.
\end{enumerate}

\begin{trapbox}
Dividir el momento entre el área geométrica de la región en lugar de la masa ponderada por la densidad.
\end{trapbox}

\vspace{0.5cm}
\hrule
\vspace{0.5cm}

\subsection*{Problema B.5: Teorema de Green para el cálculo de áreas de regiones generales \hfill \color{primaryblue}\small[Nivel 4]}
\addcontentsline{toc}{subsection}{Problema B.5: \texorpdfstring{Teorema de Green para el cálculo de áreas de regiones generales}{Teorema de Green para el cálculo de áreas de regiones generales} (Nivel 4)}

\begin{tcolorbox}[colback=white,colframe=gray!40,arc=1mm,boxrule=0.5pt,left=3mm,right=3mm,top=2mm,bottom=2mm]
\textbf{Enunciado:} Enuncie la fórmula del área mediante la integral de línea de Green: $\text{Área}(D) = \frac{1}{2} \oint_{\partial D} (x dy - y dx)$. Aplíquela para deducir el área de la región encerrada por la elipse $\frac{x^2}{a^2} + \frac{y^2}{b^2} = 1$.
\end{tcolorbox}

\begin{tcolorbox}[colback=yellow!4,colframe=accentamber!70,arc=1mm,boxrule=0.5pt,left=3mm,right=3mm,top=1.5mm,bottom=1.5mm]
\textbf{Pista Conceptual:} Parametrice la elipse mediante $\gamma(t) = (a\cos t, b\sin t)$ con $t \in [0, 2\pi]$ orientada en sentido antihorario.
\end{tcolorbox}

\paragraph{Solución Paso a Paso:}\par
\begin{enumerate}[label=\textbf{\arabic*.}]
\item \textbf{Fórmula del área por Green:}
Por el Teorema de Green en el plano:
\begin{equation*}
\oint_{\partial D} (P dx + Q dy) = \iint_D \left(\frac{\partial Q}{\partial x} - \frac{\partial P}{\partial y}\right) dA.
\end{equation*}
Eligiendo $P(x,y) = -y/2$ y $Q(x,y) = x/2$:
\begin{equation*}
\frac{\partial Q}{\partial x} - \frac{\partial P}{\partial y} = \frac{1}{2} - \left(-\frac{1}{2}\right) = 1.
\end{equation*}
Por tanto:
\begin{equation*}
\text{Área}(D) = \iint_D 1 \, dA = \frac{1}{2} \oint_{\partial D} (x \, dy - y \, dx).
\end{equation*}
\end{enumerate}

\begin{enumerate}[label=\textbf{\arabic*.}]
\item \textbf{Parametrización de la frontera elíptica:}
\begin{equation*}
\gamma(t) = (a\cos t, b\sin t), \quad t \in [0, 2\pi].
\end{equation*}
Diferenciales:
\begin{equation*}
dx = -a\sin t \, dt, \qquad dy = b\cos t \, dt.
\end{equation*}
\end{enumerate}

\begin{enumerate}[label=\textbf{\arabic*.}]
\item \textbf{Cálculo de la forma diferencial:}
\begin{equation*}
x \, dy - y \, dx = (a\cos t)(b\cos t \, dt) - (b\sin t)(-a\sin t \, dt) = a b (\cos^2 t + \sin^2 t) \, dt = a b \, dt.
\end{equation*}
\end{enumerate}

\begin{enumerate}[label=\textbf{\arabic*.}]
\item \textbf{Integración:}
\begin{equation*}
\text{Área}(D) = \frac{1}{2} \int_0^{2\pi} a b \, dt = \frac{1}{2} a b [t]_0^{2\pi} = \frac{1}{2} a b (2\pi) = \pi a b.
\end{equation*}
Queda demostrada la clásica fórmula del área de la elipse.
\end{enumerate}

\begin{trapbox}
Recorrer la curva en sentido horario, lo que invierte el signo del área a $-\pi ab$.
\end{trapbox}

\vspace{0.5cm}
\hrule
\vspace{0.5cm}
\end{document}